\chapter{Intervalo}

\lettrine[lines=2, lhang=0.1, nindent=0.2em]{L}{a} lasca se desprendió a las 9:47 y pesó ochocientos veinte microgramos. Vidal la vio caer en el pocillo de la balanza a través del visor del micromanipulador, esperó a que la cifra se estabilizara en pantalla y recién entonces asintió. La conservadora de la Autoridad ---guantes de nitrilo, pulso de cirujana--- transfirió la lasca al vial, lo cerró, lo giró para que él leyera el número de serie en voz alta.


---HA-32-011 ---leyó Vidal---. Sellado 9:49.

Ella fotografió el vial contra la escala de color. Él escribió en el registro: fecha, hora, masa, punto de extracción con coordenadas sobre la ortofoto del óstracon, nombre del operador, nombre del testigo. Su letra salía pareja incluso sobre la mesa plegable, incluso con el generador vibrando a cuatro metros. Un registro con letra dudosa es un registro que después alguien discute, y en este caso el después estaba lleno de gente esperando para discutir.

La carpa olía a lona recalentada y a yeso. Afuera, el valle de Elah era una lámina blanca de mediodía; adentro, treinta y un grados que el aire acondicionado portátil bajaba a lo largo de un gradiente medible: fresco junto al equipo, horno junto a la puerta. Vidal se había vestido igual que todos los días de campo, igual que todos los días: camisa gris, la tercera de las cinco que había traído. El polvo del sitio ya le había cambiado el tono a las dos primeras.

Faltaban dos extracciones. Residuo orgánico adherido a la cara inscrita, doce microgramos si había suerte, para radiocarbono. Y pátina del surco de una letra, la \textit{dalet} de la segunda línea, para verificar que la incisión y la superficie hubieran envejecido juntas. Esa era la pregunta que valía dinero: no cuándo se coció la cerámica, sino cuándo se escribió encima. Una falsificación moderna sobre un tiesto antiguo era el escenario que los diarios llamaban improbable y que él llamaba hipótesis número dos.

---Doctor Vidal.

El profesor Be'eri había entrado agachando la cabeza bajo la lona, con esa manera suya de ocupar espacio pidiendo permiso. Era el director de la excavación, maximalista de tercera generación, y llevaba cuatro días encontrando motivos para pasar por la carpa.

---Estamos en cadena de custodia ---dijo Vidal, sin levantar la vista.

---Dos minutos. ---Be'eri se quedó junto a la puerta, del lado del horno---. Anoche releí la fórmula. \textit{Bet} dudosa o no, la sintaxis es cancillería. Cancillería implica Estado, y Estado en este valle, en el siglo diez, tiene un solo nombre posible. Usted lo sabe.

---Sé que la cerámica es local. Del resto voy a saber un intervalo.

---Acá abajo, a un kilómetro, un pastor volteó a un gigante con una piedra. ---Be'eri sonrió---. Los intervalos también cuentan historias, doctor.

---No exactamente. Las historias las cuenta la gente. El intervalo dice siglo diez más menos cuarenta años, o dice otra cosa, y ahí termina.

Be'eri levantó las dos manos, rindiéndose sin rendirse, y salió. La conservadora, sin mirar a nadie, acercó el vial siguiente.

La segunda emboscada llegó a las cuatro, en el estacionamiento de grava, cuando Vidal cargaba la caja térmica en la camioneta. La doctora Segev enseñaba en Tel Aviv, publicaba contra Be'eri desde hacía veinte años y había manejado una hora para cruzárselo casualmente.

---No le voy a pedir nada ---empezó---. Solo que tenga presente el contexto. Ese óstracon apareció en la campaña más financiada de la década, en el año exacto en que el gobierno necesita que David tenga escritura. Las casualidades existen, pero se datan mal.

---La procedencia está documentada. Estratigrafía, fotos in situ, testigos.

---Los testigos son empleados de Be'eri.

---Por eso la muestra viaja conmigo y el análisis lo hace un modelo que ninguno de los dos entrenó. ---Cerró la puerta de la caja térmica. Verificó el precinto. Lo verificó otra vez, porque el primer control había sido con la doctora Segev hablándole---. Si el resultado le da la razón a usted, va a citar mi metodología. Si se la da a él, va a citar el reglamento europeo y va a decir que Tamiz es una caja negra. Los dos discursos ya están escritos. Yo solo elijo cuál se publica.

Segev se quedó callada un momento.

---Es una manera muy triste de tener razón ---dijo, y era lo primero que decía sin objetivo, y por eso fue lo único de toda la conversación que él siguió repasando en la ruta a Jerusalén, entre el kilómetro treinta y el cuarenta y dos.


\scenebreak


El laboratorio provisional ocupaba una sala en el sótano del edificio de la Autoridad, en Har Hotzvim: paredes recién pintadas, un piso de baldosas grises de treinta por treinta, y los tres racks de Tamiz alineados contra la pared del fondo, respirando con un zumbido parejo que Vidal habría reconocido en la oscuridad. Descargó los viales él mismo. Firmó la recepción él mismo. La técnica local, una mujer de pocas palabras llamada Ruti, cotejó los números de serie contra el manifiesto y le devolvió el registro sin comentarios, que era la clase de trato que él prefería.

La primera pasada arrancó a las 19:20. Extracción de rasgos por canal: espectro de la arcilla, mineralogía de la pátina, geometría de incisión bajo microscopía, y el canal documental masticando ochenta años de cerámica excavada del Levante meridional. Vidal abrió el resumen del corpus de referencia mientras el modelo corría, y ahí, en la tabla de cobertura, encontró el primer placer real de la semana: para cerámica judaíta del siglo diez el corpus era denso, bien fechado, con procedencias limpias ---Qeiyafa, Laquis, Tell es-Safi---, cuatrocientas doce muestras con estratigrafía publicada. Se sirvió café de la máquina del pasillo y volvió a mirar la tabla, despacio, del modo en que otros hombres miraban otras cosas.

---Buen corpus ---le dijo a Ruti, que ordenaba viales.

---Es lo que hay.

---Es más que lo que hay en casi cualquier otro lado. Con esta referencia, el intervalo va a salir estrecho y va a resistir un ataque. Eso casi nunca pasa.

Ruti asintió como quien archiva un dato meteorológico. Vidal se llevó el café a la mesa y dejó que la barra de progreso hiciera lo suyo.

Los canales convergían recién de madrugada. La arquitectura era suya línea por línea, y sin embargo había una asimetría que conocía mejor que nadie: los canales individuales se dejaban auditar ---podía extraer qué banda del espectro pesaba, qué muestra del corpus tiraba del resultado---, pero la capa de convergencia, la que fundía todo y emitía el certificado de confianza global, devolvía sus razones en un formato que ninguna herramienta de interpretación traducía del todo. Funcionaba. Se validaba contra casos conocidos con una tasa de error que avergonzaba al campo entero. Explicarse, eso no. Era el precio de la precisión y él lo había pagado con los ojos abiertos.

Esperar era la parte del trabajo sin procedimiento, así que se fabricó uno. Contó las baldosas de la sala: catorce a lo ancho, veintidós a lo largo, trescientas ocho, menos las cortadas contra la pared del rack, que eran fracciones y exigían criterio. Se podía redondear cada fracción a media baldosa o medirlas una por una. Midió tres con el pie, estableció que el corte era constante, extrapoló. Trescientas quince con un margen de dos. El número quedó fijado y la espera, organizada. Contar era eso: convertir un intervalo muerto en una serie de operaciones verificables. Lo venía haciendo desde los treinta años y le daba mejores resultados que cualquier otra cosa que la gente hiciera para esperar.

A las 21:10 golpearon el marco de la puerta abierta. La directora de distrito, Dalia Ashkenazi, todavía con la credencial colgada, entró y miró los racks como se mira un motor caro que otro maneja.

---¿Cómo vamos?

---Primera pasada limpia. Convergencia esta noche. Mañana a mediodía tengo intervalos preliminares, con la reserva de siempre.

---¿Que sería?

---Que preliminar significa preliminar.

Ashkenazi acercó una silla y se sentó del lado del escritorio que ocupan los que van a decir algo incómodo.

---Mañana a las diez lo quiero en la sede, con las dos comisiones en la sala. Be'eri, Segev, la gente del ministerio. ---Hizo una pausa que había practicado---. Doctor, le voy a hablar claro porque usted factura por claridad. Este país lleva dos años esperando ese tiesto. Hay un discurso escrito en la oficina del primer ministro con espacio para una fecha. El gobierno necesita el resultado.

---El modelo no sabe qué necesita nadie.

---Ya sé. Por eso lo contratamos a usted y no a un patriota. Pero se lo digo igual, porque el que va a estar en la sala mañana es usted, no el modelo, y las salas son mías: sé cómo se ponen.

---Depende de qué llame resultado. Si llama resultado a una datación con mecanismo trazable, se la doy y la defiendo contra quien sea, incluido su ministerio. Si llama resultado a la palabra David, eso no sale de esta sala porque nunca va a entrar. Ninguna fibra de este mundo tiene escrito un nombre propio.

---Es cerámica, no fibra.

---Es una manera de decir.

Ashkenazi se levantó. En la puerta se dio vuelta, y lo que dijo llegó por fuera del guion, con otra voz:

---Mi abuelo excavó en este país con un cepillo de dientes y una Biblia. Encontraba lo que buscaba, siempre. Todo el mundo lo quería. ---Se colgó bien la credencial---. A usted no lo va a querer nadie, doctor, y por eso le pago. Diez en punto.

Los pasos se apagaron en el pasillo. Vidal se quedó frente a la barra de progreso, que iba por el sesenta y uno por ciento, y anotó en el registro la visita, la hora, el requerimiento de la reunión. La letra le salió prolija, como siempre; la releyó igual. Después levantó la vista al piso y contó las baldosas de nuevo, desde el otro extremo de la sala, por control cruzado. Trescientas quince, margen de dos.

El número era el mismo. Con eso alcanzaba para dormir.

\clearpage

\chapter{La unidad de medida}

\lettrine[lines=2, lhang=0.1, nindent=0.2em]{L}{a} sala del tercer piso olía a café de máquina y a alfombra nueva, y el aire acondicionado estaba dos grados por debajo de lo razonable. Vidal llegó a las 9:52, conectó su equipo, verificó que la diapositiva ocho siguiera siendo la ocho. Lo era. La verificó igual, porque la noche anterior había reordenado el bloque de metodología y un archivo reordenado es un archivo que puede fallar.


A las 10:04 Ashkenazi cerró la puerta y las dos comisiones quedaron sentadas a los dos lados de la mesa como si alguien hubiera dibujado una línea en el laminado. Be'eri y los suyos contra la ventana; Segev y dos colegas de Tel Aviv contra la pared; en la cabecera, tres personas del ministerio con carpetas idénticas. Diecinueve en total, sin contarlo a él.

Expuso el diseño en veintidós minutos. Tres extracciones, dijo, todas documentadas: lasca de cerámica para composición y procedencia, residuo orgánico de la cara inscrita para radiocarbono, pátina del surco de la dalet de la segunda línea para verificar que incisión y superficie hubieran envejecido juntas. Cuatro canales de análisis, cada uno auditable rasgo por rasgo. Corpus de referencia de cuatrocientas doce muestras con estratigrafía publicada. Controles ciegos intercalados en la cola de procesamiento, de modo que ni él sabía en qué posición corría el óstracon. Los datos crudos se publicarían completos, junto con el informe, para que cualquier laboratorio del mundo pudiera intentar destruir el resultado. Esa era la parte que más le gustaba explicar y la explicó despacio: un resultado que nadie puede atacar no vale nada; un resultado que todos pueden atacar y sigue en pie es lo único que él firmaba.

Nadie discutió el método. Ni una pregunta sobre el muestreo, ni una sobre los controles. Be'eri asintió durante toda la exposición con la cabeza de quien ya escribió la conclusión.

---Excelente ---dijo cuando Vidal terminó---. Impecable. Ahora, cuando el intervalo confirme el siglo décimo, la fórmula regia queda anclada en contexto de cancillería, y eso, con la sintaxis que ya publicamos\ldots{}

---El intervalo, si es del siglo décimo, ancla la cerámica en el siglo décimo ---dijo Vidal---. La sintaxis es problema de ustedes.

---El contexto es parte de la evidencia, doctor.

---El contexto es parte de su argumento. Mi informe termina donde termina la medición.

Segev se inclinó hacia adelante.

---Ahí lo tienen ---dijo, sin dirigirse a nadie en particular---. Aunque el tiesto sea del siglo décimo, un tiesto del siglo décimo con una fórmula regia prueba que alguien escribía fórmulas regias. El salto de ahí a un reino, y del reino a un nombre, lo van a dar los diarios, no los datos. Y todos en esta mesa lo sabemos y nadie lo va a decir en la conferencia de prensa.

---Eso tampoco es mi problema ---dijo Vidal.

---Va a serlo ---dijo Segev---. Usted firma, doctor. El que firma es dueño del uso que le den a la firma. Pregúntele a los de Basilea.

Uno de los del ministerio, un hombre joven con la carpeta más gruesa, levantó una mano a medias.

---Una consulta operativa. ¿El sistema puede emitir un resultado preliminar acotado, digamos, para uso interno, y el definitivo después? Porque los tiempos de comunicación del gobierno\ldots{}

---No exactamente. Puede emitir un resultado preliminar, que es un resultado con intervalos más anchos. Lo que no puede es emitir dos resultados, uno para adentro y otro para afuera. El preliminar se publica con el definitivo, en el mismo archivo. Si difieren, se ve por qué difieren, línea por línea.

---Entiendo ---dijo el hombre, y anotó algo, y por cómo lo anotó Vidal calculó que había entendido otra cosa.

Ashkenazi miró su reloj y propuso un cuarto de hora de pausa. La línea del laminado se disolvió en grupos junto a la máquina de café.

Ella lo alcanzó en el pasillo, frente a una vitrina con lucernas bizantinas mal iluminadas. Traía dos vasos de cartón y le tendió uno. El café estaba quemado; Vidal lo sostuvo sin beberlo.

---Adentro le fue bien ---dijo ella---. Nadie le tocó el método. Es lo que esperaba.

---El método es sólido. Que nadie lo toque es lo mínimo.

---Nadie le tocó el método porque a nadie le importa el método, doctor. ---Lo dijo sin acidez, mirando las lucernas---. Escúcheme un minuto, porque esto no se lo va a decir nadie más en este edificio. Usted entrega intervalos. Es honesto, es correcto, y yo pago por eso. Pero la gente que espera afuera no vive en un intervalo. Vive en una fecha, en un nombre, en un sí o un no. En algún punto todos eligen, doctor; lo que usted vende es el derecho a elegir más tarde.

---Yo no vendo nada.

---Factura, entonces.

---Facturo mediciones. Lo que cada uno elija hacer con ellas es una operación distinta, posterior, y no figura en mi contrato.

Ashkenazi bebió de su vaso, despacio.

---Sabe qué es lo curioso ---dijo---. Que le creo. Le creo que usted de verdad piensa que esas dos operaciones están separadas.

Volvió a la sala antes que él. Vidal se quedó tres segundos frente a la vitrina. La frase de ella tenía una estructura defendible, casi un teorema: todos eligen, el intervalo solo aplaza la elección. Le buscó el error y tardó en encontrarlo, y lo que encontró era más débil de lo que esperaba: que fuera verdad para la gente de afuera y no dijera nada sobre la medición. Una afirmación puede ser verdadera e irrelevante a la vez. La mayoría lo es. Tiró el café entero en el cesto y entró.

La segunda mitad duró menos y valió menos, hasta que dejó de valer menos. Segev retomó lo de Basilea, esta vez con el expediente abierto sobre la mesa, y ya sin ironía.

---Lo digo como advertencia de procedimiento, no como ataque. En 2027 el modelo certificador de Basilea era el mejor del mercado. Certificados perfectos, trazabilidad perfecta. Y alguien entrenó treinta ídolos cicládicos contra ese modelo hasta que los treinta pasaron. Cientos de millones. Museos, dos fondos soberanos. Mi punto es este: el mejor certificado del mundo es exactamente tan bueno como el secreto de sus pesos, y en Basilea nunca se supo cómo salieron los pesos. Ni quién los sacó, ni por dónde. Seis años y ni una imputación. Entonces cuando el doctor nos dice que su resultado va a resistir cualquier ataque, yo le pregunto: ¿resiste también ese?

---Tamiz tiene un solo dueño ---dijo Vidal---. Los pesos residen en hardware propio, sin réplica en la nube, con acceso de una sola persona. Yo.

---El modelo de Basilea también tenía dueño.

---Tenía dueño corporativo, ciento cuarenta empleados y tres centros de datos contratados. La superficie de exposición era otra.

---Puede ser. ---Segev cerró el expediente---. Solo digo que alguien, en algún lugar, sacó los pesos de un sistema que se vendía como cerrado, y que hasta que sepamos cómo, la palabra cerrado significa menos de lo que todos queremos que signifique.

Vidal abrió la boca para responder y la cerró, porque la respuesta que tenía era una probabilidad, y las probabilidades, dichas en voz alta en una sala así, suenan a excusa. El punto de Segev quedó en la mesa, sin refutar. Lo anotó mentalmente en la columna de asuntos abiertos, que era una columna que prefería vacía.

Entonces habló Halevi.

Había estado callado toda la mañana, en la segunda fila del lado de Be'eri, un hombre de unos sesenta y cinco años con chaqueta de tweed en un país donde nadie usa tweed. Epigrafista. Vidal conocía sus publicaciones: cuarenta años de atribuciones por ojo, por oficio, por conocimiento del trazo. El tipo de carrera que la generación de Vidal había vuelto ornamental.

---Una observación de método, ya que estamos en método. ---La voz era suave---. El doctor nos garantiza que él revisa personalmente cada canal antes de firmar. Me tranquiliza. Porque los sistemas fallan, pero también fallan las personas, ¿verdad? A veces uno firma cosas que avaló otro. Un maestro, por ejemplo. Alguien de mucho prestigio. Y después resulta que había que haber mirado el metal más de cerca. ---Hizo una pausa breve---. ¿O era bronce, aquello? Etrusco, si mal no recuerdo.

El proyector zumbaba. Alguien movió una silla.

Vidal contó las sillas de la sala. Veintidós. Diecisiete ocupadas al empezar, dieciocho ahora, porque un asesor había entrado durante la pausa. Cuatro apiladas contra la pared del fondo, una junto a la puerta con una carpeta encima. Veintidós. La cuenta cerraba. La hizo de nuevo desde la puerta hacia la ventana y volvió a cerrar.

---¿Alguna pregunta sobre el diseño del peritaje? ---dijo.

Halevi sonrió apenas y bajó la vista a sus papeles. Fue Segev, de todas las personas de la sala, la que quebró el silencio.

---Sí, yo. ¿Cuándo tenemos los intervalos?

---El procesamiento cierra esta noche. Consolidación y controles mañana por la mañana; a mediodía entrego el preliminar a la Autoridad, por el canal que fije la dirección.

---Por mí ---dijo Ashkenazi---. Todo por mí. ---Y miró a los tres del ministerio uno por uno, sin apuro, hasta que los tres anotaron algo.

La reunión se levantó a las 12:40. Be'eri le estrechó la mano demasiado tiempo y le dijo que la ciencia y la historia iban por fin a caminar juntas; Vidal respondió que caminarían en paralelo, que era distinto, y Be'eri se rio como si hubiera sido un chiste.


\scenebreak


En el taxi de vuelta a Har Hotzvim abrió el registro de la jornada. Anotó la hora de inicio y de cierre de la reunión, los asistentes que pudo identificar, el compromiso de entrega del preliminar, la objeción de Segev sobre Basilea con la etiqueta de asunto abierto: verificar literatura sobre el vector de extracción de pesos, si existía alguna, que no existía, y esa inexistencia era en sí un dato.

Del comentario de Halevi anotó: ``intervención sin contenido técnico''. Nueve caracteres de más, incluso así. Lo dejó.

El taxi frenó en un semáforo largo de la calle Bar-Ilan. Vidal contó los segundos de la luz roja por hacer algo con la atención: cuarenta y uno. Después, al pie de la página, fuera de las columnas, sin categoría asignada, escribió la frase de la directora completa, con sus palabras exactas, porque una afirmación irrelevante bien construida merece registro aunque no merezca respuesta: \textit{todos eligen; el intervalo aplaza la elección.}

La miró un momento. El semáforo cambió.

No era falsa. Era ajena a su trabajo, que era otra cosa, y en el margen, junto a la frase, dibujó el signo que usaba para las entradas que no requerían acción: un círculo pequeño, cerrado, sin salida hacia ninguna otra celda de la planilla. El taxi arrancó y el registro quedó abierto sobre sus rodillas hasta el laboratorio, donde el procesamiento del óstracon iba, según la pantalla del rack central, por el noventa y cuatro por ciento.

\clearpage

\chapter{Ruido}

\lettrine[lines=2, lhang=0.1, nindent=0.2em]{E}{l} procesamiento cerró a las 21:56, seis minutos después de la estimación de la pantalla, y Vidal anotó la desviación antes de abrir el resultado. Seis minutos en catorce horas de cómputo. Dentro de tolerancia.


El sótano de Har Hotzvim olía a plástico tibio y a solvente de limpieza, y los tres racks emitían ese zumbido de ventiladores que a la segunda semana el oído ya descontaba. Ruti seguía en la mesa lateral, cotejando viales contra el manifiesto, aunque su turno había terminado a las nueve.

Datación primero. El residuo orgánico de la cara inscrita devolvió un intervalo de 1005 a 925 antes de la era común, con un factor de calibración documentado y una confianza de 0,94. Siglo X. Estrecho. El canal mineralógico era consistente: la pátina del surco de la dalet compartía composición y espesor con la pátina de la superficie no inscrita, dentro de un margen que el propio informe desglosaba capa por capa. Incisión y tiesto habían envejecido juntos. La hipótesis número dos, inscripción moderna sobre cerámica antigua, quedaba descartada con los números a la vista. Cualquiera podía repetir el cálculo. Eso era un resultado: una afirmación con su mecanismo adjunto, auditable de punta a punta, y Vidal lo leyó dos veces porque leer dos veces era el protocolo, no porque dudara.

La segunda pantalla era la lectura.

El canal epigráfico había resuelto la secuencia de la segunda línea con confianza alta: dalet, vav, dalet, trazo por trazo, con la geometría de incisión superpuesta a los cuatrocientos doce comparanda del corpus. Hasta ahí, letras. El problema empezaba donde empezaba siempre: en qué eran las letras además de letras. El modelo listaba tres lecturas del contexto sintáctico. Antropónimo regio dentro de una fórmula de cancillería, la lectura de Be'eri, la del discurso con el espacio en blanco: 0,61. Sustantivo común en cadena constructa, sin nombre propio de ninguna especie: 0,33. Daño de superficie en la tercera letra, con reconstrucción alternativa: 0,06. El umbral de atribución fijado en el protocolo, firmado por las dos comisiones en la reunión de la mañana, era 0,90.

No concluyente.

Vidal corrió el diagnóstico del canal epigráfico. Limpio. Revisó la asignación de comparanda, uno por uno, cuarenta minutos. Correcta. El modelo tenía datos suficientes para las letras y datos insuficientes para el rey, y esa asimetría era exactamente la que un instrumento honesto debía exhibir: la piedra dice dalet-vav-dalet; quién era dalet-vav-dalet lo dice el corpus, y el corpus, para fórmulas regias judaítas del siglo X, era una docena de fragmentos disputados y una estela de Tel Dan un siglo más tardía. Con eso el sistema podía proponer, y proponía. Con eso el sistema podía cerrar, y se negaba a cerrar.

Mañana a mediodía, en la sede de la Autoridad, iba a entregar un siglo X con intervalo estrecho y una atribución abierta, y ninguna de las dos facciones iba a irse conforme. El peritaje estaba bien hecho. Eso era otra cosa.

---¿Terminó? ---dijo Ruti. Había apilado los viales en la gradilla y tenía el manifiesto firmado bajo el brazo.

---El procesamiento, sí. La consolidación corre esta noche.

Ella miró la pantalla desde donde estaba, a tres metros, una distancia que no le permitía leer nada, y aun así preguntó:

---¿Es de David o no es?

Vidal tardó un segundo de más. En dos semanas, Ruti había cotejado ciento cuarenta números de serie sin comentar un solo dato.

---Las letras son las que dice el profesor Be'eri. Que nombren a un rey queda por debajo del umbral. La atribución no cierra.

---Ah. ---Se colgó el bolso---. Mi abuela va a decir que sí igual. Con menos letras que eso, ya decía que sí.

Apagó la luz de su mesa y se fue, y el zumbido de los racks ocupó el hueco que dejó la puerta. Vidal se quedó mirando la gradilla. Doce viales. Los contó de todos modos.


\scenebreak


En el hotel, con el escritorio contra la ventana y la calle Jaffa todavía encendida abajo, abrió la plantilla del informe preliminar. Eran las 23:10. El texto técnico salió solo durante media hora: metodología, cadena de custodia de las muestras, tabla de intervalos, desglose por canal. En la sección de atribución transcribió las tres lecturas con sus cifras y escribió que ninguna alcanzaba el umbral pactado, y que el modelo, en consecuencia, se abstenía.

Se abstenía. Releyó el verbo. Un modelo se abstiene cuando la señal existente admite más de una explicación y ningún dato disponible las separa. Eso era comportamiento correcto. Eso era, además, un comportamiento que él había visto antes, en otro caso, con otra diferencia: aquella vez la señal que admitía dos explicaciones tenía un peso de 0,04 sobre el agregado, y quien se abstuvo no fue el modelo.

Dejó el informe abierto y fue a la otra partición. El directorio no tenía nombre: una fecha, 2029-03, y un contenedor cifrado cuya clave existía en un solo lugar, que era él. El respaldo del contenedor viajaba con el resto de sus discos, opaco; podía copiarse mil veces sin que nadie leyera adentro. Tipeó la clave. Dos errores antes de acertar, y tenía la boca seca, y fue a servirse agua del minibar antes de abrir el primer archivo.

Amberes. El díptico. El log crudo de la corrida final, previo a la consolidación que se publicó. Ahí estaba, en la línea 3.407 y las once siguientes: una discordancia menor en el canal de imprimación, una firma espectral en la capa de preparación del panel izquierdo que el modelo marcaba como no explicada por la hipótesis de pastiche del siglo XIX. Peso agregado: 0,04. Todo lo demás, la madera, los pigmentos, la craquelura inducida, el canal documental, convergía en falsificación con una masa de evidencia que ningún tribunal académico había podido tocar. El panel era falso. Lo era. Cuatro años de intentos de refutación lo habían confirmado por seis vías independientes.

La clasificación de la línea 3.407 decía \textit{ruido de sensor, excluir de consolidación}. La fecha de la clasificación era del 14 de marzo de 2029, a las 2:31 de la madrugada. El campo de operador decía SV.

El modelo no había excluido nada. El modelo había levantado la mano y un hombre se la había bajado, y el hombre había tenido razón, y la publicación salió sin la línea, y las dos carreras que el informe destruyó se destruyeron sobre un resultado del que faltaba un renglón. Un renglón irrelevante. Probablemente irrelevante. Con un 0,04 de peso, la palabra probablemente era generosa consigo misma, y él lo sabía en 2029 y lo sabía ahora, y también sabía otra cosa: que la decisión correcta tomada fuera de registro deja de ser auditable, y que una decisión inauditable era, según su propia regla, dicha en congresos, impresa en sus protocolos, otra cosa que verdad.

Bajó el cursor hasta la mitad del log. Lo subió. Cerró el contenedor sin llegar a la línea 3.418, donde terminaba la secuencia discordante, y verificó dos veces que el volumen quedara desmontado.

La regla venía de antes de Amberes. Los bronces habían pasado por el estudio de su director con un dictamen de ojo y una bibliografía de tres generaciones, y Vidal, veintinueve años, tercer autor, firmó la datación sin pedir los datos crudos porque pedirlos habría sido dudar del hombre que le había enseñado a medir. El director murió en enero; el informe holandés que desarmó los bronces salió en abril. La retractación se publicó en junio con una sola firma debajo, la del autor que quedaba vivo.

Volvió al informe preliminar. Eran las doce y cinco. Quedaba la sección de síntesis, dos párrafos que las comisiones iban a leer antes que todo lo demás y en lugar de todo lo demás, y ahí cada sustantivo era una toma de posición. \textit{La inscripción} daba por resuelto lo epigráfico. \textit{El óstracon de la fórmula regia} regalaba la atribución en el nombre. \textit{El hallazgo} sonaba a comunicado de prensa. Escribió: ``El objeto presenta una datación de siglo X a. e. c. con intervalo estrecho; la lectura de la secuencia dalet-vav-dalet es firme a nivel de grafema y no concluyente a nivel de atribución.''

El objeto. Lo releyó. Era la palabra exacta para una cosa que existía con certeza y significaba por debajo del umbral: sin dueño, sin bando, sin rey adentro. Servía para este tiesto y serviría para cualquier cosa que entrara en su cola de procesamiento, una tabla flamenca, un bronce, una tela. Una palabra limpia. La dejó en las cuatro apariciones de la síntesis y la incorporó, sin anotarlo en ninguna parte, al vocabulario de base de sus informes.

Terminó a la una menos veinte. Exportó, firmó con su clave, programó el envío para las 11:50 por el canal de Ashkenazi. Después abrió el registro de jornada para el cierre del día y fue llenando las columnas: hora de resultado, desviación del cómputo, diagnósticos corridos, entrega programada. En la columna de asuntos abiertos escribió ``corpus de fórmulas regias: insuficiente, no subsanable en plazo''. Quedaba una entrada por hacer y el cursor se detuvo sobre la celda. Cuarenta minutos de esta noche correspondían a un directorio de 2029 que ningún proyecto activo justificaba. El registro tenía una categoría para eso desde octubre, desde una fila frente a una tumba sin nada adentro: tiempo perdido. Escribió la duración, 0:40, y en el campo de descripción, después de considerarlo, puso solamente ``archivo''.

Cerró la laptop. La calle seguía encendida y un camión de reparto descargaba cajones frente al hotel, a la una de la mañana, y Vidal se quedó de pie junto a la ventana mirando la operación completa: el hombre bajaba un cajón, lo apilaba, volvía por otro, con un ritmo de dieciocho, veinte segundos por ciclo. Diecinueve. Dieciocho. Le vino la frase de Ruti, la abuela que decía que sí con menos letras, y la examinó un momento con la misma atención que al camión. La abuela no estaba cometiendo un error de cálculo, porque no estaba calculando; hacía otra operación, sin intervalo, sin umbral, y esa operación quedaba fuera de su incumbencia igual que la frase de la directora en el taxi, igual que la fila de octubre.

El estómago le hizo un ruido largo. Buscó cuándo había comido por última vez y encontró el desayuno, 7:15, un café y media tostada, y nada después: catorce horas de cómputo también para él, sin registro. Llamó a recepción y pidió lo que hubiera. Había sopa. Que fuera sopa.

\clearpage

\chapter{Día sin datos}

\lettrine[lines=2, lhang=0.1, nindent=0.2em]{B}{e'eri} llegó a la presentación con el preliminar impreso y anotado en dos colores, y eso ya era un dato: lo había leído entero entre las 11:50 del envío y las 9:00 de esa mañana. Segev llegó con la laptop bajo el brazo y eligió la silla más lejana a Be'eri sin mirar dónde se sentaba Be'eri, lo que exigía haberlo calculado antes de entrar. Halevi ocupó el mismo lugar que en la reunión anterior, contra la ventana, con la misma chaqueta. Ashkenazi cerró la puerta y dijo que había café en el pasillo y que nadie iba a salir a buscarlo hasta que terminaran.


Vidal proyectó tres diapositivas. Podría haber sido una, pero tres distribuían mejor la discusión: cada facción iba a querer pelear una distinta. La primera: intervalo 1005--925 antes de la era común, confianza 0,94, pátina del surco solidaria con la superficie no inscrita. La segunda: secuencia dalet-vav-dalet resuelta con confianza alta. La tercera: atribución repartida en tres hipótesis, 0,61 la mejor, umbral de atribución firmado por ambas comisiones en 0,90.

---Mil cinco a novecientos veinticinco ---dijo Be'eri. Se puso los anteojos para leer en la pantalla lo que ya tenía subrayado en el papel---. Dígame qué otra cancillería escribe en ese intervalo, doctor. Una. Nómbreme una.

---El intervalo fecha la cerámica. La atribución es otra operación.

---Es que no se puede separar la fecha de la...

---Se puede. Es lo que acabo de mostrar.

Segev giró la laptop hacia la mesa como si la laptop fuera un testigo.

---Cero sesenta y uno no es una lectura, es una moneda cargada. Con ese número, en mi departamento, un doctorando repite el análisis.

---No exactamente. Es la hipótesis mejor puntuada de tres, con el reparto completo publicado. Un doctorando que repita el análisis va a obtener el mismo reparto. Esa es la diferencia con una moneda.

---El umbral del 0,90 lo firmamos para tranquilizarla a ella ---dijo Be'eri, y señaló a Segev con los anteojos---. Si hubiéramos sabido que la máquina iba a quedarse a tres décimas...

---Lo firmaste porque tu propio epigrafista te dijo que... ---empezó Segev.

---Firmaron los dos ---dijo Ashkenazi---. Está en acta. Siga, doctor.

Vidal siguió. Cuarenta minutos de preguntas, ninguna nueva: todas estaban contestadas en el informe, pero un informe leído en silencio no le sirve a nadie que necesite que lo vean objetar. Be'eri quería que el intervalo arrastrara la atribución; Segev quería que la atribución débil contaminara el intervalo; el trabajo de Vidal consistía en impedir ambos movimientos, que eran el mismo movimiento en direcciones opuestas. Halevi habló una sola vez, al final, cuando ya juntaban papeles.

---La dalet está bien leída ---dijo, sin dirigirse a nadie---. Su máquina lee mejor de lo que atribuye. En eso se parece a mis alumnos.

Nadie respondió y él pareció conforme con eso.

En el pasillo, Ashkenazi caminó a su lado hasta el ascensor. Llevaba la copia de Be'eri, que Be'eri había olvidado o abandonado sobre la mesa, y la iba hojeando por los subrayados ajenos.

---¿Sabe lo que hizo hoy? Dejó insatisfecha a toda la gente importante de este edificio en proporciones casi idénticas.

Dos facciones descontentas en magnitudes comparables: en el oficio de Vidal, esa era la firma de un arbitraje bien hecho, y estuvo a punto de decirlo así. La ecuanimidad no se declara; se mide en el reparto del malestar.

---El informe sostiene las dos cosas a la vez ---dijo---. La fecha cierra, la atribución queda abierta. Cualquiera de los dos que cite una sola mitad va a estar citando mal.

---Van a citar una sola mitad. ---Ashkenazi apretó el botón del ascensor y lo miró---. Hay un discurso escrito en la oficina del primer ministro con un espacio en blanco. El espacio era para una fecha. Usted acaba de darles una fecha con un cero noventa y cuatro adelante, firmada por el único perito del mundo al que ninguna de las dos comisiones puede acusar de nada. Van a decir ``la época del rey David'' y, técnicamente, van a tener razón. El intervalo alcanza para eso.

---El informe dice que la atribución no...

---El informe lo van a leer doce personas. El intervalo lo va a leer el mundo. ---Le devolvió la copia subrayada de Be'eri, como si fuera de Vidal, y en cierto modo lo era---. Quería que lo supiera antes de irse. Usted les dio menos de lo que pedían y les sirvió igual. Eso, acá, es lo más parecido al éxito que va a conseguir nadie.

El ascensor llegó. Vidal bajó solo, con la copia ajena en la mano, releyendo los subrayados de Be'eri: todos en la primera diapositiva, ninguno en la tercera.


\scenebreak


Esa noche abrió la planilla de asignación. Vuelo el sábado a las 6:40; entre el cierre de la jornada y el embarque quedaban dos días. Los contó en horas de vigilia: treinta y una, redondeando la eficiencia de sueño al promedio del mes. El viernes por la mañana tenía destino: supervisar el embalaje de los tres racks, cotejar el manifiesto de salida contra el de entrada, firmar la cadena de custodia del equipo. Cuatro horas, cinco con la aduana. El resto era una columna en blanco, y una columna en blanco era una anomalía de proceso: desde 2029 cada bloque de su tiempo tenía un código. Escribió en el jueves, de 8:00 a 18:00, ``reconocimiento urbano''. El código existía en su catálogo desde un congreso en Atenas donde había perdido un día por una huelga; designaba caminata con registro de observaciones, sin entregable. Diez horas asignadas. La columna cerró.

El jueves salió a las 8:04 con el teléfono cargado y un objetivo formulado como lo habría formulado un topógrafo: relevar la superposición.

La ciudad vieja se dejaba leer. Eso fue lo primero que registró, bajando hacia la puerta de Damasco entre carros de pan con sésamo y un olor a cardamomo quemado que salía de un local sin nombre: la ciudad se dejaba leer porque estaba construida como un corte estratigráfico puesto de pie. Bajo la puerta otomana, tres metros más abajo del nivel de calle actual, asomaba el arco romano de Adriano, y la diferencia de cota era información pura: dieciocho, veinte siglos de escombro compactado, a razón de quince centímetros por siglo, un valor razonable para una ciudad destruida con frecuencia. Caminó el cardo bizantino, seis metros bajo la rasante, con las columnas todavía en pie dentro de un pozo techado, y anotó la cota en el teléfono. Ninguna capa había borrado a la anterior. Ese era el rasgo notable. En cualquier otro sitio, construir era demoler; acá, construir era escribir encima sin tachar, y el resultado era un archivo donde cada época seguía siendo legible bajo la siguiente. La metáfora se le instaló sin resistencia, cosa rara: su cabeza rechazaba las metáforas como rechazaba los adjetivos, pero esta era técnicamente exacta. Superposición, terminus post quem, lectura por capas. La ciudad era un archivo. Podía trabajar con eso.

Lo segundo que registró fue el uso. En el muro occidental contó las hiladas desde la plaza: las siete inferiores herodianas, con el almohadillado característico, después cuatro de bloque menor, romano tardío o más probablemente omeya, después el remiendo mameluco, después lo otomano arriba de todo, y frente al muro un cartel en tres idiomas explicaba qué hilada probaba qué derecho. A doscientos metros, en la Vía Dolorosa, una placa franciscana citaba un pavimento excavado como fundamento de la estación. Más allá, un panel del waqf databa una fachada. Cuatro barrios, y en los cuatro alguien había convertido una capa del archivo en escritura de propiedad. Le resultó familiar con una precisión incómoda: era exactamente el uso que el discurso con el espacio en blanco iba a darle a su intervalo. Acá llevaban siglos haciéndolo. Las piedras eran el registro y el registro era el título, y cada confesión leía la capa que le convenía y llamaba ruido a las demás.

Al mediodía el muecín y las campanas de una iglesia se superpusieron durante cuarenta segundos, cada uno en su frecuencia, sin interferir. Dos señales simultáneas en el mismo canal, ninguna corrompida. Se quedó parado en una esquina del barrio cristiano hasta que terminaron las dos.

Los escalones del zoco estaban pulidos hasta un coeficiente de fricción que ya era un riesgo, y por ellos bajó, entre puestos de especias y de camisetas, siguiendo la pendiente, hasta que la calle desembocó en un patio de piedra clara y el teléfono confirmó lo que ya sabía desde octubre: basílica del Santo Sepulcro. Fachada cruzada del siglo doce sobre restos bizantinos sobre la plataforma del templo de Adriano: tres capas legibles a simple vista, coherentes con el patrón general del archivo. En la cornisa, bajo una ventana del segundo nivel, había una escalera de madera apoyada. Consultó el teléfono: llevaba ahí desde el siglo dieciocho, porque moverla exigía el acuerdo de las seis confesiones que compartían el edificio, y el acuerdo nunca se había producido. Vidal la miró un rato más largo del que la escalera justificaba. Un registro de cambios congelado durante doscientos setenta años porque el costo del consenso superaba el costo de la anomalía. Conocía comisiones así. Había estado con dos de ellas el miércoles.

Por la puerta única entraba y salía gente en flujo desparejo: un grupo de monjas con el paso coordinado, hombres solos, una familia con un cochecito que hubo que plegar. Adentro habría más capas, las más disputadas del edificio. Su recorrido del día estaba completo en un noventa por ciento y el interior de un templo en uso no era relevable con la metodología del resto: demasiada gente, ninguna cota accesible. Marcó la ubicación en el mapa del teléfono ---31,7785 norte, 35,2296 este, el punto cayó sobre un techo gris entre dos cúpulas--- y emprendió la vuelta por el barrio cristiano, cuesta arriba, contra el flujo de la tarde.

En el hotel volcó las observaciones: cotas, hiladas, tiempos de marcha. Veintiséis entradas. La última era el punto marcado, que quedó en la lista sin descripción, porque la descripción disponible ---basílica, tres capas, escalera--- ya estaba dicha en las entradas anteriores y el punto en sí no agregaba dato. Lo dejó igual.

El viernes por la mañana los racks entraron en sus cajas en el orden inverso al del armado, y Ruti cotejó los cuarenta y un números de serie del manifiesto de salida contra el de entrada, en voz alta, sin saltearse ninguno, y a las 12:20 la cadena de custodia estaba firmada y el laboratorio provisional era otra vez un sótano con trescientas quince baldosas. Almorzó de pie. Volvió al hotel a las 13:40.

Y ahí estaba la celda: viernes, de 14:00 a 20:00, sin código.

Escribió ``reconocimiento urbano (continuación)'' y lo miró. Un segundo relevamiento del mismo sector, con los mismos instrumentos, producía datos redundantes: eso tenía nombre en su oficio y el nombre era desperdicio. Lo borró. Probó ``preparación de viaje'': la valija de un hombre con cinco camisas iguales se preparaba en once minutos, y once minutos no llenaban una celda de seis horas. Lo borró también. Revisó el correo, que ya había revisado. El cursor quedó parpadeando sobre el blanco, un pulso por segundo, y Vidal lo dejó parpadear.

Era, hasta donde su propio registro alcanzaba, la primera celda vacía desde 2029.

Sobre el escritorio, el teléfono seguía con el mapa abierto de la tarde anterior. El punto estaba donde lo había dejado, sobre el techo gris entre las dos cúpulas, sin etiqueta, y quedaban seis horas de luz.

\clearpage

\chapter{Cuarenta y una personas}

\lettrine[lines=2, lhang=0.1, nindent=0.2em]{L}{a} escalera seguía en la cornisa, apoyada contra la ventana de la derecha, y Vidal la registró al entrar al atrio como se registra un testigo de nivel: si alguien la movía, algo se había movido en el edificio de las seis confesiones. Las 14:37. Cruzó el atrio leyendo la fachada por fases constructivas ---jambas cruzadas del XII, dintel expoliado, remiendos del terremoto del 27--- y pasó la puerta con el método de siempre, el de cualquier yacimiento en uso: primero la planta, después la circulación, después la gente.


La planta era conocida por bibliografía. La circulación era un caos administrado, seis liturgias compartiendo un volumen sin pasillos. La gente era una fila.

Empezaba en la puerta del edículo, bajo la rotonda, y daba media vuelta alrededor de la estructura antes de deshacerse en el borde, donde los que llegaban dudaban un momento y después se sumaban. Vidal se ubicó junto a una columna, a distancia de observación, y contó. Cuarenta y una personas.

Se detuvo en el número. Esa misma mañana Ruti había cotejado cuarenta y un números de serie contra el manifiesto, en voz alta, y un número usado hace pocas horas contamina el conteo siguiente: sesgo de anclaje, modo de fallo documentado de la atención humana. Volvió a contar desde la puerta, esta vez señalando con la mirada cabeza por cabeza. Cuarenta y una. Mientras terminaba se sumaron dos más al final, de modo que el dato ya era viejo al producirse, pero el dato original quedaba: cuarenta y una personas entre él y la puerta baja del edículo cuando había empezado a mirar.

Midió el avance con el reloj y una junta del pavimento como cota. Un metro cada tres minutos, con desviaciones cuando el monje de la puerta hacía pasar grupos de a tres. Y en el panel de la entrada, en tres idiomas, estaba escrito lo que había adentro: la tumba está vacía.

Ese era el dato que le sobraba. En su campo, una tumba vacía era el caso de manual de evidencia negativa: la ausencia de un cuerpo prueba una resurrección tanto como la ausencia de cerámica prueba un incendio, es decir, nada. Probó la primera categoría y la descartó: nadie en la fila parecía estar ahí para evaluar evidencia. Probó la segunda, sesgo de muestra: peregrinos, población autoseleccionada, cero grupo de control; correcto y a la vez inútil, porque el sesgo explicaba quiénes venían, y su pregunta era otra. Probó la tercera, valor simbólico: la gente hace fila ante monumentos sin contenido, tumbas de soldado desconocido, llamas votivas, y eso funcionaba casi, casi cerraba, salvo por un detalle. El soldado desconocido está adentro. Acá el objeto de la fila era la ausencia misma, exhibida, con panel en tres idiomas. Tres categorías, tres residuos. La fila seguía avanzando un metro cada tres minutos y ninguna de sus casillas la contenía entera.

Se quedó mirando la puerta.

Era baja, de mármol claro, y la gente salía de a una, doblada, y se enderezaba en la luz de la rotonda con un segundo de ajuste, como quien sale de un pozo de sondeo. Vidal eligió una para observación: una mujer de unos cincuenta años, chaqueta gris, que salió, se apartó dos pasos del flujo, y se detuvo. Esperaba de ella alguna conducta clasificable. Llanto, teatro, fotografía. La mujer recogió una mochila que había dejado con otra del grupo, se la colgó de un hombro, dijo algo breve que hizo asentir a la otra, y le quedó en la cara una expresión que Vidal conocía de su propio laboratorio y de ningún otro lugar: la de un operador cuyo resultado coincidió con la predicción. Repasó desde sus seis metros de distancia la lista de errores de método posibles. Muestreo, expectativa, instrumento contaminado. Desde donde estaba, la mujer no exhibía ninguno. Había esperado una hora para ver un lugar donde el panel mismo declaraba que su contenido era cero, lo había visto, y se iba con la cara de quien verificó algo.

Datos insuficientes a distancia. En su práctica eso tenía un protocolo, uno solo: la discrepancia que resiste la clasificación remota se inspecciona en persona. Fue hasta el final de la fila y se sumó. Puesto cuarenta y cuatro.

La espera fue una serie de mediciones porque las manos necesitaban algo. El pavimento bajo la fila estaba gastado en cóncavo, dos o tres centímetros de flecha en las losas más viejas: décadas de esa misma fila, quizá siglos, un desgaste auditable si alguien quisiera datarlo por tasa de abrasión. Olía a cera y a resina quemada, un incienso denso que bajaba de la rotonda, y desde el catolicón llegaba un canto grave en griego, sostenido, sin melodía que él pudiera retener. Cada tres minutos, un metro. El monje de la puerta administraba el ingreso con un gesto de dos dedos y una consigna repetida en inglés: two minutes. Delante de Vidal, un grupo hablaba en español, acento centroamericano, y la mujer mayor del grupo repartía velas finas de un manojo envuelto en papel. Le tendió una a él.

---Tome. Para adentro.

---Gracias ---dijo Vidal, y la vela quedó en su mano porque devolverla exigía una conversación.

La mujer lo miró un momento más de lo necesario.

---¿Primera vez?

---Sí.

---Se nota.

Se dio vuelta hacia los suyos sin esperar respuesta. Vidal revisó el reloj, que ya había revisado, y dedicó los siguientes metros de avance a la pregunta de en qué era observable, exactamente, que fuera la primera vez. Su ropa era neutra. Su conducta, contenida. Contaba losas del piso, quizá; movía los labios al contar, era posible; nunca se había filmado contando. Un hombre que se creía ilegible acababa de ser leído por una peregrina en cuatro segundos, con confianza suficiente para decirlo en voz alta, y el diagnóstico era correcto. Anotó mentalmente el episodio bajo ``calibración externa'' y la etiqueta le duró poco, porque la fila dobló la última esquina del edículo y la puerta quedó a la vista.

De cerca, la estructura era legible: mármol rosado del rebozado del XIX, limpio desde la restauración, juntas nuevas sobre fábrica vieja, todo posterior en muchos siglos a lo que decía contener. El monje le hizo el gesto de dos dedos a él y a la pareja de adelante.

Adentro había una antecámara con una lámpara y un pedestal con un fragmento de piedra bajo vidrio ---rotulado como parte de la piedra de cierre, sin cadena de custodia que él conociera---, y después la segunda puerta, más baja todavía, un metro treinta a ojo. Vidal se dobló casi en dos para pasar. Se había fijado el procedimiento en la fila: ciento veinte segundos, conteo interno, inventario de superficies, salida. Empezó a contar al cruzar el umbral.

La cámara medía unos dos metros por dos. A la derecha, la losa: mármol claro, pulido por manos, con una grieta transversal que según una fuente del XVI se había hecho a propósito para desalentar el saqueo, dato imposible de verificar como casi todo ahí. La losa era un revestimiento de 1555. Debajo del revestimiento, el banco de roca. Debajo del banco de roca, la afirmación. Cada superficie visible era posterior al evento conmemorado por quince siglos como mínimo; en términos de su oficio, el interior del edículo contenía cero información sobre su propia pregunta, y la pregunta ni siquiera estaba bien formulada, porque una tumba vacía es compatible con todas las hipótesis, incluida la trivial. Trece lámparas colgaban sobre la losa. Las contó. La pareja que había entrado con él estaba arrodillada y la mujer tenía la frente apoyada en el mármol, y el mármol, donde el codo de Vidal lo rozó sin intención, estaba tibio. Calor de manos acumulado, calorías de fila, transferencia perfectamente explicable. Golpearon la puerta desde afuera.

Salió doblado, se enderezó en la rotonda, y buscó el número del conteo interno para cerrar el procedimiento.

El número faltaba.

Sabía que había empezado a contar en el umbral. Recordaba el once, o el trece. Después el registro estaba en blanco hasta el golpe en la puerta, y el reloj decía que adentro habían pasado poco más de dos minutos, ciento treinta segundos quizá, consistente con la consigna del monje, de modo que la duración total estaba verificada por instrumento externo y solo faltaba la serie interna. Volvió a mirar el reloj. La misma hora, más doce segundos. En treinta años de conteos ---filas, baldosas, semáforos, números de serie--- la serie se había interrumpido por sueño o por interrupción ajena, con causa asignable en cada caso. Repasó las candidatas: la losa tibia, las trece lámparas, la frente de la mujer contra el mármol. No cerró en ninguna.

Abrió la nota del teléfono para volcar la observación en caliente, como correspondía. Escribió: ``edículo, interior, 2 min:". El cursor parpadeó después de los dos puntos. Estuvo un rato con el teléfono en la mano, el tiempo suficiente para que dos grupos más entraran y salieran, y al final guardó la nota así, con los dos puntos abiertos, porque borrarla habría sido falsear el registro y completarla habría requerido un dato.

Todavía tenía la vela. Delante de los candelabros de arena, donde ardían decenas iguales a la suya, la plantó derecha entre las otras y la dejó sin encender. Le pareció el único acto disponible que no afirmaba nada en ningún sentido, y recién al alejarse se le ocurrió que una vela apagada en un lecho de velas encendidas era, para cualquier observador, la afirmación más visible de todas.

Se quedó veinte minutos más de lo planificado, junto a la misma columna del principio. Contó la fila de nuevo: treinta y siete. Otros integrantes, otro número, la misma forma doblando alrededor del edículo a un metro cada tres minutos. La fila no era la gente. La gente se recambiaba entera y aquello persistía, con la estabilidad de un proceso y sin el mecanismo de ninguno. A las 17:05 salió por el atrio, la escalera seguía en su cornisa, y cuesta arriba hacia la puerta de Jaffa lo alcanzó la sirena de Shabat, larga, sostenida sobre los techos, mientras los comercios bajaban las persianas casi al unísono, como un sistema entrando en modo seguro.

En el hotel, a las 18:40, abrió la planilla de asignación. La celda del viernes seguía en blanco, 14:00 a 20:00. Probó ``reconocimiento urbano'': el código exigía registro de observaciones, y el registro de la tarde consistía en una nota con dos puntos y nada después. El código quedaba fuera por definición propia; usarlo habría sido el tipo de contabilidad indulgente que él facturaba caro cuando la encontraba en peritajes ajenos. Quedaba la otra columna, la de las entradas con círculo al margen, las que no requieren acción. Escribió: ``basílica del Santo Sepulcro, 14:37--17:05, sin entregable''. Columna: tiempo perdido. El círculo al margen.

La planilla vivía donde había vivido siempre, en una carpeta local de 2029 que nunca había migrado al esquema de respaldo; a las 20:00 el sincronizador barrería todos sus archivos hacia el servidor de Zúrich y pasaría de largo junto a ese, como cada noche, porque nadie ordena la copia de seguridad del tiempo que perdió. Las reglas de su propio sistema eran claras: un archivo sin respaldo se migra o se elimina. Miró la entrada nueva, la fecha, las horas, la letra limpia de la celda. Once kilobytes en total.

Cerró el archivo sin hacer ninguna de las dos cosas.

\clearpage

\chapter{Bibliografía hostil}

\lettrine[lines=2, lhang=0.1, nindent=0.2em]{L}{a} búsqueda llevaba abierta cuarenta minutos cuando bajó el brillo de la pantalla por segunda vez. El gesto era absurdo y lo sabía: el asiento 11A daba a la ventanilla, el de al lado estaba ocupado por una mujer que dormía con la boca entreabierta, y nadie en el vuelo de las 6:40 tenía interés en lo que un pasajero leyera sobre el Mediterráneo. Bajó el brillo igual. Después corrigió los términos de búsqueda, que era la parte que más se parecía a esconder algo: donde había escrito \textit{resurrección --- evidencia} puso \textit{sábana de Turín, historiografía crítica}, y con la sintaxis nueva la pregunta ya podía leerla cualquiera por encima de su hombro sin que significara nada de él.


Los resultados llegaron ordenados por relevancia y la relevancia, para variar, tenía razón. Primero: 1389, memorial del obispo Pierre d'Arcis al papa de Aviñón. Un artista había confesado haber pintado la tela. La confesión constaba; el nombre del artista, no constaba. Segundo: 1988, datación por radiocarbono en tres laboratorios independientes, Oxford, Arizona, Zúrich. Intervalo combinado 1260--1390. Los tres intervalos se solapaban. El del laboratorio de Zúrich se había medido a veinte minutos de tranvía de su estudio.

Suficiente. Una confesión del siglo catorce y una datación triple con solapamiento: cualquier peritaje del mundo cerraba con menos. Lo anotó así, mentalmente, con la fórmula de sus informes: \textit{evidencia convergente, dos vías independientes, caso resuelto en 1988}. Cuarenta y cuatro años de resuelto.

Compró el paquete de datos del vuelo y descargó el memorial completo.

Mientras el archivo bajaba ---lento, el satélite repartía ancho de banda entre doscientas personas--- se dijo que era higiene profesional: uno cita la fuente primaria o no cita. El documento se abrió sobre el mar, en latín, cuarenta y dos páginas escaneadas de una edición de 1900. Leyó despacio, con el traductor al costado para las abreviaturas. Y a la tercera página el perito que llevaba adentro, el que facturaba caro precisamente por esto, empezó a marcar márgenes. La copia sobreviviente era un borrador. Sin sello, sin fecha cierta, sin firma. La confesión del artista era de segunda mano: d'Arcis decía que su predecesor la había obtenido, treinta y tantos años antes. Y d'Arcis tenía un litigio abierto con la colegiata que exhibía la tela, un litigio por ofrendas, es decir por dinero. Testigo con interés. Declaración de oídas, sin cadena de custodia documental. Si un cliente le llevaba ese memorial como prueba de falsificación, él lo aceptaba como indicio y le cobraba por decirle que no alcanzaba solo.

Alcanzaba con el carbono. El carbono no tenía litigios. Cerró el memorial y dejó la mano sobre el trackpad más tiempo del necesario.

---¿Eso es la Síndone? ---dijo la mujer de al lado. Se había despertado sin ruido y miraba la pantalla, donde la página de resultados seguía abierta con la miniatura de un rostro en negativo.

---Es un documento del siglo catorce.

---Mi marido la vio en 2015. Tres horas de fila para pasarle por delante, ¿cuánto habrá sido?, medio minuto.

Vidal giró apenas la pantalla hacia ella, lo contrario de lo que había estado haciendo toda la mañana.

---¿Cuánto tiempo estuvo delante? Exacto, si se acuerda.

---Exacto no sé ---dijo la mujer---. Poco. De la tela no se acuerda: dice que era oscura y que estaba lejos. De la fila se acuerda de todo. Había un señor con un termo que convidaba café a desconocidos.

Ella volvió a acomodarse contra el respaldo y cerró los ojos. Vidal se quedó con el dato donde no tenía columna: un hombre con un termo, y nada del objeto que había cruzado el planeta para ver.


\scenebreak


El estudio de Zúrich olía a los racks recién reinstalados, ese olor de electrónica tibia y polvo de embalaje que tardaba una semana en irse. Afuera llovía desde el martes. Vidal trabajaba de día en lo pendiente ---el informe final del óstracon, dos consultas de museos, la agenda de noviembre--- y a las diez de la noche, cuando el edificio se quedaba en silencio salvo el radiador y los ventiladores, abría la otra carpeta.

Secondo Pia, mayo de 1898. Primera fotografía de la tela, placa de vidrio. Al revelar, el negativo mostró un positivo: la imagen sobre el lino ya estaba invertida, luces por sombras, como un negativo fotográfico fabricado cinco siglos antes de la fotografía. A Pia lo acusaron de retocar la placa. Vidal buscó las réplicas: 1931, otro fotógrafo, otras placas, mismo resultado. El dato era limpio y era raro, y la rareza limpia era su materia prima desde hacía veinte años. Un pintor del siglo XIV podía pintar un negativo, en teoría; lo que costaba explicar era para quién, dado que nadie iba a poder verlo derecho durante quinientos años.

Después, 1978. El STURP: más de treinta científicos, ciento veinte horas de acceso directo, espectrometría, fluorescencia, radiografía, cinta adhesiva sobre la superficie para levantar fibras. La bibliografía posterior era un delta que salía entero de ese río: cada artículo de los últimos cincuenta y cuatro años, a favor o en contra, masticaba datos de esas ciento veinte horas, porque ningún equipo había vuelto a tener nada comparable. Los sindonólogos vivos databan sus carreras con ese reloj. Había hombres de ochenta años cuya autoridad entera consistía en haber estado en esa sala una semana de octubre, y hombres de sesenta cuya amargura entera consistía en haber llegado tarde.

La conclusión oficial del proyecto cabía en una línea: la imagen no es pintura; mecanismo de formación desconocido.

Leyó la línea dos veces. Mecanismo desconocido, en un informe final, con membrete. Un informe serio decía \textit{datos insuficientes} y especificaba qué medición faltaba, con qué instrumento, a qué costo; \textit{desconocido} era una tarea pendiente redactada como conclusión, el equivalente metodológico de una falta de ortografía en la primera página. Le molestó de un modo físico, en la mandíbula. Instrumentos de 1978. Sensores de 1978, potencia de cómputo de 1978, y desde entonces todo el mundo discutiendo la interpretación de esas cintas adhesivas en lugar de volver a medir. El problema llevaba cincuenta y cuatro años abierto por la razón más barata de todas: nadie lo había medido bien.

Abrió el archivo de notas que había creado el domingo ---se llamaba, todavía, ``lectura''--- y escribió debajo de lo anterior: \textit{El mejor argumento contra la autenticidad no es el carbono. Es el silencio. Un objeto de esa importancia deja rastro: inventarios, cartas, recibos, pleitos. Trece siglos sin un recibo. La historia documental empieza en 1355 en Lirey y nadie explica de dónde salió.} Releyó el párrafo. Era lo más sólido que había escrito en tres noches, más sólido que el memorial y que los tres laboratorios, y guardó el archivo con la sensación de haber adelantado trabajo para un informe que nadie le había encargado.

También descargó el corpus fotográfico de referencia, el moderno: dieciséis mil placas de alta resolución, hilo por hilo, publicadas entre 2004 y 2019 por un químico de la Comisión de Conservación de Turín. Ferrero, O. La documentación técnica era impecable, calibración cromática incluida, y las notas al pie usaban con la misma naturalidad la palabra \textit{urdimbre} y la palabra \textit{reliquia}, sin comillas ninguna de las dos, cosa que a Vidal le produjo el desconcierto leve de un instrumento que mezcla unidades y aun así devuelve valores correctos.


\scenebreak


La ingesta la armó un viernes a la noche, como armaba todas: manifiesto de fuentes y pesos por canal. Trazabilidad de cada archivo, sin excepción. Las placas de Pia digitalizadas, las de 1931, las dieciséis mil de Ferrero, las fotografías de la intervención de 2002, los subconjuntos publicados de los espectros del STURP, los datos crudos de 1988 liberados décadas después por insistencia de terceros. Todo lo que existía en dominio público sobre la tela entró a Tamiz en una cola de procesamiento de nueve horas.

El resultado ocupaba tres líneas.

\textit{Canal de imagen: cobertura suficiente, resolución degradada por generaciones de copia; conclusiones limitadas a morfología.} \textit{Canales espectral y material: corpus de entrada incompatible; archivos sin calibración de origen; cadena de custodia no documentada.} \textit{Datación: datos de 1988 internamente consistentes; imposible evaluar representatividad de la muestra sin adquisición directa.}

Y debajo, en el campo de recomendación que el propio Vidal había diseñado para que el modelo nunca callara: \textit{se requiere adquisición directa}.

Corrió la cola de nuevo con los umbrales relajados. Nueve horas. Mismas tres líneas. Relajó los umbrales un paso más, hasta un nivel que en un peritaje real le habría dado vergüenza, y el modelo, con esa lealtad suya que era lo único parecido a un carácter, devolvió lo mismo con una línea extra: \textit{advertencia: configuración fuera de rango de validez}. Fotografías de fotografías. Espectros huérfanos. Una datación de hace cuarenta y cuatro años, sobre una esquina de la tela. Con eso el mejor sistema de autenticación existente decía lo que un becario honesto habría dicho gratis: para saber algo hay que tocar el objeto.

Buscó quién controlaba el acceso. La respuesta era corta: la Custodia pontificia, operativamente la Comisión de Conservación, Turín. Entre los miembros de la Comisión figuraba el autor del corpus fotográfico. La tela estaba plana, en atmósfera inerte, bajo el crucero de la catedral, desde 1997; el último acceso científico con instrumentos, 1978. Abrió un documento en blanco y escribió una línea: \textit{Solicitud de acceso con fines de estudio.} La miró. La solicitud correcta se presenta en persona, con protocolo adjunto. Eso también lo sabía de su oficio: por correo se rechaza fácil.

El lunes siguiente le tocaba el cierre administrativo del mes, la única rutina de su vida que odiaba sin matices. Facturación: las horas de octubre y lo corrido de noviembre, cada bloque con su código, cada código con su cliente. La Autoridad de Antigüedades, cerrado. Las consultas de museos, cerradas. Y después la planilla siguió mostrando entradas, una por noche, veintiuna en total desde el regreso, de 22:00 a la una, a las dos, una vez hasta las 3:40. El campo de código estaba en blanco en todas. Ochenta y siete horas y media.

Se quedó mirando la cifra. Verificó la suma a mano, celda por celda, aunque la planilla sumaba sola desde hacía años. Ochenta y siete horas y media. Desde 2029 cada bloque de su tiempo tenía código; la regla era suya y la había cumplido mejor que ninguna otra. Tres semanas de noches trabajadas con el rigor de un encargo ---manifiesto de fuentes, control de versiones, hasta la advertencia de rango--- y ni una línea facturable, ni un cliente, ni un entregable definido. Si un colega le mostraba esa planilla, el diagnóstico le habría llevado un segundo: eso ya no es una lectura. Eso tiene forma de proyecto.

Creó un código nuevo. El catálogo pedía tres campos. En \textit{nombre} escribió: evaluación de viabilidad. En \textit{entregable} escribió: informe interno. El tercer campo era \textit{cliente} y el cursor se quedó parpadeando ahí, sobre el blanco, a un ritmo de un parpadeo por segundo, once, doce, trece, mientras él repasaba la lista de clientes posibles y la lista devolvía vacío. Escribió su propio nombre.

\clearpage

\chapter{La ciudad de la tela}

\lettrine[lines=2, lhang=0.1, nindent=0.2em]{L}{a} tela estaba a once metros de él y no había forma de verla.


Vidal había llegado a las 10:07 para una cita de las once, porque un lugar se verifica antes de usarlo, y la catedral de Turín era, a efectos de su solicitud, el lugar. Adentro hacía más frío que afuera. Olía a cera y a piedra húmeda, y el ruido de la calle entraba amortiguado, como filtrado por la masa del edificio. Bajo el crucero, detrás de una barrera baja y de un paño oscuro que caía hasta el suelo, estaba la caja: horizontal, sellada, atmósfera inerte desde 1997. Delante había dos filas de reclinatorios y seis personas, cuatro de rodillas, dos sentadas. Las contó dos veces.

Lo que sí podía verse era una reproducción fotográfica en un marco, a un costado, iluminada desde arriba. Tamaño natural, alta resolución; con probabilidad alta, una placa del corpus de Ferrero. Una mujer mayor estaba parada frente a la fotografía con las manos juntas. A once metros de la tela real, sellada e invisible, la gente le rezaba a una copia calibrada. Vidal registró el dato y le buscó categoría durante un rato. En Jerusalén la fila esperaba para ver un lugar donde no había nada; acá el objeto existía, estaba presente, y la devoción se dirigía igual a un sustituto. La variable relevante no parecía ser la visibilidad. Anotó eso en el teléfono, sin desarrollo, y a las 10:48 salió a buscar la sede de la Comisión.


\scenebreak


---Doctor Vidal. Llega temprano, cosa que en esta ciudad se considera una virtud teologal.

Ottavio Ferrero era alto y estaba encorvado de una manera específica, la de décadas sobre un ocular, y tenía las manos manchadas de reactivos viejos, un moteado marrón que ningún jabón saca. Corbata de lana en noviembre. Le tendió la mano, le señaló una silla frente a un escritorio cubierto de carpetas en un orden que parecía tener sistema, y sirvió café de una cafetera de filtro que debía llevar horas encendida.

---Leí su protocolo anoche ---dijo---. Dos veces. La segunda con tinta roja.

---Le agradezco la lectura.

---Agradézcame la tinta roja, que es donde está el trabajo. ---Se sentó---. Antes de eso, una cortesía que no es cortesía: su análisis del díptico de Amberes es de lo más serio que produjo su campo en veinte años. Lo digo habiendo sido amigo de uno de los dos connoisseurs que su resultado jubiló. La química no tiene lealtades, y a mi edad uno aprende a agradecérselo.

Vidal dejó el café en el borde del escritorio.

---Usé su corpus fotográfico como fuente primaria en mi evaluación preliminar. Las dieciséis mil placas. La calibración cromática está documentada mejor que en ningún corpus patrimonial que yo haya procesado.

---Dieciséis mil placas y usted me habla de la calibración. ---Ferrero movió la cabeza---. Me va a caer bien, doctor, y eso me va a complicar la mañana.

Abrió la carpeta de arriba. El protocolo de Vidal, impreso, con anotaciones en tinta roja en los márgenes, densas, algunas con signos de admiración.

---Vamos a lo que importa. Punto tres. Usted pide muestras físicas.

---Tres miligramos de fibra del borde lateral, zona sin imagen, adyacente a los muestreos de 1988. El daño es del orden de lo ya infligido y queda documentado hilo por hilo.

Ferrero lo miró un momento. Después metió la mano en el bolsillo del saco, sacó un par de guantes de algodón de conservación y se los puso mientras hablaba, sin explicación, como quien se acomoda los anteojos.

---Del orden de lo ya infligido ---repitió---. En 1988 se cortó una tira de siete centímetros y le aseguro que todos los involucrados usaron esa misma aritmética: es el borde, es zona sin imagen, es poco. La aritmética era correcta. ¿Sabe qué costó?

---La ubicación de la muestra puso en duda el resultado. Contaminación, remiendos medievales, la hipótesis del borde reparado. Mi protocolo responde a eso con muestreo múltiple y caracterización de la trama previa a la extracción.

---Eso responde a la crítica técnica. Yo le pregunté el costo. ---Ferrero apoyó las manos enguantadas sobre la carpeta---. Si le doy sus tres miligramos y su modelo devuelve un intervalo, el que sea, con la precisión que usted publica, ¿quién administra ese resultado? ¿Usted? ¿Su modelo? ¿La Custodia? Porque alguien lo va a administrar, doctor, eso no es opcional. La única pregunta es si lo decidimos antes o lo improvisamos después, y en el ochenta y ocho se improvisó después.

---Mi informe se publica completo, con datos crudos y advertencias de rango. La administración posterior excede el peritaje.

---Ahí está. ---Ferrero levantó un dedo enguantado, sin agresión, casi con gratitud, como quien encuentra el precipitado que esperaba---. Esa frase. ``Excede el peritaje.'' Es una frase honesta y es la frase que me obliga a decirle que no, aunque todavía no se lo estoy diciendo. Un intervalo de confianza puede ser verdadero y hacer daño. Las dos cosas a la vez, sin contradicción. La verdad del intervalo no lo absuelve del daño.

---El daño no es una variable del peritaje ---dijo Vidal---. Si lo fuera, ningún peritaje sería confiable. Un resultado que se modula por sus consecuencias deja de ser un resultado.

---Deja de ser un resultado suyo. Pasa a ser un resultado de todos, que es exactamente lo que usted no quiere y exactamente lo que va a pasar. ---Ferrero se sirvió más café. A Vidal le rellenó la taza sin preguntar---. En el ochenta y ocho también había un intervalo estrecho, doctor. Tres laboratorios, curvas solapadas, un trabajo limpio. Lo que no había era nadie esperando del otro lado del anuncio. Se leyó en conferencia de prensa, mil doscientos sesenta, mil trescientos noventa, y después cada uno se fue a su casa. Los de la conferencia de prensa, digo. Los otros se quedaron donde estaban, con eso.

Se pasó dos dedos por la muñeca izquierda. No había nada ahí, ni reloj ni pulsera.

---Yo conocía a alguien que siguió ese anuncio como se sigue... ---Se detuvo. Ordenó el borde de la carpeta, que ya estaba ordenado---. Conocí a varios. Le ahorro las historias. Mi pregunta es de método, para que la entienda en su idioma: usted diseña controles para el error del instrumento y controles para el error del operador. ¿Dónde está, en su protocolo, el control para el error del anuncio? ¿Quién espera del otro lado del suyo?

Vidal repasó mentalmente el punto siete de su protocolo, comunicación de resultados, que ocupaba cuatro líneas de las veintidós páginas. Verificó de memoria que decía lo que recordaba que decía. Lo decía.

---El protocolo prevé revisión de pares previa a cualquier publicación.

---Pares. ---Ferrero dejó la palabra en el aire un momento antes de seguir---. Los pares de usted revisan la química, y la química va a estar bien, de eso yo no dudo, se lo digo sin ironía: si alguien vivo puede fechar ese lino con honestidad es usted. El problema es que del otro lado del anuncio no hay pares. Hay dos mil millones de personas sin comité de lectura.

Se levantó, fue hasta la ventana. Desde ahí, supuso Vidal, se veía el campanario del Duomo; él veía la espalda encorvada y las manos enguantadas cruzadas atrás.

---¿Sabe qué es lo que más rabia me da de su protocolo? ---dijo Ferrero, de espaldas---. Que el punto cuatro está bien. La caracterización hiperespectral de la zona de imagen, sin contacto. Eso yo lo quise hacer... ---se corrigió---, eso la Comisión lo evaluó en dos mil veinte y no había instrumento con la resolución. ¿Qué resolución espacial alcanza su banda corta sobre fibra?

---Doce micrones por píxel en condiciones de laboratorio. Sobre celulosa, con el contraste de deshidratación descrito en la literatura, distingo estructura a nivel de fibrilla.

Ferrero tardó en darse vuelta. Cuando lo hizo, se miró un segundo las manos enguantadas como si le hubieran aparecido puestas.

---Fibrilla ---dijo. Y después, más bajo---: Cincuenta y cuatro años. El último instrumento serio que tocó esa tela es anterior a mi doctorado entero. ---Volvió a sentarse. Se sacó los guantes, los dobló, los guardó---. Bueno. Le cuento el trámite, que es a lo que vino.

El trámite era el previsible: la solicitud se elevaba a la Comisión, la Comisión emitía dictamen técnico, la Custodia decidía sobre el dictamen, y sobre la Custodia había otro nivel que Ferrero mencionó con un movimiento de mano hacia arriba que abarcaba a la vez la jerarquía y el techo. Plazos: cuatro a seis semanas. Vidal tomó nota. Después Ferrero, hojeando el protocolo, se detuvo en el anexo de condiciones ambientales.

---Esto va a estar desactualizado para cuando le contesten, dicho sea de paso. Hay teca nueva en construcción. Control de atmósfera de otra generación, sensores continuos, telemetría. Para la ostensión del Jubileo el lienzo va a estar mejor guardado de lo que estuvo nunca, lo cual, tratándose de una tela que sobrevivió a dos incendios y a un bombero con una maza, es decir bastante.

---¿Telemetría con registro exportable?

---Continua y auditada. Los ingenieros de la donación mandaron especificaciones más completas que las que esta Comisión hubiera sabido pedir. ---Ferrero pasó la página---. Donación anónima, eso sí. Canalizada por una asociación de amigos del lienzo, que es una figura que existe y que yo mismo, después de treinta años acá, sería incapaz de explicarle del todo. Llegó en el treinta y uno, en un momento en que esta diócesis no estaba en condiciones de rechazarle una vela a nadie. Anónima y perfecta. En mi experiencia las donaciones perfectas son como los blancos de laboratorio: cuando salen demasiado limpios, uno repite la corrida. ---Se encogió de hombros---. La repetimos. Salió limpia de nuevo. Uno se cansa de desconfiar antes de que la plata se canse de llegar.

Vidal había dejado de escuchar en ``registro exportable''. Una serie temporal continua de la atmósfera de la teca era, para cualquier análisis futuro, una línea de base gratuita: humedad, temperatura, composición del gas, deriva estacional de la fibra. Anotó ``telemetría teca nueva --- solicitar acceso a histórico'' y subrayó la línea. Cuando levantó la vista, Ferrero lo estaba mirando con una expresión que tardó en clasificar y archivó como cansancio.

---Le hago la última pregunta y lo dejo ir ---dijo Ferrero---. Si su modelo devuelve siglo primero, ¿usted qué cree que pasa?

---Depende de qué llames pasar. En términos de evidencia, un intervalo compatible con el siglo primero es condición necesaria y no suficiente. La autenticidad seguiría abierta.

---Eso es lo que pasa en su informe. Le pregunté qué pasa. ---Ferrero juntó las carpetas---. No me conteste. La diferencia entre esas dos preguntas es todo lo que traté de decirle esta mañana, y usted es un hombre inteligente: le va a llegar sola, con retardo, como los productos de pirólisis. Uno huele el humo mucho después del calor.

En la puerta le dio la mano. La cortesía del apretón era de otra década.

---Voy a elevar su solicitud con nota favorable en lo técnico. Es la mejor solicitud que pasó por esta mesa desde que estoy en ella, y lo voy a escribir así, con esas palabras.

---¿Y en lo demás?

---En lo demás, la Comisión va a hacer lo que hace una institución que se quemó dos veces con el mismo fuego. ---Ferrero le sostuvo la mano un segundo más---. Le va a llegar una carta muy bien redactada, doctor. Léala entera. A veces ponemos cosas ciertas en los considerandos.

Vidal volvió por la plaza. Entró de nuevo a la catedral, sin decidirlo del todo: el camino a la estación pasaba por la puerta y la puerta estaba abierta. Bajo el crucero seguía el paño oscuro sobre la caja que nadie veía. Ahora había nueve personas en los reclinatorios. La mujer mayor seguía frente a la fotografía, o era otra mujer mayor en la misma posición; a esa distancia y con esa luz, indistinguibles.

Calculó la probabilidad de aprobación de su solicitud y le salió un número que revisó dos veces y que las dos veces salió igual. Guardó el teléfono sin anotarlo.

Del crucero a la puerta había cuarenta y dos pasos. Los contó.

\clearpage

\chapter{Denegado}

\lettrine[lines=2, lhang=0.1, nindent=0.2em]{L}{a} carta llegó un jueves de diciembre con el correo de las diez, en un sobre de papel verjurado que pesaba más que la correspondencia ordinaria. Vidal lo sopesó en la palma antes de abrirlo. Cuarenta gramos, quizá cuarenta y cinco. El remitente decía Custodia, con una dirección de Turín que él conocía porque la había escrito tres veces en el expediente de solicitud. Afuera pasó el tranvía de la Universitätstrasse, el número 6, con el chirrido de curva que el estudio recibía cada siete minutos.


La abrió con el cortapapeles, no con el dedo. El papel de adentro era del mismo gramaje que el sobre.

Ferrero le había dicho que la leyera entera, así que la leyó entera, de pie junto a la ventana, dos veces. La primera para el contenido. La segunda para la estructura, que era lo que él sabía leer mejor: cinco considerandos y una parte resolutiva de tres líneas. La resolución denegaba el acceso. Los considerandos explicaban por qué, y estaban bien escritos, mejor que bien: cada uno se sostenía solo y ninguno dependía del anterior, de modo que refutar uno dejaba cuatro en pie.

El considerando segundo citaba el dictamen técnico de la Comisión de Conservación y reproducía una frase textual, entre comillas: \textit{la mejor solicitud que pasó por esta mesa desde que estoy en ella}. Ferrero había cumplido. Lo habían citado con exactitud, con la referencia del acta, y dos considerandos más abajo la misma carta usaba esa calidad como agravante: precisamente porque el protocolo era serio, correspondía evaluar el marco en que sus resultados existirían. El considerando cuarto hacía esa evaluación. Reglamento europeo de peritaje algorítmico de 2029, artículo 9, apartado 2: los sistemas de arquitectura cerrada quedan excluidos del peritaje de patrimonio protegido; sus resultados carecen de existencia a efectos de certificación. El sistema Tamiz, decía la carta, opera bajo la excepción académica del artículo 17, que habilita investigación y excluye certificación. En consecuencia, cualquier resultado obtenido sobre la Reliquia carecería de valor legal y académico, y un riesgo material sobre el lienzo ---tres miligramos de fibra, la carta daba la cifra de su propio protocolo--- resultaría injustificable a cambio de un resultado que, a efectos del ordenamiento, no existiría.

Era un argumento limpio. Vidal lo desarmó en la cabeza buscando la soldadura floja y la soldadura floja tardó en aparecer porque la carta la había previsto: la Custodia ni siquiera estaba obligada por el reglamento. Una entidad eclesiástica, patrimonio propio, jurisdicción propia. Invocaba una norma que le era ajena, voluntariamente, porque la norma le servía. Eso era lo único atacable, y era inatacable: nadie puede reprocharle a una institución que adopte un estándar más exigente que el que la ley le impone.

El considerando quinto era el que Ferrero le había anunciado sin anunciarlo. Hablaba de ``la experiencia acumulada por esta Custodia en materia de comunicación pública de resultados de análisis científicos'' y de la necesidad de que ``toda investigación futura se inscriba en un marco que garantice la gradualidad de dicha comunicación''. Sin fecha, sin laboratorios, sin la palabra carbono. Vidal le puso la fecha en el margen, a lápiz: 1988. Después leyó la frase de nuevo. \textit{La comunicación}, decía. El defecto que la carta le imputaba a 1988 era la conferencia de prensa, la falta de custodia del relato. El resultado mismo quedaba fuera del reproche.

Anotó eso también. Le pareció un dato, todavía sin columna donde ponerlo.

Volvió al escritorio y abrió el texto consolidado del reglamento. Lo conocía; había participado como experto consultado en dos rondas de comentarios, en 2028, y había escrito entonces que el requisito de trazabilidad del razonamiento era correcto en principio e inaplicable en la práctica a las arquitecturas de mejor desempeño. Su comentario constaba en el anexo público de la consulta, con su nombre. Podía recitar el artículo 9 sin mirarlo. Lo leyó igual, apartado por apartado, siguiendo las líneas con el cursor, y después leyó el 17, y después las definiciones del artículo 3, donde ``auditable'' tenía dos criterios acumulativos y Tamiz cumplía uno. El desempeño de Tamiz era verificable hasta el hartazgo: controles ciegos, tasas de error publicadas, seis vías independientes de refutación en el caso de Amberes y las seis a favor. Ninguna de esas verificaciones contaba. El criterio era el mecanismo, no el resultado. Un perito que acierta siempre y no puede mostrar cómo, decía en sustancia el artículo 3, no es un perito: es un oráculo, y los oráculos quedan fuera del expediente.

Él había escrito una versión de esa frase. En un informe de 2031, sobre el trabajo de otro: \textit{una certeza inauditable no es una certeza; es un artefacto, y se descarta}. La carta de Turín le aplicaba su propia regla con papel mejor y sin nombrarlo a él tampoco.

Abrió el registro de jornada y escribió la entrada del día. Fecha, hora de recepción, remitente, síntesis en una línea: \textit{denegación; fundamento técnico art. 9.2; fundamento adicional considerando 5°, ver margen}. Le quedaba una línea disponible en el campo de observaciones. Escribió: \textit{criterio de exclusión = criterio propio (inf. 2031)}. Miró la frase. Era exacta. La dejó como estaba y cerró el campo sin agregar nada, porque cualquier cosa que agregara sería un comentario y los comentarios eran para los informes de otros.

Guardó la carta en la carpeta física del expediente, que ahora estaba completa: solicitud, protocolo de veintidós páginas, acuses, denegación. Un expediente cerrado en cuatro documentos. En su clasificación de archivo, los expedientes cerrados iban al armario del fondo. Este lo dejó en el escritorio, contra la lámpara, y esa desviación del procedimiento quedó sin registrar.


\scenebreak


El lunes llamaron de Múnich a las 9:20. Un tal Brandmeier, de la aseguradora que llevaba la cartera de dos colecciones privadas bávaras, con un caso que en otro trimestre habría sido interesante: una tabla atribuida a Cranach el Viejo, litigio de herencia, honorarios que el hombre mencionó en el segundo minuto de conversación, cosa que en Vidal siempre bajaba la estimación previa del interlocutor.

---Necesitaríamos el informe antes de marzo ---dijo Brandmeier---. El calendario procesal no...

---El plazo es viable.

---Perfecto. Entonces le mando el borrador de encargo hoy y...

---El plazo es viable y voy a declinar el encargo.

Hubo un silencio de línea, dos segundos, con el ruido de una oficina detrás.

---¿Puedo preguntar el motivo? Para el expediente interno, entiéndame.

---Capacidad.

---Doctor, discúlpeme la franqueza, consultamos antes de llamar. Nos dijeron que su calendario del primer trimestre estaba libre. Fue el argumento para autorizar el presupuesto.

Vidal miró el calendario abierto en la segunda pantalla. Enero en blanco, febrero en blanco, marzo con una sola reunión académica. La información de Brandmeier era correcta y la suya no, y él estaba del lado equivocado de esa asimetría por primera vez en mucho tiempo.

---No exactamente ---dijo---. Hay un compromiso previo sin fecha cerrada.

---¿Y si nosotros nos adaptamos a esa fecha?

---El compromiso es excluyente.

Brandmeier lo intentó una vez más, con delicadeza profesional, y después agradeció y cortó. Vidal se quedó con el teléfono en la mano. \textit{Compromiso previo sin fecha cerrada.} La frase era formalmente verdadera del modo en que son verdaderas las frases que uno construye para que lo sean. Abrió el registro para asentar la llamada. En el campo \textit{motivo de declinación} el catálogo ofrecía siete códigos. Conflicto de interés, no. Fuera de especialidad, no. Capacidad, no: acababa de comprobarse que no. Leyó los siete dos veces. Dejó el campo vacío y el registro le marcó la entrada en amarillo, incompleta, como estaba diseñado que hiciera.

La segunda oferta llegó por escrito el miércoles: una casa de subastas de Viena, verificación de atribución de un Klimt de procedencia argentina, exclusividad, adelanto del cuarenta por ciento. La respondió en dos líneas, declinando sin motivo, porque el correo, a diferencia de su propio registro, aceptaba la ausencia de motivo sin marcar nada en amarillo. Después se levantó, fue hasta el rack del fondo y corrió la verificación de integridad de los discos fríos, que estaba programada para el viernes y que ningún dato sugería adelantar. Integridad correcta. Volvió al escritorio.

Dos encargos declinados en tres días sumaban, a tarifa corriente, más que todo lo facturado en el semestre. El número era fácil de calcular y lo calculó, porque calcular era barato y el resultado no lo obligaba a nada. Lo que le habría costado algo era la otra operación, la de escribir en el campo amarillo la razón verdadera, y esa operación la tenía clasificada, con precisión, como pendiente.

El expediente cerrado seguía contra la lámpara. La luz de la tarde de diciembre, que en Zúrich duraba poco y venía gris, le daba al papel verjurado una textura de tela vista de cerca.

El jueves a la noche, después del sincronizador de las 20:00, abrió la carpeta local de 2029. La planilla pesaba once kilobytes, igual que la última vez; la fecha de modificación era la de octubre, Jerusalén, el hotel de la calle Jaffa. Desde entonces el archivo había viajado con él en el disco sin que él lo tocara, salteado cada noche por la rutina de respaldo, un pasajero sin registro en un sistema donde todo lo demás dejaba traza doble.

La abrió.

La columna decía \textit{tiempo perdido}. La entrada decía \textit{basílica del Santo Sepulcro, 14:37--17:05, sin entregable}, con el círculo pequeño al margen, el signo de las cosas que no requieren acción. Dos horas veintiocho. Él había estado ahí, había contado una fila de cuarenta y una personas frente a un lugar donde un panel en tres idiomas declaraba que adentro faltaba lo importante, y había clasificado esas dos horas veintiocho con el único código de su catálogo que significaba error. Ahora tenía sobre el escritorio una carta que clasificaba dos meses de su mejor trabajo con el código equivalente de otro catálogo: sin existencia a efectos del ordenamiento. Los dos asientos eran simétricos y los dos estaban bien hechos según las reglas de sus respectivos libros.

El cursor quedó un rato sobre la celda del código. Reclasificar era una operación de cuatro teclas. Conocía el catálogo de memoria: \textit{asunto abierto} habría sido defendible, incluso riguroso, dado que la nota del teléfono seguía con los dos puntos sin completar. Cuatro teclas y la entrada dejaba de ser un error para ser una pregunta.

Dejó el código como estaba y cerró la planilla, y el gesto de dejarla intacta le llevó más tiempo que el que habría llevado cambiarla.

Después tomó la carta de Turín, la última hoja, y la sostuvo contra la lámpara. A trasluz apareció la marca de agua: un escudo, nítido, centrado en el pliego. Papel de archivo, cien por ciento algodón, hecho para durar siglos. Verificó el escudo contra el membrete. Coincidían. De todo el expediente, era la única afirmación que había podido comprobar con sus propios instrumentos, y era verdadera: la carta que lo dejaba afuera estaba escrita en un papel que iba a sobrevivirlos a los dos.

\clearpage

\chapter{Maître Ansermet}

\lettrine[lines=2, lhang=0.1, nindent=0.2em]{E}{l} eufemismo estaba en la segunda línea. \textit{Acceso a material textil de interés histórico}: seis palabras que no nombraban nada y que solo podían nombrar una cosa. Vidal leyó la carta entera antes de volver a esa línea, porque el orden de lectura era parte del método, y la carta entera confirmó lo que la segunda línea ya había dicho. Membrete de la Fondation Cassiodore, dirección en Ginebra, un párrafo de cortesía sin adjetivos, la propuesta de una conversación en el estudio del abajo firmante, tres fechas posibles. Firmaba J.-M. Ansermet, \textit{avocat}. Ninguna frase de la carta proponía nada. Una conversación en un estudio de Ginebra era un acto tan legal como el papel en que estaba escrita la invitación, y Vidal registró que ese razonamiento ---clasificar la carta por lo que no decía--- era un razonamiento de abogado, y que él lo había producido sin esfuerzo, en el primer minuto, antes del café.


Verificó lo verificable. La fundación existía: registro de comercio ginebrino, fines culturales, un consejo de tres nombres que no devolvían nada útil en ninguna base. Lo demás era opaco, y la opacidad, en Ginebra, no era un dato: era un servicio.

Contestó aceptando la segunda fecha y compró el tren de las 8:02.

En Ginebra soplaba viento del lago y las calles altas de la ciudad vieja olían a piedra fría. El estudio ocupaba un tercer piso en la place du Bourg-de-Four; la placa de bronce era más chica que una tarjeta de visita y decía solo el apellido. Abrió el propio Ansermet. Sesenta y tantos años, traje oscuro sin marca visible, un apretón de manos seco y breve. Adentro había un reloj de péndulo, estanterías de expedientes encuadernados en tela gris y un escritorio sin computadora a la vista. Olía a papel viejo y a cera de piso. El radiador crujía a intervalos de unos cuarenta segundos.

---Gracias por venir, doctor. Le propongo una regla para esta reunión, y es la única que voy a proponerle: yo no voy a decir en esta habitación ninguna frase que no pudiera leerse en voz alta ante un tribunal. Le pido que evalúe todo lo que escuche con ese criterio, porque es el criterio con que fue redactado.

Hablaba un español sin acento localizable, aprendido en algún lugar caro. Las oraciones salían terminadas, con sujeto, verbo y cierre, sin una sola contracción del pensamiento.

---Depende de qué tribunal ---dijo Vidal.

---De cualquiera. Es lo que vuelve interesante el ejercicio.

Se sentaron. Ansermet cruzó las manos sobre el escritorio y expuso. La Fondation Cassiodore actuaba por cuenta de un mandante cuyo nombre no formaba parte del mandato. El mandante estaba en condiciones de hacer posible un acceso directo, sin supervisión institucional, a un material que Vidal había intentado estudiar por la vía ordinaria durante el último trimestre. El acceso duraría diecisiete días, del doce al veintinueve de abril. Los instrumentos serían los de Vidal, el protocolo sería el de Vidal, los resultados los produciría Vidal. La restitución del material a su ubicación original estaba garantizada dentro del marco de la inspección final del veintinueve de abril, y todo el asunto quedaba concluido antes del cuatro de mayo, fecha en que abría una exhibición pública que ninguna de las partes tenía interés en perturbar.

En ningún momento dijo Sábana. En ningún momento dijo Turín.

---Está hablando de la Sábana de Turín ---dijo Vidal.

---Usted acaba de decir ese nombre. Yo no. Le agradecería que, cuando reconstruya esta conversación, recuerde de quién fue cada palabra.

El reloj dio la media. Vidal esperó a que terminara.

---Lo que usted describe es un delito.

---Yo no he descrito ninguna acción, doctor. He descrito un calendario y una capacidad. Si en algún punto de esta reunión yo le propusiera cometer un delito, usted haría bien en levantarse e irse, y yo habría hecho mal mi trabajo. Ninguna de las dos cosas va a ocurrir.

Vidal apagó el teléfono. Ya estaba apagado; verificó igual, con la mano dentro del bolsillo, el botón bajo el pulgar.

Después empezó a preguntar. Era lo que hacía con los peritos de parte, con los gestores de aseguradoras, con los marchands que traían procedencias demasiado limpias: preguntas de precisión, cerradas, diseñadas para que cualquier respuesta comprometiera al que respondía con una afirmación verificable. En once años nadie había aguantado más de seis sin entregar algo: una fecha falsa, un nombre de más, una negación que después se podía medir.

---¿La Fondation fue la donante de la teca nueva de 2031?

---La Fondation no ha hecho donaciones a ninguna entidad italiana. Consta en sus estados financieros, que se depositan cada año y que usted puede pedir mañana en el registro.

Verificable. Verdadero, casi con seguridad. Y sin contenido: la donación de la teca había pasado por una asociación de amigos del lienzo, no por una fundación ginebrina. La respuesta era una puerta que se abría a una habitación vacía.

---¿Su mandante es una persona física?

---Mi mandato incluye la instrucción de no caracterizar al mandante. La instrucción es anterior a usted y no se refiere a usted.

---¿Cuántas ofertas como esta ha transmitido la Fondation?

---Es la primera vez que la Fondation le transmite una oferta a usted, doctor.

---Eso no es lo que pregunté.

---Es lo que estoy en condiciones de responder. Observe que la distancia entre su pregunta y mi respuesta también es información, y que se la estoy dando gratis.

Vidal se recostó en la silla. El radiador crujió. Llevaba nueve preguntas y tenía las manos vacías: ni una fecha comprometida, ni un nombre, ni una negación medible. Cada respuesta era verdadera, completa en su sintaxis y hueca en su referencia, y las nueve juntas dibujaban con exactitud el perímetro de lo que no se le iba a decir. Era la primera vez en años que alguien jugaba ese juego con él y terminaba las rondas intacto. Contó las respuestas otra vez, de memoria, buscando la grieta. Nueve. Ninguna.

---¿Cómo se pagaría un trabajo así? ---dijo, y eligió el condicional a conciencia.

---De un modo que no le genere a usted obligaciones de declaración ante terceros. La Fondation no opera por vía bancaria en asuntos de esta naturaleza. Los detalles del instrumento figurarían en el contrato, que también está redactado para leerse en voz alta.

---Un contrato sobre un objeto que nadie nombra.

---Un contrato sobre servicios profesionales de análisis de material textil. Los contratos nombran obligaciones, doctor, no objetos. Es una distinción que a los científicos les cuesta y a los archivistas no.

Vidal miró la estantería de expedientes grises. Fundación, abogado, mandante sin nombre, pagos que no tocaban un banco: tres capas, cada una legal en sí misma, cada una opaca hacia la siguiente. La estructura era limpia del modo en que era limpio un dato sin trazabilidad: no se le podía reprochar nada porque no ofrecía superficie donde apoyar el reproche.

Quedaba una prueba. La había preparado en el tren, entre Lausana y Ginebra, como pregunta final: la que iba a terminar la reunión, porque nadie podía responderla.

---Si esto fuera algo más que una conversación ---dijo---, yo necesitaría el expediente técnico de la ventana de acceso. Calendario real de la manipulación, personal previsto, especificaciones de la teca, condiciones de transporte y de atmósfera, tolerancias. Sin ese expediente no hay nada que evaluar, y ese expediente no existe fuera de la Comisión.

Ansermet se levantó sin apuro, fue hasta la credencia bajo la ventana y volvió con una carpeta encuadernada en la misma tela gris de los estantes. La puso sobre el escritorio, del lado de Vidal, con el lomo hacia él para que pudiera leerlo. El lomo no decía nada.

---Está indexado ---dijo Ansermet, y volvió a sentarse.

Catorce pestañas. Vidal las contó con el pulgar antes de abrir la primera. El calendario ocupaba la segunda: traslado del material a la teca nueva el once y el doce de abril, unas treinta horas de manipulación en la sacristía con personal reducido, cada turno con horario y función; del doce al veintinueve, diecisiete días; inspección final de la teca el veintinueve; apertura de la ostensión el cuatro de mayo, con la visita anunciada que hacía imposible cualquier cosa después de esa fecha. La pestaña seis era la telemetría de la teca: sensores, frecuencias de muestreo, protocolos de alarma, el histórico al que él había querido pedir acceso en noviembre y que nunca le contestaron. La pestaña once era un anexo de requisitos instrumentales para caracterización hiperespectral de superficie. Decía doce micrones por píxel.

Doce. Su número. El de su protocolo de veintidós páginas, el que la Custodia había usado de agravante.

---¿Cómo obtuvo esto?

---No redacté yo el expediente. No estoy en condiciones de describir las fuentes de un documento que no redacté.

Vidal volvió a la pestaña del calendario. Se dio cuenta de que la releía con un lápiz mental en la mano y de que la primera anotación no era una objeción de principio sino una objeción de ingeniería: el expediente asignaba cuatro horas de estabilización de atmósfera tras el transporte y cuatro horas era poco, para lino con esa historia térmica él habría exigido doce. La segunda anotación era sobre la vibración del vehículo. La tercera era una pregunta sobre la humedad relativa del laboratorio de destino. Iba por la tercera cuando registró lo que estaba haciendo: corregir el diseño. En algún punto de los últimos minutos el documento había dejado de ser una prueba en su contra y se había vuelto un borrador a mejorar, y él no había decidido ese cambio; el cambio se había hecho solo, con su letra.

Cerró la carpeta.

---¿Cómo garantizan la devolución? ---dijo, y oyó la pregunta un segundo tarde, cuando ya estaba afuera: era la pregunta de alguien que estaba adentro del problema, no enfrente.

Ansermet, si lo registró, tuvo la cortesía de no mostrarlo.

---La devolución es la única cláusula del proyecto que no admite escenario alternativo. El mandante lo sabe mejor que nadie. Si algún día lo conoce, entenderá por qué le digo que esa garantía es la más sólida de todas las que puedo ofrecerle, y también la única que le pido que acepte sin auditar.

---Yo no acepto nada sin auditar.

---Lo sé. Por eso el expediente estaba preparado antes de que usted lo pidiera.

El reloj marcó la hora. Afuera había empezado a oscurecer sobre los tejados y el radiador crujía a intervalos más cortos, o eso le pareció; no los había medido. Ansermet acercó la carpeta dos centímetros más hacia el borde del escritorio.

---Llévela. Léala como leería el trabajo de otro. No le pido una respuesta: le pido que encuentre dónde falla, y que me llame cuando tenga la lista. Si la lista es larga, no habrá pasado nada. Una conversación en Ginebra y un expediente sin membrete no constituyen nada, en ninguna jurisdicción.

---¿Y si digo que no?

---Nada, doctor. Esa es la parte de la oferta que más le va a costar creer.

Vidal guardó la carpeta en el bolso, en el compartimiento interno, entre el pasaporte y el billete de vuelta, donde llevaba las cosas que no podían perderse. Lo notó al cerrar el cierre. La dejó ahí.

\clearpage

\chapter{Debida diligencia}

\lettrine[lines=2, lhang=0.1, nindent=0.2em]{L}{a} Fondation Cassiodore ocupaba once líneas en el extracto del registro de comercio de Ginebra. Vidal las leyó por tercera vez a las diez y veinte de la noche de un martes, en su estudio de Zúrich, con el flexo bajo y el chirrido del 6 entrando por la ventana en la curva de la Universitätstrasse. Fines: promoción, estudio y conservación del patrimonio cultural europeo. Capital fundacional: cincuenta mil francos, el mínimo legal. Inscripción: marzo de 2019. Órgano de revisión: una fiduciaria de la rue du Rhône que figuraba como auditora de otras doscientas catorce entidades. Consejo: los tres nombres que ya conocía, tres abogados de tres estudios distintos, sin más huella pública que la matrícula.


Once líneas. En un peritaje, un objeto con esa documentación se habría clasificado sin deliberar: procedencia construida. Demasiado limpia para ser descuido, demasiado escueta para ser inocencia.

Aplicó el único método que tenía, que era el suyo: seguir el material. Los estados financieros depositados ante la autoridad de vigilancia declaraban ingresos por ``contribuciones de terceros afectadas a proyectos'', sin un nombre en tres ejercicios, y un patrimonio administrado por mandato a través de dos sociedades anónimas, una de Ginebra y una de Zug, cuyos objetos sociales estaban redactados con las mismas palabras exactas: servicios de trazabilidad de materias primas y certificación de origen. La coincidencia textual indicaba un mismo redactor; la tomó como el primer dato firme de la noche. El accionista único de ambas era una Anstalt de Vaduz cuyo reglamento interno, por derecho de Liechtenstein, se guarda sin depositar. El fundador fiduciario de la Anstalt remitía a una fundación privada de Willemstad, Curazao, donde el beneficiario económico es, por diseño de la jurisdicción, un dato que ningún registro público contiene. Tres jurisdicciones, cada una elegida por lo que se niega a decir. Ahí la cadena dejaba de ser cadena.

Pagó cuatrocientos francos por un informe de una base de cumplimiento que usaban los bancos. Llegó en once minutos: sin coincidencias adversas, sin personas expuestas políticamente, sin litigios. Lo único que ese resultado probaba era que la estructura la habían levantado profesionales y que funcionaba según diseño.

Contó los nombres humanos verificables en toda la cadena. Cinco. Los cinco eran fiduciarios de oficio, hombres que firmaban por cientos de estructuras ajenas, y ninguno era dueño de nada. Alguien cuyo patrimonio se administraba a través de empresas dedicadas a saber dónde estaba cada cosa y quién la había tocado había pagado, con precisión de especialista, para que nadie pudiera saber dónde estaba él. Vidal anotó la frase en la libreta de trabajo y la subrayó, porque describía una competencia, y las competencias del otro lado de una mesa conviene tenerlas medidas.

A las doce menos cuarto corrió el respaldo nocturno de sus discos, aunque el respaldo corría solo a las doce.

El miércoles a las 8:10, antes del café, sacó del cajón el cuaderno de tapas grises donde anotaba lo que su sistema de planillas carecía de columna para recibir. Había pocas páginas usadas. Escribió la fecha, la hora, y debajo, en dos líneas:

\textit{El mandante no es identificable con los medios a mi alcance.}
\textit{Acepto no saberlo.}

La primera línea era un resultado y podía defenderse ante cualquiera. La segunda era otra cosa. La releyó con la distancia con que revisaba las conclusiones ajenas y observó que la letra no cambiaba de una línea a la otra: el mismo trazo parejo, las mismas mayúsculas alineadas de sus cadenas de custodia. Un grafólogo habría dicho que las dos frases salían del mismo hombre en el mismo estado. Cerró el cuaderno y lo devolvió al cajón, debajo de los manuales de calibración.

Esa tarde llamó a Ginebra con la lista que Ansermet le había pedido: once objeciones al expediente, todas de ingeniería. Estabilización de doce horas en lugar de cuatro. Registro de vibración durante el transporte con acelerómetro propio, del que él conservara los datos. Humedad relativa del laboratorio de destino con banda de tolerancia y no con valor nominal. Ocho más del mismo orden. Ansermet las anotó, pidió una precisión sobre la sexta y aceptó las once sin negociar ninguna. Vidal colgó y se quedó mirando la lista tachada. La había armado también para medir resistencia, y la resistencia medida era cero. Cero también era una medición; todavía le faltaba saber de qué.

El borrador de contrato llegó el viernes por mensajería, en papel sin membrete, cosido con hilo. Diecinueve cláusulas. Leyó las primeras siete al ritmo de quien coteja: partes, objeto descrito por perífrasis, plazos, indemnidades, la restitución con rango de cláusula esencial. La octava se titulaba ``Archivo'' y la leyó tres veces. Los datos crudos, los resultados intermedios y definitivos y ``toda anomalía registrada durante los trabajos'' quedaban en custodia de la Fondation a perpetuidad. El experto conservaba ``el uso académico de los resultados, sujeto a autorización''.

Marcó tres puntos con lápiz. Después registró la palabra que sobraba: \textit{anomalía}. Un contrato prudente dice hallazgos, resultados, información. Anomalía es vocabulario de quien ya estuvo frente a una. Y registró la palabra que faltaba: un plazo. El documento entero era fechas ---el 12, el 29, el 4 de mayo---, y el único punto sin fecha de las diecinueve cláusulas era el que decía \textit{a perpetuidad}.


\scenebreak


El lunes tomó otra vez el tren de las 8:02. Ginebra estaba bajo una niebla que cortaba los edificios a la altura del segundo piso, y en el estudio del Bourg-de-Four el radiador crujía con el mismo intervalo de la primera visita, unos cuarenta segundos; esta vez llevó la cuenta sin proponérselo. Olía a cera de piso y a papel viejo. Ansermet le sirvió agua sin preguntarle.

---La cláusula octava ---dijo Vidal, y puso el borrador sobre el escritorio, abierto, con las tres marcas de lápiz a la vista.

---La releí esta mañana suponiendo esta conversación. Dígame las tres marcas en el orden que prefiera.

---"A perpetuidad''. ``Custodia''. Y ``sujeto a autorización'', que en la práctica es\ldots{} ---se detuvo, buscó la palabra exacta---. Es un veto. Sin plazo y sin causa.

---En los dos primeros puntos puedo ofrecerle precisión y en el tercero puedo ofrecerle procedimiento. ``Sujeto a autorización'' puede convertirse en ``previa autorización escrita, que no será denegada sin causa fundada'', con un plazo de respuesta de sesenta días. Puedo añadir que la denegación deberá motivarse por escrito y que usted tendrá derecho a una segunda presentación. Todo eso lo firmo hoy.

---La causa fundada la evalúa la Fondation.

---Correcto.

---Y el silencio de sesenta días se tiene por denegación.

---Correcto también. Veo que leyó el derecho aplicable durante el fin de semana.

Vidal apoyó el índice sobre la página.

---Entonces la propuesta es invertirla. El archivo queda en mi custodia. La Fondation recibe copia íntegra y el mismo uso que me ofrece a mí. La redacción se la dejo a usted.

Ansermet juntó las manos sobre el escritorio. Tardó en contestar lo que tardaba el radiador en crujir, y crujió.

---Le voy a ahorrar tres semanas, doctor, porque las tres semanas las perderíamos los dos. En este contrato puedo concederle todas las redacciones que pida y ninguna sustancia. La cláusula octava es sustancia. Prefiero decírselo hoy, completo y de una vez, antes que dejárselo descubrir coma por coma, porque usted lo descubriría y a mí me habría costado su confianza en todo lo demás, que sí es genuina.

---Hay una cuarta marca ---dijo Vidal---. ``Anomalía''. Un abogado escribe ``hallazgos''. ``Anomalía'' es una palabra de laboratorio, y de un laboratorio que espera encontrarla.

---Le contesto lo que me consta. La cláusula octava es la única del contrato que recibí ya redactada; las otras dieciocho son mías. La transcribí sin cambiarle una palabra, incluida esa. Si la palabra le informa algo sobre su redactor, esa inferencia es suya: yo carezco de base para confirmarla y de instrucciones para negarla.

Se quedaron callados. Afuera, una campana dio la media en algún patio de la ciudad vieja. Vidal miró la página y dejó que el cálculo corriera entero, porque los cálculos a medias eran los que después cobraban interés. Lo que la cláusula le quitaba tenía un valor de mercado conocido: cero. Un resultado de Tamiz era, por el artículo 9.2, legalmente inexistente; obtenido sobre este objeto, por esta vía, era además impublicable por origen, en cualquier revista, en cualquier siglo de su vida profesional. Estaba cediendo a perpetuidad un activo que ningún registro del mundo le reconocería jamás. El cálculo cerraba. Le molestó la velocidad con que cerraba, como le molestaban los ajustes perfectos en datos sucios. Del otro lado de la ecuación había una sola cantidad: sin esto, la pregunta quedaba abierta para siempre. Con esto, quedaba abierta la cláusula. Puesto en limpio era una permuta de incógnitas, y de las dos incógnitas él sabía cuál podía medir y cuál iba a medirlo a él.

Eligió la pregunta.

---Redacte las tres correcciones ---dijo---. La cláusula queda.

---Quedará con sus tres correcciones. ¿Algo más para esta versión?

---Sí, pero excede el contrato. Protocolo técnico, composición del equipo, mecánica de la devolución. Condiciones, no cláusulas. Las traigo por escrito.

---La Fondation puede recibirlo el viernes. Si el viernes le resulta incómodo, cualquier día que usted proponga será el correcto.

Ansermet lo acompañó hasta la puerta. En el descanso de la escalera, con el sobretodo ya puesto, Vidal se dio vuelta.

---Una cosa. Cuando dijo que su confianza en todo lo demás era genuina, ¿eso también podría leerlo ante un tribunal?

---Podría ---dijo Ansermet---. Sería el único momento de la audiencia en que el tribunal me creería sin esfuerzo.


\scenebreak


En el tren de vuelta, pasado Lausana, sacó el cuaderno de tapas grises. Buscó la página del miércoles y leyó las dos líneas. Debajo, con fecha y hora, agregó:

\textit{Cláusula octava: aceptada. Obtuve tres redacciones. La sustancia quedó donde estaba.}

Después, en el renglón siguiente, escribió la palabra \textit{anomalía}, sola, y le puso al lado el signo con que marcaba en sus planillas las entradas sin acción pendiente: un círculo pequeño. Miró el círculo un rato. Lo dejó, aunque la clasificación era discutible.

Las tres entradas de la página tenían la misma letra, pareja, sin una corrección, la letra con que había firmado cuatrocientas cadenas de custodia. Verificó la fecha contra la pantalla del teléfono antes de cerrar el cuaderno. La sabía de memoria. La verificó igual.

\clearpage

\chapter{Firma}

\lettrine[lines=2, lhang=0.1, nindent=0.2em]{M}{arzo} llevaba tres semanas de redacciones cruzadas por correo ---Ansermet no había exagerado--- cuando volvieron a sentarse frente a frente. El viernes final, el radiador del estudio de Ansermet crujió tres veces mientras Vidal leía su lista. Intervalos de cuarenta segundos; la lectura le llevó dos minutos y diez. Había traído las condiciones en una hoja, a un solo espacio, numeradas, y las leyó en el orden en que las había escrito, que era el orden de importancia y no el de calendario.


---Primera. La devolución es innegociable. Si en algún punto de los diecisiete días surge un conflicto entre completar el análisis y cumplir la ventana de restitución, el conflicto se resuelve a favor de la restitución. Sin excepción, sin consulta, sin voto. Aunque el análisis quede trunco.

Ansermet tenía las manos cruzadas sobre el escritorio vacío.

---Usted subordina el objeto de su interés profesional al cierre de la operación ---dijo---. Quiero que conste que lo formuló usted, porque es la instrucción más antigua que tengo en este asunto. La recibí antes de conocer su nombre, y es la única que me fue dada dos veces. La Fondation acepta.

---Segunda. El protocolo técnico lo apruebo yo. Punto por punto, por escrito, con fecha. Ningún procedimiento que toque la tela se ejecuta sin mi firma previa, y me reservo el veto sin causa sobre cualquiera de ellos. Si alguien del equipo considera que un veto mío pone en riesgo el calendario, la discrepancia se registra y el veto se mantiene.

---Aceptado en esos términos. El registro quedará en papel, en el laboratorio, en un solo ejemplar.

---Tercera. El equipo. Necesito conservación textil de primer nivel. No un restaurador general: alguien con experiencia documentada en lino arqueológico, en gran formato, en tela con historia térmica. La manipulación es el punto de falla de todo el plan y no pienso...

---Permítame ahorrarle la redacción del resto del párrafo ---dijo Ansermet. Era la primera vez en tres reuniones que lo interrumpía---. El equipo está seleccionado.

Abrió el cajón izquierdo del escritorio y puso sobre la mesa dos hojas sueltas, sin membrete, mecanografiadas. Perfiles, sin nombres. El primero: conservación textil, formación completa en el Opificio delle Pietre Dure, especialización en tejidos planos de gran formato, experiencia de escuela directa en la intervención de 2002 sobre el mismo objeto. El segundo: seguridad y logística, historia laboral verificable de dieciocho años, sin observaciones.

Vidal leyó el primer perfil dos veces. Después puso su propia hoja al lado y comparó su condición tercera con el párrafo mecanografiado. «Gran formato.» «Historia térmica.» La frase sobre la manipulación como punto de falla estaba, con otro orden de palabras, en el tercer renglón del perfil.

---¿Cuándo se redactó esto?

---Los perfiles llevan fecha de impresión de noviembre.

La denegación de la Custodia había llegado a Zúrich en diciembre. En noviembre su solicitud todavía estaba en trámite, elevada con nota favorable, plazo estimado de cuatro a seis semanas. Alguien había armado el equipo de una operación cuyo requisito previo era que a él le dijeran que no, semanas antes de que se lo dijeran.

---En noviembre mi solicitud seguía abierta.

---Es correcto.

---Entonces la fecha de estos perfiles indica una de dos cosas. Previsión o información.

---Carezco de elementos para decirle cuál de las dos ---dijo Ansermet---. Y de instrucciones para negarle cualquiera.

El radiador crujió. Vidal contó hasta el crujido siguiente: treinta y ocho segundos. Volvió a leer su lista, que ya sabía de memoria, buscando la condición que ellos no hubieran anticipado. La lista tenía tres puntos. Los tres estaban resueltos desde antes de que él los escribiera. Había venido a negociar y estaba recitando; la diferencia entre las dos cosas era exactamente la diferencia que su oficio existía para detectar, y la había detectado tarde, con el documento ya sobre la mesa.

---Falta un perfil ---dijo---. Ansermet había hablado de un equipo de tres.

---El tercer integrante se presentará en Collegno, el primer día. El fundador considera que esa incorporación se explica mejor en persona que en papel, y en este caso comparto el criterio.

---Ese criterio no es auditable.

---No lo es. Lo dejo constar.

Los tres ejemplares del contrato estaban en el borde del escritorio, cosidos con hilo, con las diecinueve cláusulas y sus tres correcciones a la octava incorporadas al texto. Junto a los ejemplares había una lapicera.

Era de laca negra, sin marca visible, con el capuchón puesto. Vidal la levantó. Dieciocho, veinte gramos. La laca estaba comida en la zona de agarre, un desgaste parejo, de décadas, del tipo que produce una sola mano sosteniendo siempre igual; el plumín, cuando lo destapó, mostraba el bisel lateral de un dueño único que escribe inclinado. Un objeto con trazabilidad completa. Procedencia ininterrumpida, un solo custodio, uso continuo verificable en el propio material. Lo único en toda la operación que cumplía sus estándares era el instrumento con que iba a firmarla.

Firmó los tres ejemplares. Ansermet firmó después, secó cada firma con un papel, y guardó dos ejemplares en el cajón. El tercero se lo entregó sin comentario.

---Hay algo más ---dijo Ansermet. Se levantó, fue hasta el estante de los expedientes de tela gris y bajó un paquete envuelto en papel madera, atado con cordel---. De parte del fundador. Me pidió transmitirle una sola frase: que este objeto, a diferencia del otro, no requiere devolución.

Vidal sostuvo el paquete. Pesaba como un libro grande.

---¿Debo entender la frase como una cortesía o como una cláusula?

---Debe entenderla como fue dicha. Yo la transmití completa.


\scenebreak


En el tren, pasado Nyon, cortó el cordel con la llave del estudio y abrió el papel. Era un ejemplar de las actas del STURP, la edición de los proceedings de 1981, tela editorial gastada en el lomo, quinientas y pico de páginas. Lo abrió por el índice y una nota a lápiz le salió al paso antes de llegar al primer artículo: junto al título de la ponencia sobre química de la imagen, una mano pequeña y pareja había escrito, en francés, «conclusión correcta, método insuficiente: ver p. 214».

Fue a la página 214. La nota tenía razón.

Hojeó el resto como se hojea un log ajeno. El libro estaba anotado margen por margen, sin una página limpia en las secciones técnicas. La mano discutía con los autores: corregía una estimación de temperatura en el modelo de chamuscado («gradiente no considerado»), reordenaba las hipótesis de formación de imagen con una numeración propia que anticipaba, comprobó Vidal con incomodidad creciente, la taxonomía que la literatura recién adoptaría en los noventa. Junto a un pasaje sobre la densidad de la trama, un cálculo al margen, hecho a mano, con el resultado subrayado; Vidal lo rehízo de cabeza y le dio lo mismo. Las minas eran distintas. Al menos tres durezas de lápiz, es decir al menos tres épocas de lectura, la más vieja ya gris de decenios, la más nueva todavía brillante. Alguien había vuelto a este libro durante toda una vida, como se vuelve a un expediente que no cierra.

Buscó un nombre. Guardas, portadilla, colofón. Ni ex libris, ni firma, ni inicial. La única marca de propiedad era la letra misma, que no figuraba en ningún registro que él pudiera consultar y que sin embargo era más identificatoria que una firma: una mano entrenada en química, que leía el inglés técnico sin esfuerzo y anotaba en francés, que discutía con los mejores especialistas de 1978 y les ganaba con veinte años de ventaja. Vidal conocía el género de documento. Era una cadena de custodia intelectual, completa, sin titular.

Cerca del final, en la sección sobre difusión de resultados y prensa, encontró la única anotación sin argumento del libro. Junto a un párrafo que recomendaba coordinar los anuncios públicos entre los equipos, la mano había trazado una doble raya vertical al margen. Nada escrito. La mina era más blanda que las otras y el trazo más hondo: el papel estaba marcado hasta el reverso de la hoja. Vidal pasó el índice por el relieve, del otro lado, como quien verifica una incisión. En quinientas páginas de discusión, el único lugar donde el lápiz había apretado era el único donde no decía nada.

Anotó en el teléfono: «actas STURP, ed. 1981, anotado, s/nombre, 3 minas mín.». Miró la nota. Le agregó: «la doble raya, p. 487». Después miró por la ventanilla un rato, el lago del lado equivocado del tren, y cuando volvió al libro ya estaban entrando en Lausana y no lo abrió de nuevo.


\scenebreak


Llegó al estudio de la Universitätstrasse a las siete y media, con el tranvía 6 chirriando la curva abajo. El contrato fue a la caja de los contratos, entrada nueva en el registro, código y fecha, dos minutos. El libro no fue a ninguna parte durante una hora.

El inventario del estudio tenía una entrada por cada objeto físico, ciento noventa y cuatro en total, cada una con procedencia, propietario y disposición final prevista. Vidal abrió una ficha nueva. En «objeto» escribió: actas STURP, anotado. En «procedencia» empezó a escribir «Fondation Cassiodore» y se detuvo, porque la Fondation era la vía y él acababa de pasar dos semanas estableciendo que la vía no era la fuente. En «propietario» estuvo más tiempo. El fundador había dicho que el objeto no requería devolución, lo cual describía una obligación ausente, la disposición de una cosa, y omitía con precisión de mandato la única palabra que la ficha pedía. Un regalo transfiere propiedad. Un préstamo sin plazo la retiene. La frase transmitida no era ninguna de las dos cosas y era verificablemente ambas.

Dejó los dos campos vacíos. El sistema pintó la entrada de amarillo.

Había otra entrada amarilla en el inventario, una sola, de marzo de 2029: «soporte, misceláneo», el disco donde vivía la copia primaria del contenedor 2029-03, el directorio de Amberes, la ficha que él había redactado imprecisa a propósito y nunca corregido. Abrió el cajón donde estaba el disco, debajo de los manuales de calibración, al lado del cuaderno de tapas grises. Puso el libro encima del disco y el cajón cerró igual, con dos milímetros menos de juego.

Ciento noventa y cinco objetos en el estudio. Ciento noventa y tres clasificados. Los dos restantes compartían cajón, y compartían otra cosa que Vidal formuló de pie, con la mano todavía en la manija: eran los dos únicos objetos de su vida profesional que sabían algo de él que él no había dicho. El disco sabía lo que había excluido en Amberes. El libro sabía quién lo había leído a él antes de elegirlo, porque un hombre que anota así los márgenes de un libro anota así los márgenes de todo.

En doce días empezaba el traslado. Faltaba conocer al equipo que otro había elegido, en un laboratorio que otro había equipado, dentro de un calendario que otro había escrito. Lo que quedaba de él en la operación era el protocolo, la firma y el veto. Se dijo que era suficiente. Contó los objetos del escritorio, que eran seis, y los volvió a contar.

\clearpage

\chapter{Collegno}

\lettrine[lines=2, lhang=0.1, nindent=0.2em]{A}{bril,} ocho días antes del traslado. La sala blanca de Collegno mantenía veinte grados con dos décimas de oscilación y Vidal llevaba cuarenta minutos tratando de encontrarle un defecto.


El edificio, por fuera, era un taller de mantenimiento electromecánico en la segunda fila de un polígono industrial: chapa gris, un cartel con un nombre que existía en el registro de comercio, dos camionetas con el mismo nombre en la puerta. Por dentro, alguien había construido una caja dentro de la caja. Piso flotante desacoplado de la losa, esclusa de doble puerta, presión positiva, línea de nitrógeno con caudalímetro propio. La mesa de sujeción medía cinco metros con veinte, gran formato real, con sujeción por depresión regulable en dieciséis zonas independientes. El lienzo medía cuatro con cuarenta. Sobraban ochenta centímetros por diseño, no por descuido: los certificados de la mesa citaban el largo del lienzo con dos decimales.

Vidal revisó los certificados de calibración uno por uno. El banco hiperespectral traía firma de un laboratorio de metrología de Ivrea, fecha 14 de marzo, trazabilidad al patrón nacional. Lo verificó igual: montó su tarjeta de referencia, corrió la secuencia, midió. Doce micrones por píxel. En una nave industrial de Collegno, con el climatizador encendido, doce micrones. En Zúrich, con todo a favor, él llegaba a doce en días buenos.

---El registro de vibración ---dijo.

---Acelerómetro instalado bajo la mesa desde el lunes ---dijo Ledda---. Los datos se vuelcan a papel cada mañana. Copia única. La conserva usted.

Su objeción número siete del expediente, cumplida antes de que llegara a pedirla. Ledda tenía una tablilla con hojas sueltas y contestaba en el mismo formato en que Vidal preguntaba: sustantivo, plazo, responsable. En veinte minutos de recorrida no había usado un solo adjetivo. El formulario de auditoría de equipos que Vidal traía impreso tenía treinta y una líneas y a las 11:40 estaban todas tildadas. Volvió al principio de la lista y la recorrió de nuevo buscando la línea que hubiera tildado por inercia. La secuencia cerró dos veces.

Ninguna universidad le había dado esto. Ningún museo, ninguna aseguradora, ningún gobierno. Se lo había dado un nombre que él había aceptado no conocer, y el orden correcto de esa frase era el problema, así que la dejó sin ordenar.

La puerta de la esclusa hizo su doble ciclo y entró una mujer baja con el pelo atado con un hilo de algodón. Traía dos rollos de entretela bajo el brazo izquierdo y los apoyó en la mesa auxiliar antes de acercarse.

---Chiara Fabbri ---dijo---. Vos sos el de la máquina.

---Sebastián Vidal.

---Ya sé. ---Miró el banco hiperespectral, después el escáner de mano que Vidal tenía sobre la mesa, el portátil de dos kilos con que hacía los barridos preliminares---. ¿Ese es el que llega a doce micrones?

---En banco. En mano, veinte.

---¿Me dejás?

Vidal se lo alcanzó. Ella lo tomó con las dos manos, una bajo el cuerpo del instrumento y otra en el lateral, y lo sostuvo así el tiempo entero que lo estuvo mirando. Dos kilos. Un agarre bimanual era innecesario para ese peso. Vidal lo registró como registraba las cosas que funcionaban fuera de especificación sin fallar: un dato sin casilla. Tenía una cicatriz fina en el índice izquierdo, de bisturí, vieja.

---La mordaza de abajo ---dijo ella, girando el escáner--- te marca la tela si trabajás a menos de cinco milímetros. Este borde. ¿Ves el canto vivo?

---El instrumento no toca la superficie.

---El instrumento no. La mano que lo sostiene ocho horas, sí. ---Se lo devolvió, también con las dos manos---. Le hacemos un perfil de fieltro y listo. No es crítica, es barata.

Vidal pasó el pulgar por el canto. Vivo, en efecto. Once años usando ese modelo y el borde nunca había figurado en ninguna de sus listas, porque sus listas medían el instrumento y el instrumento cumplía. Anotó «perfil de fieltro, mordaza inferior» en la columna de pendientes, y le llamó la atención que la primera entrada de esa columna en todo el día viniera de una persona y no de una máquina.

El perfil sin nombre de noviembre decía Opificio delle Pietre Dure, escuela directa de la intervención de 2002. La mujer coincidía con el perfil. Coincidencia verificada.

A las doce y diez la esclusa cicló de nuevo y Vidal supo quién era antes de que terminara de abrirse la segunda puerta, por la altura y por la corbata de lana, y empezó a contar. Uno. Ferrero se detuvo a dos metros de la mesa. Dos. Tres. Miró el banco hiperespectral, la mesa de cinco metros, la línea de nitrógeno. Cuatro. Cinco. Seis. Al séptimo segundo dijo:

---Usted necesita mis ojos. Yo necesito que esto no se haga sin alguien que sepa lo que puede romper.

---Sus dieciséis mil placas ya están en mi corpus.

---Las placas son lo que el lienzo era el día que las tomé. Los ojos saben lo que fue haciendo desde entonces. ---Se acercó a la mesa y apoyó una mano en el borde, sin tocar la superficie de trabajo---. Buena mesa. En 2002 lo sostuvimos con pesas de tela rellenas de arena, como en 1978. Cincuenta años de progreso para llegar a lo que cualquier tapicero le habría dicho gratis.

Chiara se rio corto, desde la mesa auxiliar. Ferrero la saludó con una inclinación de cabeza y un «Chiara» sin apellido; ella contestó «dottore» con una cortesía tan exacta que Vidal la archivó también, otro dato sin casilla. Entre noviembre y hoy había habido una carta de cinco considerandos, un dictamen técnico favorable citado como agravante, una firma de Ferrero elevando la solicitud con nota de mérito. Nada de eso entró a la sala. Vidal esperó a ver si Ferrero lo nombraba, contó hasta donde llegó la pausa, y la pausa terminó en otra cosa.

---¿Cuándo llega su modelo? ---dijo Ferrero.

---Los racks viajan conmigo el nueve.

---Entonces tenemos cuatro días para pelearnos por el protocolo con la sala vacía. Es el orden correcto. Después del doce, cada pelea cuesta horas de lienzo.

Almorzaron en la antesala, sándwiches de una bolsa de papel que había traído Ledda, y a la una y media se sentaron los cuatro a la mesa de reuniones, que era una puerta sobre caballetes con las veintidós páginas del protocolo en cuatro copias. Ferrero se puso los anteojos, leyó el punto 3 completo en silencio, y cuando llegó a «tres miligramos de fibra del borde lateral» sacó del bolsillo del saco un par de guantes de algodón de conservación y se los puso. Las manos enguantadas quedaron quietas sobre el papel. La voz salió firme.

---Custodia compartida de toda muestra. Toda. Desde el segundo en que la fibra se separa del lienzo hasta el segundo en que se reintegra o se documenta su consumo. Una caja de dos cerraduras, una llave suya, una llave mía. Ningún traslado, ninguna preparación, ningún montaje en portaobjetos con un solo par de manos presente.

---El protocolo ya prevé registro de cada manipulación ---dijo Vidal.

---Un registro dice lo que pasó. Dos llaves impiden que pase. ---Ferrero dio vuelta la página sin esperar respuesta---. Usted va a tomar tres miligramos y para usted van a ser tres miligramos. Yo necesito poder decirle a quien me lo pregunte, dentro de veinte años, dónde estuvo cada microgramo cada minuto. Porque me lo van a preguntar. A usted no. Usted se vuelve a Zúrich. ¿Quién paga la diferencia entre su registro y mi respuesta?

Vidal revisó la cláusula contra su propio esquema de cadena de custodia. La doble llave agregaba entre cuatro y siete minutos por manipulación. Multiplicado por el plan de muestreo, costaba menos de tres horas en diecisiete días.

---Aceptado. Con un ajuste: si una de las dos llaves no está disponible en ventana crítica, el veto lo resuelvo yo y la discrepancia queda registrada.

---Registrada con mi firma en contra, cuando corresponda.

---Eso es lo que dije.

---No exactamente ---dijo Ferrero, y por primera vez algo parecido a humor le pasó por la cara---. Usted dijo «registrada». Yo dije «con mi firma en contra». Son dos documentos distintos el día que esto llegue a un tribunal.

Ledda esperó a que los dos terminaran de anotar y puso sobre la mesa una sola hoja escrita a máquina.

---Reglas de operación. Sin excepciones desde hoy hasta la devolución. Efectivo para todo gasto; el circuito de billetes ya está armado y no se toca. Papel para todo registro; nada digital sale de esta sala, nada digital entra sin inventario. Teléfonos en la caja de la esclusa. Y personal con historia limpia real: cada persona que cruce ese portón tiene dieciocho, veinte, treinta años de vida verificable, empleadores que existen, domicilios que existen. Ninguna leyenda, ninguna cobertura.

---Una cobertura profesional resistiría un control de rutina ---dijo Vidal.

---Una identidad falsa no resiste 2033. Reidentificación en cámaras, cruce de contratistas por los Carabinieri antes del evento papal, historial de pagos de veinte años consultable en horas. Lo que resiste es lo contrario. ---Ledda alineó la hoja con el borde de la mesa---. Una identidad verdadera con motivos ocultos, sí. Esa resiste todo. Los sistemas verifican historia. Los motivos quedan fuera de alcance del procedimiento.

Chiara tenía la taza de café sostenida con las dos manos y la bajó sin ruido hasta la mesa. Ferrero se quitó un guante para firmar la hoja de reglas. Vidal la firmó después, con su letra de cadena de custodia, y mientras firmaba hizo la cuenta que venía postergando desde las once de la mañana: en esa sala había cuatro personas y él había elegido a cero. Había redactado el requisito de conservación textil; otro había elegido a Fabbri. Había exigido seguridad con historia verificable; otro había elegido a Ledda. A Ferrero no lo había pedido nadie que él conociera, y era, de los tres, el único cuya presencia Vidal podía explicar por mérito técnico puro y el único que había firmado un dictamen en un expediente que terminó en denegación.

La primera concesión, en enero, le había costado tres semanas, dos viajes a Ginebra y una anotación con fecha y hora en el cuaderno de tapas grises. Esta la hizo entre la firma de Ferrero y la suya. Treinta segundos, ninguna anotación. El precio bajaba con el uso, como en cualquier serie.

---El jueves a las nueve ---dijo Ledda, juntando las copias--- llega el lino de ensayo. Cuatro metros y medio, sarga de espina, envejecido térmico. Práctica de manipulación completa antes del doce.

---El lino lo recibo yo ---dijo Chiara---. Nadie lo toca hasta que yo diga cómo.

Nadie discutió. Vidal anotó «jueves 9:00, ensayo» y debajo, por procedimiento, «verificar procedencia del lino de prueba», y tachó la segunda línea antes de cerrar el cuaderno, porque la procedencia del lino de prueba era exactamente el tipo de pregunta que el contrato le había enseñado a dejar sin ordenar.

El piso franco quedaba en Turín, tercer piso de un edificio de los años sesenta cerca de la estación Dora, escalera con olor a cera vieja y una cocina donde todo funcionaba y nada combinaba. Sobre la mesa había un sobre con efectivo para dos semanas, contado y fajado, un juego de llaves de la nave de Collegno y la hoja del día siguiente, mecanografiada.

Vidal contó el efectivo, coincidió, y leyó la hoja mientras el agua se calentaba. 9:00, taller textil, ensayo de manipulación sobre lino de prueba. Responsable: Fabbri. Debajo, la lista de asistentes: Ferrero, O.; Ledda; Vidal, S.

Leyó su apellido en tercer lugar de una lista alfabéticamente incorrecta. Buscó la columna de función, la que en todos sus proyectos decía quién dirigía y quién asistía.

La hoja no tenía esa columna.

\clearpage

\chapter{Dos manos}

\lettrine[lines=2, lhang=0.1, nindent=0.2em]{E}{l} lino de ensayo llegó a las 9:07, en un tubo de PVC de dos metros con tapas termoselladas, y Chiara lo hizo esperar.


---Cuarenta minutos ---dijo---. Viene de otra humedad.

Apoyó el tubo en dos caballetes forrados, sin abrirlo, y colgó un higrómetro del caballete izquierdo. La hoja del día decía «Responsable: Fabbri», y Vidal había venido preparado para verificar la promesa del perfil de noviembre: manos de primer nivel, escuela directa de 2002. Cuarenta minutos de espera eran una buena práctica de taller. La anotó como tal.

Ledda confirmó los plazos desde la esclusa, teléfono ya en la caja: práctica completa hoy, repetición el sábado, perímetro toda la tarde, regreso a las cinco. Nadie le preguntó qué perímetro. Ferrero llegó a las nueve y media con la corbata de lana y se paró a dos metros de la mesa, las manos atrás.

Chiara usó la espera. Del borde del tubo, donde el proveedor había dejado flecos francos, cortó un hilo de quince centímetros y se lo dio a Vidal.

---Tirá. Despacio, de las dos puntas.

El hilo aguantó, aguantó, y cedió sin aviso, con un chasquido que Vidal sintió más en los dedos que en el oído.

---¿Cuánta fuerza hiciste?

---No la medí. ---Calculó de memoria---. Menos de un kilo.

---Ese hilo tiene un envejecido térmico de laboratorio. El otro tiene cuatro siglos de pliegues, un incendio y costuras encima. Dividí tu kilo por lo que te parezca prudente y después dividilo otra vez. ---Recogió el hilo roto y lo guardó en un sobre, como si también eso tuviera cadena de custodia---. El número que importa no está en tu protocolo. Está acá. ---Levantó las dos manos, sin drama, y se puso a preparar el fieltro.

A las 9:47 abrieron el tubo. Olor a pan tostado, tenue, el perfume del envejecido térmico; abajo, el siseo continuo de la presión positiva y el clic del caudalímetro de nitrógeno cada tanto, sin período fijo. Vidal lo había cronometrado el martes: entre once y diecinueve segundos. Hoy no lo cronometró.

El despliegue llevó una hora y diez. Chiara dirigía con frases de cuatro palabras: la tela no se levanta, se acompaña; el rodillo va a la tela, la tela nunca al rodillo; el peso se pone, no se apoya. Las pesas de arena ---Ferrero había dicho que en 2002 se usaron iguales, y en 1978--- se distribuían de a dos, simétricas, y la depresión de la mesa se activaba por zonas, de adentro hacia afuera, dieciséis válvulas en un orden que ella dictó de memoria y que coincidía con el que Vidal había escrito, salvo en las zonas seis y once, invertidas.

---El protocolo dice seis antes que once ---dijo Vidal.

---El protocolo supone que la tela está seca pareja. Esta esquina quedó cuarenta minutos más cerca del caño de clima. ---Tocó el lino con el dorso de un dedo enguantado, dos centímetros afuera del borde de trabajo---. Once primero. Corregilo si querés.

Lo corrigió. Escribió «inversión 6/11, causa: gradiente de humedad local» y le llamó la atención que la causa fuera verificable y que él hubiera necesitado que se la dijeran.

A las once ensayaron el plegado de transferencia. Vidal ejecutó el suyo siguiendo su propia hoja, ángulo y radio dentro de especificación. Chiara ejecutó el espejo en el otro extremo. Después llevó la lupa de conteo a los dos pliegues, sin comentario, y lo dejó mirar a él primero.

Sobre la línea de su pliegue: siete fibras quebradas por centímetro, blancas, erizadas contra la trama como pasto pisado. Sobre la línea de ella: una. Quizás dos, si la segunda era rotura y no torsión.

---El radio era el mismo ---dijo Vidal.

---El radio sí.

Ferrero se acercó a mirar por la lupa. Tardó en hablar, y cuando habló lo hizo hacia el lino, con el usted de siempre.

---Siete por centímetro en un ensayo, doctor. En abril, cada una de esas fibras tiene dueño hace dos mil años o hace setecientos, según usted gane o pierda. En cualquiera de los dos casos, ¿a quién se las explica? ---Se enderezó---. La Comisión revisa el cronograma del traslado a las dos. Conviene que del lado de la puerta que se abre con nombre y apellido haya alguien de la casa.

Se fue al mediodía. Vidal repitió el plegado dos veces más antes de parar. Cinco fibras, después cuatro. La curva mejoraba y el residuo seguía siendo ella: lo que las manos de Chiara sabían y las suyas todavía no, una variable sin unidad, sin certificado de Ivrea, sin trazabilidad. La dejó anotada como «diferencial de ejecución» porque el registro exigía un nombre y ninguno de los disponibles era exacto.

Antes de cortar para comer, Chiara ordenó las tiras de fieltro sobrantes, alineándolas por largo contra el canto de la mesa, aunque iban a la basura igual. Recién cuando estuvieron paralelas preguntó:

---¿Cuántas horas de tela tenés vos? Manos, digo. Sin escáner en el medio.

Vidal empezó a calcular y el cálculo le devolvió un número que decidió no redondear hacia arriba.

---Las suficientes para saber cuándo llamar a alguien como usted.

---Va bene. ---Tiró el fieltro---. Por lo menos eso lo medís bien.

Comieron en la cocina de la nave, donde todo funcionaba y nada combinaba, igual que en el piso franco: la misma mano había equipado los dos lugares. Ledda comió en nueve minutos, pan con algo, de pie, y salió al perímetro. Chiara sirvió sopa de un termo, cruzó apenas las manos sobre la mesa, bajó la cabeza y dijo dos frases en italiano, en voz baja, sin mirar a nadie; la segunda terminaba en un nombre. Después cortó el pan y le pasó la mitad a Vidal como si las dos operaciones fueran contiguas. Vidal miró el reloj de la pared. Andaba bien.

---Le hago una consulta técnica ---dijo él, porque el silencio pedía relleno y el relleno, ya que existía, podía rendir---. De las once hipótesis de formación de imagen, la que mejor pondera para un origen medieval es el bajorrelieve calentado. Explica el negativo, explica la gradación, y es reproducible con medios del siglo catorce. Si yo tuviera que apostar por un falsificador, apostaría ahí.

---Perdés.

---Depende de qué llames perder. Es la hipótesis con mejores antecedentes de\ldots{}

---Un chamuscado fluoresce. ---Sopa, sin apuro---. Las quemaduras del treinta y dos fluorescen rojizo bajo ultravioleta. Está fotografiado desde el setenta y ocho, Miller y Pellicori, podés bajarlo esta noche. La imagen no fluoresce. Nada. Si el mecanismo fuera calor, tendrías la misma firma en los dos lugares.

---La intensidad del calentamiento podría\ldots{}

---Segundo problema. El calor entra. Yo cosí linos quemados, tres, de dos iglesias y una colección privada. El chamuscado tuesta el hilo hasta el alma: cortás una fibra y está marrón adentro. Las fibras de imagen están coloreadas en la superficie y blancas adentro. Eso también está publicado. ---Dejó la cuchara---. Y tercero: el chamuscado fragiliza. Una tela con quinientos años de imagen-chamuscado se te quiebra en las líneas de imagen. Esa tela se plegó y desplegó siglos y las líneas aguantaron. Tu falsificador tendría que haber calentado sin fluorescencia, sin penetración y sin fragilizar. Ya no es un bajorrelieve. Es un milagro al revés.

Vidal cotejó de memoria. Las tres citas estaban en su corpus; las tres auditaban limpias; el argumento era el que él habría aceptado de un revisor de primer nivel, y lo acababa de dar una mujer que treinta segundos antes rezaba sobre una sopa de termo. Tenía dos carpetas para ella desde noviembre: creyente, con sesgo esperado a favor de la autenticidad, y técnica, con datos descontables del sesgo. El argumento era técnico, la conveniencia era devota, las dos flechas apuntaban al mismo lado y ninguna contaminaba a la otra. Probó la carpeta «caso mixto». La etiqueta duró lo que la sopa.

---¿Usted preferiría que la hipótesis fuera cierta o falsa? ---dijo, y era una pregunta con trampa, de las suyas, de las que en once años nadie esquivaba sin entregar algo.

---¿Vos descartás hipótesis según quién te las trae?

---No.

---Entonces esa pregunta no te sirve para nada. ---Juntó los platos---. A las dos seguimos. El rodillo, esta vez entero.

La corrida de la tarde salió completa: mesa, rodillo, mesa, cuatro metros y medio acompañados de punta a punta, dieciséis zonas en el orden corregido, cero pliegues nuevos bajo la lupa. Vidal manejó el escáner de mano sobre el tramo central, ya con el perfil de fieltro que ella le había hecho el primer día, y cuando terminó se lo pasó para que verificara la mordaza.

Chiara lo recibió con las dos manos. Dos kilos de instrumento, pero también había recibido así su birome esa mañana, y el martes una tarjeta de calibración que pesaba menos que un naipe. Vidal ya tenía tres observaciones de la conducta y pidió el dato.

---¿Eso es procedimiento?

---Costumbre. ---Giró la mordaza, revisó el fieltro con el pulgar---. Todo lo que tocás es de alguien que no está.

---¿Regla de quién?

---De la persona que me enseñó. ---Le devolvió el escáner, con las dos manos también---. Estuvo del lado de adentro en 2002. Descosió los parches. Aspiró el reverso. Sabe cómo pesa esa tela con parches y cómo pesa sin parches, y son dos telas distintas. ---Una pausa del ancho de una puntada---. Después le pasaron cosas que no vienen al caso.

Cerró el tema con el estuche del escáner. Vidal lo dejó cerrado. Esa noche, en el cuaderno, lo asentaría bajo «observaciones de taller»: regla mnemotécnica, transmisión oral, folclore de oficio. Útil para predecir conducta de manipulación. Nada que modelar.

A las seis y veinte el lino de ensayo quedó enrollado en su tubo, el higrómetro colgado, la hoja del día firmada. Chiara desconectó el escáner y se puso a enrollar el cable, vuelta sobre vuelta contra el codo, pareja, sin mirarlo, con el cuidado de quien guarda algo que no es suyo. Vidal registró el patrón del día completo: primero las manos ordenaban algo, después llegaba una pregunta que no figuraba en ninguna hoja.

---¿Y si el resultado te dice que es del siglo XIII? ---dijo Chiara.

---Entonces es del siglo XIII.

---Te pregunto qué hacés vos. No qué hace el modelo.

Vidal revisó una calibración que ya había revisado.

---Publico ---dijo.

La respuesta había salido en presente, automática, de un archivo anterior a enero, anterior a la cláusula octava, y los dos la dejaron sobre la mesa sin tocarla.

---¿Y si te dice que no puede explicarlo?

---Los modelos no dicen eso.

Ella terminó con el cable, lo colgó del gancho, y recién entonces lo miró.

---Los tuyos tampoco, ¿no?

Se sacó la banda elástica, se volvió a atar el pelo con el mismo movimiento, y juntó su campera del respaldo de la silla. En la esclusa se dio vuelta, la mano ya en la caja de los teléfonos.

---El sábado a las nueve. El rodillo lo armás vos solo, yo miro. Si rompés más de dos fibras por metro, lo repetimos hasta que no.

---¿Y si no llego a dos?

---Llegás. ---Sacó el teléfono, lo prendió, esperó la señal---. Vos aprendés rápido lo que se puede medir.

La puerta exterior cerró con su doble golpe de burlete y Vidal se quedó en la sala blanca con el siseo de la presión positiva y el clic del caudalímetro. Contó los intervalos. Trece segundos, diecisiete, doce. Dentro de rango.

Fue hasta la mesa y pasó la lupa una vez más por el pliegue de la mañana, el primero, el de las siete fibras. Seguían ahí, blancas, erizadas. Un ensayo se documenta con sus fallas; eso lo hacía desde los treinta años. Midió la línea de rotura con la reglilla de la lupa, anotó la densidad exacta ---6,8 por centímetro, había redondeado mal a la mañana--- y corrigió el registro con la fecha del ajuste al lado, como correspondía.

Después buscó en la memoria del día la frase de ella sobre el modelo y la sometió al procedimiento estándar para afirmaciones ajenas: ¿era verificable? Los modelos no dicen «no puedo explicarlo»; dicen intervalos anchos, confianzas bajas, banderas de datos insuficientes. Decir «no puedo explicarlo» exigiría una capa que evaluara la propia evaluación, y esa capa, en Tamiz, era la de convergencia, y la de convergencia era la que devolvía sus razones en un formato que ninguna herramienta traducía del todo. De modo que la respuesta técnica correcta a la pregunta de Chiara era: mi modelo puede llegar a un estado que equivale a «no puedo explicarlo» y carece de las palabras para decírmelo. Y la pregunta de ella no había sido sobre el modelo. Había sido la otra, la primera, la que él contestó con «publico» usando un verbo que la cláusula octava había vuelto condicional dos meses atrás.

Nadie le había hecho esa pregunta antes. Ashkenazi le había preguntado qué diría el informe. Ferrero le preguntaba quién pagaba los resultados. Ansermet no preguntaba. El fundador, el mandante, la firma detrás de la firma, compraba las respuestas antes de que existieran. Chiara le había preguntado qué hacía él, con el modelo puesto entre paréntesis, y él había tardado un cable entero en no poder contestar.

En el piso franco, esa noche, el cuaderno de tapas grises quedó abierto en la mesa de la cocina más tiempo del que llevaba escribir lo que escribió. «Ensayo 1: completo. Diferencial de ejecución: alto, decreciente. F. objeta densidad de rotura: fundada.» Y abajo, en la misma letra pareja: «Fabbri: sin carpeta.»

Miró la última línea. Era imprecisa, y su sistema no admitía entradas imprecisas: una persona sin carpeta era una variable sin unidad, otra más, la segunda del día. La tachadura habría sido lo correcto. Dejó la línea intacta, cerró el cuaderno y lo guardó, y mientras el agua se calentaba leyó la hoja del día siguiente, mecanografiada, que alguien había deslizado bajo la puerta: 9:00, reunión de arquitectura de la operación. Responsable: Ledda. Tema: ventanas.

\clearpage

\chapter{Ventanas}

\lettrine[lines=2, lhang=0.1, nindent=0.2em]{L}{edda} había pegado el plano de la catedral a la pared de la cocina con cinta de papel, cuatro tiras por lado, y encima había clavado una hoja de calco con líneas rojas que no cruzaban ninguna pared. Eran trayectos de personas.


---Treinta horas ---dijo---. Del once a las siete de la mañana al doce a la una de la tarde. Personal reducido: cuatro de conservación, dos de la Comisión, un electricista de la empresa de la teca, dos de seguridad diocesana. Los turnos están en el expediente que usted ya leyó. Lo que no está en el expediente es esto.

Marcó con el dedo un rectángulo pequeño a la izquierda del crucero.

---La sacristía. El lienzo sale de la caja actual acá y entra a la teca nueva acá, y entre una cosa y la otra hay un intervalo que en el cronograma oficial se llama «verificación de soporte y limpieza de superficie de apoyo». Duración estimada: cincuenta minutos. Duración real de las últimas tres operaciones equivalentes documentadas: entre setenta y ciento diez.

Vidal buscó el número en su copia. Estaba, en la página once, con la palabra «estimada» y sin la palabra «real».

---¿De dónde saca los ciento diez?

---Informes de intervención de 2002 y dos traslados de textiles de gran formato en Florencia, 2019 y 2028. Los tres los firmó gente formada en la misma escuela que Fabbri.

---Eso no es un dato del once de abril.

---No. Es la distribución. Yo trabajo con distribuciones ---Ledda no levantó la vista del plano---. La sustitución ocurre dentro de esa cola. Necesitamos cuarenta minutos con dos personas nuestras adentro y el resto del personal fuera de línea de vista. Fabbri manipula. Ferrero certifica el estado ante la Comisión, porque su firma en esa acta es lo que hace que nadie vuelva a mirar la tela hasta el veintinueve.

Ferrero, sentado del otro lado de la mesa con el pocillo vacío entre las manos, se tocó la muñeca izquierda.

---Y la salida.

---Tubo de transporte idéntico al de ingreso, misma rotulación, mismo peso bruto. Sale con la carga de material de embalaje. Camioneta rotulada, chofer con historia laboral verificable de once años en la empresa, sin observaciones. El chofer no sabe nada y no tiene que saber nada: transporta descartes de obra, y eso es exactamente lo que transporta según el remito.

---El remito es falso ---dijo Ferrero.

---El remito es incompleto.

Chiara, contra la pared, tenía el termo en las dos manos y no había servido a nadie todavía.

---¿Y la vuelta? ---preguntó.

Ledda pasó a la segunda hoja de calco.

---Veintinueve de abril. Inspección final de la teca antes de la ostensión: acceso técnico autorizado, presencia de la Comisión, presencia de la Custodia, presencia probable de un monseñor que todavía no está confirmado en la lista. Dos ventanas. La primera, once minutos, durante el cambio de guardia de las diecinueve horas: relevo de dos puestos y traspaso de bitácora, con solapamiento nominal de cero. La segunda, seis minutos, durante la revisión programada del sistema de clima, que exige despresurizar la campana y aislar los sensores de la línea de alarma.

---Once es casi el doble de seis ---dijo Vidal.

---Once tiene doce personas moviéndose en un radio de cuarenta metros. Seis tiene dos técnicos que ya están adentro por procedimiento y una alarma silenciada por procedimiento. Yo recomiendo la de seis. Menos variables humanas.

---¿Y el margen? ---Vidal contó la línea de tiempo en la hoja: seis minutos, en la escala del dibujo, medían siete milímetros---. La operación de reinserción, con la tela plana, sin doblar, con el marco de la teca en su sitio.

---Cuatro minutos con veinte segundos en el ensayo del martes, con Fabbri y una asistente. ---Ledda hizo una pausa que no significaba nada---. Ese es el número que tenemos.

---Entonces el margen es de un minuto cuarenta.

---Uno con cuarenta.

Vidal lo escribió. Después dividió mentalmente: un margen del veintiocho por ciento sobre el tiempo de ejecución. En cualquier informe suyo, un intervalo con esa cola habría llevado una nota al pie.


\scenebreak


Salieron a las tres y media, separados, con veinte minutos de diferencia entre uno y otro, y se encontraron en la piazza San Giovanni haciendo lo que hace cualquiera: mirar la fachada. Llovía poco, apenas mojaba el empedrado, y el olor a piedra fría venía del atrio.

Ledda caminó primero por el lado norte, sin señalar nada, hablando bajo y de frente, nunca hacia una cámara.

---Cinco fijas en la fachada, tres en la plaza, dos en el acceso de servicio. Las de la plaza son municipales, con reidentificación por marcha y silueta desde el veintinueve. No importa si usted se pone anteojos. Importa cómo camina.

---¿Cuánto guarda el municipio?

---Noventa días. Los Carabinieri, para un evento con visita papal, piden extractos de doce meses y los cruzan con los registros de contratistas. Cada persona que entra a esa obra tiene nombre, número fiscal, cobertura médica y una fecha de alta. Lo que buscan no son caras conocidas. Buscan repeticiones.

Doblaron por vía XX Settembre. Un andamio cubría el flanco del palacio arzobispal y dos operarios bajaban tablones con una cuerda; abajo, un tercero contaba los tablones en voz alta, en dialecto.

---Enumere lo que nos mata ---dijo Vidal.

---Tres cosas. Un pago rastreable: una transferencia, una tarjeta, un pasaje comprado con nombre. Por eso todo es efectivo y por eso el efectivo viene de un circuito armado hace catorce meses, con billetes que no salieron juntos de ningún cajero. Segunda: una cara que aparece dos veces donde no corresponde. La misma persona en el perímetro el nueve y el veintinueve. Los sistemas no reconocen a nadie: reconocen coincidencias. Tercera, y es la que nos mata de verdad, la casualidad.

---Defina casualidad.

---Un incendio de un tacho a doscientos metros. Una manifestación que corre el cordón. Un tipo que se descompone en la fila y hace que la seguridad reasigne dos puestos a las diecinueve horas del veintinueve. Contra las primeras dos yo trabajo. Contra la tercera solo tengo la ventana de seis minutos, que es corta justamente porque las casualidades necesitan tiempo para llegar.

Vidal contó las cámaras del acceso de servicio: dos, con carcasa nueva, cableado embutido, sin la mugre que dejan cuatro inviernos. Instalación reciente.

---Esas son del proyecto de la teca ---dijo Ledda, sin que él preguntara---. Vigilancia de obra. Las paga la donación.

Se detuvo. Después siguió caminando, porque detenerse en el perímetro era una de las cosas que Ledda había prohibido en la primera reunión.

Las especificaciones de la teca las habían aportado los ingenieros de la donación. La telemetría de la teca estaba en la pestaña seis del expediente, con protocolos de alarma e histórico, y él la había leído en Ginebra pensando que era generosidad de un anexo técnico bien armado. El calendario de las treinta horas de sacristía estaba en la pestaña dos porque alguien que había financiado el traslado sabía cómo iba a organizarse el traslado. Cuatro horas de estabilización de atmósfera: un plazo insuficiente para lino con esa historia térmica, y un plazo que él, al objetarlo, había corregido a doce, y las doce horas eran ahora doce horas de lienzo fuera de la caja.

No compraron la reliquia. Compraron la ventana.

Lo revisó en orden, como se revisa una cadena de custodia. 2031: la Comisión, endeudada la diócesis, acepta una donación anónima canalizada por una asociación de amigos del lienzo. La donación construye la teca. La teca exige un traslado. El traslado exige manipulación. La manipulación exige personal reducido, y el personal reducido exige una lista, y la lista la propone quien puso las especificaciones. Después, la inspección final, que también está en el proyecto de la teca, porque un equipo nuevo se inspecciona antes de exponerlo a dos millones de personas. Las dos ventanas del veintinueve no eran una debilidad del sistema. Eran el sistema.

Dos años. Alguien había escrito esto dos años antes de que a él le denegaran una solicitud de veintidós páginas.

Y era limpio. Esa era la parte que le hizo bajar el ritmo del paso hasta que Ledda tuvo que acomodarse. Ninguna identidad falsa, ningún soborno, ninguna puerta forzada: un donante que dona, una comisión que agradece, un cronograma que se cumple. Toda la operación viajaba adentro de procedimientos verdaderos. Él había pasado once años desarmando falsificaciones y no había visto nunca una construcción de este orden, y la sensación de estar frente a un trabajo mejor que el suyo le duró una cuadra entera antes de que se le ocurriera la pregunta correcta.

Quién lo pensó.

No Ansermet: Ansermet ejecutaba. Ledda había entrado hacía meses. La donación era de 2031, y en 2031 nadie sabía que Vidal iba a pedir acceso, y en noviembre de 2032 ya había dos perfiles mecanografiados con la fecha de impresión encima.

---Ledda.

---Dígame.

---El expediente de la ventana de acceso. Las catorce pestañas. ¿Quién lo armó?

---Me llegó armado.

---¿De quién?

---De la misma persona de la que le llega a usted. ---Ledda cruzó al otro lado de la vereda para evitar el cono de una cámara sin cambiar la cadencia---. Yo pedí tres correcciones. Me las aprobaron el mismo día. Doctor, en dieciocho años nunca me aprobaron nada el mismo día. Eso también es un dato.


\scenebreak


En Collegno, a las siete y diez, Vidal aprobó por escrito el punto 2 del protocolo: sustitución durante el intervalo de verificación de soporte, ventana operativa de cuarenta minutos dentro de una cola estimada entre setenta y ciento diez, Fabbri responsable de manipulación, Ferrero responsable del acta ante la Comisión. Chiara leyó el punto entero antes de firmar y pidió una modificación: la gemela entraba a la sacristía ya desplegada y aclimatada, no en tubo. Cuarenta minutos con un desenrollado adentro eran treinta y uno de trabajo real y ella no operaba con treinta y uno.

---¿Y cómo la metemos desplegada? ---preguntó Ledda.

---En la caja de apoyo de la mesa de conservación. Va plana debajo del tablero, con el falso fondo. La caja entra el once a las siete con el resto del equipamiento y nadie la abre porque el que la abre soy yo.

Ledda anotó. Ferrero levantó una mano a media altura.

---Y si el electricista de la empresa de la teca decide que ese día quiere ayudar a mover la mesa, ¿quién lo saca de la sacristía?

---Yo ---dijo Chiara---. Con un motivo textil que él no va a entender.

---Ese es el punto: que no lo entienda. Cuando esto se investigue, y se va a investigar, alguien le va a preguntar a ese hombre qué le dijeron y por qué salió. Y él va a decir la verdad, que es lo peor que puede decir.

---Va a decir que la conservadora le pidió que se retirara mientras se manipulaba el textil. ---Ledda seguía escribiendo---. Es el procedimiento normal y consta en los tres traslados que le mencioné. La declaración de ese hombre nos protege.

Ferrero se tocó la muñeca izquierda otra vez, no dijo nada más, y firmó donde correspondía firmar.

El punto 3, devolución, quedó abierto.

Vidal lo escribió con la letra de siempre, mayúsculas alineadas, la misma que usaba en las cadenas de custodia: «Punto 3. Restitución. Ventana A: 11 min (cambio de guardia, 19:00). Ventana B: 6 min (revisión sistema de clima). Tiempo de ejecución medido: 4 min 20 s. Recomendación de seguridad: B. Estado: pendiente de decisión.»

Ledda leyó la hoja al revés desde su lado de la mesa.

---Yo la recomendé.

---Sé que la recomendó. La aprobación es mía y todavía no la tengo.

---¿Qué le falta?

---Tres repeticiones más del ensayo, con el marco de la teca real y con la presión negativa. Y el histórico de la telemetría de los últimos noventa días, para ver cuántas veces esa revisión de clima se corrió de horario.

---Se lo pido mañana.

---Pídalo hoy.

Ledda salió. Ferrero se puso el saco y se quedó un momento con la mano en el respaldo de la silla, mirando la hoja del protocolo desde arriba, y después se fue también, sin decir buenas noches, que era su forma de decirlas.

Chiara juntó los pocillos y empezó a enrollar el cable del hervidor. Lo enrolló con más vueltas de las necesarias.

---«Pendiente de decisión» ---dijo.

---Es lo que es.

---Los otros dos puntos los aprobaste hoy.

---Los otros dos tienen los datos completos.

Terminó con el cable, lo colgó del gancho, y recién entonces se dio vuelta.

---El once entramos. El veintinueve salimos. Si el veintinueve no sale, ¿el once qué fue?

Vidal miró la hoja. La frase «pendiente de decisión» le había quedado en el renglón de abajo, sola, con espacio en blanco a los dos lados, y no había en la página ninguna otra entrada sin resolver.

---Depende de qué llames\ldots{}

Se detuvo. Chiara esperó, porque ella no interrumpía. Afuera, un camión entró al polígono en primera y el ruido subió por el piso flotante hasta la mesa.

---Voy a pedir el histórico yo ---dijo él, y buscó el teléfono en la caja de la esclusa, aunque el teléfono no servía para eso y ella lo sabía.

\clearpage

\chapter{La gemela}

\lettrine[lines=2, lhang=0.1, nindent=0.2em]{L}{a} gemela estaba plegada en cuarenta y ocho dobleces sobre la mesa del taller y el olor a lino chamuscado seguía ahí con el extractor al máximo. Vidal lo medía en minutos: cincuenta y dos desde la última aplicación y la nave todavía olía a eso. Chiara trabajaba con la máscara colgada del cuello, porque para lo que estaba haciendo necesitaba la nariz.


---Piano ---dijo, aunque nadie la estaba ayudando.

Las quemaduras de 1532 se habían replicado por el mismo mecanismo que las había hecho: plegar la tela, apoyar metal fundido, dejar que el calor atravesara las capas y saliera del otro lado en la misma posición. Ferrero había mandado la curva de temperatura por escrito, tres páginas con la aleación, los tiempos de contacto y la advertencia de que la plata a novecientos sesenta grados perdona menos que el fuego. No había venido a verlo. Vidal anotó la ausencia en la hoja del día, en la columna de asistentes, con la raya que usaba para las inasistencias justificadas.

El lino era de un telar manual de Bélgica: sarga de espina tres a uno, treinta y ocho hilos de urdimbre por centímetro, veintiséis de trama, cuatro metros cuarenta por uno trece. Dos meses de tejido. Después, ciento veinte horas de envejecimiento térmico a noventa y ocho grados con ciclos de humedad, cuatro lavados para sacarle el apresto, y un amarillo que ahora se sostenía dentro de un delta E de dos contra las placas calibradas del corpus de referencia. Ese número lo había verificado él, dos veces, con el colorímetro apoyado en catorce puntos distintos de la tela. Dos no era cero. A un metro cuarenta, con el cristal de la teca en medio y la iluminación de la teca, dos era invisible; lo tenía medido con ocho observadores del polígono a los que Ledda les había pagado en efectivo por mirar una tela durante quince segundos y decir cuál de las dos era vieja.

Cuatro habían acertado. Estadísticamente, azar.

Las manchas de agua le habían llevado a Chiara nueve días y eran, según ella, lo más difícil, porque el agua de 1532 había entrado por los pliegues y había arrastrado lo que había en la tela hacia los bordes de cada mancha, y eso no se pinta: se hace pasar. Los parches de las clarisas fueron otra cosa. Desde 2002 los parches no estaban; lo que quedaba en el lienzo eran las perforaciones de las agujas y un halo más claro donde la tela había estado cubierta cuatrocientos setenta años. Chiara las había contado sobre las placas de Ferrero, ampliación por ampliación, y las estaba reproduciendo a mano con una aguja del calibre que correspondía al hilo de 1532.

---Cuatrocientas ochenta y dos en el parche grande de arriba ---dijo---. Voy por trescientas nueve.

---¿Cuánto por hora?

---Cuarenta. Y bajando, porque me tiembla.

Trabajaba de pie, con la lámpara a cuarenta y cinco grados, y cada vez que dejaba la aguja la apoyaba en el cartón y no en la tela. Cuando levantaba el orillo lo levantaba con las dos manos. Se lavaba las manos entre cada tramo de veinte perforaciones aunque la tela era una tela hecha en Bélgica el año anterior y no era de nadie.

La palabra que le venía a Vidal para eso lo incomodaba. Reverente no era una magnitud. La sustituyó por consistencia de procedimiento y siguió.

---¿Por qué la tocás así? ---preguntó.

Chiara terminó la perforación antes de contestar, lo que en ella era una forma de contestar.

---Porque el once mis manos no van a tener tiempo de decidir. ---Enderezó la espalda---. Si acá la agarro de cualquier manera, allá la agarro de cualquier manera.

Aceptable. Verificable, incluso: memoria motora, tasa de error bajo presión, todo eso estaba publicado.

Ella juntó el resto de hilo de lino de la mesa y empezó a enrollarlo en dos dedos.

---¿Vos podrías reconocerla?

---A un metro cuarenta, con ese cristal y esa luz, no.

---Yo sí.

Vidal esperó el mecanismo.

---No por la imagen ---dijo Chiara---. Por cómo cae. La otra tiene seiscientos años de plegados en el mismo lugar y ninguna tela finge eso. ---Guardó el ovillo en el bolsillo del delantal---. Pero yo el veintinueve no voy a estar del lado del vidrio, así que da igual.


\scenebreak


La celda de ensayo la había armado Ledda contra la pared del fondo: bandeja de apoyo con las mismas cuatro células de carga que la teca, sonda de humedad relativa, sensor de oxígeno residual, la misma frecuencia de muestreo, el mismo protocolo de alarma. Los modelos los había conseguido por el proveedor de la teca, que había sido pagado por la donación y facturaba a quien le pidiera un repuesto.

Cuarenta horas de registro. Vidal las leyó en papel, a las seis y cuarto de la mañana, con el café enfriándose al lado.

La masa seca cerraba: cuatro gramos de diferencia sobre mil doscientos veintiséis, dentro de banda. El problema empezaba en la tercera columna. A cuarenta y cinco por ciento de humedad relativa la gemela pesaba seis coma cuatro gramos más que la referencia, y a cincuenta y cinco, once coma dos. La pendiente era distinta. La tela nueva chupaba agua con otra curva, y en la teca esa curva quedaba registrada minuto a minuto durante dieciocho días, con telemetría continua y auditada, en un archivo que la Comisión iba a poder abrir cuando quisiera durante los próximos veinte años.

Vidal repitió la lectura de las cuatro células. Después repitió el cálculo de pendiente, que ya había hecho.

---El respaldo ---dijo Chiara, mirando la hoja por encima de su hombro---. Es la batista. Está tratada.

---El respaldo son once metros de costura.

---Once metros ochenta.

Ledda llegó a las siete y escuchó de pie.

---Tela nueva, sin apresto, dos lavados con agua desmineralizada, envejecimiento corto ---dijo Chiara---. Descoser, medir, coser. Tres días si duermo cinco horas.

---Tres días ---repitió Ledda---. La aclimatación del conjunto baja de cinco a dos.

---Sí.

---El once no se mueve.

---Sé que el once no se mueve.

Vidal cruzó los números: dos días de aclimatación contra cinco de protocolo, sobre un conjunto de dos telas con historias térmicas distintas, y la pendiente de adsorción medida recién después de coser, sin margen para una segunda corrección. Y las repeticiones de manipulación que Chiara tenía asignadas para esos tres días pasaban a cero.

---¿Cuántas te quedan hechas del plegado de transferencia?

---Nueve.

---El protocolo pide quince.

Chiara ya estaba levantando la orilla de la batista con la hoja del descosedor.

---Entonces el protocolo va a tener una discrepancia registrada. ---Sin levantar la vista---: Ponela con mi nombre.


\scenebreak


Ferrero llegó a las siete de la tarde con el saco puesto y del bolsillo derecho le asomaba un puño blanco de algodón, un centímetro de tela doblada. Se sentó frente a la hoja del punto 4 y leyó los tres miligramos como si fuera la primera vez, aunque el número era el mismo que Vidal le había puesto por escrito en Turín en noviembre.

---Tres miligramos del borde lateral, zona sin imagen, adyacente al corte del ochenta y ocho ---dijo Vidal---. Es una fracción de tres diezmilésimas de por ciento de la masa. En el ochenta y ocho cortaron una tira de siete centímetros.

---Y en el ochenta y ocho, doctor, ¿le fue bien a alguien?

---Depende de qué llames\ldots{}

---Le fue bien al laboratorio. ---Ferrero se tocó la muñeca izquierda---. Los otros no publicaron nada.

---El daño es medible. Puedo darte la superficie exacta, la pérdida de resistencia local, el número de fibras comprometidas por milímetro. Eso es lo que hay para discutir.

---Eso es lo que usted tiene para discutir. ---Ferrero sacó los guantes del bolsillo y se los puso, primero el izquierdo, estirando cada dedo hasta el fondo---. Usted toma tres miligramos y para usted es una muestra. Para dos mil millones de personas es una herida en el cuerpo de otra persona. No le pido que lo crea. Le pido que calcule con esa unidad.

---Esa unidad no tiene dimensión.

---Ninguna de las que a usted le importan tiene dimensión, y las usa todos los días. ---Se acomodó el puño del guante---. ¿Y a quién se lo explica usted, después?

Vidal miró la hoja. Tres miligramos era el número que salía de la aritmética: tres alícuotas independientes, ochocientos microgramos útiles cada una, el resto en pérdida de preparación. Con dos coma cuatro las alícuotas bajaban a dos. Dos alícuotas fechan igual de bien y ensanchan el intervalo unos veinticinco años, irrelevante para separar el siglo I del XIII. Lo que se pierde con la tercera es la capacidad de distinguir una señal de una contaminación: con tres, una discordancia se localiza; con dos, una discordancia queda flotando y no hay a qué atribuirla.

Lo escribió al costado, en números, para verlo.

---Dos coma cuatro ---dijo.

Ferrero no agradeció. Preguntó qué perdía, y Vidal le contestó exactamente eso, la tercera alícuota, la localización de discordancias, y Ferrero asintió una vez y dijo que le parecía un precio razonable, sin aclarar de quién.

Se fue a las ocho menos veinte con los guantes todavía puestos.


\scenebreak


El acta la redactó a las nueve y cuarto, en la mesa de la sala blanca, con el caudalímetro haciendo clic cada once o diecinueve segundos.

Antes había hecho una cosa más. Cortó dos milímetros del orillo de la gemela, los procesó y le dio la muestra a Tamiz sin etiqueta. Once minutos. Lino cultivado en Bélgica, cosecha reciente, firma de fertilizante nitrogenado moderno, intervalo 2030--2033, confianza 0,99. El modelo tampoco tuvo que esforzarse en la microscopía: a cincuenta aumentos, la pared celular de una fibra de setecientos años y la de una de dos no se parecen a nada.

«Punto 5. Límites de la unidad de sustitución. Reproduce: apariencia a distancia de inspección igual o mayor a 1,40 m, con cristal e iluminación de la teca; masa seca y masa en equilibrio entre 45 y 55 \% HR dentro de ±1,5 g una vez rehecho el respaldo; respuesta de células de carga, higrometría y oxígeno residual dentro de bandas de alarma. NO reproduce: microestructura de fibra a partir de 50×; edad radiométrica; firma espectral de la zona de imagen. Corolario: cualquier examen por contacto o toma de muestra sobre la unidad la identifica en minutos. En consecuencia, la restitución del original no es opcional.»

Lo leyó en voz alta a los tres y lo firmaron los cuatro. Ferrero escribió debajo de su firma, con letra chica: «Conforme con el texto. Disconforme con el hecho.»

Cuando se quedó solo abrió su registro y anotó la modificación del punto 4: toma reducida de 3,0 a 2,4 mg. En el campo de motivo escribió «margen de seguridad adicional en la manipulación del borde», releyó la frase y la dejó, porque describía un resultado cierto y él no tenía otra manera de nombrar lo que había pasado a las siete y media.

Al margen le puso el círculo chico. No requiere acción.

\clearpage

\chapter{Protocolo}

\lettrine[lines=2, lhang=0.1, nindent=0.2em]{E}{l} protocolo interno quedó consolidado en dieciocho puntos y Vidal los firmó en el orden en que los había numerado. Diez de abril, 9:20, la mesa larga de la nave de Collegno, el caudalímetro haciendo clic contra el silencio a intervalos de once a diecinueve segundos. Cuatro sillas, una carpeta, copias numeradas en papel para cada uno, según la regla de Ledda. Cada punto llevaba la aprobación de Vidal al margen, con fecha, y los que tocaban muestras llevaban además la firma de Ferrero: custodia doble, caja de dos cerraduras, ninguna manipulación con un solo par de manos presente. Eso ya estaba acordado desde marzo. Lo que se firmaba hoy era el cuerpo entero, con el cronograma adjunto.


El cronograma cubría las cuatrocientas ocho horas entre la recepción y la restitución, hora por hora. Día uno: ingreso del tubo a las 13:40, apertura de esclusa, aclimatación de doce horas con registro de humedad cada quince minutos. Día dos: despliegue sobre la mesa de sujeción, responsable Fabbri, activación de las dieciséis zonas de depresión en el orden corregido, de adentro hacia afuera con la seis y la once invertidas. Días tres a seis: barrido hiperespectral completo, doce micrones por píxel, ochenta y dos franjas con solapamiento del quince por ciento. Día siete: toma de muestra, 2,4 miligramos, dos alícuotas, borde lateral, zona sin imagen. Días ocho a catorce: procesamiento, microscopía, corridas de control ciego. Quince y dieciséis: reembalaje y ensayos de restitución con marco real. Día dieciocho, veintinueve de abril: la ventana de seis minutos. Entre el día catorce y el quince había un bloque de dieciocho horas marcado «margen», sin asignación. Ledda lo había pedido así, sin explicar, y Vidal lo había aprobado porque un cronograma sin holgura es una promesa, y las promesas no se auditan.

Leyó los dieciocho puntos en voz alta, uno por uno. Ferrero pidió una coma en el punto once. Ledda confirmó los horarios de relevo contra su propia hoja. Chiara siguió la lectura con el dedo sobre su copia de trabajo y asintió en cada punto que la nombraba.

Firmaron los cuatro. La tinta de Ferrero tardaba en secar y él sopló sobre su firma con una paciencia de otro siglo.

---¿Quién recibe los datos crudos? ---dijo Chiara.

Tenía las manos quietas sobre la mesa. Sin cable, sin hilo, sin nada entre los dedos. La pregunta llegó así, desnuda, y por eso Vidal tardó un segundo más de lo normal en ubicarla: las preguntas difíciles de Chiara venían siempre precedidas de algo que enrollar.

---Cláusula octava del contrato marco ---dijo---. Los datos crudos, los resultados intermedios y definitivos y toda anomalía registrada quedan en custodia de la contraparte, a perpetuidad. Uso académico del experto sujeto a autorización escrita, que no será denegada sin causa fundada. Plazo de respuesta, sesenta días.

---Eso dice el papel. Yo pregunté quién.

---La Fondation Cassiodore.

---La Fondation es una dirección en Ginebra. ¿Quién?

---No es identificable con los medios a mi alcance ---dijo Vidal---. Lo verifiqué. Dos sociedades de mandato, una Anstalt en Vaduz, una fundación en Curazao. El rastro muere en tres jurisdicciones.

Chiara se miró las manos un momento, como si le faltara la herramienta.

---¿Y a vos te parece normal que al que paga no le importe el resultado sino el archivo?

---El resultado está dentro del archivo. Quien exige el archivo completo exige el resultado con su contexto, que es lo que exigiría cualquiera que sepa leer datos. Es la cláusula de un rigor que yo firmaría.

---La firmaste.

---Con correcciones. ---Se oyó a sí mismo y ajustó---: Con tres correcciones de redacción.

Ella asintió despacio. Ledda había empezado a juntar las hojas del cronograma; Ferrero seguía con la vista fija en su firma, esperando que secara.

---Desconfías de ellos desde el primer pago ---dijo Vidal---. Te lo he visto hacer en cada reunión donde aparece el nombre. Muy bien. Dime el mecanismo. Un dato, uno, medible, que sostenga la desconfianza.

---No lo tengo.

---Entonces lo que tienes es una corazonada con el vocabulario cambiado. Desconfiar sin mecanismo es superstición con vocabulario de prudencia. Es exactamente lo mismo que confiar sin mecanismo, con el signo invertido. Y yo no trabajo con ninguna de las dos.

Chiara no discutió. Levantó su copia de trabajo con las dos manos, aunque la copia era suya, la alineó contra el borde de la mesa y volvió al respaldo de batista, once metros ochenta de costura que le quedaban por delante. La aguja retomó su ritmo, un chasquido corto contra el dedal, más regular que el caudalímetro.

Vidal anotó en el registro la hora de cierre de la firma, 10:05, y debajo dejó una línea en blanco. La línea en blanco no correspondía a nada del formato. La miró, decidió que era un error de espaciado, y siguió.

El café lo hizo Ledda a las once, en la cocina donde nada combinaba, y lo tomaron de pie alrededor de la mesa chica porque la grande ya era del protocolo. Olía a café quemado y a cinta de embalar. Ferrero contaba los días que faltaban para el ingreso con un método propio: no los contaba, los llamaba por su nombre litúrgico, y el once de abril era para él la vigilia de algo que se guardó de precisar.

Fue ahí, con las cuatro tazas todavía en la mano, cuando Chiara habló de nuevo.

---Ottavio. Para la línea de base de contaminación. ---Dejó la taza. La dejó con las dos manos---. El inventario de la intervención del dos mil dos. Lo que se aspiró del reverso, los hilos que se retiraron con los parches. ¿Existe completo en algún lado? Publicado, digo. Con destinos.

El caudalímetro hizo clic. Pasó el intervalo entero, y el clic siguiente, antes de que nadie se moviera. Ledda miró a Chiara, después a Ferrero, con la expresión neutra de quien registra un cambio de presión sin conocer la fuente. Vidal contó: dos clics, entre veintidós y treinta y ocho segundos de silencio, y en ese lapso Ferrero terminó su café sin apuro, apoyó la taza, y se acomodó el puño de la camisa sobre la muñeca izquierda.

---El acta de cierre es del catorce de marzo de dos mil tres ---dijo---. El material removido quedó registrado bajo la nomenclatura de la Comisión, categoría de material de conservación en curso, que por definición no se publica hasta que la conservación concluye. Para su línea de base le alcanzan mis placas del reverso, que están calibradas, y el informe de aspirado, que le puedo facilitar.

---Gracias ---dijo Chiara.

---De nada.

La cortesía entre esos dos siempre había sido perfecta. Vidal la tenía medida desde noviembre: un trato de usted invertido, él tuteándola desde la autoridad del oficio, ella devolviendo una amabilidad sin una sola arruga. Lo de ahora había sido lo mismo con un decimal de más. Una respuesta técnicamente completa a una pregunta que pedía otra cosa, aunque Vidal no habría sabido decir qué otra cosa pedía. Registró la temperatura del intercambio como registraba las derivas de un sensor: valor anómalo, causa no asignada, y ninguna columna donde ponerlo. Pensó un momento en el cuaderno de tapas grises, que estaba en Zúrich, en un cajón, debajo de los manuales. Volvió a su taza.

Chiara ya estaba otra vez en la batista. Ferrero se quedó junto a la mesa chica, de espaldas a la sala blanca, mirando por la ventana alta el techo de la nave vecina.

A las cuatro de la tarde, cuando el protocolo firmado ya estaba en la caja de dos cerraduras junto con el lino de ensayo documentado, Ferrero pidió la palabra con un gesto que no necesitaba pedirla: se puso de pie.

---Una condición más. Es la última y la pongo hoy porque después del doce no pienso negociar nada con la tela sobre la mesa. ---Sacó del bolsillo del saco los guantes de algodón. Los sostuvo, sin ponérselos, doblados en la mano izquierda---. Derecho de veto. Sobre cualquier anuncio, comunicación, adelanto o insinuación de resultados, de cualquiera de nosotros, por cualquier vía, a cualquier destinatario, mientras dure la operación. Veto sin causa, sin apelación y sin plazo de respuesta.

Vidal dejó la lapicera.

---Nadie planea anunciar nada. El resultado es impublicable por método desde dos mil veintinueve y por origen desde que existe esta nave. Su cláusula prohíbe algo que ya es imposible dos veces.

---En el ochenta y ocho el protocolo también preveía coordinación de anuncios. Yo lo leí, doctor. Estaba bien escrito. Tres laboratorios, comunicación conjunta, gradualidad. Después alguien juntó a la prensa en una sala y el intervalo salió por los diarios en una tarde. ---Los guantes giraron una vez en su mano---. Las cláusulas contra lo imposible son las únicas que me interesan, porque lo imposible es lo único que en este asunto tiene antecedentes. Si un dato de esta sala llega afuera antes del veintinueve, ¿quién lo paga? No lo pagamos nosotros. Nosotros pagamos con cárcel, que es barato. Lo pagan los que están del otro lado del anuncio, los que no firmaron nada. Escríbalo.

---El veto le da poder sobre mi informe final.

---El informe final es posterior a la operación. Léame la frase: mientras dure la operación. Después del veintinueve, usted hace con sus datos lo que su contrato le deje hacer, que según entiendo no es mucho.

Era exacto. La cláusula, tal como estaba formulada, expiraba con la devolución y en el intervalo prohibía un conjunto vacío. Vidal la sopesó igual, punto por punto, buscando el uso que Ferrero podía darle que él no estuviera viendo, y encontró uno solo: la palabra insinuación era ancha, y ancha a favor de Ferrero. Podía vetarle un correo, una consulta externa, una llamada a un colega. Podía, en teoría, vetarle una frase dicha en la cocina.

---Insinuación se reemplaza por «comunicación a terceros ajenos al equipo» ---dijo Vidal---. Lo demás queda como lo dictó.

---Aceptado.

Lo escribieron como punto diecinueve. Ferrero lo leyó dos veces, guardó los guantes en el bolsillo y firmó con la tinta lenta de siempre. Ledda firmó sin comentario. Chiara dejó la aguja clavada en la batista, vino hasta la mesa, leyó el punto entero con el dedo bajo cada línea y firmó cuarta.

---Nadie planeaba anunciar nada ---dijo, todavía inclinada sobre el papel---. Ahora está firmado que nadie planeaba. Qué oficio raro el nuestro.

Ferrero soltó una risa corta, de una sola sílaba, la primera del día.

Vidal cerró la carpeta. Dos asuntos del día quedaban fuera de ella: una desconfianza sin dato y un silencio de dos clics de caudalímetro. Ninguno tenía formato de registro. Los dejó donde estaban.

Ledda recogió las tazas y habló desde la pileta, sin darse vuelta, con el tono que usaba para los horarios de relevo.

---Punto fuera de protocolo. Hoy es el último día sin material sensible en ningún sitio. Veinte y treinta, piso franco, cena. Asistencia recomendada, no obligatoria.

Se quedaron mirándola. En cuatro semanas Ledda no había propuesto nada que no tuviera responsable y duración estimada. Fue Ferrero el que rompió el silencio, colgándose el saco del brazo.

---Cocino yo ---dijo---. Alguien tiene que controlar al menos una variable esta semana.

\clearpage

\chapter{Cena en Turín}

\lettrine[lines=2, lhang=0.1, nindent=0.2em]{E}{l} brasato llevaba tres horas y Ferrero lo vigilaba con un termómetro de cocina que había calibrado contra el de la nave, según dijo, porque un instrumento sin contraste era una opinión con pilas. Tenía un repasador al hombro y la corbata de lana metida entre dos botones de la camisa, para protegerla de las salpicaduras. La cocina del piso franco olía a vino reducido y a laurel, y en la mesa había cinco platos de tres vajillas distintas, porque en ese departamento todo funcionaba y nada combinaba.


Vidal llegó a las 20:31. Traía una botella que le habían recomendado en el negocio de la esquina de Collegno y que Chiara recibió con las dos manos, la giró para leer la etiqueta y la aprobó con un movimiento de cabeza, como si la botella hubiera pasado una inspección de trama.

---Nebbiolo ---dijo---. Te asesoraron bien.

---Pedí el que fuera con carne estofada. La variable la eligió el vendedor.

---Entonces brindamos por el vendedor.

Ledda ya estaba sentado a la mesa con un vaso de agua por la mitad. Sería el primero de tres. Vidal los contó sin proponérselo y después dejó de contar, porque la cuenta cerraba sola: Ledda, agua, la puerta a la vista.

Comieron. El brasato estaba en su punto y Ferrero lo sirvió anunciando la temperatura interna de la pieza como quien lee un acta. Chiara hizo la primera pregunta personal de la noche a los diez minutos, con la naturalidad de quien pasa el pan.

---Ledda, ¿vos de dónde sos?

---Dato sin uso operativo.

---Era conversación, no interrogatorio.

---Lo sé. Mismo criterio.

Ferrero apoyó los cubiertos.

---Probemos otra cosa. ¿Vacaciones?

---Mayo. Después del cuatro.

---¿Familia?

---Tema archivado en 2019.

---¿Y el archivo se puede consultar?

---Hay plazos de reserva.

Chiara se rio con la copa en la mano y a partir de ahí el juego quedó instalado: cada pregunta personal a Ledda se contestaba con un plazo, y la mesa empezó a hacerle preguntas solo para escuchar el plazo. Ledda no se defendió del juego ni lo alentó: lo administró, con la misma cara con que administraba las ventanas de acceso. En algún momento Ferrero giró hacia Vidal.

---¿Y usted, doctor? ¿Qué hace cuando no trabaja?

---Está previsto para mayo ---dijo Vidal.

La mesa entera se rio, Ledda incluido, medio segundo de exhalación por la nariz que en esa cara equivalía a una carcajada. Vidal revisó la frase que acababa de decir. Era un chiste. Lo había construido con material ajeno, sintaxis de Ledda, y había funcionado. La mecánica exacta del efecto se le escapaba, pero el resultado era replicable: Chiara todavía se estaba riendo.

Con el segundo vaso de vino, Ferrero contó lo del incendio.

---Abril del noventa y siete ---dijo---. Yo tenía veintisiete años y el sueño pesado, así que me despertó mi madre, no el teléfono. Desde el balcón se veía el resplandor sobre el Duomo y mi madre dijo: se quema la capilla del Guarini, andá. Como si yo fuera bombero. Llegué a las tres de la mañana con media Turín en la plaza, todos mirando para arriba, nadie rezando todavía porque rezar es lo que se hace después, cuando ya se sabe qué se perdió. Y adentro había un bombero con una maza.

Hizo una pausa para servirse salsa. La administraba como el suspenso.

---El relicario estaba detrás de un cristal blindado. Certificado. Antibalas, antimartillo, anti todo lo previsible. La casa que lo fabricó tenía los ensayos documentados: tantos julios, tantos impactos, cero penetración. Y este hombre le pegó con la maza hasta que el cristal dijo basta. Después le preguntaron cómo lo había logrado, porque técnicamente era imposible, y el hombre dio una respuesta que no figura en ninguna norma de materiales. Yo prefiero la versión técnica: el certificado medía impactos de laboratorio y el bombero traía otra unidad. ---Levantó el tenedor---. La capilla estuvo cerrada veintiún años. El cristal, en cambio, cedió en minutos. De las dos estructuras, la que estaba diseñada para resistir fue la que falló primero. Eso, en mi pueblo, es una comedia. Con un muerto adentro, como todas las buenas.

---¿Quién es el muerto? ---preguntó Chiara.

---La certificación ---dijo Ferrero, y se limpió la boca con la servilleta---. Que en paz descanse.

Vidal masticó despacio. El razonamiento tenía un hueco, el cristal había fallado por fatiga acumulada de impactos sucesivos, no por una unidad de medida desconocida, y estuvo a punto de decirlo. Lo dejó pasar. Corregir un chiste era un procedimiento sin producto.

Fue Ledda quien tendió el puente, sin querer, preguntando si el protocolo de emergencia de la teca nueva contemplaba un escenario de incendio. Ferrero contestó que sí, que estaba mejor escrito que el del noventa y siete, y que los protocolos siempre mejoraban después, nunca antes, porque el papel aprende de lo que la gente paga. Y de ahí, sin cambiar el tono, pasó al otro año.

---Con el anuncio del ochenta y ocho también hubo protocolo ---dijo---. De prensa. Se convocó una conferencia, se leyó un intervalo, se agradeció a los laboratorios. Todo correcto. Lo que no se previó fue la recepción. Hubo gente a la que ese anuncio le llegó sin nadie al lado. Se dio la noticia y no se preguntó quién la recibía, ni en qué habitación, ni con qué recursos. Eso tampoco figura en ninguna norma de materiales.

Hablaba a la fuente de brasato, no a la mesa. La mano derecha le subió a la muñeca izquierda y se quedó ahí, dos dedos sobre el punto donde se abrocha una correa, y no había correa, ni reloj, ni marca; solo el gesto, ajustado a nada.

---Ottavio ---dijo Chiara---, ¿la nuez moscada la pusiste al principio o al final?

Ferrero parpadeó.

---Al final. Al principio se evapora lo volátil y queda lo amargo.

---Mi maestra decía lo mismo del elogio.

---Tu maestra tenía razón en más cosas de las que le reconocieron.

---En casi todas ---dijo Chiara, y le acercó la fuente para que sirviera lo último.

El empalme duró menos de dos segundos. Vidal lo repasó mientras Ferrero explicaba la diferencia entre rallar y moler, ya de vuelta en su registro de químico con delantal. El corte había sido limpio: una pregunta trivial, insertada en el punto exacto donde la frase anterior todavía no terminaba de caer, y la conversación había cambiado de vía sin sacudón, sin que el propio Ferrero registrara el desvío. O registrándolo y agradeciéndolo, eso era indecidible desde afuera. Lo que Vidal tenía era la maniobra, ejecutada con precisión de bisturí, y ninguna hipótesis sobre el criterio que la disparaba. Chiara sabía dónde estaba el punto blando de esa tela. El cómo lo sabía quedaba fuera de su alcance.

Ledda se fue primero, a las once y cuarto, después de lavar su propio vaso y anunciar que a las seis, antes del ingreso del tubo, había circuito de billetes que cerrar. Ferrero se fue veinte minutos más tarde con el saco al brazo y una porción de brasato envuelta en papel de aluminio que Chiara le impuso en la mano. Desde la puerta, ya en el pasillo, se dio vuelta.

---Doctor. La cena no va en ninguna hoja del día, ¿verdad?

---No hay columna para esto.

---Mejor ---dijo Ferrero---. Que algo quede sin acta. Aunque sea la nuez moscada.

Quedaron los platos. Chiara lavaba y Vidal secaba, porque el escurridor era chico y alguien tenía que resolver el cuello de botella. Por la ventana entreabierta llegaba el ruido de una formación de carga por las vías de Dora, largo, con dos frecuencias superpuestas, y el agua salía con un golpe de aire en la canilla cada tanto, un defecto de presión que a esa hora sonaba casi doméstico.

---¿Puedo preguntarte una cosa? ---dijo ella.

---Depende de qué llames una cosa.

---Si el resultado dice siglo primero. ¿Vos qué hacés?

---Verifico contra los controles ciegos, reviso contaminación de borde, y el intervalo se contrasta con el corpus, que para lino de esa cuenca es pobre, así que el primer sospechoso sería el propio corpus. Después se repite la corrida con semilla distinta y se documenta la varianza. Si todo eso cierra, el resultado se sostiene.

---Te pregunté qué hacés vos.

---Eso hago yo. El modelo corre; las decisiones de las que te hablé las tomo yo.

Chiara le pasó la olla del brasato, pesada, con el fondo todavía tibio.

---Ya ---dijo.

Y nada más. Enjuagó los cubiertos, los fue dejando de a uno en el escurridor, y el silencio que siguió no era incómodo ni era cómodo: era un dato de otra clase, y Vidal lo dejó donde estaba porque su sistema entero carecía de un campo para eso. Secó la olla dos veces. La segunda pasada era innecesaria y la hizo igual.

---Le cambiaste el tema a Ferrero ---dijo él después---. Con la nuez moscada.

---Sí.

---Fue precisa. Quiero decir: el punto de inserción. ¿Por qué ahí?

Chiara cerró la canilla. Se secó las manos con el repasador de Ferrero, que había quedado colgado del respaldo de una silla, y lo dobló en cuatro antes de contestar, con el cuidado que le ponía a los cables.

---Cuando una tela tiene un tirón, no lo estirás delante de todos. Lo llevás a la mesa de trabajo, con luz buena, cuando el dueño quiere.

---¿Y si el dueño no quiere nunca?

---Entonces la tela se guarda con el tirón. Se puede vivir así. Vos deberías saberlo.

Lo dijo sin filo, ya juntando su abrigo, y le deseó buenas noches en italiano, dos palabras gastadas de uso, y se fue. Vidal escuchó la puerta del edificio tres pisos más abajo y después el tren de nuevo, o era otro.

En su cuarto, el bolso estaba contra la pared, con el registro del día adentro, sin abrir. La rutina nocturna tenía cuatro pasos: volcado de notas, cotejo del cronograma de mañana, revisión de calibraciones pendientes, cierre. Llevaba años sin saltearla; la última interrupción documentada había sido una fiebre en 2027, y aun esa noche había hecho el cierre desde la cama. Se sentó en el borde del colchón con la intención de empezar por el volcado y se encontró repasando, en cambio, la mecánica del chiste de mayo, el arco de la maza contra el cristal certificado, la mano de Ferrero ajustando una correa inexistente, la olla tibia.

Ninguna de esas entradas tenía formato. Todas eran de esta noche y ninguna era auditable, y mañana no quedaría de ellas más soporte que cuatro memorias con sesgos distintos, la suya incluida.

Apagó la luz con el registro todavía en el bolso. La omisión quedó registrada por el único sistema que seguía despierto, que la anotó como omisión, sin causa asignada, y a la mañana siguiente ni siquiera eso.

\clearpage

\chapter{Listas cruzadas}

\lettrine[lines=2, lhang=0.1, nindent=0.2em]{T}{rece} días antes de la cena, con la gemela todavía con el respaldo descosido sobre la mesa, el veintiocho de marzo a las 8:05, Ledda entró a la nave de Collegno con una hoja doblada en tres y la abrió sobre la mesa de embalaje antes de decir nada. Era una lista impresa, cuarenta y tantos renglones, con dos marcados a lápiz.


---El cruce de contratistas se adelantó ---dijo---. Estaba agendado para la tercera semana de abril. La prefectura lo movió al primero de mes por la visita del cuatro de mayo. La lista provisoria circuló el lunes.

Apoyó el índice en el primer renglón marcado.

---Bertoglio, distribución técnica, Moncalieri. Figura en dos registros. Subcontratista de la empresa de la teca en la obra de la catedral, julio a noviembre del año pasado. Y tres entregas en este polígono. Febrero, febrero, marzo.

Vidal leyó el renglón dos veces. Las tres entregas eran las sondas de humedad de la celda de ensayo, el juego de repuestos y el material de la línea de nitrógeno. Los remitos estaban en la carpeta de la esclusa, con sus fechas.

---¿Qué ve un inspector cuando cruza eso?

---Primero documentación. Remitos, domicilios de entrega, personal de reparto. Si un domicilio se repite donde no corresponde, pasa a inspección física. Este polígono no resiste una inspección física.

---¿Plazo?

---Lo fija la prefectura. Fuera de nuestro control.

Ledda había enumerado tres modos de falla en la reunión de ventanas: pago rastreable, cara repetida donde no corresponde, casualidad. Esto era el segundo, llegando antes de fecha y en papel, sin cara todavía. Vidal contó los renglones de la lista mientras Ledda hablaba. Cuarenta y seis. Dos marcados. Volvió a contarlos y dieron lo mismo.

---Bertoglio queda cortado desde hoy ---siguió Ledda---. La deuda se salda en efectivo por vía indirecta, con recibo genérico. Lo que estaba pendiente de entrega se compra de nuevo por otra distribuidora, sin vínculo con obra de catedral, facturación en efectivo. Plazo de reposición: doce a quince días.

---El cronograma tenía catorce días de margen.

---Tenía.

Lo pendiente eran las dos células de carga de repuesto y el segundo juego de sondas. Sin repuestos, la celda de ensayo funcionaba con un solo juego montado, y cada prueba completa contra los sensores exigía dos días de ciclo más uno de reconfiguración. Vidal hizo el cálculo tres veces, con lápiz, en el margen de la lista de Ledda, y las tres veces el resultado fue el mismo: donde el protocolo decía tres pruebas completas de la unidad de sustitución, la aritmética dejaba lugar para una.

---¿Y si pido las células por un canal de Zúrich? Instrumental de laboratorio, uso declarado en Suiza, tránsito normal.

---Aduana registra el tránsito y el destino final queda a cargo nuestro. Es un rastro nuevo para ahorrar doce días. Recomendación: no.

Ledda plegó la lista en tres, por los mismos dobleces.

---El circuito de billetes absorbe la re-facturación ---dijo---. Eso ya está resuelto. Lo que se quemó es margen. El margen no se re-factura.

Se fue a las 8:40. Vidal se quedó junto a la mesa de embalaje escuchando el caudalímetro de la línea de nitrógeno, que hacía clic a intervalos de entre once y diecinueve segundos, según su propio cronometraje del martes anterior. Midió cuatro intervalos con el segundero. Trece, diecisiete, doce, quince. Dentro de rango. La línea funcionaba igual que ayer y él ya lo sabía antes de medir.

Esa tarde llevó el cálculo a la mesa de costura. Chiara estaba en el segundo día del respaldo de batista: once metros ochenta descosidos, medidos, a medio coser de nuevo, con el hilo tensado en pasadas que llevaban su ritmo propio y ninguna prisa visible. Vidal le puso delante la hoja con la aritmética.

---Tres pruebas se vuelven una ---dijo---. La versión dos va contra los sensores una sola vez, con todo montado, antes del once. Si querés, lo registro con la discrepancia a tu nombre, como la otra.

---Registralo a nombre del cronograma. Yo coso lo mismo con una prueba que con tres.

Terminó la pasada, clavó la aguja en el acerico y juntó el hilo sobrante, que empezó a enrollar sobre dos dedos, despacio, vuelta por vuelta.

---¿Y si la única prueba da mal?

---Depende de qué llames mal. Si es la pendiente de adsorción otra vez, con el respaldo nuevo tendríamos un día para... ---Se detuvo. Rehízo el cálculo mentalmente y le dio lo que ya le había dado en papel---. No tendríamos un día.

---Eso te pregunto.

---Entonces el once la caja de apoyo entra con el falso fondo vacío. La sustitución se cancela y esto se desarma.

Chiara terminó de enrollar el hilo, lo dejó junto al acerico, un anillo prolijo, y retomó la aguja.

---Allora que dé bien ---dijo, y fue todo.

Vidal registró esa noche, en la hoja del día: pruebas completas de unidad de sustitución contra sensores: una; planificadas: tres; motivo: exposición logística sobrevenida. Releyó el renglón. Era exacto y no decía nada de lo que importaba, y esa combinación, en su sistema, contaba como un renglón bien escrito.

Ferrero llegó dos días después, a las siete de la tarde, con una carpeta bajo el brazo. Adentro venía una tabla de emisividad corregida para la calibración hiperespectral de las zonas de chamuscado: dieciocho valores, cada uno con su fuente, dos de ellos mejores que los que Vidal tenía cargados. Trabajo fino, hecho sin que nadie se lo pidiera. Vidal lo incorporó al banco delante de él y Ferrero verificó la carga por encima de su hombro, señalando con el dedo sin tocar la pantalla.

Después se sentó del otro lado de la mesa y se tocó la muñeca izquierda, donde no llevaba nada.

---Doctor, la situación de la distribuidora cambia el balance de su protocolo. Le propongo una reducción de alcance. El equipo de fluorescencia ultravioleta de campo completo sale del plan. Es un alquiler, viaja con generador propio, y el canal por el que venía ahora tiene una bandera encima. Menos instrumentos, menos huella. Si un inspector encuentra un generador de esa potencia en un taller electromecánico que factura mantenimiento, ¿quién se lo explica? ¿Y con qué remito?

---De acuerdo.

Ferrero siguió hablando.

---Considere además que el resultado ultravioleta ya existe. Miller y Pellicori lo fotografiaron en el setenta y ocho: las quemaduras fluorescen, la imagen no. Usted replicaría un dato publicado asumiendo un riesgo nuevo, y si esto sale mal por un generador alquilado, el costo no lo paga usted solo, lo paga... ---Se detuvo a mitad de la subordinada---. ¿Dijo que sí?

---Dije que sí. El argumento logístico alcanza. El otro es discutible: yo no replico datos ajenos, los verifico. Pero el generador sale del plan.

Ferrero se acomodó en la silla como quien llega a una puerta abierta con el hombro por delante. Tardó un momento en soltar el resto.

---Ya que estamos en reducción ---dijo---, las ochenta y dos franjas del barrido podrían bajarse a sesenta. El solapamiento del quince por ciento es generoso; con diez por ciento cubre igual, y son dos días menos de exposición del lienzo a la mesa.

---No. Las ochenta y dos quedan. El solapamiento es lo que me deja auditar la costura entre franjas; sin eso, cualquier discordancia en un borde de franja es indecidible.

---Que conste mi propuesta.

---Consta.

Cuando Ferrero se fue, Vidal se quedó con la tabla de emisividad en la pantalla y las dos propuestas en la hoja del día, una aceptada y una vetada. Las miró un rato. La tabla mejoraba la lectura de las zonas de chamuscado. Las dos propuestas recortaban canales de la zona de imagen: la fluorescencia distinguía chamuscado de imagen, el solapamiento cosía el mapa de la imagen entera. Puesto en columnas, el patrón era visible: la mano que perfeccionaba trabajaba en un sector de la tela y la mano que recortaba trabajaba en el otro.

Lo pesó como pesaba cualquier correlación con n igual a dos: insuficiente. Un hombre que llevaba cuarenta años protegiendo ese objeto tenía motivos de sobra para querer menos instrumentos alrededor, y el generador era, de verdad, el bulto más ruidoso del inventario. Prudencia. La palabra cerraba el caso y la dejó cerrada, con la incomodidad menor de las palabras que cierran demasiado rápido, esa que él sabía descontar.


\scenebreak


La noche del diez de abril se despertó a las cuatro menos veinte sin alarma, en el piso franco, con el registro del día todavía dentro del bolso. La rutina salteada seguía salteada; nadie la había hecho por él. Encendió la lámpara, sacó el registro y la hizo entera, fuera de horario y fuera de orden: primero el cotejo del cronograma ---a las siete empezaba el traslado, y con él, según la pestaña dos del expediente, el ingreso del equipamiento con la caja de apoyo---, después el volcado de notas de la cena, reducido a asistentes y horarios porque el resto seguía sin formato, y al final las calibraciones.

Las calibraciones estaban al día. Lo sabía. Las copias en papel del acelerómetro, una por mañana desde el lunes, estaban en la funda del bolso, sumadas y firmadas por él mismo. Sacó la del martes y volvió a sumar la columna de picos, a mano, en el reverso de una hoja usada. Dio igual que la suma original. Sacó la del miércoles. Dio igual. El certificado del banco hiperespectral decía doce micrones por píxel con trazabilidad al patrón nacional y fecha catorce de marzo, y él lo había verificado dentro de la nave con su propia tarjeta de referencia; releyó el número de serie del certificado contra su anotación y coincidían, como habían coincidido cuando los cotejó la primera vez.

Afuera pasó el primer tren de la mañana, o el último de la noche. La ventana clareaba por el borde.

Quedaba una hoja: el cronometraje del caudalímetro, del martes anterior al ensayo. Intervalos de once a diecinueve segundos, veinte mediciones, media de catorce coma nueve. Nada en el traslado dependía del caudalímetro. El caudalímetro quedaba en Collegno, a doce kilómetros de la sacristía, y su régimen de clics era irrelevante para todo lo que iba a pasar hoy entre las siete de la mañana y la una de la tarde de mañana. Sumó la columna igual. Media: catorce coma nueve.

Guardó las hojas en su orden, cerró el bolso, y se vistió con la camisa que tocaba. A las seis en punto sonó la alarma que había programado hacía una semana y que ya no hacía falta. La dejó sonar hasta el final, de pie junto a la cama, con el bolso colgado del hombro.

Nueve pitidos. Los contó.

\clearpage

\chapter{Ensayo general}

\lettrine[lines=2, lhang=0.1, nindent=0.2em]{L}{a} alarma seguía sonando cuando abrió el registro por última vez antes de salir. No buscó la hoja de hoy: buscó atrás, la semana entera, porque en unas horas el once iba a borrarla como fecha de trabajo y la quería completa antes de que eso pasara. Seis de abril. Ensayo general. Tres corridas.


La sacristía ocupaba, sobre el piso de la nave de Collegno, cuarenta y un metros cuadrados de cinta azul. Ledda la había trazado durante la noche con la pestaña dos del expediente abierta sobre un caballete: cada cota transcrita, cada vano marcado con dos tiras cruzadas, la puerta simulada por un marco de listones que se abría hacia adentro porque la de verdad abría hacia adentro. La mesa de conservación era la real, la misma que entraría el once a las siete. La caja de apoyo, con su falso fondo, era la real. Lo único falso del rectángulo era la tela.

Ferrero recorrió el perímetro antes de empezar, comparando la cinta con una memoria de treinta años.

---El armario de ornamentos está veinte centímetros más cerca de la mesa ---dijo---. En 2002 nos costó un codo. El codo de la doctora Marini, para ser exacto.

Ledda corrió la tira sin comentario. Anotó la corrección en su hoja, con hora.

A las 8:10 Chiara se ató el pelo con un hilo de algodón y ocupó su posición junto a la caja. El lino de ensayo, cuatro metros y medio de sarga envejecida, esperaba desplegado sobre la mesa haciendo de original. Ferrero se paró donde iba a pararse la Comisión, tres pasos atrás del borde de cinta, las manos cruzadas a la espalda, y desde ahí miró el lino de ensayo de un modo que Vidal registró sin encontrarle categoría: como se mira una radiografía propia. Cuando Chiara ajustó la tensión de una esquina, Ferrero adelantó medio paso y se detuvo en el borde azul, exactamente sobre la cinta, sin pisarla.

---Ventana ---dijo Ledda, y apretó el cronómetro. Mecánico, de cuerda. Nada digital entraba en la sala.

Chiara abrió la caja de apoyo, liberó el falso fondo, y la gemela apareció plana, ya aclimatada, con las quemaduras de 1532 en su sitio y las cuatrocientas ochenta y dos perforaciones del parche grande esperando a un ojo que supiera contarlas. El intercambio tenía nueve fases y Vidal las conocía de memoria porque las había aprobado por escrito una por una. Rodillo, plegado de transferencia, tubo, gemela al soporte, verificación de posición, cierre. Chiara trabajaba sin apuro y sin pausa, con ese ritmo que en el primer ensayo de marzo le había parecido lento y que después, cronometrado, resultó más rápido que cualquier apuro. El único sonido era el clic del caudalímetro, a intervalos de once a diecinueve segundos, media catorce coma nueve, y el roce del lino contra el fieltro del rodillo.

---Treinta y tres quince ---dijo Ledda cuando la caja quedó cerrada---. Margen: seis cuarenta y cinco.

Dentro de la ventana de cuarenta. Vidal lo anotó. La segunda corrida empezó a las 9:05.

En el minuto dieciocho de la segunda corrida, Ledda abrió el marco de listones sin aviso.

---Puerta. Dos personas de la Comisión. Sin anuncio previo.

Nadie le había comunicado esa variable a Vidal. Chiara tampoco levantó la vista: soltó el rodillo en el soporte, desplegó la tela de servicio sobre la mesa con un solo movimiento y se volvió hacia la puerta imaginaria con las manos vacías y una pregunta de taller en la cara, algo sobre humedad, algo aburrido. Ferrero, desde su marca, dijo en voz de acta:

---Verificación de soporte en curso. Ruego no ingresar hasta el cierre de la fase.

Vidal contó. Cuarenta y un segundos hasta que Ledda declaró despejada la puerta, ciento diez más de recomposición antes de que Chiara retomara la fase interrumpida en el punto exacto. La corrida cerró en treinta y siete cincuenta.

---Dentro de ventana ---dijo Ledda---. Con dos minutos diez de margen. Es el escenario malo. El peor no se ensaya: se cancela.

---¿Y quién decide la cancelación? ---dijo Ferrero---. Porque si la decide el cronómetro, muy bien. Si la decide una persona con las manos en la tela, esa persona carga con las dos mitades del error: cancelar de más y cancelar de menos.

---Umbral en el minuto treinta y cuatro ---dijo Ledda---. Pasado el umbral con fase abierta, se cierra y se aborta. Lo decide el número.

Ferrero se tocó la muñeca izquierda, donde veinte años hubo un reloj, y volvió a su marca.

La tercera corrida, limpia, dio treinta cuarenta. Vidal pasó los tres tiempos al registro y les sumó la desviación. Todo cabía en las columnas. Esa era, revisada dos veces, la propiedad general de la mañana: todo cabía en las columnas.

A las once hicieron la lectura del registro en papel: las cuarenta y ocho horas que la gemela versión dos había pasado dentro de la celda de ensayo, cerradas anoche antes de sacarla para el simulacro de esta mañana, contra los sensores replicados de la teca. Era la única prueba completa que el cronograma recortado había dejado en pie. Ledda leyó las bandas en voz alta desde la copia en papel. Masa seca, dentro de uno coma uno gramos. Masa en equilibrio a cuarenta y cinco por ciento de humedad, dentro de cero coma ocho. A cincuenta y cinco, dentro de uno coma tres. Células de carga, higrometría, oxígeno residual: sin alarma en cuarenta y ocho horas. Vidal cotejó cada cifra contra el punto cinco del protocolo, firmado por los cuatro, y escribió «conforme» junto a cada línea.

---A un metro cuarenta, con el cristal, ni la madre que la tejió ---dijo Chiara. Después torció la boca---. La madre que la tejió soy yo. Mala frase.

---La frase es auditable ---dijo Vidal---. Ocho observadores, cuatro aciertos. Azar.

Chiara fue a guardarla. Y ahí ocurrió la observación que Vidal estuvo un rato largo sin poder asentar: la levantó con las dos manos abiertas, por debajo, el peso repartido en los antebrazos, la cabeza inclinada sobre la tela, exactamente igual que levantaba el lino de ensayo, exactamente igual ---supuso, y el supuesto era el problema--- que iba a levantar el otro. La gemela tenía dos años de existencia y ningún dueño ausente. Era de ellos; peor: era de nadie. La regla de su maestra, esa que Chiara citaba ---todo lo que tocás es de alguien que no está---, tenía acá un valor de entrada nulo, y las manos de Chiara la aplicaban igual, sin excepción, como si la regla midiera a quien toca y la tela diera lo mismo.

Vidal abrió el registro, apoyó la lapicera en la columna de observaciones y escribió «manipulación conforme». Quedaba un resto sin palabra. Dejó el resto donde estaba.

Ferrero se acercó a la mesa cuando ya estaban cerrando. Sacó del bolsillo derecho del saco los guantes de algodón de conservación y se los puso con dos tirones parejos, aunque su parte del día había terminado y no quedaba nada que tocar. Con los guantes puestos apoyó dos dedos sobre el borde del lino de ensayo, que valía lo que una lona, y le corrigió un pliegue que nadie más había visto.

---Repitamos lo que sabemos ---dijo. Tenía la voz que usaba para leer actas, con otra cadencia debajo, más antigua---. La gemela engaña al ojo. Engaña a la balanza. Engaña a los sensores de la teca. Bajo un microscopio, a cincuenta aumentos, es una tela belga con fertilizante del siglo veintiuno en las fibras, y cualquier estudiante lo ve en una tarde. La imagen de la zona central es una impresión y su espectro lo dice a gritos. Y si alguien corta un hilo y lo data, el intervalo es dos mil treinta con confianza cero noventa y nueve, porque ese número lo produjo usted mismo, doctor, con su propia máquina.

---Dos mil treinta a dos mil treinta y tres. Cero noventa y nueve ---dijo Vidal---. Correcto lo demás.

---Entonces la aritmética es esta: diecisiete días. Ni uno más. Cada día dieciocho que ese objeto pase bajo el cristal es un día en que la operación entera depende de que a nadie, en toda la diócesis, se le ocurra mirar de cerca lo que dos mil millones de personas van a venir a mirar de lejos. ---Se alisó los guantes, dedo por dedo---. Si esto pasa del veintinueve, ¿quién lo paga? Se lo digo yo: lo paga gente que hoy ni siquiera sabe que existe una pregunta.

---El veintinueve está en el cronograma con hora y responsable ---dijo Ledda---. Ventana de seis minutos. Ensayada.

---Los cronogramas del ochenta y ocho también tenían horas ---dijo Ferrero---. Lo que no tenían era a nadie del otro lado.

Se quitó los guantes recién en la esclusa, de espaldas, y los dobló dentro del bolsillo antes de salir.

A las siete de la tarde sonó el teléfono fijo de la oficina del taller, en la hora fijada, sobre el escritorio con olor a grasa fría y facturas de mantenimiento electromecánico verdaderas. Vidal atendió al segundo timbre.

---Buenas tardes, doctor. Tengo tres cosas para comunicarle y las tres me constan por documento. Primera: la nómina del operativo del once y el doce se mantiene en nueve personas, las mismas nueve que figuran en la pestaña dos, sin altas ni bajas al día de hoy. Segunda: el intervalo de verificación de soporte y limpieza figura en el cronograma oficial con la duración estimada que usted conoce, y ninguna comunicación posterior lo modifica. Tercera: los cruces de la prefectura concluyeron sobre la lista provisoria y ningún renglón vinculado a este asunto fue observado.

---¿Eso constituye una garantía?

---Constituye una descripción. Las garantías las emite quien controla los hechos; yo describo documentos, y los documentos describen la ventana. La ventana existe en los términos que usted evaluó.

Vidal sostuvo el auricular con el hombro, abrió el registro sobre las facturas y escribió, en el punto de ejecución del protocolo: «Confirmación recibida 19:04. Nómina y ventana vigentes. Ejecución: adelante.» Firmó. La letra le salió pareja, mayúsculas alineadas, la misma de las cuatrocientas cadenas de custodia.

---Registrado ---dijo---. Adelante.

---Queda una cuarta cosa ---dijo Ansermet---, que me fue instruida esta mañana y que le transmito completa porque completa me fue dada. El fundador lo recibe pasado mañana a las once, en Cologny. Le dicto la dirección cuando usted me indique. El objeto de la reunión no me fue comunicado, y si me hubiera sido comunicado, carecería de instrucciones para transmitírselo. El automóvil corre por cuenta de la Fondation desde la estación de Ginebra.

Vidal miró el marco de la puerta de la oficina, donde la pintura se descascaraba en una franja del ancho de un dedo, y contó las capas visibles. Tres.

---Faltan cinco días para el traslado.

---Tres días antes del traslado, para ser exactos, a las once de pasado mañana. El fundador conoce el cronograma en el mismo nivel de detalle que usted. La superposición, doctor, me permito suponer que la ha considerado.

---Dicte.

Ansermet dictó: una ruta, un número, una comuna. Vidal lo escribió en la última línea del día, dentro de las columnas, porque afuera de las columnas ya había dejado demasiado esta semana. Asunto: reunión con el mandante, Cologny. Duración estimada: sin dato. Después la lapicera quedó un momento quieta sobre el campo que faltaba, el que su sistema exigía para cerrar cualquier entrada, y lo llenó con las únicas palabras que tenía preparadas desde enero.

Entregable: identidad del mandante.

\clearpage

\chapter{Cologny}

\lettrine[lines=2, lhang=0.1, nindent=0.2em]{O}{cho} de abril, tres días antes del traslado. La grava del patio estaba rastrillada en líneas paralelas y las líneas se interrumpían donde alguien había caminado esa mañana, un solo par de pisadas, del garaje a la casa. Vidal las contó mientras el automóvil de la Fondation daba la vuelta y se iba. Once. El hombre que las había dejado pisaba largo y sin apuro.


Nadie salió a recibirlo. La puerta principal estaba entornada y detrás había un pasillo de piedra clara que olía a cera y, más al fondo, a tierra mojada. Siguió el olor. El pasillo desembocaba en un ala de vidrio adosada al lado sur de la casa: una estructura de hierro pintado de verde oscuro, paneles curvos arriba, condensación en los bordes bajos. Un jardín de invierno. Adentro, la temperatura subía seis o siete grados y la humedad más, y entre dos mesas de trabajo cargadas de macetas un hombre grande, de espaldas, apretaba tierra alrededor de un tallo con las dos manos desnudas.

---Doctor Vidal ---dijo el hombre, sin darse vuelta todavía---. Le agradezco la puntualidad. En esta casa es una cortesía que casi nadie me hace.

Se dio vuelta despacio. Grande de una manera que el traje disimulaba mal, o que el traje ni intentaba disimular: una chaqueta de lana gastada en los codos, un pantalón con rodillas de jardinero. Vidal había tasado ropa para tres aseguradoras y reconoció el género. Esa chaqueta que parecía vieja costaba lo que un coche. Las manos que se limpiaba ahora con un trapo eran anchas, con tierra en las uñas, tierra de verdad, incrustada, de años.

---Emeric Sandoz ---dijo, y le tendió la mano.

Vidal se la dio. El apellido cayó en su sistema y encontró casillero solo: trazabilidad de materias primas, certificación de origen, infraestructura de datos. Los objetos sociales idénticos de las dos sociedades, la de Ginebra y la de Zug, redactados con las mismas palabras. Era la primera verificación en tres meses que cerraba al primer intento, y le llamó la atención el orden: el dato había llegado cuando ya era inútil, cuando el protocolo estaba firmado y la gemela cosida. Le habían dado la respuesta después del examen.

---Siéntese ---dijo Sandoz---. Acá, donde el vidrio no gotea.

Había una mesa baja entre dos sillones de mimbre, y en la mesa un servicio de té: dos tazas, una tetera, un colador de plata. Sandoz sirvió una taza, la giró para que el asa quedara del lado de Vidal, y la empujó tres centímetros hacia él. Su propia taza quedó boca abajo sobre el platito. Ni la tocó.

---¿Usted no toma?

---Yo sirvo ---dijo Sandoz, y se sentó. El mimbre crujió con su peso---. Usted vino con una lista, doctor. La conozco sin haberla leído porque yo habría traído la misma. Empecemos por el punto que más le importa, si me lo permite, que no es mi nombre. Mi nombre ya lo procesó. Es la cláusula octava.

Vidal dejó la taza donde estaba.

---La cláusula octava está negociada.

---Está corregida, que es distinto, y las correcciones son buenas. Pero su objeción de fondo sigue en pie, y como usted es demasiado educado para formularla en mi jardín, permítame formularla yo. ---Sandoz apoyó las manos en las rodillas---. Usted sostendría lo siguiente: que el archivo es la respuesta; que quien custodia el archivo es dueño de la respuesta; que una respuesta con dueño deja de ser una respuesta y pasa a ser un instrumento; y que por lo tanto usted, Sebastián Vidal, va a trabajar diecisiete días con el mejor instrumental de su vida para producir exactamente la única cosa que su método le prohíbe producir: un resultado que existe y que nadie puede examinar. Una expropiación, diría usted, no de la tela. De la respuesta. ¿Es fiel el resumen?

Era mejor que el original. Vidal lo había pensado en fragmentos, a las dos de la mañana, en los márgenes de otras cosas; nunca lo había armado entero. Este hombre se lo devolvía armado, pulido, con las juntas a la vista.

---Es fiel ---dijo---. Le falta la conclusión.

---La conclusión es suya y no me corresponde. Lo que me corresponde es esto. ---Sandoz se inclinó apenas---. Usted escribió una vez, en un informe sobre el trabajo de otro, que una certeza inauditable es un artefacto y se descarta. Es una buena regla. La aplicó toda su vida a los objetos de otros. Lo que le propongo es la ocasión de aplicársela a algo que le importe, y ver si la regla sobrevive.

El vidrio del techo goteaba en algún punto a la izquierda, a intervalos regulares. Vidal empezó a contarlos y se obligó a parar en cuatro.

---Mi informe de 2031 tuvo circulación restringida.

---Tuvo la circulación que tienen las frases verdaderas ---dijo Sandoz---. Se citan sin permiso.

Sobre la mesa de trabajo más cercana, entre una bandeja de plantines y un pulverizador de cobre, había tres libros apilados con un lápiz cruzado encima. El lápiz estaba gastado hasta la mitad, afilado a cuchillo, con las facetas irregulares de la madera cortada a mano. El libro de arriba estaba abierto y el margen bajaba lleno de una letra pequeña, en francés, inclinada a la derecha, con la mina apretada en las palabras que importaban.

Vidal conocía esa letra. La había leído durante tres noches, margen por margen, en quinientas y pico de páginas de las actas del STURP, edición de 1981, envueltas en papel madera. Tres durezas de mina, tres épocas de lectura, una taxonomía de hipótesis renumerada quince años antes de que la literatura la adoptara. Y una doble raya vertical, sin texto, en la página cuatrocientos ochenta y siete.

---El ejemplar que me envió ---dijo---. Las anotaciones son suyas.

---Todas.

---En la sección de difusión de resultados hay una marca doble, junto al párrafo que recomienda coordinar los anuncios entre equipos. Es la única anotación del libro sin palabras.

Sandoz tardó un momento. Cuando habló, el tono era el mismo, y eso fue lo que Vidal registró: el mismo tono exacto.

---Hay párrafos que uno subraya y párrafos con los que uno discute ---dijo---. Esa raya es una discusión que perdí. La marqué mucho antes de saber contra quién la había perdido. ---Señaló la taza con la cabeza---. Se le enfría.

Vidal tomó un sorbo. Era el té que él tomaba en Zúrich, la misma casa, el mismo grado de tueste. Anotó el dato y lo dejó sin columna.

---Las actas ---dijo---. Usted las corrige. En el modelo de chamuscado hay un gradiente de temperatura enmendado a lápiz que la publicación original tiene mal.

---Lo tiene mal por transcripción, no por medición. Pellicori midió bien. ---Sandoz cruzó las manos---. Usted habrá notado, doctor, que ese libro es el documento más honesto que produjo este asunto en casi siete siglos. Ciento veinte horas de acceso, decenas de instrumentos, y la conclusión oficial cabe en una línea: la imagen no es pintura, mecanismo desconocido. Un equipo entero que viaja hasta Turín para volver con las manos vacías y lo dice. Eso está más cerca de su método que de mi Iglesia, si me permite el posesivo, que uso por costumbre y sin derecho. Su solicitud de acceso del año pasado era la continuación natural de ese libro. La Custodia la rechazó por las razones que usted conoce y por una que usted subestima: en 1978 nadie era dueño de lo que se iba a encontrar, y diez años después la institución aprendió el precio de eso en una conferencia de prensa. Yo decidí no repetir esa lección, doctor. La diferencia entre la Custodia y yo es una diferencia de método, no de propósito.

---Eso es un argumento a favor de mi objeción, no en contra.

---Lo sé ---dijo Sandoz---. Por eso se lo doy. Le doy los mejores argumentos de su lado porque los conozco mejor que los del mío. ---Levantó la tetera y volvió a llenarle la taza, que estaba a dos tercios---. ¿Quién es su cliente en este asunto, doctor?

La pregunta llegó sin cambio de tono, entre el ruido del té y el crujido del mimbre. Vidal sostuvo la taza con las dos manos.

---Depende de qué llame cliente.

---Llamo cliente a quien figura en el campo correspondiente de su registro. Usted registra todo. Alguien figura.

En el campo cliente de la evaluación de viabilidad, creada un domingo de noviembre después de veintiuna noches sin código, figuraba su propio nombre. Vidal lo había escrito él y lo había leído cien veces y nunca lo había dicho en voz alta.

---Yo ---dijo.

---Ahí está ---dijo Sandoz, y se recostó, y no lo repitió, y ese fue el gesto: la respuesta quedó en el aire sin que la usara.

El goteo de la izquierda seguía. Vidal dejó la taza.

---Los motivos ---dijo---. Ansermet describe documentos. Los documentos describen la ventana. Ninguno describe por qué existe la ventana. Un análisis de este costo, con este riesgo penal, sin destino de publicación posible. La estructura del gasto es inauditable.

---La estructura del gasto está a su disposición, factura por factura, si la pide. Lo inauditable son los motivos, y en eso tiene razón, y va a seguir teniéndola. ---Sandoz se paró, despacio, y fue hasta la mesa de plantines; enderezó un tallo con dos dedos, sin mirar a Vidal---. No le pido confianza, doctor. La confianza es una moneda que usted no emite. Le pido diecisiete días y le doy lo que ninguna institución de la Tierra va a darle nunca. Los motivos, si le sirven de algo, se los cuento cuando esto termine.

---Cuando esto termine ---repitió Vidal.

---Es un plazo, no una evasiva. Usted trabaja con plazos. ---Sandoz se volvió, y por primera vez en la hora sonrió, poco, con algo que a Vidal le costó clasificar porque se parecía a la paciencia---. Pascal escribió que hay que apostar, que uno ya está embarcado y que abstenerse también es apostar. Yo desconfío de esa página, le aviso: es demasiado buena, y las cosas demasiado buenas se auditan. Pero en un punto el hombre midió bien. Usted ya apostó, doctor. Apostó en noviembre, cuando le escribió a Turín. Lo que hacemos acá es discutir las condiciones de una apuesta hecha.

Lo acompañó hasta el patio de grava él mismo. En la despedida le estrechó la mano con las dos suyas, un segundo más de lo protocolar, y Vidal sintió la aspereza de la tierra seca contra el dorso. El automóvil ya esperaba con el motor en marcha.


\scenebreak


El tren de la tarde bordeó el lago con el sol bajo. Vidal abrió la libreta a la altura de Nyon y empezó por donde empezaba siempre: la lista de afirmaciones, una por línea, cada una con su casillero de verificación al margen.

El apellido: verificable, verificado. El gradiente de Pellicori: verificable contra la publicación, lo comprobaría esa noche, y ya sabía el resultado. La cita de Pascal: exacta, incluida la desconfianza, que era de Sandoz y estaba declarada como suya. La estructura del gasto ofrecida factura por factura: una oferta, comprobable si la ejercía. La discusión perdida de la página cuatrocientos ochenta y siete: sin contenido afirmado, imposible de falsar. El plazo de los motivos: un plazo. Un plazo se cumple o se incumple; hasta el cuatro de mayo era indecidible por definición.

Repasó la lista dos veces. En once años de peritajes, todo interlocutor le había entregado, antes de la sexta pregunta, un dato falsable: una fecha corrida, un porcentaje redondeado a favor, una procedencia embellecida. Era el mecanismo de la gente, la grieta por donde se los medía. Buscó la grieta de Sandoz afirmación por afirmación y llegó a Lausana con la página limpia. El hombre había hablado una hora y cuarto y hasta sus opiniones venían con la incertidumbre declarada, como un resultado bien informado.

Al pie de la lista, en el renglón de síntesis, escribió: «No dijo nada verificablemente falso.»

Miró la frase. Era una descripción de ausencia, la clase de resultado que su propio informe de 2031 mandaba descartar, y la dejó, porque era lo único que tenía.

Cuando guardaba la libreta vio la mota de tierra en el puño derecho de la camisa, donde las dos manos del jardinero habían apretado la suya. Gris parduzca, dos milímetros, tierra de maceta. Levantó la otra mano para cepillarla y la detuvo a mitad de camino, sin causa asignable, y el puño siguió así, con su muestra adherida, todo el resto del viaje hasta Zúrich.

\clearpage

\chapter{Treinta horas}

\lettrine[lines=2, lhang=0.1, nindent=0.2em]{E}{l} primer nombre entró a las 7:01 por el acceso de servicio y Vidal lo tachó de la lista con lápiz. La cámara de obra daba un plano alto, sin audio, con la fecha corriendo en la esquina; la imagen refrescaba cada dos segundos y en cada refresco las figuras saltaban medio paso, como fotogramas de un registro antiguo. Nueve renglones en la pestaña dos del expediente. Cuatro de conservación, dos de la Comisión, un electricista, dos de seguridad diocesana. Fue tachando a medida que las siluetas cruzaban el umbral: reconoció a Chiara por la forma de cargar ---el cuerpo entero detrás de la caja de apoyo, las dos manos abajo, otro par de manos ajenas en la otra punta--- y a Ferrero por la altura y la curva de la espalda, adelante de todos, sin cargar nada, con el paso de quien entra a un lugar que conoce mejor que su casa.


A las 7:09 los nueve estaban adentro. Cerró la lista. Nueve tachados, cero pendientes.

Los dos estaban ahí por derecho propio y ese era el diseño entero. La nómina decía Fabbri, C., conservación textil, Opificio delle Pietre Dure, escuela directa de la intervención de 2002; decía Ferrero, O., Comisión de Conservación, autor del corpus fotográfico de referencia. Cada palabra era verificable y verdadera. Si los Carabinieri cruzaban esa nómina contra doce meses de registros, encontrarían dos carreras impecables que tenían todo el derecho del mundo a estar en esa sacristía esa mañana. Lo único que la nómina no podía listar eran los motivos, y los motivos, había aprendido Vidal en Ginebra, eran la única categoría que en 2033 seguía sin dejar traza.

La nave de Collegno estaba en silencio salvo el caudalímetro. Sobre la mesa del puesto remoto había tres cosas: el monitor con la cámara del acceso, el terminal de telemetría de la teca ---la vieja y la nueva, dos columnas de números que el proveedor de la donación volcaba por el mismo canal que figuraba en la pestaña seis--- y la planilla de papel donde Vidal asentaba una verificación cada quince minutos. Humedad, temperatura, oxígeno residual, células de carga. La telemetría tardaba entre dos y cuatro segundos en refrescar cuando él pedía el volcado. Empezó a cronometrar esa latencia sin decidirlo: tres segundos, dos, cuatro, tres. A la cuarta verificación ya tenía una media y a la octava un desvío. En treinta horas iban a ser ciento veinte verificaciones. Anotó la primera serie en el margen de la planilla, donde el formato tenía lugar para observaciones y él escribió números.

El teléfono fijo sonó a las 10:22.

---Nómina ---dijo Ledda---. Renglón diez. Incorporación de la Comisión, comunicada hoy a las nueve y cuarto. Ceruti, monseñor, conservador. Presencia parcial: una hora. Horario a confirmar.

Vidal miró la lista de nueve tachados.

---¿A confirmar por quién?

---Por él. La Comisión transmite, no fija. ---Una pausa de papel dándose vuelta---. Recalculo: si la hora cae en la cola de verificación de soporte, la fase se congela hasta que salga. Si cae de noche, improbable pero sin garantía, el umbral de aborto se adelanta al minuto treinta. Tercer efecto: el relevo de Fabbri se cancela.

---Explicá el tercero.

---La caja de apoyo la abre ella y nadie más. Con nueve adentro, la caja podía quedar cuatro horas bajo llave con Fabbri en descanso. Con un observador de horario propio, la caja tiene que tener a la responsable a menos de tres metros todo el tiempo que él esté en el edificio, y como el horario no existe, todo el tiempo es todo el tiempo. Fabbri pasa de dieciocho horas de presencia a veintiséis. La manipulación fina le va a tocar en la hora veinte de pie.

Vidal hizo la cuenta sobre el papel aunque ya la tenía hecha. Veintiséis horas. La ventana de sustitución en la madrugada, las manos de Chiara sobre la tela con un cansancio acumulado que ningún ensayo de Collegno había incluido como variable. El plan tenía dieciocho puntos firmados y ninguno contemplaba a un monseñor sin horario.

---La hora sin asignar es mala ---dijo Ledda.

Vidal tardó un segundo en contestar, y el segundo se le fue en registrar que era el primer adjetivo que le escuchaba en cuatro semanas.

---¿Alternativa?

---Ninguna que mejore. Se absorbe. Confirmación de Fabbri en la rotación de comida, catorce cero cinco, por esta línea. Cuarenta segundos, guion acordado.

Cortó. Vidal se quedó con el tubo en la mano un momento más de lo funcional y después lo colgó y volvió a la telemetría, que refrescó en tres segundos, dentro de la media.

A las 14:05 el fijo sonó de nuevo.

---Soy yo ---dijo Chiara. Se oía la calle detrás, un motor, una campana lejos---. Ya sé lo del proveedor nuevo. Lo vi entrar a saludar a las once. Saluda bien.

---Ledda recalculó. Tu relevo\ldots{}

---Me lo dijeron. Veintiséis.

---La manipulación te queda en la hora veinte. En los ensayos nunca pasaste de la hora nueve. Es una variable sin medir.

---Sebastián. ---Una pausa; algo que se acomodaba del otro lado, quizás el pelo, quizás nada---. La tela no sabe cuántos somos.

Él abrió la boca para corregir el razonamiento y la volvió a cerrar porque el razonamiento era correcto: la variable relevante era el estado de la tela y el estado de las manos, y las manos eran de ella, y ella acababa de decirle que respondía por sus manos. Todo lo demás era exposición de ella, con su nombre verdadero en el renglón cuatro de la nómina, su historia limpia de veinte años puesta como fianza de un plan que ella había firmado sabiendo el precio. Si la operación caía, el primer apellido de la lista era Fabbri. Ese cálculo Vidal lo tenía hecho desde enero y hasta esa mañana había sido un cálculo.

---¿Comiste? ---dijo ella.

---Son las dos de la tarde.

---Eso contesta la hora. ¿Comiste?

---Hay una verificación a las catorce y quince.

---Comé antes. Se me acaba el tiempo del guion. ---Y cortó ella, dos segundos antes de los cuarenta.

Vidal se quedó mirando la planilla. En la columna de observaciones de las 14:05 escribió «confirmación Fabbri, conforme» y al lado, sin decidirlo, «40 s: usó 38». Después fue hasta la cocina de la nave, comió de pie algo frío, y volvió al puesto antes de las 14:15.


\scenebreak


La tarde fue una serie. Verificación, latencia, columnas, margen. Tres segundos, tres, dos, cuatro. Entre una verificación y la siguiente el monitor mostraba la puerta de servicio cerrada y quieta, y la quietud, refrescada cada dos segundos, tenía un temblor mínimo de compresión que Vidal aprendió a descontar. A las 16:40 cronometró el caudalímetro de nitrógeno de la nave, veinte intervalos, media de quince coma uno. La media histórica era catorce coma nueve. Repitió la serie completa. Catorce coma ocho. Promedió las dos series, obtuvo el valor de siempre, y guardó el cronómetro sabiendo que el instrumento medido había sido otro.

No había nada que auditar en Turín. Las treinta horas corrían dentro de un edificio al que sus instrumentos llegaban por dos hilos ---una cámara de obra, una telemetría prestada--- y todo lo demás era gente. Once años construyendo un sistema donde ningún resultado dependiera de confiar en nadie, y el punto crítico de la operación más importante de su vida era el pulso de una conservadora en su hora veinte de pie y la firma de un químico de sesenta y tres años que había pasado la vida entera a un cristal de distancia de esa tela. Vidal revisó la calibración del terminal de telemetría, que estaba calibrado. La revisó de nuevo. Anotó las dos revisiones con la hora, porque su sistema exigía anotar hasta lo redundante, y la palabra «redundante» se quedó un rato en el borde de la página sin que la escribiera.

A las 19:07 la cámara de acta se activó.

Era el segundo hilo del canal del proveedor: un plano fijo, alto, sobre la mesa de conservación de la sacristía, montado por la Comisión para documentar la apertura como se había documentado todo en 2002. Sin audio, doce cuadros por segundo, luz pareja. Vidal vio entrar a los de conservación con la caja larga, vio el despeje de la mesa, vio a Chiara del lado izquierdo con otra conservadora enfrente. Vio a Ferrero ponerse los guantes.

Eran los de algodón. Los suyos, los del bolsillo del saco, los de siempre.

La apertura llevó veintitrés minutos contra los veinte estimados y Vidal los contó todos. Cuando la tapa cedió, las figuras se ordenaron alrededor de la mesa según el protocolo de la Comisión y hubo un intervalo en que nadie se movió, seis segundos que el plan no asignaba a nada. Después Chiara y la otra conservadora tomaron el soporte y el traslado a la mesa se hizo en el orden previsto, y a Ferrero le tocó lo que el acta le pedía: sostener el borde lateral, zona sin imagen, mientras las otras manos ajustaban la posición. La maniobra requería sostener y soltar. Sostener, según los ensayos con el lino de práctica, insumía entre dieciocho y veintidós segundos.

Vidal cronometró treinta y uno.

Las manos enguantadas de algodón quedaron sobre el borde nueve segundos después de que la posición estuviera ajustada, después de que la conservadora de enfrente hubiera retirado las suyas, después de que el motivo técnico de tenerlas ahí hubiera terminado. En el monitor eran dos formas blancas sobre una franja clara, refrescándose doce veces por segundo, quietas. Cuarenta años de placas, dieciséis mil, calibración cromática documentada, y todas tomadas desde el otro lado de algo: un cristal, una teca, un mandato. Las formas blancas se abrieron por fin, se levantaron dos centímetros, y se quedaron suspendidas un cuadro más antes de retirarse del plano.

Vidal tenía el lápiz sobre la columna de observaciones. Escribió la hora de inicio y fin de la fase, la duración total, la desviación contra el estimado. El renglón siguiente lo empezó ---«sostén del borde: 31 s, exceso de»--- y lo dejó ahí, sin sustantivo, porque los sustantivos disponibles eran todos de otra disciplina.

A las 21:50 el fijo sonó una vez y cortó antes del segundo timbre. Protocolo de Ledda para «sin novedad, próxima fase en curso». Vidal anotó la hora en la planilla y volvió al monitor, donde la puerta de servicio seguía cerrada y quieta, refrescándose cada dos segundos, y donde, de los nueve nombres de la lista, él era hasta esa hora el único que no había cruzado el umbral.

\clearpage

\chapter{La sacristía}

\lettrine[lines=2, lhang=0.1, nindent=0.2em]{E}{l} overol venía doblado en una bolsa de plástico y olía a fábrica. Vidal se lo puso encima de la camisa gris en la parte de atrás de la camioneta, con las rodillas contra la caja de herramientas, y Ledda le alcanzó la credencial de la cinta antes de que terminara de subir el cierre.


---Léala.

Plástico rígido, foto de tres meses atrás, banda magnética. AMBIENTI CONTROLLATI S.R.L. Debajo, VIDAL, S. Debajo, VERIFICACIÓN DE INSTRUMENTOS.

---Es mi nombre.

---Es su nombre. Contrato marco del catorce de febrero, dos facturas emitidas, aportes al día. Usted verificó las sondas de la celda de ensayo en Collegno en marzo y cobró por eso. La empresa lo dio de alta en el anexo de personal autorizado de la diócesis el mismo día. Los carabineros cruzaron el anexo el primero de abril. Su renglón salió sin observaciones.

---El anexo existe porque la empresa tiene contrato con la diócesis.

---Desde el treinta y uno.

Vidal pasó el pulgar por el canto de la credencial. La empresa estaba adentro de la catedral porque alguien había pagado la teca dos años atrás, y él iba a cruzar esa puerta con un contrato real, un trabajo real y una factura real. Lo único no auditable era por qué.

---Dos cosas ---dijo Ledda---. Una: el monseñor se fue a las diez y le dijo al sacristán que vuelve para la última fase. No dio hora.

---Entonces la ventana de las tres diez tiene una variable sin acotar.

---Sí.

---Se puede correr.

---No se corre la ventana. Se corre el aborto. Umbral en el minuto treinta y cuatro, lo decide el número. ---Ledda miró el reloj del tablero---. Dos: adentro hay cinco. Con usted, seis. Un cuerpo más son cuarenta gramos de agua por hora en el aire de la sacristía. Ferrero lo va a decir. Déjelo decirlo.

Cruzaron el acceso de servicio a las 22:18, entre la ronda del vigilante y el cambio de la bitácora. Dos cámaras nuevas sobre el dintel, instalación de obra, ángulo bajo. Las había pagado la misma donación. Vidal levantó la caja del banco de sondas con las dos manos, sostuvo la credencial hacia el lector con el mentón, y el lector hizo un pitido corto y verde.

Adentro el aire estaba seco y le secó la nariz en cuatro respiraciones. Cuarenta y un metros cuadrados; él conocía las cotas de memoria, transcritas en cinta azul sobre el piso de una nave de Collegno, y la sacristía real tenía el armario de ornamentos veinte centímetros más adentro de lo que la cinta decía, exactamente como Ferrero había avisado. Olor a cera y a madera vieja, y debajo un olor químico de la manguera flexible del equipo de clima portátil, que soplaba contra el zócalo con un zumbido de doce ciclos por minuto. La luz era fría, de paneles montados sobre trípodes. En el fondo, la teca nueva, abierta, con el marco desmontado y las cuatro células de carga a la vista.

Y en el centro de la sala, sobre la mesa de conservación, una superficie plana cubierta con papel de conservación y una sábana de algodón crudo.

Vidal fue a la teca. Su trabajo esa noche era ese: conectar el banco portátil a la telemetría, verificar la respuesta de las células de carga, dejar firmada la planilla de la empresa. Tardó veintidós minutos en lugar de los quince estimados porque el conector de la sonda de humedad estaba montado del lado incómodo y tuvo que trabajar de rodillas. Ferrero pasó detrás de él dos veces sin decirle nada, con los guantes de algodón puestos, y a la tercera se detuvo.

---Seis personas ---dijo---. En el noventa y siete, cuando el bombero rompió el cristal, adentro había dos y sobraba uno. ¿Usted sabe cuánta agua transpira un hombre despierto?

---Entre treinta y cincuenta gramos por hora, según carga térmica.

---Y la banda de la sonda es de dos puntos de humedad relativa. Haga la cuenta usted, que la hace más rápido.

---La hice en la camioneta.

Ferrero se tocó la muñeca izquierda con la palma de la mano derecha, un gesto de dos segundos, y siguió camino hacia el armario de ornamentos.

A las 22:51 Chiara le hizo un gesto con la barbilla desde la mesa. Vidal dejó el banco de sondas en el piso y se acercó, y ella retiró la sábana de algodón doblándola en tres hacia sí, y después el papel, y no dijo nada.

Lino. Sarga de espina tres a uno, urdimbre visible a cuarenta centímetros, el tono entre marfil y paja. Dos líneas paralelas de agujeros a lo largo de todo el eje mayor: las quemaduras de 1532, con los bordes carbonizados y las aureolas de chamuscado alrededor, simétricas porque el lienzo estaba plegado en cuarenta y ocho capas cuando la plata fundida entró. Los rombos de las manchas de agua, más claros, con el halo de arrastre en el borde. Donde habían estado los parches de las clarisas hasta 2002, la tela mostraba las perforaciones de aguja en hileras y un halo más claro que dibujaba el contorno de lo que ya no estaba; cuatrocientas ochenta y dos agujeros en el parche de arriba, contados sobre placas, reproducidos a mano, cuarenta por hora. Un pliegue transversal en el tercio inferior con la fibra vencida hacia adentro.

De la imagen, a cuarenta centímetros, no había nada. Fibrillas amarillentas discontinuas, algunas sí y las de al lado no, sin contorno, sin línea, sin acumulación de materia en las intersecciones de la trama. Nada que un ojo pudiera integrar.

Vidal retrocedió.

A dos metros y medio, entre la mesa y el armario, se formó. La figura frontal, los brazos cruzados sobre la pelvis, la banda oscura de las manchas del costado. Y arriba el rostro, que él había visto en dieciséis mil placas calibradas, en las de Pia, en las réplicas de 1931, en pantallas de calidad clínica a tres monitores, y que ahí, en una sacristía con luz de trípode, a dos metros y medio, era un rostro.

Contacto: no, porque la deformación de proyección exigiría. No siguió. Vapor: la difusión daría un gradiente que. Tampoco. Las once hipótesis estaban en su cabeza numeradas y ordenadas, con sus objeciones publicadas y sus réplicas, y ninguna arrancó. Se quedó donde estaba.

Después Chiara volvió a poner el papel, y él fue a la caja de herramientas y verificó por segunda vez el par de apriete del conector de la sonda, que ya había verificado.

---Ocho minutos ---dijo ella, más tarde, en voz de taller: baja, sin aire, la voz que no rebota en las paredes---. Estuviste ocho minutos parado ahí.

---No tengo ese número.

---Yo sí. Lo miré porque tenía que cerrar la fase.

Se habían sentado en el piso, contra el flanco de la mesa, del lado del armario, donde el vigilante del pasillo no llegaba con la vista. Chiara tenía el pelo atado con una banda elástica de las que vienen en los mazos de fichas. Le pidió la linterna de bolsillo y la tomó con las dos manos, como tomaba todo, y la apoyó apagada sobre la rodilla.

---Repasame las nueve fases ---dijo Vidal---. Desde el rodillo.

---Rodillo desde el borde de los pies, seis vueltas, tensión constante, yo marco el ritmo con la respiración y vos no me hablás. Plegado de transferencia, radio ocho, un solo movimiento. Tubo. Cierre del tubo lo hace Ledda. Después la gemela al soporte, verificación de posición con las dos referencias del marco, y ahí sí me hablás: me cantás las cotas.

---En la fase cuatro, cuando levantás el objeto para pasarlo al tubo, la mano izquierda va---

---La tela.

---¿Qué.

---Cuando levanto la tela. ---Chiara movió el pulgar sobre el cuerpo de la linterna, de arriba abajo---. Vos decile como quieras en la planilla. Adentro de la fase decime la tela, porque si me decís el objeto yo tengo que traducir, y traducir son dos décimas.

---La mano izquierda va debajo del borde a doce centímetros de la esquina, no a diez.

---A doce. Sí. ---Levantó las dos manos, las abrió, las volvió a bajar---. Con las dos, siempre. Nunca una sola aunque parezca que aguanta. Eso lo repito hace veinte años y lo voy a repetir mañana también.

El equipo de clima cambió de tono y volvió. Del otro lado de la puerta, el vigilante corrió una silla.

---Es la primera vez que la toco ---dijo Chiara.

---Tu maestra estuvo en 2002.

---Mi maestra estuvo en 2002. Yo tenía veintiuno y estaba pegando etiquetas en Florencia. ---Se rio con la boca cerrada, un sonido de nada---. A las siete y media, cuando la sacamos de la caja, pensé: pesa menos. Toda la vida escuchando cuánto pesa y pesa menos. Es una sábana. Después me acordé de que tenía que respirar.

Vidal miró la mesa por encima del hombro. El papel de conservación tenía una arruga larga en el tercio derecho.

---¿Vos qué viste? ---dijo ella.

---Fibrillas coloreadas en superficie, sin penetración aparente a esa distancia, discontinuidad compatible con lo publicado. Nada que agregue.

---Te pregunté qué viste.

---Eso es lo que vi.

Chiara le devolvió la linterna con las dos manos, la puso en la palma de él y esperó a que la cerrara antes de soltarla.

---Vale ---dijo, y se levantó apoyándose en la rodilla, no en la mesa.

A las 23:14 Ferrero estaba sentado en la banqueta baja junto al armario de ornamentos, con la espalda contra la puerta de madera, los guantes de algodón puestos y las manos sobre las rodillas. Tenía los ojos abiertos. Vidal pasó por delante con el banco de sondas.

---Cuatro horas ---dijo Ferrero.

---Tres horas cincuenta y seis.

---Usted, si le sacan el decimal, se queda sin conversación. ---Ferrero acomodó la espalda contra la madera---. En vigilia se espera y ya está, doctor. Es la única liturgia que usted practica desde que nació.

Vidal armó el banco al pie de la teca. Conectó la sonda, esperó el arranque, y en la pantalla aparecieron las cuatro células de carga con la lectura en cero y el cuarto decimal moviéndose: 0,0003, 0,0000, 0,0002, por la corriente de aire del equipo de clima contra el zócalo. Tolerancia declarada, ±1,5 gramos. El cuarto decimal era ruido y el ruido estaba adentro de banda, y él ya lo había verificado en Collegno con el mismo instrumento, y lo verificó otra vez ahí, de rodillas sobre la piedra, mirando cómo el número se movía sin ir a ninguna parte.

\clearpage

\chapter{Sustitución}

\lettrine[lines=2, lhang=0.1, nindent=0.2em]{A}{ las 2:58 la sacristía olía a cera y a piedra fría, y de los nueve nombres de la lista quedaban cinco cuerpos adentro: los cuatro del equipo y la asistente de conservación que Chiara había pedido para el sostén del borde. Un sexto, el electricista de la empresa de la teca, cubría el perímetro dormido en la camioneta desde la una, con autorización escrita de su encargado. El monseñor seguía afuera. Había dicho que volvía para la última fase, sin dar hora, y la última fase del cronograma oficial empezaba a las siete. Ledda había hecho el cálculo en voz alta a medianoche: cuatro horas de margen contra una palabra de sacerdote. Después había repartido las posiciones.}


---Ventana a las 3:10 ---dijo---. Fases uno a nueve, responsable Fabbri. Umbral de aborto, minuto treinta y cuatro. Lo decide el reloj. Preguntas.

Nadie tenía.

---A la hora, entonces.

Vidal ocupó la mesa auxiliar contra la pared norte, a dos metros de la mesa de conservación, con la hoja de servicio de Ambienti Controllati sujeta al portapapeles de la empresa. Verificación de instrumentos, teca de exhibición, turno nocturno. Su función legítima cubría la otra, la real: purgar el tubo de transporte con nitrógeno seco hasta cero coma dos por ciento de humedad relativa interior, y mantenerlo abierto el tiempo exacto que la fase crítica iba a durar, ni un segundo más, porque un tubo abierto de más era una segunda anomalía montada sobre la primera. Ya había medido las cuatro células de carga, la sonda de humedad, el sensor de oxígeno residual; las cifras estaban en la hoja, correctas, auditables una por una. Faltaba el renglón de cierre y la firma. Cada palabra del documento era verificable. El documento entero era falso. Escribió «sin novedades en el instrumental» con la letra de sus cuatrocientas cadenas de custodia, pareja, de mayúsculas alineadas, y comprobó que la letra ni se enteraba.

A las 3:10 Chiara dijo «fase uno» en el mismo tono con que pedía la sal.

El rodillo entró por el extremo de los pies. Vidal conocía las nueve fases de memoria y las había firmado una por una, y aun así el orden de los hechos lo tomó de costado: había aprobado un procedimiento y lo que avanzaba sobre la mesa era una tela de cuatro metros cuarenta con seiscientos años de plegados en los mismos lugares, cediendo al rodillo con un sonido que el ensayo con el lino belga había imitado casi bien. Casi. El lino de ensayo susurraba más agudo. Lo anotó en la cabeza y en ningún papel.

Chiara trabajó todo el tiempo con las dos manos, incluso donde el protocolo admitía una. La asistente sostenía el borde; Ferrero, a la cabecera, cantaba los tiempos con la voz de acta que había ensayado para las interrupciones, aunque nadie interrumpía. En el plegado de transferencia Vidal contó las fibras que el pliegue iba a costar y le dio una por centímetro, quizá menos, la marca de ella. Su propio pliegue, en el ensayo, había costado siete.

---Fase cuatro ---dijo Chiara, y las dos manos entraron debajo del borde a doce centímetros exactos.

---Cota confirmada ---dijo la asistente.

Por un segundo largo las dos telas estuvieron sobre la mesa a la vez: la verdadera a medio camino del tubo, la gemela todavía plegada en la caja de apoyo. Desde sus dos metros Vidal no pudo distinguirlas por la vista, solo por el sonido que no se repetía: la verdadera se dobló sobre sí misma con un ruido seco, de papel viejo, que el lino de ensayo nunca había imitado igual.

---Transferencia, en curso ---cantó Ferrero.

La tela cruzó el aire, cuatro metros cuarenta cediendo en un solo pliegue continuo, y entró al tubo de transporte que Vidal sostenía abierto con las dos manos desde hacía nueve minutos de purga. En cuanto Chiara liberó el borde, cerró la tapa: el sello de goma resistió un instante y después cedió con un chasquido que sintió en los dedos antes de oírlo. Bloqueó los tres seguros en el orden que había practicado ciento veinte veces con un tubo vacío y nunca, hasta esa noche, con algo adentro. Contuvo el aire sin decidirlo hasta el tercer seguro.

---Fase cinco ---dijo Chiara, sin que la voz le cambiara un tono entre una tela y la otra. La gemela salió del falso fondo de la caja de apoyo, ya plana, ya aclimatada, y subió al soporte de la teca con el mismo ritmo de rodillo que había marcado toda la noche.

---Cota confirmada ---dijo la asistente otra vez, idéntica a la primera; solo el número distinto en la hoja de Vidal decía que acababa de ser otra tela.

De las nueve fases, las cuatro que tuvieron a las dos telas comprometidas a la vez duraron doce minutos. Vidal los tenía cronometrados de los ensayos: doce y medio, la mejor corrida. Esta dio doce y ocho segundos. El tercer seguro del tubo le costó dos segundos más que en los ensayos: las manos menos firmes de lo que había calculado. Firmó la hoja de servicio a las 3:24 porque tenerla abierta ya era una anomalía.

---Fase seis, verificación de posición ---dijo Ferrero, con la voz de acta---. Referencia norte.

---Coincide ---dijo la asistente, sin levantar la vista.

---Referencia sur.

---Casi. ---La asistente corrigió con dos dedos, un milímetro---. Coincide.

La verificación de posición cerró al minuto veintinueve. Chiara bajó el marco, Ledda confirmó el sellado, y las cuatro células de carga tomaron la masa de la gemela.

La curva subió, tocó banda y siguió.

---Nadie toca nada ---dijo Ledda.

Vidal miró la pantalla del banco de sondas. Masa en equilibrio, dos gramos y una décima por encima de tolerancia y todavía en ascenso. La debilidad conocida: la pendiente de adsorción. Cinco cuerpos durante seis horas en cuarenta y un metros cuadrados, unos cuarenta gramos de agua por hora cada uno; la humedad relativa de la sacristía estaba dos puntos arriba de la de cualquier ensayo, y la gemela bebía. Lo habían corregido en el respaldo, lo habían medido en la celda, y la celda nunca había tenido cinco personas respirando alrededor. El cálculo cabía en un renglón y él lo hizo en medio segundo, y el cálculo servía exactamente de nada, porque el protocolo de alarma disparaba con desviación sostenida treinta segundos y la única variable disponible era el tiempo.

La asistente hizo un gesto hacia el marco, dos dedos adelantados, como quien va a sostener algo que empieza a ceder.

---Quieta ---dijo Ledda, sin subir la voz, y la mano volvió atrás.

Contó. Con las manos quietas sobre el portapapeles, contó los segundos como los había contado toda la vida, y la serie aguantó entera. Al nueve la curva perdió pendiente. Al once cruzó de vuelta hacia banda. Al doce estabilizó, un gramo y tres décimas adentro de tolerancia, y ahí se quedó.

Ledda esperó otros veinte segundos completos antes de mover la cabeza.

---Sigue.

---¿Cuánto quedó? ---preguntó Chiara, sin mover las manos del soporte.

---Un gramo y tres décimas adentro de tolerancia. Estable desde hace cuatro segundos.

Ella asintió una vez, sin preguntar el resto, y volvió al rodillo.

El pico quedaba escrito: doce segundos fuera de banda en una telemetría continua y auditada que el veintinueve alguien iba a revisar entera. Una corrección al registro deja firma de corrección; el pico, sin tocar, solo queda como el pico. Quedaba, entonces. Le habría puesto el círculo de no requiere acción, si el riesgo hubiera estado en un papel suyo.

Ferrero firmó el acta de estado a las 3:47, ante la cámara de la Comisión, con los guantes de algodón puestos. La pluma le patinó en la primera curva del apellido y la firma salió más gruesa que la de los documentos que Vidal le conocía. Estado de conservación conforme, superficie de apoyo verificada, sin observaciones. El renglón de observaciones, vacío, con la rúbrica debajo: a partir de esa firma, nadie volvía a mirar la tela hasta el veintinueve. Ferrero tapó la lapicera y se quedó con las dos manos en el borde de la mesa, la cabeza baja, sin moverse durante unos segundos que Vidal no llegó a contar.

---Diecisiete días ---dijo, sin levantar la cabeza---. Que sean los mismos diecisiete.

Chiara pasó junto a la teca al recoger el rodillo. Se inclinó apenas hacia el cristal.

---Portati bene ---le dijo a la gemela, en voz baja, y siguió enrollando.

Vidal lo anotó como lo que era: una instrucción impartida a un objeto sin capacidad de cumplirla. Le llamó la atención estar de acuerdo con la instrucción.


\scenebreak


Vidal volvió al puesto remoto de Collegno a las cuatro y media, cruzando el acceso de servicio entre dos rondas, la credencial de la empresa en el bolsillo, su nombre verdadero en la credencial. A las 6:50 la nave estaba a diecinueve grados; Chiara dormía sentada contra la esclusa, el pelo atado con un hilo de algodón ---su firma y la de Ferrero en el acta de conservación habían cerrado su parte antes del amanecer, y la asistente quedaba en la sacristía para las firmas de cierre de las trece, en representación de los dos, cláusula prevista desde el protocolo original---, y Ledda revisaba el caudalímetro de nitrógeno, anotaba el valor, y pasaba al siguiente instrumento sin cambiar el ritmo de la mano.

A las 7:40, cuarenta minutos después de lo previsto, el monseñor entró a la catedral para la última fase y bendijo una teca que contenía lino belga con intervalo de datación 2030-2033. A las 8:20 el tubo salió por el mismo acceso, entre bolsas de escombro real, en la camioneta rotulada, con el chofer de once años de historia laboral limpia y un remito que declaraba descartes de obra sin detallar cuántos. En el monitor, el vigilante de la ronda se acercó a la puerta trasera abierta, miró adentro con la linterna un segundo largo, y firmó el remito apoyado en la carrocería sin pedir que bajaran nada. Vidal contó los segundos que tardó la camioneta en salir del encuadre desde que el vigilante se apartó ---dieciocho--- y no volvió a contar nada más esa mañana.

Ferrero llegó a Collegno a las nueve. Traía los guantes puestos desde la puerta, o se los había puesto en el auto; no dijo cuál, y Vidal no preguntó.

En el inventario de la nave, el ingreso previsto para las 13:40 figuraba ya de puño y letra de Ledda: material de calibración, un bulto.

A la 13:38 el portón lateral sonó dos veces.

Entraron el tubo entre Ledda y el chofer, lo depositaron en el soporte junto a la mesa de sujeción de cinco metros veinte, y el chofer se fue con el remito firmado y sin haber mirado nada más de tres segundos seguidos.

---Cuatro paradas de la ruta real, ninguna de las nuestras ---dijo Ledda, cerrando el portón---. Cinco horas. Sin desvío. El acta cerró a la una en punto. Nueve firmas, ninguna novedad.

La línea de nitrógeno hizo su clic, catorce coma nueve segundos de media, y el tubo quedó horizontal bajo las luces de la sala, con sus tapas termoselladas y su rotulación idéntica a la del ingreso a la catedral.

Los cuatro quedaron alrededor del tubo sin que nadie lo decidiera. Chiara tenía las manos cruzadas adelante, quietas por primera vez en dos días. Ferrero, a los pies del tubo, con los guantes de algodón que ya no se sacó en toda la tarde, miraba el cilindro de PVC sin moverse, las manos quietas a los costados, más tiempo del que cualquier inspección visual necesitaba. Ledda consultó su reloj una vez y guardó la muñeca. Doce horas de aclimatación por delante: hasta la 1:40 de la madrugada, el contenido del tubo era una hipótesis sellada, y todos los presentes conocían la cifra y ninguno la dijo.

Vidal abrió el registro sobre la mesa auxiliar. Fecha, 12 de abril. Evento, ingreso de material, inicio del período de análisis, cuatrocientas ocho horas. Hora, escribió, 3:10.

Se detuvo. Tachó con una sola línea horizontal, la primera tachadura de ese registro desde que el registro existía, y escribió 13:40. La mano había ido sola a la otra hora, la de la sacristía, la de las telas cruzándose sobre la mesa; algún proceso suyo consideraba que los diecisiete días habían empezado ahí y llevaba diez horas corriendo por su cuenta, sin código asignado. Miró la corrección. En cualquier cadena de custodia suya, una tachadura sin inicial al margen invalidaba el renglón entero. No la inicialó.

Después buscó la columna. Tenía el catálogo entero disponible: evento, responsable, duración, entregable, observaciones, el círculo del margen para lo que se documenta y se suelta. Recorrió los campos dos veces con la vista y en ninguno cabía lo que había que asentar, porque lo que había que asentar carecía de sustantivo en su sistema: un valor alto en una variable que ninguna de sus planillas medía, presente desde las 3:22 de la madrugada según su propia estimación, sostenido, sin señal de decaer. Dejó observaciones en blanco. Le puso el círculo.

Cuando levantó la vista, los otros tres seguían alrededor del tubo, en silencio, y él tenía la palma derecha apoyada en el plástico, a la altura del primer tercio, donde adentro estaría el extremo de la cabeza. Retiró la mano. El plástico le había dejado el frío en la palma y ningún registro suyo decía desde cuándo la tenía ahí.

\clearpage

\chapter{Diecisiete días}

\lettrine[lines=2, lhang=0.1, nindent=0.2em]{E}{l} tubo se abrió a las 8:09 y lo primero que salió fue el olor. Seco, con un fondo orgánico apenas dulce, distinto del lino de ensayo, distinto de cualquier textil que hubiera pasado por las manos de Vidal en veinte años de peritajes. Un olor con historia térmica, habría escrito si el registro tuviera columna para eso. El caudalímetro de nitrógeno hizo clic dos veces mientras Chiara aflojaba la segunda tapa, a catorce y a dieciséis segundos, dentro de rango.


---Zona once antes que la seis ---dijo Chiara---. Como en el ensayo. Ferrero, usted el borde norte. Sebastián, vos el rodillo, y a mi ritmo, no al tuyo.

El despliegue llevó una hora y cuatro minutos. Vidal ejecutó el rodillo como lo había practicado el sábado bajo la mirada de Chiara, con las muñecas bajas y el peso repartido, y cuando el último tramo cedió sobre la mesa la tela hizo el ruido que el lino belga nunca había imitado: seco, de papel viejo, un sonido que venía de adentro de la trama. Después la sala quedó en silencio, salvo el clic irregular del caudalímetro, y los cuatro estuvieron un momento sin hacer nada frente a cuatro metros cuarenta de lino extendido bajo las lámparas, y fue Vidal el que rompió el intervalo porque el intervalo carecía de asignación en el cronograma.

---Inventario de estado ---dijo---. Cuarenta minutos.

Tardó ochenta y cinco. Lo supo porque se cronometró, como siempre, y dejó correr el reloj sin corregir el desvío. Recorrió la superficie con la lupa frontal y la voz en volumen de dictado, y Ferrero cotejaba cada punto contra sus propias placas, extendidas en la mesa auxiliar en el orden del corpus. Las quemaduras de 1532, simétricas porque la tela estaba plegada en cuarenta y ocho capas cuando la plata la alcanzó, medidas una por una. Las perforaciones de aguja donde estuvieron los parches de las clarisas, con su halo más claro alrededor. El faltante de 1988 en el borde lateral, siete centímetros, los hilos del corte todavía francos. En la zona de imagen la voz se le puso más lenta. Dictó fibrillas amarillentas, discontinuas, coloración superficial, sin acumulación de materia en las intersecciones, y entre un rasgo y el siguiente se le abrían pausas que la descripción técnica no explicaba. Nadie las señaló.

---La mancha de agua grande del lado izquierdo ---dijo Ferrero de memoria, sin mirar sus placas--- termina a cuatro hilos del pliegue axial.

---Siete ---dijo Vidal, con la lupa puesta.

---Cuatro. La fotografié en 2011.

Vidal corrió la regla graduada. Siete hilos. Ferrero fue hasta la mesa auxiliar, buscó la placa 9.114, la suya, la estudió medio minuto con la cabeza inclinada.

---Siete ---dijo---. La placa era mía y la memoria también. Quédese con la placa.

Sacó los guantes de algodón del bolsillo del saco y se los puso, aunque el inventario era visual y las manos de Ferrero llevaban toda la mañana lejos de la tela.

Del despliegue quedaron tres fibras sueltas sobre la superficie de la mesa, dentro de la proyección del orillo. Chiara las levantó con pinza de punta blanda y ahí empezó el procedimiento que Ferrero había impuesto en marzo y que Vidal había aceptado con salvedad registrada: cada fragmento pesado dos veces, dos operadores, la balanza verificada contra pesa patrón antes y después. Fibra uno, ciento cuarenta microgramos; la segunda pesada dio ciento cuarenta y uno y se asentó el promedio. Vial, precinto, caja de dos cerraduras. Ferrero giró su llave, Vidal la suya, y las firmas quedaron cruzadas en el acta de custodia, la de Ferrero arriba en la primera muestra, la de Vidal arriba en la segunda, alternadas según el protocolo. Los dos coma cuatro miligramos del punto cuatro esperaban al día siete y ese mismo circuito. Vidal midió el tiempo del procedimiento completo: nueve minutos por muestra. En cualquier otro peritaje habría tachado el paso por redundante. Acá firmó las tres veces sin objeción y comprobó que la letra le salía pareja.

Cerró el acta de despliegue al mediodía. Escribió: «Objeto desplegado sin incidencias. Estado conforme al corpus fotográfico de referencia. Inventario adjunto.» Chiara firmó debajo y agregó, de su mano: «La tela quedó plana. Tensión conforme en las dieciséis zonas.» Los dos sustantivos quedaron uno arriba del otro, en el mismo papel, y el papel se archivó así.

El barrido de reconocimiento arrancó a las 12:40. Resolución gruesa, doscientos cuarenta micrones por píxel, la superficie completa en una sola pasada para fijar el mapa de referencia antes de las ochenta y dos franjas finas de los días siguientes. El cabezal corrió sobre el riel sin una vibración fuera de banda ---el acelerómetro de la mesa, que había volcado a papel esa mañana él mismo, en su copia única, venía limpio--- y a las 16:10 Vidal tenía en pantalla lo que el STURP había perseguido durante ciento veinte horas con los instrumentos de 1978: la tela entera, caracterizada, georreferenciada hilo por hilo, antes de que terminara la tarde del primer día de trabajo. Comió de pie frente al monitor algo que Chiara le puso al lado, y cuando quiso registrar la pausa de almuerzo descubrió que se le había pasado la hora de registrarla.

Para calibrar la ingesta hizo falta el corpus de referencia, y el corpus de referencia estaba en la cajonera metálica del fondo, donde había estado desde antes de que él pisara la nave. Treinta y una muestras de lino arqueológico datado, mediterráneo oriental, del siglo III antes de la era común al II de la era común, cada una con procedencia documentada, informe de laboratorio o datación publicada, nueve de ellas de contextos de Judea del siglo primero. Vidal revisó los expedientes uno por uno, buscando el defecto, porque su propio corpus para lino judeo-romano era pobre ---doce muestras, y un colega honesto atacaría cualquier resultado por ahí--- y este lo llevaba a treinta y una sin que pobre dejara de ser la palabra, pero con un espesor que ninguna universidad que él conociera podía exhibir. Ni Zúrich. Ni el Instituto de Jerusalén, que llevaba cuarenta años excavando el material en su propio suelo. Los expedientes estaban en regla. Las cadenas de custodia cerraban. Abrió el registro para asentar la integración y la mano le empezó una línea que decía «procedencia del corpus:» y la línea se quedó en los dos puntos, igual que otra nota suya de octubre, y arriba escribió en cambio «corpus de referencia local: verificado, apto, integrado a calibración». Quién juntaba treinta y una muestras de lino datado del Mediterráneo oriental y las dejaba en un galpón de Collegno era una pregunta que cabía en un renglón del mismo registro.

No la hizo.

A las cinco y veinte la humedad relativa de la sala subió dos puntos y medio. Causa asignable: los ciclos de esclusa de la tarde, cuatro entradas en dos horas, la carga de agua de los cuerpos. Chiara detuvo el barrido con la mano levantada, sin tocar nada.

---La tela está tomando ---dijo---. Se corta acá y se deja reequilibrar.

---¿Cuánto?

---Dos horas. Puede que dos y media.

Vidal abrió el cronograma para calcular la cascada: dos horas hoy corrían la primera franja fina de mañana, que corría la segunda, que comía margen del día tres. Tenía el argumento armado, el mismo que había usado toda su vida contra todo retraso, y lo que dijo fue:

---De acuerdo.

Chiara ya estaba ajustando la consigna del clima y ni registró la respuesta. Vidal sí. Se quedó con el cronograma abierto, mirando el punto donde su objeción habría empezado, y anotó: «Detención por reequilibrio, 2 h, responsable Fabbri, aprobada.» La palabra aprobada le quedó rara en la página, no por falsa sino por rápida: la aprobación había salido antes que el cálculo, y en su sistema el orden era siempre el inverso. Usó las dos horas para auditar las costuras del mapa de reconocimiento, franja contra franja, y encontró dos zonas donde el solapamiento del quince por ciento le permitía corregir una deriva de registro que con apuro se le habría pasado. Cuando el barrido se reanudó a las 19:40, con la tela reequilibrada, la señal venía más limpia que a la tarde. Lo verificó dos veces. La segunda verificación sobraba y la hizo igual, pero por primera vez en memoria la hizo por gusto.

A las nueve de la noche Ledda pegó en la pared del fondo el plan completo: cuatrocientas ocho horas en columnas, del ingreso a la restitución, hora por hora, cada bloque con responsable y entregable, el bloque de dieciocho horas sin asignación entre el día catorce y el quince, la ventana de seis minutos al final como una celda igual a las otras. Once hojas con cinta en las cuatro esquinas. Ferrero se paró delante con las manos atrás, todavía con los guantes de algodón puestos, y lo leyó entero, de arriba abajo, dos veces. Se tocó la muñeca izquierda por encima del puño de la camisa.

---En el setenta y ocho fueron ciento veinte horas ---dijo--- y los americanos durmieron por turnos en el piso del salón. Acá hay cuatrocientas ocho y una columna que dice quién duerme cuándo. ---Hizo una pausa---. En los hospitales cuelgan estos planes donde el paciente los vea. Para que sepa que alguien tiene un plan. ¿Y si el día nueve se enferma uno de nosotros, quién lo tacha?

---El margen del catorce absorbe hasta dieciocho horas ---dijo Vidal.

---Eso contesta otra pregunta. ---Ferrero se abotonó el saco---. Buenas noches. Mañana traigo la actualización de la tabla de emisividad que dejé en marzo. Encontré dos valores más para revisar.

Se fue. Ledda cerró la caja de teléfonos de la esclusa, confirmó el circuito de la mañana con una frase de siete palabras y se fue detrás. Chiara tardó más: apagó las lámparas de trabajo, dejó la de guardia, y extendió la tela de servicio sobre la mesa con el mismo cuidado con que a la mañana la había retirado, las dos manos abiertas, alisando desde el centro. Debajo de la tela de servicio, la depresión de las dieciséis zonas siguió sonando en su frecuencia baja, apenas audible entre los clics del nitrógeno.

Vidal hizo el volcado de notas, el cotejo del cronograma, la revisión de calibraciones pendientes. Antes del cierre volvió a la ventana de la esclusa y miró la mesa un momento: el bulto largo y bajo bajo la tela gris, cuatro metros cuarenta documentados hasta el hilo. Mañana a las ocho, primera franja fina. Doce micrones por píxel sobre una superficie que once hipótesis publicadas llevaban un siglo sin explicar. Sintió la boca hecha agua, literal, salival, y le buscó causa: había cenado mal y rápido. La causa cerraba a medias.

Chiara apareció a su lado con el bolso al hombro y las llaves en la mano.

---¿A qué hora venís mañana?

---A las siete.

---El plan de la pared dice ocho.

Vidal miró la pared del fondo, las once hojas, la columna del día tres con su nombre en el primer bloque.

---A las siete ---repitió.

Chiara se colgó mejor el bolso y salió primera por la esclusa, y si sonrió, lo hizo ya de espaldas, donde el registro de él llegaba sin resolución.

\clearpage

\chapter{La esquina del ochenta y ocho}

\lettrine[lines=2, lhang=0.1, nindent=0.2em]{D}{ía} seis, diecisiete de abril: la discordancia apareció en la consolidación nocturna de la franja setenta y uno, la penúltima del lote, y Vidal la encontró a las 7:12, con el abrigo todavía puesto. El canal de imagen la marcaba en amarillo: subcanal de textura, borde lateral, ángulo inferior respecto de la cara frontal. Una región de once centímetros por seis donde la distribución de diámetros de fibra dejaba de ser una campana y se partía en dos.


Lo primero fue el instrumento. Recalibró el banco contra la tarjeta de referencia, doce micrones por píxel, dentro de especificación. Volvió a correr la franja setenta y uno completa, cincuenta minutos de barrido sobre una zona que ya tenía barrida. La discordancia volvió a salir en el mismo lugar, con los mismos contornos, y el caudalímetro de nitrógeno hizo clic dos veces mientras la pantalla la dibujaba.

Chiara llegó a las ocho menos cuarto y él la llamó a la pantalla sin preámbulo.

---Decime qué ves.

Ella miró el mapa de falso color un rato largo, con las dos manos apoyadas en el borde de la mesa, sin tocar nada.

---¿Esto es el borde del ochenta y ocho?

---La zona adyacente. La tira que cortaron salía de acá. ---Señaló la línea recta que interrumpía el mapa: siete centímetros de ausencia, documentados en todas las placas desde 1988---. Lo que marca el modelo es lo que quedó alrededor.

---¿Y qué marca?

---Dos poblaciones de fibra. Una coincide con el resto del lienzo. La otra tiene diámetro medio menor, torsión distinta, y esto. ---Cambió de capa: el subcanal químico---. Celulosa de algodón entre el lino. Trazas de rubia y de goma. Alguien tiñó hilo nuevo para que igualara al viejo.

Chiara se acercó a la mesa. Pidió la lupa frontal, pidió luz rasante, y estuvo veinte minutos sobre esos once centímetros con la cara a un palmo de la tela, sin apoyar nada, las manos cruzadas atrás como quien se las ata para no usarlas. Cuando se enderezó tenía la voz más baja que de costumbre.

---Hay empalmes. Hilo a hilo, en bisel, con la torsión copiada. ---Hizo una pausa---. Yo pasé por esa esquina en el inventario. La conté. Treinta y ocho hilos de urdimbre, como todo el resto. Quien hizo esto reprodujo la cuenta exacta.

---¿Vos lo habrías visto?

---A ojo, no. Con lupa, sabiendo dónde mirar, quizás. ---Se corrigió---: Ahora sí lo veo. Ayer estuve a diez centímetros y era tela.

El resto de la mañana fue procedimiento. Tamiz cruzó el subcanal químico contra el corpus documental de técnicas de reparación textil europea: el empalme en bisel con teñido de igualación estaba descrito, tenía nombre de taller y tenía siglos. La química del tinte y la hilatura del algodón cerraron el intervalo de la intervención: 1290 a 1420, confianza 0,93, canal por canal, cada rasgo auditable hasta el espectro de origen. Vidal recorrió la cadena entera dos veces. No había capa opaca en el resultado: era el modelo trabajando como él lo había construido, razones a la vista, apelable en cada eslabón.

Y el eslabón final era este: el corte de 1988 caía completo dentro de la zona remendada. El límite entre tejido original y tejido de reparación pasaba a cuatro milímetros de la línea de corte, del lado de afuera. Los tres laboratorios habían medido bien. Habían medido un zurcido.

Ferrero llegó a las once y media, y Vidal le puso el mapa delante sin decirle qué era. Lo dejó leer. Ferrero se inclinó hacia la pantalla, pidió la escala, pidió las coordenadas contra su propio corpus de placas, y a medida que entendía se fue inclinando más, hasta que el nudo de la corbata de lana le colgó sobre el teclado. Preguntó la resolución del subcanal de textura. Preguntó si el bisel de los empalmes tenía dirección consistente. La tenía. Entonces Ferrero se rió: una risa corta, de dos golpes, que se cortó sola, y que Vidal registró porque en cuatro semanas de operación era la primera.

---Cuarenta años ---dijo Ferrero---. Cuarenta años discutiendo protones y nadie miró la costura. ---Se enderezó. Le cambió la cara, y las manos le fueron al bolsillo derecho del saco, donde asomaban los guantes de algodón---. ¿Usted entiende las consecuencias de esa imagen, doctor? Del ochenta y ocho vive media literatura y de la otra media vive el que la refuta. Las consecuencias.

Se puso los guantes y no dijo más.

Chiara seguía junto a la mesa. Habló mirando la esquina, no a ellos.

---No fueron las clarisas. Las clarisas cosían para que se viera: parche, refuerzo, puntada honesta. ---Tocó el aire sobre la zona, a dos centímetros de la tela---. Esto lo hizo alguien que cosía para que nadie lo viera nunca. Un siglo antes que ellas, si tu intervalo está bien. Otra escuela. ---Levantó la vista---. Me hubiera gustado conocerle las manos.

Vidal abrió el registro y empezó a redactar el hallazgo. La primera versión decía «zurcido invisible de factura extraordinaria, indetectable para cualquier examen previo a 2033». La leyó y tachó «invisible», que era redundante con «indetectable», y «extraordinaria», que era una opinión. La segunda versión decía «reparación por empalme de alta calidad técnica». Tachó «de alta calidad técnica»: la calidad estaba en los datos, no hacía falta adjetivarla. La tercera quedó así: «Zona de reparación por empalme con teñido de igualación, 11 x 6 cm, borde lateral. Intervalo de intervención 1290--1420, confianza 0,93. El muestreo de 1988 cae íntegramente dentro de la zona.» La leyó en voz alta, para la sala. Sonaba a lo que era.

---Habría que avisarle a Zúrich, a Oxford y a Arizona que midieron bien ---dijo Chiara---. Les va a dar pena y alivio al mismo tiempo.

---El punto diecinueve sigue vigente ---dijo Ferrero, con los guantes puestos---. Nadie avisa nada a nadie hasta la devolución.

---Era un chiste, Ottavio.

---Ya sé. Contesto igual.

Ledda, que había escuchado todo desde la mesa de logística sin intervenir, hizo una sola pregunta:

---¿Esto cambia algo del cronograma?

---Nada ---dijo Vidal---. La toma del día siete sigue donde estaba. Adyacente a la zona, del lado del tejido original. Ahora con más razón: sabemos exactamente qué lado es cuál.

Se hizo un silencio de trabajo, de los buenos, y Vidal se dio cuenta de que llevaba un rato con los hombros bajos, una postura que no se conocía en horario de operación. Chiara había vuelto a la lupa. Ferrero copiaba coordenadas a su libreta con una letra apretada de químico viejo. Durante quizás veinte minutos, la nave de Collegno fue un lugar donde cuatro personas habían encontrado algo cierto.

Fue Vidal el que lo terminó.

---Hay que decir la otra mitad ---dijo---. Antes de que la primera se instale sola.

Chiara levantó la cabeza de la lupa.

---¿Qué otra mitad?

---Esto demuele el ochenta y ocho. No exactamente: demuestra que el ochenta y ocho dató otra cosa. Es distinto y es peor. ---Fue hasta la pantalla y puso el mapa completo del lienzo, cuatro metros cuarenta, con la zona amarilla en la esquina como una estampilla en una sábana---. La datación medieval queda explicada. La tela queda sin fecha. Once centímetros tienen edad conocida y son los once que menos importan. El resto volvió a no tener ninguna.

---Tiene la que tenía ayer ---dijo Chiara.

---Ayer tenía una datación publicada en contra. Hoy tiene cero dataciones. La pregunta está más abierta que esta mañana, no menos. ---Miró la estampilla amarilla---. Ganamos un resultado y perdimos el único dato duro que existía sobre este objeto. Los dos enunciados son verdaderos y el segundo es el que importa.

---Hasta el día siete ---dijo Chiara.

---Hasta el día siete, la tela no tiene edad. Y si el día siete sale mal, tampoco después.

Ferrero se había quedado quieto, con la libreta abierta. Habló despacio, eligiendo.

---En el ochenta y ocho hubo tres laboratorios, conferencia de prensa y un número. El número era correcto y midió un remiendo. ---Se tocó la muñeca izquierda, por encima del guante---. Guarde esa secuencia, doctor. Le va a servir cuando tenga su propio número.

Nadie contestó eso. La depresión de las dieciséis zonas siguió con su zumbido bajo y el clic del nitrógeno cayó donde caía siempre, entre once y diecinueve segundos, sin patrón.

A las seis de la tarde Ferrero salió a la esclusa a revisar su teléfono en la caja, como cada día a esa hora, por si la Comisión lo necesitaba. Tardó más que de costumbre. Cuando volvió traía el saco abotonado y la cara neutra que Vidal ya sabía leer al revés.

---Novedad administrativa ---dijo---. Monseñor Ceruti pidió a la Comisión los registros del sistema de clima de la teca nueva. El histórico completo desde la puesta en servicio.

---¿Quién es Ceruti? ---preguntó Chiara, que lo sabía.

---El conservador que la Comisión sumó durante el traslado. Estuvo el once a la mañana, saludó, se fue. ---Ferrero se sentó---. El pedido entró hoy por mesa de entradas, con nota formal. Verificación de condiciones previas a la ostensión. Rutina preostensión.

---¿El pedido dice clima o dice telemetría? ---preguntó Ledda.

---Dice clima. Humedad, temperatura, oxígeno residual.

---¿Con qué firma se extrae y quién la prepara?

---Un técnico de la empresa de la teca prepara el volcado. La Comisión lo visa antes de entregarlo. Plazo ordinario: dos semanas, aunque con la ostensión encima puede ser menos.

Ledda abrió su cuaderno, escribió el apellido, el objeto del pedido y la fecha, y subrayó el renglón una vez, con regla. Guardó el cuaderno. En ningún momento dijo que fuera rutina.

Vidal hizo el cálculo sin proponérselo, con la velocidad con que hacía todos los cálculos que hubiera preferido no hacer. El pico de dos gramos y una décima vivía en la tabla de células de carga, con marca de tiempo 3:10 y segundos, doce de abril. Las células de carga eran otra tabla. Clima era clima: humedad, temperatura, oxígeno. El pedido de Ceruti, tal como estaba redactado, pasaba a cuatro milímetros del corte. Del lado de afuera. Por ahora, la frontera caía del lado de afuera.

---Si amplía el pedido a telemetría completa ---dijo Ledda---, quiero saberlo el mismo día. Antes del visado, no después.

---Lo va a saber ---dijo Ferrero.

Chiara empezó a levantar la mesa para el cierre de jornada. Vidal volvió al registro del día. Bajo el hallazgo de la esquina, redactado tres veces y firmado, abrió un renglón nuevo en la columna de asunto abierto: «Ceruti --- logs sistema de clima teca --- alcance del pedido: clima. Rango temporal: histórico completo. Ampliaciones: verificar.» Dejó el margen sin círculo, porque el círculo era para lo que no requería acción.

Antes de guardar releyó la página entera. Arriba, un resultado con intervalo estrecho y cadena auditable, el mejor de su carrera sobre el peor problema de su campo. Abajo, un apellido con una nota formal. Los dos renglones eran correctos. Los leyó dos veces, y las dos veces el de arriba pesó menos.

\clearpage

\chapter{Mecanismo desconocido}

\lettrine[lines=2, lhang=0.1, nindent=0.2em]{L}{a} consolidación del canal de química de superficie cerró a las 6:47 del día once. Vidal leyó la salida de pie, con el zumbido de la presión positiva de fondo y el clic del caudalímetro a sus intervalos irregulares, que había dejado de cronometrar porque ya tenía la media.


Los datos eran los mejores de la historia del problema y lo eran de un modo verificable, línea por línea. Doce micrones por píxel sobre la zona de imagen completa, ochenta y dos franjas con solapamiento del quince por ciento, cuarenta y ocho bandas espectrales por punto. Ninguna partícula de pigmento por encima del límite de detección, que estaba dos órdenes de magnitud por debajo del de 1978. Ningún rastro de direccionalidad: el mapa de orientación de rasgos daba isotropía dentro del error en toda la superficie frontal, y una mano ---cualquier mano, con cualquier herramienta--- deja dirección. La coloración era una oxidación selectiva de la capa externa de las fibras, con deshidratación de la celulosa, confinada a una profundidad que el corte virtual resolvía en fracciones de micrón: la pared primaria alterada, el interior blanco. Fibrilla por fibrilla, además, la alteración era discontinua: una fibrilla coloreada junto a otra intacta en el mismo hilo, sin gradiente entre ambas, sin capilaridad, sin acumulación de materia en las intersecciones de la trama. Todo lo que el STURP había medido en ciento veinte horas estaba ahí, confirmado y refinado hasta un decimal que ellos no habían podido soñar.

Ferrero llevaba veinte minutos frente al monitor de la mesa auxiliar, con las manos cruzadas a la espalda, recorriendo el corte virtual de una fibrilla como se recorre una radiografía propia.

---Pellicori midió esta fluorescencia con un instrumento que ocupaba una mesa entera ---dijo---. Y le salió bien, dentro de lo que su época le dejaba. ¿Sabe cuánto tardamos en escribir el párrafo de conclusiones, en el setenta y ocho? Yo lo sé porque lo tengo estudiado. Meses. Cuarenta personas, meses, para tres líneas.

---Todavía falta la clasificación ---dijo Vidal.

---Eso digo.

A las 9:15 alimentó a Tamiz con el paquete consolidado y las once hipótesis de formación de la literatura, cada una parametrizada contra rasgos medidos. El modelo tardó lo previsto. La salida entró al registro a las 15:03.

Contacto: descartado. El contacto exige correlación entre intensidad y presión, y la intensidad de la imagen variaba donde la topografía de un cuerpo no puede variar. Vapor: descartado; la difusión deja gradiente y capilaridad, y no había ni una cosa ni la otra. Pigmento: descartado por ausencia de materia. Bajorrelieve calentado: descartado; el chamuscado tuesta el hilo hasta adentro y fluoresce rojizo, y la imagen era superficial y muda bajo ultravioleta, tal como Chiara recitaba de memoria desde los ensayos. Ácido diluido: descartado por perfil químico. Radiación, en sus tres variantes publicadas: descartadas por ausencia de perfil de penetración. Las otras tres, descartadas cada una por un rasgo medido, con el rasgo citado.

Y después, en el campo de clasificación, tres palabras: «fuera de corpus». Y debajo, en el campo de detalle: «mecanismo no representado».

Vidal corrió el diagnóstico de integridad. Limpio. Verificó la calibración del banco contra la tarjeta de referencia, que había verificado esa mañana. Doce micrones. Volvió a la consola.

En noviembre, en Zúrich, había leído las conclusiones publicadas del STURP ---antes de que Ansermet le hiciera llegar, en enero, el ejemplar anotado envuelto en papel madera--- y frente a la frase «mecanismo desconocido» había escrito en su archivo de lectura una sola línea: cuarenta personas, ciento veinte horas, y entregan una abdicación. Un perito que concluye ``desconocido'' está declarando el límite de sus instrumentos, había pensado, y confundiéndolo con un límite del mundo. Ahora los instrumentos eran suyos, la resolución era de 2033, el corpus de rasgos era el más denso jamás construido sobre el problema, y su propio modelo acababa de firmar la misma frase con mejor ortografía.

Cincuenta y cinco años de progreso instrumental para volver, con intervalos más estrechos, al mismo renglón.

---Está donde estábamos nosotros ---dijo Ferrero. Sonaba a diagnóstico, sin triunfo.

---No exactamente.

---¿En qué difiere?

Vidal abrió la boca para citar la resolución, las bandas, el límite de detección. Cada número que tenía para ofrecer hacía el renglón final más pesado, no más liviano. Cerró la ventana de detalle.

---En que yo puedo demostrar todo lo que descarto ---dijo.

---Nosotros también podíamos. Esa nunca fue la parte difícil.

A las 15:40 empezó a atacar el resultado. Lo hizo con el método que usaba contra los resultados ajenos, el que había desarmado el díptico de Amberes y doscientos objetos menores: si la clasificación no cede ante la presión, es sólida; si cede, era un artefacto de umbral. Primera pasada: relajó los umbrales de descarte de las once hipótesis, uno por uno, hasta el borde de lo defendible. El modelo devolvió las mismas once exclusiones con residuales mayores y la misma clasificación. Cuatro horas y diez de cómputo. Segunda pasada: forzó la categoría, obligando a Tamiz a elegir la hipótesis menos mala aunque estuviera fuera de banda. El modelo eligió contacto, como Vidal había previsto, y adjuntó un mapa de residuales que encendía la zona frontal entera: un ajuste que ningún revisor del mundo dejaría pasar, y Vidal era el revisor más duro que conocía. Cuatro horas y media. Tercera pasada, ya de noche: hipótesis compuestas, mecanismos combinados de a dos, veintiocho combinaciones. Ninguna bajó del umbral de residuales. Cinco horas.

Ferrero se quedó hasta la segunda. En algún momento, sin que nadie se lo pidiera, dictó de memoria los parámetros publicados de la variante ácida ---concentraciones, tiempos, el perfil de fluorescencia esperado--- para que el descarte fuera inatacable.

---Si la va a matar, mátela bien ---dijo---. Una hipótesis mal descartada resucita en el peor momento, y después alguien tiene que salir a enterrarla en público. Pregúntele a cualquiera que haya vivido el ochenta y ocho.

Se puso los guantes de algodón para irse. Era temprano para el frío.

Poco antes del amanecer, con la tercera pasada ya cerrada, Vidal hizo la cuenta que venía postergando desde la tarde. Las tres pasadas habían consumido trece horas y cuarenta de un presupuesto de cómputo que compartía cola con los controles ciegos de la datación, y los controles ciegos eran innegociables porque él los había hecho innegociables, por escrito, con su firma al margen. En la pared del fondo, el plan de Ledda seguía pegado en once hojas, hora por hora, y las celdas del día doce que Vidal había pensado usar para el análisis fino de la franja de imagen central estaban ahora hipotecadas por sus propias repeticiones. Una cuarta pasada mordía el bloque de dieciocho horas rotulado «margen». Ese bloque tenía un dueño y el dueño lo había pedido sin explicación. Vidal escribió la orden de la cuarta pasada, la miró en la pantalla, y la dejó en cola sin ejecutar.

Había otra cosa, y la anotó porque anotar era lo que correspondía aunque no tuviera columna. Los canales individuales se comportaban como siempre: espectro, química, geometría, cada uno devolvía «no concluyente» con residuales publicables, auditables hasta el dato de origen. La capa de convergencia, en cambio, acompañaba cada clasificación con un índice global de confianza, y ese índice había hecho algo que Vidal revisó tres veces antes de asentarlo: había subido entre pasadas. 0,84 en la primera. 0,89 en la segunda. 0,93 en la tercera. Cuanto más presionaba él las hipótesis, más segura estaba la convergencia de que ninguna alcanzaba, y ningún canal individual justificaba la cifra, y las herramientas de interpretación, aplicadas al índice, devolvían ruido. Los canales decían no sé con la voz de siempre. La convergencia decía no sé con una firmeza creciente que no le rendía cuentas a nadie. Escribió: «Convergencia: comportamiento divergente del resto de la arquitectura. Caracterizar.» Le puso fecha y hora, y ningún círculo al margen. Antes de cerrar el registro revisó la entrada de Ceruti: sin novedad desde el visado, alcance sin ampliar. La dejó abierta igual.

Medir mejor. Durante veinte años esa había sido su respuesta cuando alguien le preguntaba por qué el problema de la imagen seguía abierto: porque nadie había medido bien, porque 1978 era prehistoria instrumental, porque un rasgo sin resolver es un rasgo sin resolución suficiente. La respuesta tenía la elegancia de ser comprobable, y acababa de comprobarla. Había medido mejor que nadie en la historia del problema. El problema estaba intacto. Más limpio, mejor delimitado, con intervalos estrechos alrededor de un centro vacío.

Chiara entró a la sala pasadas las seis, con la bata de trabajo todavía puesta, a preparar el banco para la sesión de la mañana. Miró la pantalla desde lejos, sin acercarse a leerla.

---¿Cenaste?

---Hay tres pasadas completas. Las tres devuelven lo mismo.

---Eso no era la pregunta.

Descolgó el cable del escáner de mano y empezó a enrollarlo, vuelta por vuelta, con el pulgar guiando el arco para que ninguna espira montara sobre otra. Vidal conocía ese enrollado. Lo había visto antes de cada pregunta que después le costaba semanas.

---El modelo evaluó once hipótesis ---dijo él, porque llenar el silencio con datos era más barato que esperar---. Las descartó todas con el rasgo citado. Sin proponer una duodécima. Frente a datos corruptos exhibe exactamente este comportamiento, así que corrí integridad dos veces. Limpio las dos.

---O sea que no sabe.

---Sabe todo lo que se puede saber con estos datos. Lo que dice es que el mecanismo está fuera de su corpus.

---Sebastián. ---Terminó una vuelta, empezó otra---. Cuando el modelo te dio 1290-1420 en la esquina, con el zurcido, ¿dudaste?

---Tenía cadena auditable. Canal por canal.

---Y esto tiene tres pasadas y dos diagnósticos.

Chiara colgó el cable del gancho, alineó el conector hacia abajo, y se quedó con las dos manos apoyadas en el borde del banco, como sostenía las cosas que eran de otro.

---¿Al modelo le creés cuando no sabe, o solo cuando sabe?

Vidal miró el frente del rack central. Los conectores estaban en cuatro columnas de seis. Los contó de arriba abajo, veinticuatro, y el número era el de siempre, y la pregunta seguía donde ella la había dejado. Chiara apagó la luz del banco y se fue a dormir sin esperar respuesta, que era su manera de dejar claro que la respuesta no era para ella.

En la pantalla, la clasificación seguía abierta. Fuera de corpus. Mecanismo no representado. Debajo, en la cola de ejecución, la cuarta pasada esperaba una tecla.

Vidal empezó a contar los conectores de nuevo, por columnas esta vez. Seis, doce, dieciocho.

\clearpage

\chapter{Cero coma noventa y siete}

\lettrine[lines=2, lhang=0.1, nindent=0.2em]{D}{ía} doce. La noche siguiente a haberla dejado en cola, lanzó la cuarta pasada a las 0:07. Antes hizo el cálculo dos veces: la cola de controles ciegos retomaba a las seis, la pasada pedía cuatro horas y diez según la estimación del propio modelo, y entre el final de una cosa y el principio de la otra quedaba una hora cuarenta y tres de margen. El bloque de Ledda, las dieciocho horas sin asignar entre el catorce y el quince, no se tocaba. En el registro de papel asentó la hora de lanzamiento y el motivo: «verificación adicional de estabilidad de resultado», y debajo dejó la línea de resultado en blanco, para llenarla cuando hubiera resultado.


Los ventiladores del rack central subieron medio tono. Era el ruido de carga plena, el mismo que en Jerusalén, el mismo que en Zúrich, y durante un rato la sala quedó reducida a ese zumbido y al clic del caudalímetro de nitrógeno, que llegaba a intervalos de once a diecinueve segundos con media de catorce coma nueve. Vidal se sirvió el café que había quedado del turno de la tarde, frío, con el borde de la taza marcado por el café anterior, y se sentó a esperar con la espalda contra la pared del fondo, donde el plan de las cuatrocientas ocho horas colgaba en once hojas con cinta en las cuatro esquinas.

Repasó las costuras entre franjas mientras esperaba. Solapamiento del quince por ciento, ochenta y dos franjas, ocho mil y pico de junturas auditadas una por una durante los días tres a seis. Todas cerraban. Ese era el trabajo que sabía defender: cada píxel con cadena, cada cadena con origen, cada origen con certificado de calibración fechado en Ivrea. Se durmió veinte minutos sin proponérselo y lo despertó el cambio de tono de los ventiladores, que bajaban.

La pantalla decía 4:03.

Índice global de convergencia: 0,97. Objeto de la evaluación: perfil integrado de degradación del soporte textil. Clasificación: compatible con lino mediterráneo oriental, siglo primero. Debajo, los canales, uno por uno, con sus salidas individuales de siempre: fibra, no concluyente por corpus insuficiente; química de superficie, compatible con múltiples trayectorias, sin discriminación; geometría de trama, dentro de tolerancia sin rasgo distintivo; historia térmica, intervalo ancho; documental, sin datos anteriores a 1355. Ningún canal, tomado solo, pasaba de 0,7. Ninguno afirmaba nada que un perito serio firmaría. La capa de convergencia tomaba exactamente esos mismos datos, los integraba de un modo que ninguna herramienta traducía, y devolvía 0,97.

Corrió integridad. Limpio. Lo corrió de nuevo. Limpio. Aplicó las herramientas de interpretación al índice y devolvieron ruido, el mismo ruido que en las tres pasadas anteriores, un ruido que ya empezaba a conocer de memoria, con su textura y todo. Después hizo lo que quedaba: ablación por canal. Sin química de superficie, el índice caía a 0,71. Sin geometría de trama, 0,74. Sin los dos, 0,66. Reponía cualquiera de ellos y el índice volvía a la zona alta sin escala intermedia. El número entero era mayor que cualquier suma de sus partes, y la parte que hacía la suma era, por diseño, la única habitación de la casa sin ventanas.

Escribió la serie en el margen del registro: 0,84, 0,89, 0,93, 0,97. Cuatro pasadas, cada una más hostil que la anterior, y el índice subiendo con cada ataque. En cualquier otro contexto esa curva se llamaba robustez. En este, él todavía no sabía cómo se llamaba.

Quedaba una vía, y la conocía desde antes de lanzar la pasada, y había esperado que el resultado la volviera innecesaria. Instrumentar la capa. Abrir la convergencia, insertar sondas en las activaciones, volcar los pesos, reconstruir el camino del dato al número. Los pesos estaban a tres metros, en hardware propio, sin réplica, con acceso de una sola persona que era él. Nadie en el mundo podía impedírselo. Podía hacerlo esa misma noche.

Hizo la cuenta despacio, porque la cuenta merecía despacio. El artículo 3 del reglamento definía «auditable» con dos criterios acumulativos: mecanismo trazable y desempeño verificable. Tamiz incumplía el primero desde siempre; vivía del segundo. Y el segundo pertenecía a una configuración congelada: once años de resultados, cuatrocientas cadenas de custodia, Amberes, el óstracon, cada acierto documentado contra esa versión exacta de los pesos y de la ejecución. Una capa instrumentada era otra ejecución. Otro modelo. Un modelo de cero horas de historial, sin un solo resultado validado, cuyas salidas no heredaban nada del anterior, ni siquiera este 0,97, porque el 0,97 era del testigo y el testigo habría dejado de existir en el momento de abrirlo. Él había escrito eso, casi con esas palabras, en las rondas de comentarios de 2028: el requisito de trazabilidad es correcto en principio e inaplicable en la práctica a las arquitecturas de mejor desempeño. Constaba en el anexo público, con su nombre. Cinco años después la frase volvía por la puerta de servicio y le cobraba en su propia moneda: para auditar el número tenía que matar al único sistema del mundo capaz de haberlo producido.

Se le había secado la boca. Fue hasta el banco de Chiara, verificó la calibración de la sonda de humedad de sala, que había verificado a las 22:00 y que daba lo mismo que a las 22:00, y volvió a sentarse.

La asociación llegó sola, sin que la llamara, mientras miraba el directorio de los pesos sin tocarlo. Alguien había hecho eso una vez. Alguien, en algún punto anterior a 2027, había accedido a los pesos de un modelo cerrado de peritaje, un modelo corporativo con ciento cuarenta empleados y tres centros de datos, y había entrenado treinta ídolos contra él hasta que pasaron todos. Nunca hubo imputación. Nunca hubo literatura sobre el vector. Lo que sí hubo fue un muerto: el modelo. Cada certificado que había emitido en su vida pasó a valer el papel en que estaba impreso, y los museos y los dos fondos soberanos que habían comprado sobre esa firma cobraron en pérdidas la diferencia entre un número y un número con cadena. Abrir un modelo cerrado tenía un precedente, y el precedente terminaba en un instrumento del que ya nadie podía fiarse y en un hombre al que nadie había encontrado.

Anotó la asociación en el registro porque el registro estaba para eso. «Asociación surgida durante análisis: Basilea 2027, extracción de pesos de modelo cerrado. Pertinencia con el caso presente: nula.» Releyó la segunda oración. La pertinencia era nula: distinto objeto, distinto agente, distinta finalidad. Le puso al margen el círculo de no requiere acción y el círculo le salió menos redondo que de costumbre.

A las 5:20 abrió el árbol de consolidación del día doce. El procedimiento estándar era uno y lo había ejecutado mil veces: el resultado se integra al consolidado, el consolidado se cierra con hora y firma, y a la mañana el equipo lo lee. En el consolidado, el 0,97 iba a estar a la vista de Ferrero, que hablaba en consecuencias, y de Chiara, que iba a hacer una sola pregunta de una línea, y tarde o temprano en el archivo de la cláusula octava, que decía «toda anomalía registrada durante los trabajos» y que era el único punto sin fecha de todo el contrato. Toda anomalía registrada. Depende de qué llames registrada. Un resultado sin cadena auditable, según la regla que él mismo había escrito en un informe de 2031 y aplicado a la carrera de otros, era un artefacto, y un artefacto se descarta. Descartarlo era un acto de método. Tenía once años de precedentes internos, incluido uno del que se acordaba con fecha y hora.

No lo borró.

Cortó el archivo de salida del árbol de trabajo, creó en su disco personal un directorio nuevo y lo llamó 2033-04. Pegó. El archivo pesaba catorce kilobytes. Verificó la copia contra el original byte a byte, borró el original de la cola, y en el registro de papel llenó la línea que había dejado en blanco a las 0:07: «Cuarta pasada: sin cambios respecto de consolidado vigente. Estabilidad confirmada.» Cada palabra de esa línea era verificable por separado. Se oyó reír, una vez, un ruido corto que la sala devolvió sin eco, y lo dejó sin clasificar.

Eran casi las seis menos cuarto y por el tragaluz de la nave el negro empezaba a aflojar en los bordes. Vidal apagó la pantalla, pasó por la esclusa interna y entró a la sala blanca, donde la luz de guardia dejaba la mesa de sujeción en una penumbra pareja, de museo cerrado. El lienzo estaba desplegado en sus cinco metros veinte de mesa, con las dieciséis zonas de depresión sosteniéndolo sin un pliegue, y el aire olía a nada, que era el olor que costaba más caro de toda la operación.

Se acercó hasta los cuarenta centímetros de trabajo. A esa distancia la imagen se deshacía, como en todos los barridos: fibrillas amarillentas discontinuas, sin contorno, sin acumulación en las intersecciones de la trama, un rasgo por vez y ninguno que sumara figura. Ahí abajo estaba todo lo que él sabía medir, y todo lo que él sabía medir decía no concluyente. Buscó con los ojos, durante un rato que su registro de jornada nunca iba a tener, el rasgo que justificara el número. La superficie devolvió lo de siempre.

Retrocedió. Un paso, dos, tres, contándolos porque eran pasos y estaban ahí. A los dos metros y medio el rostro se armó de golpe, entero, con los ojos cerrados y la asimetría del pómulo que Ferrero tenía fotografiada en dieciséis mil placas.

Ninguno de sus instrumentos trabajaba a esa distancia.

\clearpage

\chapter{La traza}

\lettrine[lines=2, lhang=0.1, nindent=0.2em]{D}{urmió} tres horas en el catre de la esclusa, lo justo para que la cola de controles ciegos retomara sola a las seis según lo previsto. La marca amarilla apareció en la consolidación de las 6:52, día trece, en el canal de residuales de química de superficie. Región de nueve por cinco centímetros, cara frontal, cuadrante inferior izquierdo, lejos del zurcido del borde lateral y fuera de la imagen, a cuatro centímetros del borde de una mancha de agua. Vidal amplió. El canal reportaba partículas discretas, entre dos y ocho micrones, con firma espectral de sulfuro de mercurio.


Bermellón.

Contó las detecciones porque estaban ahí y porque el número importaba: cuarenta y tres partículas sobre el umbral, distribuidas sin patrón dentro de la región, ninguna fuera de ella. Junto al bermellón, el canal marcaba una segunda población: granos de laca roja orgánica sobre sustrato de alúmina, degradación compatible con siglos de antigüedad, y trazas de material proteico que podían ser aglutinante o podían ser piel humana de seiscientos años de manipulaciones. La composición conjunta era una paleta. Vidal la conocía de tres peritajes de tabla gótica: bermellón, laca de rubia o de quermes, temple. Siglo XIV, con margen hacia los dos lados.

Se quedó mirando la pantalla el tiempo que tardó el caudalímetro en hacer cuatro clics. En 1389 un obispo le había escrito al papa que un artista confesó haber pintado la tela. Vidal había leído ese memorial en un avión, cuarenta y dos páginas de latín, y lo había degradado a indicio: borrador sin sello, confesión de oídas, litigio de por medio. Lo había desarmado con criterios de peritaje y el desarme seguía siendo válido. Ahora el artista sin nombre entraba al laboratorio por la otra puerta, la de los datos: cuarenta y tres partículas sobre el umbral, firma espectral de sulfuro de mercurio.

Verificó la calibración del canal contra los patrones del día. Estaba en banda. La verificó de nuevo.

Para caracterizar la laca hacía falta más que espectro: hacía falta masa, y la masa exigía levantar las partículas. Una toma adhesiva de superficie, localizada, destructiva para lo levantado. Vidal escribió el procedimiento antes de las ocho y lo dejó sobre la mesa de reuniones con el campo de firmas en blanco.

Ferrero llegó a las nueve y diez con los guantes de algodón ya puestos. Ni sacándolos del bolsillo ni ajustándoselos: puestos desde la puerta, como si la conversación hubiera empezado en el auto.

---Leí su procedimiento ---dijo---. Le pido que lo suspenda.

---Es una toma de superficie. Micrográmica. La zona está fuera de la imagen.

---Sé leer un procedimiento, doctor. Le pido que lo suspenda igual. ---Se sentó del otro lado de la mesa, las manos enguantadas una sobre otra---. Piense en la secuencia. Usted ya demolió el ochenta y ocho: el carbono midió un zurcido. Muy bien. Ahora su modelo encuentra pigmento del siglo catorce y usted lo confirma con masa. ¿Qué titular sale de esa secuencia? ``La ciencia desmiente el carbono y encuentra la pintura.'' Dos resultados de la misma máquina, en direcciones opuestas, y ningún adulto entre ellos y el público. Si esa zona entra al informe, ¿quién la explica? ¿Usted? ¿Con qué autoridad, si su informe no existe legalmente?

---El informe no se publica. Eso está firmado.

---El punto diecinueve expira con la devolución. Después de mayo esa traza es un arma cargada en un archivo que no controlamos ni usted ni yo. ---Hizo una pausa y bajó la voz otro grado---. En el ochenta y ocho mi... ---Se detuvo. Alineó los guantes con el borde de la mesa---. En el ochenta y ocho nadie pensó la secuencia. Pensaron el resultado.

Vidal registró la frase interrumpida y no le asignó columna.

---Usted no me está pidiendo que suspenda una toma ---dijo---. Me está pidiendo que elija qué resultados produce este laboratorio. Que mida lo que ayuda y deje de medir lo que estorba. Se llama selección de resultados, y usted se lo enseñó a dos generaciones de químicos.

---Eso no es lo que le pedí. Le pido que...

---Eso es lo que dijo. Lo que dijo es ``suspéndalo''. La palabra que usó fue esa.

Ferrero no había descruzado las manos enguantadas. Vidal tenía la boca seca y recalculaba de memoria, por tercera vez, la calibración del canal de residuales. El caudalímetro hizo un clic entre una frase y la siguiente.

---Si esto sale mal ---dijo Ferrero---, y con ``esto'' no me refiero a la toma, me refiero a todo, ¿quién lo paga? Porque le aseguro que no lo paga usted. Usted vuelve a Zúrich con su máquina. Lo pagan los que esperan del otro lado del anuncio.

---Yo no anuncio nada. Yo mido.

---Esa frase la dijeron tres laboratorios en el ochenta y ocho. Palabra por palabra.

A las 11:40, con la discrepancia de Ferrero registrada y firmada en contra según el mecanismo de custodia compartida que ellos mismos habían pactado, Ferrero giró su llave de la caja sin que se lo pidieran, en el instante justo, con el mismo pulso de siempre. Chiara marcó el borde de la región con dos referencias de cinta de enmascarar y esperó la señal de Vidal.

---Ahora ---dijo él.

Ella apoyó el cuadrado adhesivo con las dos manos, sin presión de más, y contó en voz baja, tan baja que era apenas un movimiento de labios: uno, dos. Vidal la siguió con el cronómetro y no con la voz. Al llegar a treinta y dos levantó la mano y ella despegó de un solo tirón continuo, sin pausa a mitad de camino, el mismo gesto que había practicado cien veces con lino de ensayo y que ahora hacía sobre lo otro. El adhesivo quedó rotulado, sellado, y el vial cerró con el chasquido de siempre. Ninguno de los tres dijo una palabra que el registro necesitara, y el silencio que siguió duró más que los treinta y dos segundos.


\scenebreak


El espectrómetro devolvió la masa a las 16:20. Laca de rubia sobre alúmina, con trazas de quermes. Aglutinante proteico, probablemente huevo. Bermellón de proceso seco, molienda irregular. Paleta completa, coherente, y para un taller europeo, con los métodos de época: siglo XIII a XV, centro de masa en el XIV. Tamiz consolidó y devolvió el intervalo con confianza 0,91, canal por canal, cada firma trazable a su patrón de calibración. Era el resultado más limpio de los diecisiete días.

Chiara lo leyó de pie, con las dos manos en la hoja aunque la hoja era de él. Después fue hasta el carro de instrumental, tomó la cinta de sujeción que había quedado suelta desde la mañana y empezó a enrollarla, despacio, pareja.

Vidal esperó. Había aprendido a esperar eso.

---¿Sabés cómo se consagraba una copia de la Síndone? ---dijo ella.

---No es mi campo.

---Se apoyaba sobre el original. ---Terminó una vuelta, empezó otra---. Está documentado. Siglos de eso. Los talleres pintaban la copia, la llevaban a Turín o a Chambéry, y la copia se ponía en contacto con la tela, extendida encima, a veces horas. Hay copias que lo dicen en la inscripción: \textit{extractum ab originali}. Tocada al original. Eso era lo que valía, más que la pintura. Mi maestra restauró dos, una española y una de Savoia. Las dos con la leyenda.

---¿Copias pintadas? Al temple.

---Bermellón, laca, lo que hubiera en el taller. Y pintura de esa época seca mal y suelta grano durante años. ---Colgó la cinta enrollada del gancho del carro y recién entonces lo miró---. Tu traza es un pintor del siglo catorce que la tocó. O un pintor del siglo catorce que la hizo. ¿Tu máquina distingue?

Vidal hizo el cálculo antes de contestar, aunque ya lo había hecho mientras ella hablaba. Cuarenta y tres partículas localizadas en una región de contacto plausible. Densidad compatible con transferencia por presión. Densidad compatible también con restos de ejecución, si el ejecutor había trabajado por zonas y aquella era marginal. Ninguna estratigrafía posible: las partículas estaban sueltas sobre la fibra, sin capa, sin orden. La razón de verosimilitud entre las dos hipótesis daba uno, con error grande. El dato no distinguía nada.

---No ---dijo.

---¿Y se puede repetir la toma? ¿Buscar más grano, otra distribución, algo que incline?

Vidal abrió el mapa de residuales en la pantalla del banco. La región de nueve por cinco estaba gris. Las cuarenta y tres detecciones habían sido todo lo que el lienzo tenía para dar en esa zona, y ahora estaban disueltas en el tren de vacío del espectrómetro, convertidas en un archivo de picos perfectamente auditables. Lo que quedaba sobre la tela estaba por debajo del umbral. Bajó el umbral hasta el piso de ruido, a sabiendas de que a ese nivel el canal inventaba. Ruido.

---La zona quedó agotada ---dijo---. Lo que había, lo medimos. Lo que medimos, lo consumimos. No hay segunda pasada.

Chiara se quedó mirando el mapa gris un momento.

---Entonces tenés un dato perfecto que no dice nada ---dijo---. ¿Eso también lo firmás?


\scenebreak


Esa noche, cuando la nave quedó en luces de guardia y el único ruido era el clic del caudalímetro a sus intervalos de once a diecinueve segundos, Vidal arrastró la pizarra blanca desde el rincón de la sala de reuniones hasta debajo de la lámpara y la limpió con el borrador seco hasta que el blanco fue parejo. Destapó el marcador. Escribió con la letra de sus cadenas de custodia, mayúsculas alineadas, cuatro renglones:

1988 --- el carbono dató una reparación. Intervalo correcto, objeto equivocado. Sin valor sobre el soporte original.

Imagen --- once hipótesis descartadas con rasgo medido. Clasificación: fuera de corpus. Mecanismo no representado.

Convergencia --- 0,97 sobre perfil de degradación, compatible siglo I. Ningún canal individual lo sostiene. Inauditable.

Traza --- paleta siglo XIV, confianza 0,91, auditable. Compatible con contacto y con ejecución; el dato no elige entre las dos y ya no queda material para que elija.

Retrocedió un paso y leyó el tablero completo. Cuatro resultados, cada uno el mejor dato de su tipo que ese problema había producido en un siglo de intentos.

Quedaban cinco días de cronograma y los iba a administrar: controles ciegos, redacción, reembalaje, ensayos de restitución. Reordenó mentalmente el índice del informe y donde decía ``Conclusiones'' puso ``Estados''. Anotó en el margen inferior de la pizarra los bloques restantes, día por día, con el marcador moviéndose casi solo mientras él repasaba la cola de cómputo, y cuando terminó la lista la mano siguió, dos centímetros más abajo, trazando algo que la vista tardó en procesar porque venía de la mano y la mano estaba fuera del registro.

Trazos verticales, cortos, con un punto encima de cada uno, uno al lado del otro, orientados hacia un rectángulo pequeño dibujado al final de la serie.

Una fila de gente esperando frente a una puerta.

Vidal miró el dibujo el tiempo de dos clics del caudalímetro. Después contó las figuras, porque estaban ahí: catorce. Levantó el borrador, lo apoyó sobre la primera y las barrió de un solo pase, de izquierda a derecha, en el orden en que habrían entrado. Sobre el blanco quedó una sombra gris, baja, del ancho exacto de lo borrado, y por más que pasó el borrador dos veces más la sombra se movió sin irse, como se mueven los residuos que están por debajo del umbral pero encima de cero.

Borró las cuatro líneas de arriba también, una por una, hasta que la pizarra volvió a estar pareja bajo la lámpara. El registro en papel podía tachar sin inicialar y quedar inválido; una pizarra no dejaba ese rastro, y esa era exactamente la diferencia que la hacía segura. Dejó el marcador en la bandeja, destapado.

\clearpage

\chapter{Espejo}

\lettrine[lines=2, lhang=0.1, nindent=0.2em]{E}{l} puerto cuatro del conmutador tenía tráfico saliente y el puerto cuatro tenía que estar muerto.


Vidal lo vio a las 2:40 del día catorce, mientras revisaba la cola de los controles ciegos. Los otros diodos parpadeaban al ritmo que él conocía: cómputo interno, cabina de almacenamiento, nada más. El cuatro parpadeaba con cadencia propia, regular, baja, la cadencia de una transferencia que administra su ancho de banda para que nadie la mire dos veces.

Abrió los contadores del conmutador. El puerto cuatro salía hacia la línea de fibra del taller, la conexión comercial de la razón social que servía de cáscara a la nave, la que en teoría existía para que un taller electromecánico pudiera facturar. Tráfico acumulado desde el día uno: cuarenta y un gigabytes con decimales. Verificó el número contra los volúmenes diarios de consolidación. La curva calcaba la suya con un tres por ciento de sobrecarga, el costo de cifrado y encapsulado, noche por noche, una ventana fija a las 4:30.

Cada espectro. Cada log. Las ochenta y dos franjas, el zurcido, la traza agotada del cuadrante inferior. Sus notas de trabajo digitales, las que precedían al papel.

Buscó el servicio. Estaba donde tenía que estar un servicio legítimo: declarado, con manifiesto propio, con hashes por archivo. Recorrió el manifiesto del día doce con el dedo en la pantalla, porque la vista sola ya le había fallado una vez esa noche. 4:31. Un archivo de catorce kilobytes. El hash coincidía con el que él había verificado byte a byte antes de cortarlo del árbol de trabajo y llevárselo a su disco.

La cuarta pasada. El 0,97. Había salido hacia Ginebra veintiocho minutos después de nacer y una hora antes de que él lo borrara de la cola creyendo que borraba la única copia ajena.

Se le secó la boca. Contó los diodos del conmutador, dos veces, ocho y ocho, y el resultado fue el mismo y no informaba nada.

Quedaba una hipótesis decente: intrusión. Un canal montado sin conocimiento del equipo era un incumplimiento, y un incumplimiento era una palanca. Fue hasta el armario de documentación y sacó el expediente de catorce pestañas, el que Ansermet le había entregado en enero y él había desarmado durante tres semanas hasta producir once objeciones, todas aceptadas. Pestaña nueve: infraestructura de datos del laboratorio. Tercera hoja. «Replicación de datos hacia custodia documental. Sincronización continua. Canal cifrado dedicado. Operador: Fondation Cassiodore.»

En el margen inferior, sus iniciales. Su letra de cadenas de custodia, mayúsculas alineadas, la misma para un resultado técnico y para una concesión personal.

Lo había leído en enero. Lo había leído como redundancia, porque un espejo de datos es buena ingeniería y él objetaba ingeniería mala. Custodia documental: dos palabras que en enero significaban respaldo y que a las tres de la mañana del día catorce significaban lo que siempre habían significado para quien las escribió.

El caudalímetro hizo clic. Otro clic, dieciséis segundos después. Vidal cerró la pestaña nueve, la volvió a abrir, verificó que las iniciales siguieran siendo suyas. Seguían.

A las 8:05 usó la línea fija de la nave. Ansermet atendió al segundo tono.

---Doctor. Buenos días.

---El laboratorio espeja los datos crudos en tiempo real hacia ustedes. Desde el día uno.

---Correcto.

---Yo entendí que el archivo se transfería al cierre de los trabajos.

---Le leo la cláusula octava, si me permite, porque prefiero que discutamos sobre el texto y no sobre el recuerdo del texto. ---Un papel se acomodó del otro lado, sin apuro---. «Los datos crudos, los resultados intermedios y definitivos, y toda anomalía registrada durante los trabajos, quedan en custodia de la Fondation a perpetuidad. El uso académico por parte del experto queda sujeto a autorización previa y escrita, que no será denegada sin causa fundada.» Las tres correcciones son suyas, doctor: la autorización escrita, el plazo de sesenta días, la denegación motivada con derecho a segunda presentación. Las tres regulan su uso del archivo. Ninguna regula el momento en que la custodia comienza.

---La custodia de algo que todavía se está produciendo no es custodia. Es captura.

---Es una lectura posible. Le propongo el ejercicio inverso, que es el que usted aplica a los peritajes ajenos: busque en el texto la palabra que fije la custodia al final de los trabajos. Si la encuentra, léamela, y suspendo la sincronización mientras se dirime la discrepancia. Tengo el contrato delante y tiempo.

Vidal tenía su ejemplar sobre la mesa desde las siete y media. Diecinueve cláusulas. La octava era la única sin fecha. Lo había asentado en enero, con estas palabras: el único punto sin fecha del contrato. Lo había clasificado como descuido de redacción en un documento que llevaba ninguno.

---No hay tal palabra ---dijo.

---No hay tal palabra. Entonces no hay incumplimiento que yo pueda subsanar ni usted invocar. Hay cumplimiento. Comprendo que eso no mejora su mañana; mi función tampoco incluye mejorarla. Sí incluye lo siguiente, y se lo ofrezco: un inventario certificado de todo lo recibido por la custodia hasta la fecha, archivo por archivo, con hash y hora. Se lo hago llegar por la vía habitual, si lo quiere.

Un recibo. Se lo estaban ofreciendo con hash y hora, y él lo quería. Lo iba a cotejar línea por línea contra sus manifiestos y las líneas iban a coincidir.

---¿Alguien leyó esos datos?

---La custodia registra recepción e integridad. Sobre lecturas carezco de registro, de instrucciones y de base para afirmar o negar. Le señalo, porque es verificable, que la clave de su modelo sigue existiendo en un solo lugar, que es su memoria. Lo que la custodia tiene son sus salidas, no su instrumento.

Sus salidas. Once años construyendo la máquina que nadie más podía operar, y resultaba que la máquina nunca había sido el producto. El producto era lo que la máquina escupía, y eso viajaba a Ginebra cada madrugada a las 4:30 mientras el dueño del instrumento dormía tres horas en un catre.

---Mándeme el inventario ---dijo Vidal, y cortó.

Se quedó con la mano sobre el teléfono. En el registro de enero, el que llevaba en papel, había una entrada del cuaderno de tapas grises que ahora podía citar de memoria: acepto no saberlo. Había creído que aceptaba no saber un nombre. Había aceptado bastante más, con la letra pareja, y el contrato era la única pieza de toda la operación que funcionaba sin desviación alguna: ni un gramo fuera de banda, ni un segundo de retardo, cumplimiento perfecto, auditable de punta a punta. Lo auditable lo ataba. Él lo había exigido así.

A las nueve, Chiara estaba en el sector textil, desmontando el tendido provisorio de la lámpara de inspección, y él le contó el tráfico del puerto cuatro, la pestaña nueve, la cláusula sin fecha, en el orden correcto y con los números correctos, porque no sabía contarlo de otra manera.

Ella escuchó todo sin interrumpir. Después levantó el cable de la lámpara y empezó a enrollarlo sobre el codo, vuelta por vuelta, pareja.

---¿Ahora sí podés explicar por qué desconfiabas vos?

Vidal abrió la boca con una respuesta de informe ya armada: la cláusula, la fecha ausente, la sobrecarga del tres por ciento. La respuesta describía el mecanismo del espejo. La pregunta era anterior al espejo. Él había desconfiado antes de encontrar el puerto cuatro; había subido a revisar contadores a las 2:40 sin ningún dato que lo justificara, y si alguien le hubiera exigido esa noche el mecanismo de su sospecha habría tenido que entregar lo mismo que ella le entregó a él el día de la firma del protocolo, cuando dijo que el primer pago le olía mal y él le contestó que una desconfianza sin mecanismo se descuenta como ruido.

---No ---dijo.

Chiara terminó el rollo, lo aseguró con el propio cable y lo colgó del gancho. Con las dos manos, aunque el cable era de la nave y pesaba nada.

---Entonces ya sabés cómo se siente ---dijo, y volvió a la lámpara.

Nada más. Ni el reproche, que estaba disponible y era legítimo, ni la pregunta siguiente, que él veía venir y no vino. Conocía esa configuración: un resultado verdadero sin mecanismo que mostrar. La emitía su propio modelo.

Ledda apareció a las 11:20 con el cuaderno de tapa dura bajo el brazo. Vidal le contó el puerto cuatro antes que nada.

---¿Desde qué día? ---preguntó Ledda.

---Desde el uno.

---Eso es de enero, no de la operación. ---Cerró el cuaderno sobre el pulgar, marcando la página---. Nada digital sale sin inventario: mi regla es sobre lo que entra y sale de esta sala. Un canal cableado que ya estaba en la pared cuando llegamos no es una regla que rompí, es una que no supe que hacía falta poner. Lo reviso hoy.

Se paró junto a la mesa de reuniones y esperó a que Vidal y Chiara estuvieran a distancia de voz baja antes de abrir el cuaderno.

---Ceruti amplió el pedido. Nota formal, mesa de entradas, ayer a última hora. Ya no es solo clima: historial completo de todos los sensores de la teca desde la puesta en servicio. ---Pasó el dedo por el renglón subrayado con regla---. Incluye células de carga.

---¿Plazo?

---Ordinario, dos semanas. El volcado lo prepara el mismo técnico de la empresa de la teca, la Comisión visa. El técnico ya tiene el pedido en bandeja.

Vidal hizo la cuenta sin papel. Doce de abril, 3:10 y segundos. Un pico de dos gramos y una décima sobre tolerancia, nueve segundos de subida, retorno a banda al segundo once, estabilización al doce. Doce segundos escritos en una telemetría continua y auditada, sin corrección posible porque una corrección dejaba firma. Él mismo lo había decidido así, esa madrugada, y la decisión seguía siendo correcta: era la única decisión correcta disponible y ahora venía a cobrarse. En un historial de dos años, doce segundos eran una aguja. Pero Ceruti no iba a leer dos años: iba a leer las noches del traslado, porque para eso pedía el historial, y en las noches del traslado la aguja estaba a las 3:10 de la madrugada del doce, con nueve personas habilitadas en la sacristía y una hoja de servicio firmada por un técnico de Ambienti Controllati cuyo nombre real figuraba en el anexo de personal autorizado.

Chiara había vuelto al sector textil. Ledda cerró el cuaderno y lo alineó con el borde de la mesa.

---¿Cuánto tarda un visado de la Comisión si alguien lo empuja? ---preguntó Vidal.

---Depende de quién empuje ---dijo Ledda---. Averiguo hoy quién empuja.

Vidal abrió el registro para asentar la novedad de Ceruti y se detuvo antes del campo de observaciones. Para la respuesta hacía falta Tamiz y hacía falta él. Para lo que se iba cada madrugada a las 4:30 hacía falta solamente que él trabajara: que midiera, que consolidara, que atacara sus propios resultados hasta dejarlos inatacables. Dejó el campo en blanco y no le puso círculo.

\clearpage

\chapter{Inventario}

\lettrine[lines=2, lhang=0.1, nindent=0.2em]{A}{ la una y veinte de la madrugada había luz en el taller textil. Vidal la vio desde la esclusa, al volver de revisar una cola de controles que ya estaba revisada, y cruzó la nave contando los clics del caudalímetro de nitrógeno. Llegó a la puerta en cuatro.}


Chiara estaba sentada a la mesa de trabajo con el estereoscopio apagado y la lámpara articulada en su ángulo más bajo. A su derecha, en fila, cinco portamuestras del corpus de referencia, los de Judea, con los códigos hacia arriba. A su izquierda, dos cuadernos abiertos. Papel cuadriculado, tinta, esquemas de fibra hechos a mano alzada con cotas al margen. La letra tenía ligaduras de otra generación, y en el ángulo superior de la página abierta se leía una fecha: junio de 2002.

Ninguno de esos cuadernos figuraba en el inventario de ingreso. Vidal lo sabía porque el inventario de ingreso lo había firmado él.

---¿No dormís? ---dijo Chiara sin levantar la vista.

---La cola de controles corre sola hasta las seis.

---Entonces no dormís.

Ella cerró los cuadernos con las dos manos, uno sobre el otro, y los guardó en su bolso de lona antes de que la pregunta de él terminara de formarse. El gesto fue el mismo con que tomaba cualquier objeto ajeno, y eso era lo que no cerraba: eran de ella. O estaban con ella. Vidal miró los portamuestras de Judea. Los contó. Cinco a la vista, nueve en el corpus. Los otros cuatro estaban en su caja, en orden. Verificó el orden aunque el orden era evidente.

---Es papel mío ---dijo Chiara---. De antes de esto.

Podía preguntar de quién era la letra. Podía preguntar qué comparaba una conservadora, a la una de la madrugada, entre un corpus de lino arqueológico con procedencia documentada y un cuaderno de campo de hacía treinta años. Las dos preguntas eran legítimas y las dos iban a rebotar, y él había aprendido en cuatro semanas que Chiara contestaba las preguntas simples y esquivaba las compuestas. Se quedó parado junto a la mesa con las manos vacías, que era su versión de esperar.

Ferrero apareció en la puerta con una carpeta bajo el brazo. Se había quedado hasta tarde redactando la ampliación de su disconformidad sobre la traza, para el archivo, y venía a dejarle a Chiara las placas del cuadrante inferior que ella había pedido a la mañana. Las apoyó en la punta de la mesa, alineadas con el borde, y ya se iba.

Chiara había empezado a enrollar un resto de hilo de la canilla de repuesto. Despacio, vuelta por vuelta, sin mirarlo.

---Ottavio.

Ferrero se detuvo con la mano en el marco.

---Los tres contenedores del cierre de 2002 ---dijo ella---. El refuerzo dorsal. Los hilos carbonizados. El polvo del aspirado. ¿Dónde están?

El taller quedó en un silencio que el caudalímetro cortó dos veces. Ferrero se metió la mano en el bolsillo derecho del saco y sacó los guantes de algodón. Se los puso ahí, de pie, con la meticulosidad de quien se prepara para manipular algo que puede romperse, aunque en la mesa lo único frágil ya estaba guardado en un bolso de lona.

---Existieron ---dijo---. Los catalogué yo. Contenedor cuarenta y uno: el refuerzo dorsal, descosido entero, enrollado sobre alma de tubo. Cuarenta y dos: los hilos carbonizados de las zonas de quemadura, sesenta y tres hilos, cada uno en tubo individual. Cuarenta y tres: el polvo del aspirado del reverso. Veintiún gramos. ---Hizo una pausa y se corrigió, porque su cabeza no le permitía otra cosa---: Veintidós con los portamuestras.

---¿Y después?

---Después firmé el acta de cierre. Catorce de marzo de 2003, sesenta y siete contenedores, categoría ``material de conservación en curso'', que no se publica hasta que la conservación concluye. Mi firma es la tercera de cinco. ---Se miró las manos enguantadas---. En el inventario de verificación del año siguiente conté sesenta y cuatro.

Chiara terminó de enrollar el hilo. Lo dejó sobre la mesa, un anillo perfecto.

---¿Y qué hiciste?

---Los reclasifiqué. ``En tránsito de estudio.'' Es una categoría que existe para no decir nada, y yo sabía que existía para eso cuando la usé. ---La voz de Ferrero no subió ni bajó---. Después callé. Veintinueve años.

---¿Por qué?

Vidal esperó una de las oraciones largas de Ferrero, una de esas subordinadas que desembocaban en una factura. Llegó una sola palabra.

---Consecuencias.

Y ahí se quedó. Otras veces la palabra había sido un argumento, un peaje que Ferrero cobraba en cada reunión de protocolo: si esto sale mal, quién lo paga. Esta vez la dijo mirando la mesa, con los guantes puestos, y no cobraba nada. Pagaba.

---En 2004 la Comisión todavía sangraba por la intervención ---siguió, porque ahora que había empezado el mecanismo corría solo---. Una denuncia de material faltante habría quemado a todos los que salían en la foto de 2002. Y la primera de la fila ya estaba en el piso. Su maestra. Pensé que el silencio la protegía. ---Se tocó la muñeca izquierda, donde hacía veintitrés años no había reloj---. Pensé mal, pero lo pensé. Eso no me lo puede sacar nadie: lo pensé de verdad.

Chiara asintió como se asiente ante un dato de calibración.

---A ella la protegió tanto ---dijo--- que lleva años escribiendo cartas a la Comisión pidiendo que revisaran el inventario. Nadie contestó ninguna. En el ambiente no la acusaron de nada, porque no hacía falta: ella había hecho el aspirado, el polvo del aspirado faltaba, la cuenta la hacía cada uno solo en su casa. La dejaron de invitar a los congresos primero y a los cafés después. Limpia estandartes municipales con la misma técnica con que tocaba el lienzo. Sigue con la sospecha puesta. Sin limpiar. Ella, que limpiaba todo.

Lo contó en un tono que Vidal conocía. Plano, sin adjetivos, con las cifras adelante y las personas atrás, cada frase lijada por años de repetición hasta que ya no levantaba nada al pasar. Era el tono exacto con que él decía, cuando no quedaba más remedio que decirlo: figuraba como tercer autor, firmé sin pedir los datos crudos, la retractación salió en junio con una sola firma. Un informe pericial sobre el propio incendio. Se reconoció en la voz de otro y le resultó un instrumento mal calibrado: la medición era correcta y la escala estaba corrida en un valor que ninguno de los dos declaraba.

Se dio cuenta de que llevaba un rato contando las vueltas del anillo de hilo sobre la mesa. Veintinueve.

---¿Quién los tiene? ---preguntó.

Ferrero levantó la vista.

---Nadie lo sabe. No hay venta documentada. Esas cosas no se venden, doctor: se agradecen. Cambian de manos entre gente que se hace favores y el favor no deja recibo. Yo busqué, a mi manera, los primeros años. El último eslabón documentado es mi propia firma en una categoría que inventé para taparlos. Después, pared.

---Un contenedor sellado no se evapora. Salió de un depósito con acceso restringido. Alguien abrió una puerta, alguien cargó un bulto, alguien lo recibió. Cada uno de esos verbos tiene un sujeto y la mayoría de los sujetos deja rastro, aunque el rastro tenga veinte años. ---Vidal oyó su propia voz tomar la pendiente del peritaje, la enumeración de vectores, y en algún punto de la enumeración apareció, sin que la llamara, la traza. Cuarenta y tres partículas de bermellón en una región de nueve por cinco. Contaminación de manipulaciones previas no registradas: si el material de 2002 había circulado por manos sin inventario, su resultado más limpio quedaba apoyado sobre una historia de custodia rota, y rastrear los contenedores era, técnicamente, blindar la traza. La justificación se armó sola, completa, con su pertinencia y su casillero.

La miró un momento y la descartó. La había construido después de decidir, no antes. Conocía ese orden. Era el orden de las malas firmas.

---Dame lo que tengas ---le dijo a Chiara---. Códigos, fechas, pesos si los hay. Una cadena de custodia trunca sigue siendo una cadena: tiene un último eslabón y alguien lo tocó.

Chiara lo miró por primera vez desde que Ferrero había entrado.

---¿Para qué lo querés?

Vidal abrió la boca para contestar como informe. Verificación de integridad del corpus, descarte de contaminación cruzada, cierre de un asunto abierto. Todo era verdad y todo era falso, y en el medio segundo que tardó en notarlo se le secó la boca.

---Todavía no sé ---dijo.

---Buena respuesta ---dijo ella.

Ferrero recogió su carpeta. En la puerta se detuvo, de espaldas, y habló hacia el pasillo oscuro de la nave.

---Si encuentra el último eslabón, doctor, va a tener que decidir qué hace con mi firma. Está en el camino. No se puede llegar a los contenedores sin pasar por encima de ella. ---Una pausa---. Le pido que no la esquive. Ya la esquivé yo por los dos.

Se fue con los guantes puestos. Los pasos se apagaron hacia la esclusa y el caudalímetro hizo tres clics, irregulares, entre once y diecinueve segundos, como siempre.

Chiara sacó del bolso uno de los cuadernos. Lo abrió sobre la mesa, buscó una página con el índice y copió a mano, en una hoja suelta, tres líneas: número de contenedor, contenido, peso de catálogo. Cuarenta y uno, cuarenta y dos, cuarenta y tres. Al cuarenta y tres le agregó, entre paréntesis, las dos cifras: veintiuno, veintidós con portamuestras. Después cerró el cuaderno y lo devolvió al bolso; los cuadernos se quedaban con ella, eso no se discutió porque no hacía falta discutirlo.

---Ella anotaba así ---dijo Chiara, empujando la hoja por la mesa---. Primero el daño, después la causa. Decía que al revés se mide lo que uno quiere encontrar.

La hoja quedó bajo la lámpara. Tres renglones de una letra redonda, sin ligaduras, la letra de la discípula copiando el catálogo de la maestra. Vidal había firmado unas cuatrocientas cadenas de custodia y ninguna le había llegado por esta vía: sin membrete, sin sello, sin más respaldo que dos personas destruidas por no poder probar lo que sabían.

La tomó con las dos manos.

Se dio cuenta después, camino a la esclusa, con la hoja ya en el bolsillo de la camisa: las dos manos, como quien agarra algo de alguien que no está. Se detuvo en mitad de la nave, a oscuras, e hizo la cuenta de cuándo había adquirido ese gesto. La cuenta no le dio ningún resultado, y por primera vez en mucho tiempo dejó una medición abierta y siguió caminando.

\clearpage

\chapter{Presupuesto de fibra}

\lettrine[lines=2, lhang=0.1, nindent=0.2em]{E}{l} remanente cabía en dos viales y en una sola columna de suma.


Vidal la hizo a las 7:40 del día quince, en papel, con la letra de las cadenas de custodia. Alícuota B, sellada desde el día siete: un miligramo doscientos. Las tres fibras del despliegue, promediadas y selladas en la caja de dos cerraduras: cuatrocientos veintidós microgramos entre las tres. Ya en custodia, listo para radiocarbono y química de degradación: un miligramo seiscientos veintidós. Para las otras dos vías hacía falta fibra fresca, del orillo, y esa parte del presupuesto todavía no existía: había que pedirla. Sumadas las dos partidas, la solicitud completa se llevaba todo el lino de esa tela que iba a existir fuera de la teca en lo que quedaba del siglo.

Quedaban seis días de operación hasta el desmontaje y tres con la tela en la mesa. El veintinueve, la restitución. El cuatro, la ostensión, y después del cuatro la teca no volvía a abrirse para nadie de esa sala, ni para sus hijos si los hubieran tenido. El presupuesto era ese: una mano de fibra y una ventana que se cerraba con calendario litúrgico.

Diseñó la batería en dos horas. Datación por radiocarbono en microescala, tres fracciones con pretratamientos independientes, para que una contaminación no pudiera esconderse en las tres a la vez: consumía la alícuota B entera. Química de degradación en paralelo, como testigo interno, sobre las tres fibras del despliegue. Para las otras dos vías necesitaba fibra que todavía estaba en la tela: cinética de despolimerización de la celulosa, calibrada contra el corpus de las treinta y una muestras, que era pobre y era el mejor del mundo, y perfil isotópico de estroncio y oxígeno, porque el agua donde se creció ese lino había dejado firma, y la firma tenía cuenca. Calculó el plan de extracción del orillo en un miligramo seiscientos veintidós, el mismo margen que ya tenía aprobado para el resto, por simetría de presupuesto más que por necesidad exacta. Cuatro vías. Si convergían, el intervalo se cerraba con un margen que ningún colega honesto podría reabrir. Si el mundo entregaba sus mecanismos en alguna parte, los entregaba ahí: cuatro relojes distintos leyendo la misma fibra, y la fibra sin manera de mentirles a los cuatro con la misma mentira.

Costo: el remanente completo. Sin reserva, sin segunda alícuota, sin tercera oportunidad en este siglo.

Escribió esa frase en el borrador del procedimiento, la miró, y la dejó, porque era exacta.

A las nueve los tenía a los dos en la sala de reuniones. Ferrero llegó con su carpeta y los guantes de algodón asomando del bolsillo; Chiara venía del taller, con hilo de embalar todavía en la muñeca, porque el reembalaje del día quince corría en paralelo y ella iba y venía entre las dos mesas. Vidal repartió el procedimiento. Dos hojas. Ferrero leyó la primera despacio y la segunda dos veces.

---Repáseme el saldo, doctor ---dijo---. Después de esto, ¿cuánta muestra queda en custodia?

---Cero.

---Cero. ---Ferrero dejó las hojas alineadas sobre la mesa---. Usted está por gastar el cuerpo de otra persona, entero, hasta el último microgramo, en una pregunta que ya le contestó tres veces que no. El carbono del ochenta y ocho dató un zurcido: no. La imagen: fuera de corpus, no. Su número inauditable: no, en un idioma que ni su máquina habla. Y la respuesta de usted a tres negativas es vaciar la caja. Si la cuarta también es no, ¿qué queda? Ni muestra, ni intervalo, ni excusa. ¿Quién paga eso? Porque el que firma acá no es el que lo paga.

---No exactamente. Las tres negativas son sobre la imagen y sobre la convergencia. La datación del soporte nunca se cerró: se acotó. Esto no es repetir la pregunta, es hacerla bien por primera vez desde 1988.

---Eso mismo dijeron en 1988.

Vidal alineó su lapicera con el borde de la hoja. La volvió a alinear.

---Puede ser ---dijo---. Pero no vine hasta acá para irme con un intervalo abierto.

Era la única frase honesta que tenía y la escuchó salir como se escucha un dato: verificable, incómodo, suyo. Ferrero se quedó mirándolo un momento largo, con las manos quietas sobre la carpeta.

Chiara había empezado a enrollar el hilo de embalar que le colgaba de la muñeca. Vuelta sobre vuelta, prolijo, sin mirar a ninguno de los dos. Vidal ya sabía leer eso. Esperó.

---Voto que sí ---dijo ella---. Con una condición. La toma la hago yo. Del borde, fibra por fibra, con lupa, al ritmo que la tela aguante. Y si una fibra me dice que no, es no. Esa fibra se queda.

---¿Qué llamás que una fibra te diga que no?

---Que está trabajando. Que sostiene algo. Vos lo ves en la tensión, no en un número. ---Terminó el rollo de hilo y lo apoyó en la mesa---. Si querés te lo escribo en el protocolo, pero no te lo puedo poner en microgramos.

Ferrero habló antes de que Vidal encontrara la objeción.

---Acepto esa versión ---dijo---. Si la mano que decide es la de ella, firmo. Si es un cronograma, no firmo. ---Se puso de pie, recogió la carpeta y agregó, ya sin mirar a Vidal---: Usted también la prefiere, doctor. Ya aprobó el criterio antes de encontrarle el defecto. Le vi la cara.

Se fue hacia la esclusa. Vidal bajó la vista al procedimiento y escribió, en el campo de condiciones de detención, una línea que su sistema no admitía: \textit{criterio de detención de la toma: juicio de la responsable de manipulación, sin parámetro, inapelable.} En once años ningún protocolo suyo había llevado una cláusula sin umbral. Firmó al margen, con la letra pareja de siempre, y le llamó la atención que la letra no cambiara para eso tampoco.

La toma empezó a las 11:40 en la sala blanca, con el caudalímetro haciendo sus clics de once a diecinueve segundos y el barrido de reconocimiento apagado. Chiara trabajó sentada, la cara a diez centímetros de la lupa binocular, la pinza de puntas de teflón en la derecha. Vidal sostuvo la lámpara. Fue la función que le tocó: la única tarea del laboratorio que no dejaba firma, y la ejecutó de pie, cuarenta minutos, moviendo el haz dos grados cuando ella lo pedía con el mentón.

Del orillo salieron las fibras de a una. Chiara las nombraba antes de cortarlas: suelta, suelta, esta viene sola. Cada una al vial, cada vial a la balanza, cada pesada dos veces. El aire de la sala olía a nitrógeno y a nada, y lo único que se escuchaba entre clic y clic era la respiración de ella, corta, contada, la respiración de quien trabaja apoyando el pulso contra la mesa.

En la fibra catorce se detuvo.

---Esta no.

Vidal miró por el segundo ocular. Una fibra como las otras, en el límite de la zona muerta del orillo.

---El diámetro está dentro de rango ---dijo.

---Está cargada. Si la saco, el hilo se corre y el tirón llega a la trama. ---Retiró la pinza---. Te dije las condiciones.

Rechazó esa y una más, la diecinueve. Al final la balanza cerró en un miligramo cuatrocientos ochenta: ciento cuarenta y dos microgramos menos que el plan. Vidal rehízo el cálculo de propagación ahí mismo, de pie, en el margen del procedimiento. El intervalo se ensanchaba dieciocho años en el peor escenario. Dieciocho años era nada contra los doce siglos que había que separar. Anotó el ajuste, lo inicialó, y cuando se dio vuelta para pedir la segunda llave de la caja, Ferrero ya la tenía puesta en la cerradura, girada, con los guantes de algodón en las manos y la vista en otra parte.

Los viales salieron de custodia a las 13:55. A las 15:10 estaban fraccionados. Vidal cargó la cola de instrumentos con las cuatro vías intercaladas con controles ciegos, verificó la calibración del espectrómetro, la verificó de nuevo, y a las 17:05 apretó la tecla que ya no tenía marcha atrás. La pantalla devolvió una estimación: cuarenta horas de instrumento. Consolidación final: día diecisiete, 9:10.

En la pared del fondo, esa tarde, apareció la batería escrita en el plan de las once hojas, con la letra de Ledda, en una celda igual a las otras.

La primera noche Vidal se quedó junto a la centrífuga. El aparato hacía ciclos de dieciséis minutos: aceleración, meseta, frenado, y en el frenado el tono bajaba por la escala como algo que se rinde despacio. Contó los ciclos. A las once llevaba veintidós. Chiara pasó a la una, de camino al taller donde el rodillo de reembalaje esperaba su turno, y le dejó un termo sin preguntar nada; Ferrero se había ido a las diez, con los guantes puestos, después de mirar la cola de procesamiento por encima del hombro de Vidal y decir una sola cosa:

---Reza rápido, su máquina.

Vidal tardó un segundo de más en clasificar la frase y la dejó sin contestar.

El día dieciséis pasó en dos mesas: en una, Chiara plegaba el lienzo para el tubo con la asistente, ensayo tras ensayo, dos fibras rotas por metro, una, cero; en la otra, las fracciones giraban y se calentaban y entregaban sus curvas parciales a una cola que Vidal se prohibió abrir. Los parciales sesgaban. Un dato a medias era peor que ninguno: uno empezaba a defenderlo antes de tenerlo. Cerró la ventana de monitoreo tres veces. Las tres veces la había abierto sin registrar la decisión de abrirla.

Esa noche tampoco usó el catre. Se sentó donde la centrífuga quedaba a la izquierda y los racks de Tamiz al frente, con sus veinticuatro conectores en cuatro columnas de seis, y contó. Ciclos de centrífuga, clics de caudalímetro, las horas restantes de a quince minutos. Alrededor de las cuatro perdió la serie de los ciclos: iba por el ochenta y uno o el ochenta y tres y de golpe iba por ninguno. Se quedó quieto. La boca seca. Después empezó de nuevo desde uno, sin anotar la interrupción, y la parte de él que llevaba el registro de sus propios registros tomó nota de que era la primera vez que hacía eso.

A las 6:40 pasó un tren de carga por detrás del polígono y la nave vibró medio milímetro en las juntas del piso flotante. El acelerómetro de la mesa lo iba a mostrar en el volcado de la mañana: un pico chico, inofensivo, con su hora exacta. Todo lo de esa sala tenía hora exacta. Todo era trazable, todo auditable, todo estaba haciendo lo que él le había pedido durante once años al mundo entero que hiciera. En algún lugar de las curvas que todavía giraban estaba el número que cerraba la pregunta para siempre, y él lo iba a leer con las mismas herramientas con que había desarmado el díptico de Amberes, y el número iba a aguantar, porque los números que convergen por cuatro vías aguantan.

Eso era lo que el método prometía. Él había construido su vida sobre esa promesa como otros la construían sobre otras.

A las 7:20 llegó Chiara con pan y dos cafés del bar del polígono. Le puso uno adelante, sobre la mesa, lejos del papel.

---¿Cuánto falta?

---Una hora cincuenta.

Ella se sentó del otro lado, sin decir más, y se puso a rehacer el atado del pelo con un hilo de algodón del taller. El café humeaba entre los dos. En la pantalla, la barra de consolidación avanzaba sin apuro, indiferente a la escala de lo que consolidaba.

Vidal envolvió el vaso con las dos manos y se dispuso a esperar la hora cincuenta de a un minuto por vez.

\clearpage

\chapter{Inescrutable}

\lettrine[lines=2, lhang=0.1, nindent=0.2em]{L}{a} barra llegó al final a las 9:10 y la pantalla pasó del gris de espera al formato de consolidado sin ninguna marca que registrara la escala de lo que cerraba. Vidal dejó el café a un costado y leyó en el orden de siempre, del canal más duro al más blando.


Radiocarbono primero. La fracción de celulosa pura devolvía 45--210 de la era común, con química de pretratamiento limpia, blancos dentro de banda, curva de combustión sin anomalías. La fracción de extracción alcalina devolvía 620--890, con la misma limpieza. La fracción sin pretratamiento, 1030--1310. Tres dataciones de laboratorio impecables, cada una defendible ante cualquier comité del mundo, cada una incompatible con las otras dos por márgenes que ningún error de procedimiento explicaba, porque el procedimiento venía adjunto y era perfecto.

Cada fracción era una muestra intachable. Las tres juntas eran imposibles.

Siguió con la cinética de despolimerización. Las tres fibras del despliegue, medidas contra las treinta y una muestras del corpus, caían en tres puntos distintos de la curva: una se comportaba como lino con novecientos años de envejecimiento equivalente, otra como lino con dos mil doscientos, y la tercera quedaba fuera de la nube entera, por arriba, en una región donde el corpus tenía exactamente cero miembros. El detalle de canal lo decía con precisión de forense: historia térmica no homogénea a escala subcentimétrica. Una fibra había vivido a cuarenta centímetros de una quemadura de 1532, plegada en cuarenta y ocho capas bajo plata a novecientos sesenta grados; su vecina había vivido bajo una mancha de agua que arrastró y depositó lo que el agua arrastra y deposita; una tercera había pasado seis siglos siendo tocada, plegada, exhibida, cosida, descosida y aspirada. La tela tenía cuatro metros cuarenta de largo y en términos de degradación era un archipiélago: cada centímetro con su propio siglo.

El perfil isotópico cerró la lectura. Estroncio y oxígeno compatibles con la cuenca del Mediterráneo oriental, banda ancha. El único canal que no contradecía a los demás, porque era el que menos afirmaba.

Abajo de todo, la capa de convergencia. Campo de clasificación: \textit{muestra inescrutable}. Campo de detalle: \textit{no existe población de referencia}.

Corrió el diagnóstico de integridad sobre la consolidación. Limpio. Revisó los cinco controles ciegos que habían corrido intercalados en la misma cola, con los mismos instrumentos, la misma noche: los cinco dentro de intervalo, incluido un lino copto del siglo IV que el sistema clavó con sesenta años de margen. El espectrómetro estaba sano. La cola estaba sana. El método, que era lo único que él había traído de suyo a este edificio, funcionaba mejor que nunca.

Volvió a la frase y esta vez la leyó entera, con el peso técnico de cada palabra. \textit{No existe población de referencia.} Un intervalo de datación es una comparación: la muestra contra una población de historias compartidas. Para que exista el intervalo tiene que existir la población. Esta tela había dejado de pertenecer a cualquier población posible en algún momento entre el incendio, el agua, los remiendos y los seiscientos años de manos; su biografía la había sacado del conjunto de las cosas comparables. El problema no era falta de datos. Tenían más datos sobre este objeto que sobre ningún textil en la historia de la disciplina. El problema era exceso de biografía: cada dato era cierto y el conjunto no describía nada, porque ya no había un ``algo'' del cual el conjunto fuera la descripción.

Vidal hizo la cuenta que le quedaba por hacer y la hizo despacio, porque era la última. Un método futuro, mejor que el suyo, con resolución diez veces mayor, iba a necesitar lo mismo: una población de referencia. La población no iba a existir en cincuenta años ni en quinientos, porque su ausencia era una propiedad del objeto y las propiedades del objeto ya estaban fijadas. La respuesta definitiva no estaba aplazada. Nunca había existido. Él había pasado un año persiguiendo un número que el universo no tenía archivado en ninguna parte.

Y sobre ese naufragio flotaba una sola cosa: el 0,97. Todos sus mecanismos habían devuelto números que no cerraban; el único número que cerraba había llegado sin mecanismo. Un número sin mecanismo sobre un mar de mecanismos sin número. Se quedó mirando la pantalla y se le ocurrió que esa frase pertenecía al cuaderno de tapas grises, que estaba en Zúrich, en un cajón, debajo de los manuales de calibración.

Ferrero había llegado a las nueve menos diez y llevaba veinte minutos parado detrás de él, leyendo por encima de su hombro con el permiso tácito de las últimas semanas. Cuando la pantalla mostró la clasificación, Vidal lo oyó buscar en el bolsillo del saco. El roce del algodón. Los guantes.

Después, nada. Ningún ruido. Vidal giró la silla para pedirle la referencia de un valor de emisividad y lo encontró llorando.

Sin sonido, sin cambio en la respiración, con la cara quieta de un hombre que espera un tren. Las lágrimas le bajaban por los surcos de la cara y caían al piso de la sala de datos, y él las dejaba caer, las manos juntas adelante, enguantadas, porque secarse la cara con guantes de conservación era algo que Ferrero no iba a hacer ni ese día.

---Está leyendo corrido, doctor ---dijo, con la voz perfectamente firme---. El seiscientos veinte es la fracción alcalina. La exportación corrió la columna.

Tenía razón. Vidal verificó el mapeo, corrigió las etiquetas, y la corrección movió el desastre dos renglones y lo dejó intacto. Lo asentó igual, porque un dato bien rotulado era un dato bien rotulado.

---En el ochenta y ocho ---dijo Ferrero, y se detuvo. Se acomodó los guantes, dedo por dedo, como quien gana tiempo con un procedimiento---. En el ochenta y ocho el resultado se anunció en conferencia de prensa. Tres laboratorios, un intervalo, mil doscientos sesenta a mil trescientos noventa, y un cardenal leyendo el papel con la cara de quien lee un parte médico. Yo estaba en Turín, en la facultad, oyendo a gente discutir la calibración antes de terminar de escuchar. Mi hermano, menor que yo, estaba en la casa de mis padres, con una radio.

Chiara, que había vuelto de la sala blanca con los viales del reembalaje, se quedó junto a la puerta sin entrar del todo.

---Era seminarista ---siguió Ferrero---. Un intervalo de confianza es un objeto químicamente inerte, doctor; puede quedarse en un papel cien años sin reaccionar con nada. Se vuelve activo cuando alguien lo respira solo. Él lo respiró solo, en un cuarto, por radio, y yo tardé tres días en llamarlo, porque estaba ocupado con la calibración, que dicho sea de paso era criticable, y eso también me ocupó. Cuando llamé, ya contestaba distinto. Tardó tres años en terminar de irse: dejó el seminario, dejó la Iglesia, dejó de contestar llamadas. Murió más tarde, a los treinta y uno, en un accidente que la familia nunca quiso mirar de cerca.

Se tocó la muñeca izquierda por encima del guante. Vidal conocía el gesto; hasta ese momento era un dato de conducta con frecuencia registrable y causa sin asignar.

---Llevé su reloj veinte años. Lo guardé en 2010 porque una mañana me descubrí mirándolo para saber la hora, y me pareció una falta. ---Hizo una pausa que duró lo que un ciclo del caudalímetro---. Usted habrá hecho su tabla de mis obstrucciones. El equipo de fluorescencia que hice bajar del camión. Las franjas que quise recortarle. Los miligramos que le discutí uno por uno como un prestamista. Póngale ahora la columna de motivo, si le sirve: yo no vine a impedir la ciencia, doctor. Vine a estar al lado del anuncio.

Vidal repasó la serie completa, entrada por entrada, y las entradas cambiaron de columna sin que él tocara nada. El generador que volvió al proveedor. La propuesta de bajar el solapamiento. La firma en contra, exigida por escrito, en cada discrepancia. Dosificación, no sabotaje. Un hombre midiendo el caudal de una noticia para que nadie la respirara solo, sin poder decir que eso hacía, porque decirlo habría exigido decir el resto.

---Y hoy no hay anuncio ---dijo Vidal.

---Hoy no hay anuncio. ---Ferrero miró la pantalla con los ojos todavía mojados y algo que en otra cara habría sido el principio de una sonrisa---. Cuarenta años esperando el resultado y el resultado es que el Señor tiene mejor abogado que todos nosotros. Inescrutable. Fíjese qué palabra fue a elegir su máquina.

---El campo lo genera una plantilla ---dijo Vidal---. La palabra viene de la taxonomía de fallos.

---Ya sé de dónde viene, doctor. Le estoy diciendo adónde llega.

Chiara terminó de entrar. Dejó los viales en la mesa ---con las dos manos, como todo lo que no era suyo--- y se quedó parada frente a la pantalla el tiempo de leer dos veces. Después agarró el cable de la lupa binocular, que ya estaba enrollado, y lo desenrolló y lo volvió a enrollar.

---¿Esto es un no, o es un no sé?

---Ninguna de las dos ---dijo Vidal---. Un ``no sé'' admite mejora: más datos, mejor método, y el intervalo se estrecha. Acá el intervalo no se estrecha con nada. Es un ``no se va a saber''. Con esta tela. Con cualquier instrumento que se construya después de nosotros.

---¿Y si medís de nuevo?

---¿Con qué? ---Señaló la línea del remanente con el mentón---. Tres miligramos ciento dos, consumidos en cuatro vías. Las cuatro funcionaron. Ese es el punto que no estoy logrando que se entienda: funcionó todo.

Chiara colgó el cable en el gancho. Miró a Ferrero, después a él.

---Entonces terminamos.

---El análisis terminó ---dijo Vidal---. Queda la devolución.

Era una frase de cronograma y la dijo como frase de cronograma, y sin embargo la oyó sonar distinto en la sala. Quedaba una tela de cuatro metros cuarenta que había que plegar sin quebrar más de dos fibras por metro, un tubo que purgar a 0,2 por ciento de humedad, una ventana de seis minutos el veintinueve, y un equipo de cuatro personas que a partir de las 9:10 de esa mañana ya no trabajaba para una respuesta. Trabajaba para una restitución. Lo que él tenía en las manos, hecho el balance del año: un procedimiento pendiente, una tela que era de otros, y tres personas que sabían hacerlo.

Ferrero se quitó los guantes con cuidado, los dobló y se fue a la sala blanca a empezar el acta de estado del reembalaje. Chiara lo siguió con los viales. Vidal se quedó ordenando la consolidación para el archivo, etiquetando cada canal con su desastre correspondiente, y estaba en eso cuando Chiara asomó de nuevo por la puerta.

---¿Y a Ginebra quién le avisa?

Vidal abrió la boca para contestar ``nadie'' y la cerró antes de emitir el dato, porque el dato estaba mal formulado. Nadie tenía que avisar. El consolidado de las 9:10 ya estaba en la cola de replicación del puerto cuatro. Esa madrugada, a las 4:30, la ventana fija iba a abrirse como todas las madrugadas, el manifiesto iba a listar los archivos con su hash y su hora, y el paquete entero ---las tres fracciones que se contradecían, la nube sin miembros, la palabra de la taxonomía de fallos--- iba a cruzar a la custodia documental de la Fondation Cassiodore, verificado byte a byte, íntegro, perfecto.

---Ya está avisado ---dijo---. Se avisa solo.

Ella lo miró un segundo de más y volvió a la sala blanca. Vidal se quedó frente al conmutador del taller, donde el diodo del puerto cuatro parpadeaba con el tráfico de mantenimiento, y por primera vez armó el pensamiento con todas las letras, sin abreviarlo, como una línea de registro que nadie le pedía.

Alguien había gastado una fortuna en catorce pestañas, un telar belga, un circuito de billetes y el mejor perito vivo, y esa madrugada iba a recibir, con firma criptográfica y cadena de custodia intacta, la noticia de que la pregunta no tenía respuesta y no iba a tenerla nunca. Y si la operación entera había sido diseñada por un hombre que jamás decía nada verificablemente falso, había que auditar también esta hipótesis: que el paquete de las 4:30 llegara a Ginebra exactamente como se esperaba. Que la no-respuesta, certificada, sellada, inapelable, fuera el producto.

\clearpage

\chapter{Ruido, otra vez}

\lettrine[lines=2, lhang=0.1, nindent=0.2em]{L}{a} hoja del día apareció bajo la puerta a las siete y cinco, mecanografiada como todas: sin columna de función, solo los apellidos en el orden de siempre. Ferrero, O.; Fabbri; Ledda. Vidal la leyó de pie en el pasillo, descalzo sobre el parquet frío, y la leyó una segunda vez buscando el renglón que faltaba. Su apellido había estado en cada hoja desde el nueve de abril, ``Vidal, S.'', al final del orden. Hoy la hoja tenía tres nombres y ninguno era el suyo.


La dejó sobre la mesa de la cocina, junto a la moka de anoche, y abrió su planilla. El bloque del día estaba en blanco. El catálogo ofrecía sus códigos de siempre ---consolidación, redacción de informe, mantenimiento de instrumental, hasta reconocimiento urbano, el que había inventado después de Atenas---, y cualquiera de ellos habría cerrado la celda en cuatro segundos. Sostuvo la lapicera encima del papel el tiempo suficiente para que la tinta amenazara con caer sola. Después la tapó. La última celda vacía de su vida era de octubre, un viernes en Jerusalén, de dos a ocho de la tarde. Esta era la segunda.

Del otro lado de la persiana pasó un tren de maniobras, largo, hacia Dora. El departamento olía a café frío y a madera vieja. En Collegno, a esta hora, Chiara ya habría abierto la sala blanca. Vidal se sirvió el café de anoche sin calentarlo y se lo tomó de pie, mirando las franjas de luz que la persiana recortaba sobre el piso.

A las nueve y veinte trajo la notebook a la mesa de la cocina y escribió, de memoria, la clave del contenedor 2029-03. La primera vez le rebotó; había metido un carácter de más. La segunda abrió. El respaldo estaba repartido entre sus discos y la reconstrucción tardó once minutos; los pasó sentado, quieto, con las manos a los costados del teclado.

Había abierto ese archivo dos veces en cuatro años, las dos para ir directo a las líneas 3.407 a 3.418 y volver a cerrarlo. Nunca lo había leído desde la línea uno.

Lo leyó desde la línea uno.

La ingesta del panel izquierdo del díptico, canal por canal. Las capas de pintura resueltas en orden estratigráfico, el craquelado, la firma del aglutinante, el corpus documental de tres generaciones de connoisseurs desmontado atribución por atribución. Cuatro mil seiscientas líneas de un trabajo impecable. Y en la 3.407, la discordancia: una firma espectral en la capa de imprimación, sin correlato en el corpus, peso agregado 0,04. El modelo la había visto, la había pesado y la había puesto sobre la mesa sin poder decir qué era. En la línea 3.418, la resolución. \textit{Ruido de sensor, excluir de consolidación.} Campo de operador: SV. Catorce de marzo de 2029, 2:31.

La entrada anterior del log era de las 2:29.

Dos minutos. Eso constaba. A las 2:29 el modelo le había entregado algo que veía sin explicar, y a las 2:31 él lo había clasificado con una causa ---sensor--- que jamás verificó. Ninguna corrida de aislamiento, ninguna repetición con el instrumento en frío. La palabra afirmaba un mecanismo y el mecanismo era una hipótesis sin prueba, elegida a las dos y media de la mañana porque cerraba el expediente. Una palabra higiénica. El panel era falso; seis vías independientes lo confirmaron después; la clasificación resultó inocua. Tener razón había sido una cuestión de dirección, no de método. La operación, en sí, había sido esta: quedarse con la parte del objeto que podía auditar y nombrar el resto con una palabra que lo eximiera de mirarlo.

Conocía esa operación. La había estudiado durante meses, en otro expediente, con fecha de 1988. Tres laboratorios midieron bien siete centímetros de tela y el número nombró el objeto entero. \textit{Medieval.} También esa palabra era higiénica, también cerraba el expediente, también resultó ser la descripción exacta de la parte equivocada. El mundo le había hecho a la tela lo que él le hizo a la línea 3.407. Con la diferencia de que el mundo tardó una conferencia de prensa y él tardó dos minutos.

Se levantó, enjuagó la taza, volvió a sentarse. El archivo seguía abierto.

Peso agregado 0,04. Un número minúsculo. Pero la especie era la misma que la del 0,97 de días atrás: una señal que la arquitectura emitía sin entregar el mecanismo. La suya había sido chica y la borró con una palabra. La de ahora era enorme y llevaba días tratando de abrirla con herramientas que devolvían ruido. Ruido. La palabra otra vez, ahora del otro lado, ahora contra él.

La máquina se había construido para algo preciso y ese algo tenía nombre de función, no de proyecto. Después de los bronces ---después de firmar como tercer autor sin pedir los datos crudos, después del informe holandés de abril y la retractación de junio con una sola firma--- había levantado once años de arquitectura para que confiar dejara de ser un paso del procedimiento. Para que nunca más hubiera un momento en que un hombre, solo, decide darle crédito a otro. Y a las 2:31 de una madrugada de marzo la máquina lo había puesto exactamente ahí: solo, de noche, frente a una señal sin mecanismo, obligado a decidir. Decidió con una palabra. Como cualquiera.

La directora de distrito se lo había dicho en Jerusalén, en una oficina con mapas, y él lo había archivado como retórica de gestión. \textit{El umbral lo firmamos todos, doctor, pero el umbral no decide. En algún punto todos eligen.} Seis meses. Le había llevado seis meses y un lienzo entender que la frase era una descripción técnica.

Cerró el contenedor. Verificó que el cierre no hubiera alterado el hash. Lo verificó de nuevo.


\scenebreak


A la tarde intentó ponerle orden al día por el método de siempre. Las baldosas de la cocina primero: hexagonales, de un gris con vetas, apretadas contra el zócalo. Llegó a treinta y ocho y un tren de maniobras pasó largo hacia Dora y cuando el ruido terminó ya el treinta y ocho podía ser cuarenta o treinta y cinco, la serie estaba muerta. La empezó de nuevo desde el rincón. Se perdió a la altura de la heladera, sin tren, sin causa.

Probó con las tablillas de la persiana. Eran veintitantas; sabía que eran veintitantas; nunca supo cuántas, porque a la quinta o sexta la vista se le iba a las franjas de luz y no a las tablillas, y las franjas se movían.

De noche probó con los autos. Sentado a oscuras, contando los barridos de luz que los faros dibujaban en el techo cada vez que un auto doblaba la esquina. Uno cada tanto, un ritmo cómodo, imposible de perder. Lo perdió al séptimo. O al noveno.

Treinta años contando. Filas, baldosas, segundos de semáforo, pitidos de alarma, ciclos de centrífuga. Siempre lo había registrado como método: contar era verificar, y verificar era lo que él hacía. Pero las baldosas no requerían verificación. La fila del Sepulcro tampoco: cuarenta y una personas eran cuarenta y una personas dijera él lo que dijera. Ninguna de esas series había tenido nunca un uso, un destino, un campo donde asentarse. Números emitidos hacia nada, durante treinta años, con la constancia de un hombre que reza sin haberse enterado.

---Cuarenta y uno ---dijo, en voz alta, a la cocina vacía.

La cocina no hizo nada con el dato.

Se quedó donde estaba. La heladera arrancó, zumbó un rato, se apagó. En algún departamento de arriba corrieron una silla. Pasó un auto y el barrido de luz cruzó el techo y él lo dejó pasar sin número, y eso ---dejarlo pasar sin número--- le produjo una sensación física concreta, un vacío chico y exacto detrás del esternón, como el hueco que queda en una fila de dientes.

A las 21:50 sonó el teléfono fijo. Un timbre, cortado antes del segundo.

Sin novedad. Próxima fase en curso.

La próxima fase era la devolución. Hoy, a la hora que fijara la revisión de clima, seis minutos, cuatro minutos veinte de ejecución medida, un margen del veintiocho por ciento. Vidal se encontró de pie junto al teléfono sin recordar haberse levantado, y la fase pendiente se le acomodó en el pecho con un peso que reconoció: era lo único de ese día que todavía tenía forma.

Ferrero lo había dictado en la nave, meses atrás, con la mano abierta sobre la mesa, cuando discutían miligramos y ventanas. \textit{Todo lo demás del protocolo se puede hacer mal y corregir, doctor. La devolución no tiene versión corregida. Si esa fase falla, ¿quién se lo explica a la gente que va a hacer fila en mayo, y con qué palabras?} Él lo había aprobado entonces como se aprueba una redundancia: la restitución ya prevalecía sobre el análisis por su propia cláusula, firmada en Ginebra antes de conocer a nadie. Esta noche la frase de Ferrero era otra cosa. Esta noche era la única línea de su vida con destinatario.

La tela tenía que volver. Esta misma noche, en seis minutos, a una caja de cristal frente a la cual la gente iba a arrodillarse sin saber que la tela había salido. Eso era todo lo que quedaba en pie del protocolo, del informe, del análisis definitivo, de los diecisiete días. Ninguna respuesta. Una devolución. Un objeto que era de otros y que había que devolver intacto, y tres personas que estaban ahora mismo, del otro lado de la ciudad, preparándolo para el viaje de vuelta.

Fue al cuarto y sacó una camisa gris de la percha. Se vistió sin encender la luz del techo, con la del velador, y se anudó los cordones dos veces cada uno. En la mesa de la cocina quedaron la hoja del día con sus tres renglones, la planilla con la celda en blanco y la notebook cerrada sobre el contenedor cifrado, que volvía a ser lo que había sido cuatro años: un archivo bajo custodia de un hombre que ya sabía qué guardaba.

Agarró las llaves del auto. La escalera del edificio olía a cera vieja y en el segundo piso una radio hablaba bajito detrás de una puerta. Afuera, la calle estaba mojada sin que hubiera llovido, ese brillo de abril que dejan las máquinas de limpieza municipal, y el auto arrancó al primer intento.

Manejó hacia Collegno por la avenida vacía, con la ventanilla apenas abierta, sin contar nada en todo el camino.

\clearpage

\chapter{La tela no sabe}

\lettrine[lines=2, lhang=0.1, nindent=0.2em]{E}{ran} las 22:35 cuando Vidal entró al playón de la nave, con una sola camioneta estacionada, y en la caja de la esclusa, al dejar su teléfono, había un solo teléfono más: el de Chiara, con la funda gastada en las esquinas. La hoja diaria, pegada a la puerta interior como todas, solo tenía los apellidos de siempre, sin columna de función; Vidal sabía igual, por la caja de teléfonos, que Ferrero se había retirado a las nueve y diez y que Ledda estaba en ronda de perímetro hasta las tres. El caudalímetro de nitrógeno hizo clic dos veces mientras esperaba el ciclo de la doble puerta.


La sala blanca estaba a media luz. Sobre la mesa de sujeción, una lámpara rasante barría el lienzo desde el lado izquierdo y Chiara trabajaba inclinada en el tercio central, lupa frontal levantada sobre el pelo, que esa noche llevaba atado con hilo de algodón. Tenía al lado la carpeta del inventario de estado del día dos, abierta en la primera hoja.

Levantó la vista lo justo.

---Hoy la hoja vino con tres nombres ---dijo.

---Ya sé.

---Bueno. ---Señaló la lámpara con el mentón---. Agarrala con las dos manos. Cuarenta grados, desde ahí. Si te cansás, avisá antes de que se te caiga el ángulo, no después.

Eso fue toda la conversación durante la primera hora. Ella avanzaba por la cara frontal en el mismo orden del inventario de ingreso, punto por punto, y dictaba el estado; él sostenía la luz y leía en voz alta la cota de referencia de la carpeta. Quemadura tres, borde interno, a cuarenta y un milímetros del pliegue axial: cuarenta y uno. Conforme. Mancha de agua grande, lado izquierdo, término a siete hilos del pliegue: siete, contra la placa 9.114 de Ferrero. Conforme. Perforaciones del parche superior, muestreo de veinte sobre cuatrocientas ochenta y dos: sin variación. La zona de la traza, agotada, por debajo del umbral, exactamente igual a sí misma. La esquina del zurcido ---Chiara le decía \textit{la furba}, la única palabra italiana que se permitía dentro del registro y que después no anotaba--- con sus treinta y ocho hilos de urdimbre haciéndose pasar por lo que no eran desde hacía siete siglos. Ella escribía «conforme» con letra pareja, línea por línea, y cada tanto le pedía dos milímetros más de ángulo, y él corregía el ángulo, y la luz volvía a quedar donde tenía que estar.

Limpieza pasiva, lo había llamado ella en el protocolo. Nada tocaba la tela. Se documentaba que la tela estaba entera, que los diecisiete días habían pasado por encima sin dejar marca, que lo que volvía era lo que había salido. Un inventario al revés: el mismo papel, las mismas cotas, la firma en la otra columna.

En algún momento Vidal calculó cuánto hacía que sostenía la lámpara y le dio setenta minutos, con un margen amplio, porque llevaba más de una hora sin cronometrar los clics del caudalímetro. Los clics seguían ahí, irregulares entre once y diecinueve segundos. Habían dejado de pedirle que los midiera.

Tenía la boca menos seca que en el departamento. Los hombros más bajos. Registró ambas cosas sin asignarles causa y siguió leyendo cotas.

Cuando el inventario cerró ---ochenta y una líneas, ochenta y una veces «conforme»--- Chiara fue al panel de la mesa y empezó a soltar la depresión de las dieciséis zonas, de afuera hacia adentro, el orden inverso del despliegue. La tela se asentó por sectores con un ruido seco y breve, de papel viejo, ese ruido que el lino de ensayo belga nunca había imitado bien. En la zona once Chiara esperó veinte segundos de más, por el gradiente del caño de clima, sin mirar el protocolo. Después extendió el papel de seda desde la cabecera, media superficie, lo justo hasta el plegado de medianoche, y alisó el borde con el canto de la mano a dos centímetros de la tela, sin tocarla.

Se quedó parada al pie de la mesa un momento.

---¿Sabés qué es lo bueno de las telas? ---dijo.

Vidal apagó la lámpara rasante. La sala quedó en la luz general, más pareja y más pobre.

---Que no saben ---siguió ella---. La cuidás diecisiete días con guantes, le medís la humedad cada quince minutos, y ella no se entera. En unas horas vuelve a una caja de cristal y tampoco se va a enterar. Lo hacés igual.

Fue hasta el banco hiperespectral, desconectó el escáner de mano y empezó a enrollar el cable. Lo hacía con el cuidado de quien guarda algo que le prestaron: vuelta pareja, paso constante, la mano izquierda de guía. Vidal conocía la secuencia. Sabía lo que venía antes de que ella abriera la boca.

---¿Y ahora qué hacés vos? ---dijo Chiara, sin mirarlo.

Era la tercera vez que le hacía esa pregunta. La primera, la primera semana en Collegno, enrollando el mismo cable, él había contestado «publico», y la respuesta ya era falsa cuando la dijo, aunque tardó semanas en poder fecharle la falsedad. La segunda, semanas después, había contestado con la cláusula octava, que era una manera de contestar con la voz de otro.

---Depende de qué llames\ldots{} ---empezó, y la fórmula se le quedó a la mitad, colgando, y la escuchó colgar, y la soltó---. No sé.

Chiara terminó la vuelta que tenía en la mano. Asintió una vez, con el mismo movimiento seco con que aprobaba una pesada dentro de banda, y colgó el cable del gancho.

---Era hora ---dijo. Y después, más bajo, casi para el gancho---: Ecco.

Trajo dos bancos altos de la zona del taller y los puso a un costado de la mesa, lejos de la tela, y sacó de su bolso un termo abollado y una sola taza, y le sirvió a él en la taza y ella tomó del vaso del termo. El café estaba tibio y demasiado dulce. Vidal lo tomó igual.

---Mi maestra limpia estandartes municipales ---dijo Chiara---. Los de las fiestas patronales, los que se cuelgan un día al año y se guardan mojados porque el empleado tiene apuro. Los lava con agua desmineralizada. La compra ella, porque el municipio dice que con la canilla alcanza. ---Giró el vaso entre las manos---. Reza poco, mi maestra. Menos que yo. Una vez le pregunté si todavía creía. Después de todo lo que le hicieron, digo. Me miró así, como si le hubiera preguntado si todavía respiraba, y me dijo: creer es conservar algo que no es tuyo. Después me corrigió la tensión de un bastidor y no tocó más el tema. Nunca más.

Conservar algo que no es tuyo. Vidal dejó la taza sobre el banco, entre las rodillas.

Él tenía un contenedor cifrado en un directorio llamado 2029-03, custodiado cuatro años con una clave que existía solo en su memoria, y hasta esta semana ni siquiera lo había leído entero. Tenía un directorio nuevo, 2033-04, catorce kilobytes verificados byte a byte, con un número adentro que ninguna herramienta de la Tierra iba a explicarle nunca. Y tenía once kilobytes más viejos que todo eso, una planilla en una carpeta local que el sincronizador de las 20:00 salteaba desde octubre, con una entrada de dos horas y veintiocho minutos y una columna que decía tiempo perdido. Tres archivos sin cierre posible. Tres cosas que mantenía ---respaldos, hashes, integridad verificada dos veces--- sin poder terminarlas jamás, sin que ninguna le devolviera nada, sin que ninguna supiera siquiera que la estaban cuidando.

Su sistema entero tenía una palabra para eso y la palabra era archivo. Se le ocurrió, por primera vez en cuarenta y seis años, que a lo mejor la palabra estaba mal elegida.

---¿Y ella sabe que vos\ldots{}? ---empezó, y se corrigió---: ¿Le contaste dónde estás?

---Le dije que tenía un trabajo grande. De lino. ---Chiara se encogió de hombros---. Me dijo que me lave las manos antes y después. Es su bendición. Es la única que da.

Se quedaron un rato sin hablar. El caudalímetro hizo clic, y otro clic, y Vidal los oyó pasar sin ponerles número, y el café siguió demasiado dulce hasta el fondo de la taza. Del otro lado del polígono pasó un tren de carga y el piso flotante lo tradujo a un temblor de medio milímetro que mañana iba a estar en el volcado del acelerómetro, con hora exacta, para nadie.

Chiara bajó del banco. Enjuagó el vaso del termo en la pileta del taller y lo secó, y cuando volvió tenía otra cara, una que él había visto dos veces en diecisiete días y las dos veces antes de una frase difícil.

---Sebastián.

---Sí.

Ella miró hacia la estantería del fondo, donde los portamuestras del corpus de Judea pasaban la noche en su gradilla, y después hacia la caja de dos cerraduras, y después a él.

---Cuando esto termine, hay una cosa que te\ldots{} ---Se detuvo. Se pasó el dorso de la muñeca por la frente, aunque la frente estaba seca---. Que te quiero mostrar. No es de ahora. Es de antes. De\ldots{}

Se calló ahí. Fue hasta el banco hiperespectral y agarró el segundo cable, el del cabezal fijo, que ya estaba enrollado en el gancho desde la tarde, con las vueltas parejas y el paso constante, terminado. Lo descolgó, lo desenrolló hasta la mitad y lo volvió a enrollar desde el principio, vuelta por vuelta, más despacio que la primera vez.

Vidal miró la operación completa. Un cable que ya estaba enrollado. Treinta y cuatro segundos de trabajo sin objeto sobre un objeto sin falla. En su registro eso tenía nombre técnico, valor anómalo, causa no asignada, y tenía procedimiento: aislar la anomalía, formular la pregunta, obtener el dato o clasificar el silencio. Once años de método decían que la pregunta se hacía ahora, con la testigo presente y la memoria fresca. Sabía además, con umbral de confianza alto, de qué lado de la nave apuntaba la respuesta: a la gradilla, a los nueve de Judea, a dos cuadernos de papel cuadriculado que él había visto una madrugada y no constaban en ningún inventario firmado por él.

Dejó la pregunta donde estaba.

Se levantó del banco, fue hasta ella y le sacó el cable de las manos ---con las dos suyas, sin pensarlo, y recién al sentir el peso repartido se dio cuenta de qué manera lo había agarrado--- y lo colgó del gancho, junto al otro.

Chiara lo miró hacer. Miró el gancho, los dos rollos uno al lado del otro, las vueltas de distinto paso. Algo se le aflojó en los hombros, un milímetro, dos.

---A medianoche plegamos ---dijo---. El rodillo lo armás vos. Yo miro.

---A medianoche.

Ella apagó la luz general y quedó solo la piloto de la esclusa, azul, baja. En la mesa, plana bajo el papel de seda, la tela siguió sin enterarse de nada: ni de los diecisiete días, ni del número que había fabricado, ni de la mujer que acababa de despedirla línea por línea, ni del hombre parado a oscuras que por primera vez en su vida adulta acababa de dejar una pregunta sin hacer y la sentía, en el pecho, con la forma exacta de una respuesta.

\clearpage

\chapter{Seis minutos}

\lettrine[lines=2, lhang=0.1, nindent=0.2em]{L}{edda} ya había tachado la celda cuando Vidal entró a la nave.


Eran las siete menos diez y el aire de Collegno estaba a diecisiete grados, con ese olor seco a nitrógeno y cartón que dejaba la sala blanca cuando llevaba horas sin gente adentro. En la pared del fondo, las once hojas del plan seguían pegadas con cinta en las cuatro esquinas. La celda de las 19:00 del veintinueve ---la ventana A, once minutos, cambio de guardia--- tenía ahora dos diagonales de lápiz y una palabra al margen: \textit{anulada}. Ledda estaba de pie frente a la hoja nueve con la regla en la mano, subrayando el renglón de abajo.

---Ventana A, muerta ---dijo---. Aviso interno de la Comisión, circulado ayer a las 18:40. Monseñor Ceruti reprogramó las inspecciones previas a la ostensión. Estaban repartidas en el veintisiete, el veintiocho y el veintinueve. Ahora están las tres el veintinueve.

---¿Motivo declarado?

---Concentrar la presencia de la Comisión en una sola jornada. ---Ledda bajó la regla---. Efecto medible: entre las 18:20 y las 20:00 hay once personas en tránsito por el crucero y la sacristía. El relevo de puestos sigue siendo a las 19:00 y sigue teniendo solapamiento nominal cero. La diferencia es que ahora el hueco está lleno de gente que no forma parte del hueco.

Vidal se acercó a la hoja. El plan tenía cuatrocientas ocho horas y a esa altura quedaban pocas celdas sin marcar. Las contó rápido, celda por celda, y volvió a contarlas, y el número le dio dos veces lo mismo.

---Queda B ---dijo.

---Queda B. Revisión del sistema de clima de la teca. Prevista 16:05 del veintinueve. Duración media medida sobre las nueve intervenciones registradas desde la puesta en servicio: seis minutos doce. Máxima: siete cero cinco. Mínima: cinco cuarenta. La ejecuta un técnico de la empresa de la teca con la Comisión mirando desde el otro lado de la barrera, y durante esos minutos el marco queda abierto y las células de carga en modo de servicio, sin alarma.

Chiara había llegado sin que Vidal la oyera. Traía el termo de metal de la cocina y dos vasos, y agarró la hoja nueve con las dos manos por los bordes de abajo para leerla mejor, aunque estaba pegada a la pared.

---Seis minutos ---dijo---. Y si sale mal, ¿hay segunda pasada?

---No hay segunda pasada. ---Ledda anotó algo en su cuaderno y lo subrayó con la regla---. Hay tres modos de falla. El primero es tiempo: si la revisión se acorta, la fase queda abierta. El segundo es reconocimiento: cara repetida donde no corresponde, y ese día los Carabinieri tienen los extractos de doce meses cruzados con contratistas. El tercero es Ceruti. Pidió el histórico de clima el diecisiete y lo amplió a células de carga el veinticuatro. Está en el edificio a horas que no le corresponden.

Ferrero estaba en la puerta de la esclusa desde hacía un rato, con el saco puesto y los guantes de algodón asomando del bolsillo derecho. Vidal esperó. Había un orden previsible para lo que venía: primero la objeción de principio, después el cálculo de costo, después la pregunta de a quién se lo explica uno. Contó los segundos por costumbre. Llegó a veinte y a los veinte Ferrero se acercó a la mesa de sujeción, que estaba vacía ---cinco metros veinte de vacío, con el tubo de transporte apoyado en sus caballetes al costado, precintado desde la medianoche---, y apoyó una mano en el borde.

---Seis minutos son ciento veinte segundos más de los que hacen falta ---dijo---. Con manos buenas.

---Podríamos pedir un día más ---empezó Vidal---, correr la restitución al treinta y...

Nadie lo miró. Ledda seguía con la regla sobre el renglón; Chiara servía el café; Ferrero tenía la mano quieta en el borde de la mesa vacía. La frase se le fue quedando sin nadie que la sostuviera del otro lado, y la dejó ahí, a la mitad, con las otras cuatro variantes que había preparado en la cabeza mientras manejaba y que ya no hacía falta decir. El caudalímetro de la línea de nitrógeno hizo clic detrás del vidrio, y otra vez, y otra vez, a intervalos que él conocía de memoria: media catorce coma nueve, dispersión de once a diecinueve. Chiara le dio un vaso de café a Ledda.

---¿Cuántas manos entran? ---preguntó.

---Eso lo resolvemos hoy ---dijo Ledda---. A las dos.

Vidal se sentó en el banco de la sala de reuniones con el cuaderno cuadriculado de tapas de cartón naranja que había comprado en el quiosco de la via Cigna la semana anterior, porque el de tapas grises se le había quedado en Zúrich y hacía tres semanas que le faltaba una columna. Escribió la fecha y la hora, como siempre. Después escribió: \textit{Vinimos por una respuesta. Nos vamos con una restitución.} Miró el renglón. Era exacto y le sobraban las dos primeras palabras, así que las dejó.

A las once y veinte imprimió el borrador del informe final. Diecinueve páginas, índice con la palabra «Estados» donde durante once años había ido «Conclusiones», y en la página catorce, tercer bloque, cuatro líneas que había escrito a las cuatro y media de la mañana y reescrito dos veces:

\textit{Índice global de la capa de convergencia: 0,97. Ningún canal individual supera 0,7. La ablación por canal no produce escala intermedia. Las herramientas de interpretación devuelven ruido al aplicarse al índice. Clasificación: señal no auditable. La señal no se descarta como ruido de sistema. Se asienta sin clasificar.}

Ferrero lo leyó de pie, con los guantes puestos, pasando las hojas por el ángulo superior. Tardó nueve minutos en llegar a la catorce. Cuando llegó, la leyó dos veces y dejó el borrador sobre la mesa.

---Esto no se anuncia.

---No es un anuncio. Es un registro.

---Usted escribe un número que nadie puede explicar, lo firma con su nombre, y en algún momento eso sale de esta nave. ---Ferrero se tocó la muñeca izquierda, dos dedos, un segundo---. En el ochenta y ocho también había un renglón. El renglón tenía intervalo, laboratorios y método. Lo que no tenía era nadie del otro lado juntando lo que se rompía. ¿Y a quién se lo explica usted, doctor, cuando el archivo llegue a Ginebra y el renglón viaje solo?

---Depende de qué llames explicar.

---Llamo explicar a estar en la misma habitación que la persona que lo lee.

Vidal apoyó las dos manos abiertas sobre el borrador, una a cada lado, y no las movió.

---El punto diecinueve dice comunicación a terceros ajenos al equipo. Ese es el anuncio. Yo te propongo la otra definición, y la escribimos ahora: registro es el asiento de un dato con fecha, autor y método, en el soporte de papel de la operación, sin destinatario designado.

---Sin destinatario designado ---repitió Ferrero---. Hoy.

---Hoy y en el papel. Lo digital ya no es nuestro: se replica a las cuatro y media y no lo decido yo. Esto ---tocó las diecinueve hojas--- es lo único que queda dentro de la caja de dos cerraduras.

Ferrero se quedó mirando la página catorce un rato largo. Del polígono llegó el ruido del tren de maniobras, ese arrastre bajo que hacía vibrar medio milímetro las juntas del piso flotante y quedaba escrito en el acelerómetro con hora exacta.

---Hay una diferencia entre las dos cosas que usted no puso ---dijo al fin---. Un registro tiene testigo. Si no tiene testigo, es un borrador, y un borrador lo corrige cualquiera después. Escriba la línea a mano, con fecha, y firme. Yo firmo debajo.

Vidal levantó la vista.

---Vos estás en contra de que esa línea exista.

---Estoy en contra. Y quiero que conste que existía y que yo la vi hoy, con esas palabras y no otras. ---Se sacó el guante derecho para agarrar la lapicera, cosa que Vidal no le había visto hacer en cuatro semanas---. Las dos cosas son ciertas, doctor. No hace falta elegir una.

Firmaron a las 12:04. La letra de Ferrero salió normal, sin el grosor de la primera curva que le había salido en la sacristía. En marzo de 2029, a las 2:31 de la madrugada, un campo de operador con las iniciales SV había puesto \textit{ruido de sensor, excluir de consolidación} junto a una discordancia de peso agregado 0,04, sin correr aislamiento y sin repetir con el instrumento en frío. El panel era falso. Seis vías independientes lo confirmaron después. La decisión, igual, había sido suya.

Vidal cerró la carpeta y la puso en la caja. Giró su llave. Ferrero giró la suya sin que se lo pidieran.

A las dos, Ledda desplegó sobre la mesa de la sala de reuniones el plano de la sacristía con las cotas de la pestaña dos y una hoja aparte, mecanografiada, con la lista de personas autorizadas para el veintinueve.

---El intercambio inverso necesita dos pares de manos adentro ---dijo---. Manipulación de tela, y gemela más contenedor. La asistente de conservación del once no está en la lista del veintinueve, así que el sostén del borde y la verificación de posición los hace la misma persona que manipula. La lista la cerró la Comisión el lunes y no admite altas.

---¿Quién ejecuta la revisión de clima? ---preguntó Vidal.

---Un técnico de la empresa de la teca.

---Un técnico de Ambienti Controllati Sociedad de Responsabilidad Limitada, contrato marco del catorce de febrero, alta en el anexo de personal autorizado de la diócesis del treinta y uno de enero, dos facturas emitidas, aportes al día. ---Vidal apoyó el índice sobre el renglón correspondiente de la hoja---. Función declarada: verificación de instrumentos. La ventana no es un hueco en el que nos metemos. La ventana la abre el técnico que hace la revisión, y el técnico tiene nombre y contrato desde el catorce de febrero.

Chiara dejó de enrollar el metro de tela que tenía en la mano.

Ledda no contestó enseguida. Miró la hoja, después a Vidal, y siguió mirándolo el tiempo que le llevaba a él dar un plazo: cuatro, cinco segundos. Vidal se descubrió repasando en la cabeza el par de apriete del cabezal fijo, que había verificado la noche anterior y que estaba verificado.

---Usted entra con su nombre real ---dijo Ledda---. Si alguien lo reconoce, no hay versión alternativa. No existe un doctor Vidal que casualmente sea técnico de clima. Existe el que es.

---Ya entré una vez con mi nombre real.

---Entró de noche, con seis personas y un electricista dormido en la camioneta. El veintinueve hay once en tránsito, dos de la Custodia y un monseñor que pidió las tablas de células de carga. ---Ledda se pasó la regla de una mano a la otra---. Y en el momento en que usted esté midiendo la sonda de humedad va a haber alguien a cuatro metros mirándole las manos.

---Voy a estar midiendo la sonda de humedad.

Ferrero, desde el otro lado de la mesa, dijo algo en dialecto que Vidal no entendió y que Chiara tampoco tradujo. Ledda destapó su lapicera, tachó un renglón de su cuaderno y escribió otro debajo.

---Mañana a las siete ---dijo---. Ensayo en seis, con la maqueta y el marco real. Tres personas y cronómetro. Si a las once no bajamos de seis minutos con margen, se reabre la discusión de la ventana y le aviso a Ginebra en el día. ---Y después, sin cambiar el tono---: Mándeme su talle de overol antes de las cinco.

Fue hasta la pared del fondo, buscó la hoja once, la celda de las 16:05 del veintinueve ---seis minutos, una celda igual a todas las otras, del mismo alto que las que decían \textit{aclimatación} o \textit{barrido franja 40}--- y escribió adentro, apretando el lápiz porque el papel estaba pegado sobre revoque, tres apellidos. Fabbri. Ferrero, O., con una flecha hacia afuera del recinto. Y el tercero.

Vidal se quedó parado a un metro, leyendo su propio apellido en letra ajena, dentro de un rectángulo de dos centímetros de ancho.

---Está mal la inicial ---dijo---. Va S, no G.

Ledda borró con el canto de la goma, sopló el revoque y la corrigió.

\clearpage

\chapter{Ensayo en seis}

\lettrine[lines=2, lhang=0.1, nindent=0.2em]{A}{ las siete y un minuto Ledda apretó el cronómetro y nadie se movió durante los primeros dos segundos, porque la fase uno empezaba con las manos quietas: mirar el marco, confirmar posición, recién entonces tocar. El rectángulo de cinta azul seguía pegado al piso de la nave desde el ensayo del seis, cuarenta y un metros cuadrados con los vanos marcados en cruz, y adentro del rectángulo, sobre dos caballetes nivelados, el marco real de la teca esperaba con la tela que hacía de gemela plana bajo el cristal. La tela de prueba, en el papel del original, dormía enrollada en el tubo a dos pasos. Hacía frío de nave a esa hora y el caudalímetro de nitrógeno hacía su clic a intervalos que Vidal conocía de memoria: media de catorce coma nueve segundos, veinte mediciones.}


La primera pasada dio siete minutos cuarenta.

Ledda leyó los parciales de su cuaderno sin levantar la vista. Fase uno, doce segundos, en tiempo. Fase dos, apertura del marco, cuarenta y uno contra treinta y cinco. Fase tres, extracción de la gemela, un minuto cincuenta y ocho contra un minuto y medio. Fase cuatro, tubo, cincuenta y cinco contra cuarenta. Y así hasta el cierre. Vidal escuchó la lista completa y le hizo la aritmética antes de que Ledda la terminara: de los cien segundos de exceso, más de setenta estaban en fases donde las manos eran las suyas.

---El desvío está localizado ---dijo Vidal---. Soy yo.

---Es un dato ---dijo Ledda---. Se corrige o se reasigna.

Chiara se acercó al marco y le pidió que repitiera la fase tres sola, sin cronómetro. Vidal levantó el borde de la tela con el soporte rígido, verificó el asiento del canto contra el fieltro, lo verificó otra vez, giró, apoyó, verificó el apoyo, y cuando levantó la vista ella tenía la cabeza inclinada, mirándole las manos y no la tela.

---¿Qué viste la primera vez que miraste el canto?

---El asiento correcto contra el fieltro.

---¿Y la segunda?

---Lo mismo.

---¿Y qué ibas a hacer si la segunda te daba distinto?

Vidal abrió la boca y la cerró. Repetir una tercera. Eso era lo que habría hecho: una tercera lectura para desempatar, y una cuarta si la tercera no coincidía con ninguna. Once años de método en cuatro palabras que ella esperó sin apuro.

---Tu segunda mirada no ve más ---dijo Chiara---. Solo tarda más. La tela no cambió entre tus dos miradas. Cambiaste vos.

Le mostró con sus propias manos, una sola vez, sin ceremonia: mirar, decidir, mover, y la mano ya estaba en la fase siguiente mientras la cabeza terminaba de soltar la anterior. El gesto entero duraba lo que a él le duraba la primera verificación. Vidal lo repitió tres veces con el marco vacío. A la tercera le salió sin costura y sintió en las muñecas algo parecido a lo que sentía cuando una calibración cerraba en el primer intento: la sospecha inmediata de que había un error escondido.

Fue hasta la mesa, abrió el protocolo de la ventana y escribió con su letra de cadenas de custodia: \textit{Verificación única por punto de control. Segunda verificación prohibida dentro de la ventana.} Firmó al margen. Un procedimiento que ordenaba confiar en la primera lectura, con su firma debajo. Lo leyó una vez y cerró la carpeta antes de leerlo de nuevo.

La segunda pasada dio seis cincuenta.

---Mejor ---dijo Ledda---. Insuficiente.

El exceso residual estaba en la entrega del tubo: Vidal lo recibía de Chiara, lo giraba para leer el rótulo y recién después lo asentaba en el carro. Ledda propuso rotular las dos tapas para que cualquier orientación fuera legible. Dos minutos de cinta y marcador. Después juntó a los tres alrededor de la hoja once y repasó el reparto con la lapicera apoyada en cada apellido, como si los nombres necesitaran presión física para quedarse en su celda.

---Fabbri: manipulación de la tela, del tubo al marco, sin otra función. Vidal: la gemela y el contenedor, del marco al carro, y la hoja de servicio del instrumental. Ferrero: afuera, en el crucero, del lado de la barrera donde está la Comisión.

---Del lado donde están todos ---dijo Ferrero. Se tocó la muñeca izquierda, encima de la manga---. Once personas en tránsito y yo soy el que las mira a los ojos. Si alguien se acerca a la sacristía, ¿quién lo frena? Yo lo freno. ¿Con qué? Con cuarenta años de acta. Es el único material del plan que no se compra ni se ensaya.

---Se ensaya ---dijo Ledda---. Ahora.

Ferrero lo miró como se mira a un alumno que preguntó algo razonable en el momento equivocado, se puso de pie, caminó hasta el marco de listones que hacía de puerta y se plantó delante con las manos cruzadas atrás. Ledda hizo de monseñor sin cambiar la cara: dio dos pasos hacia la puerta simulada.

---Dottore, ¿puedo\ldots{}?

---Verificación de soporte en curso, ruego no ingresar hasta el\ldots{} ---Ferrero se detuvo. Bajó la cabeza medio segundo---. Eso era abril. Eso fue la otra vez. ---Se acomodó la corbata de lana---. Revisión del sistema de clima en curso. El técnico necesita el recinto estable. Cinco minutos, monseñor, y le explico las tablas mientras esperamos.

---Otra vez ---dijo Ledda.

La segunda vez le salió corrida, sin la costura. La tercera, con una variación: agregó una queja sobre el calendario de la Comisión que sonaba tan verdadera que Vidal tardó en entender que era parte del muro. Ferrero volvió a la mesa y se sirvió café frío.

---Cuarenta años diciendo la verdad delante de actas ---dijo--- y resulta que la verdad era el entrenamiento. Si esto sale mal, doctor, quiero que conste que mi mejor herramienta la fabricó la honestidad.

Vidal anotó los tiempos de recomposición y no encontró columna para la frase.

---Usted coordina desde el auricular ---le dijo a Ledda---. ¿Y si hace falta un tercer par de manos adentro?

---Adentro hay dos pares y esos dos alcanzan porque lo medimos. ---Ledda tapó la lapicera---. Yo soy la persona que si la ven, ya perdimos. Un técnico con credencial tiene explicación. Una conservadora tiene explicación. Yo tengo dieciocho años de historia limpia que se ensucian con estar parado en el lugar equivocado. Coordino desde afuera del edificio.

Corrieron la tercera pasada a las diez menos cinco.

Cinco cincuenta y uno. Vidal escuchó el número y siguió con las manos en el cierre del marco, terminando la fase, porque el protocolo nuevo decía que la cabeza suelta la fase anterior mientras la mano ejecuta la siguiente, y el número era la fase anterior. La tela que hacía de gemela quedó en el carro, rotulada por las dos tapas. La tela de prueba quedó bajo el cristal, plana, con las dieciséis zonas de apoyo asentadas. Ledda leyó la celda de ensayo del fondo: las cuatro células de carga estabilizadas en banda, masa en equilibrio a cincuenta y dos por ciento de humedad relativa, un gramo con una décima dentro de la tolerancia de más menos uno y medio. Tercera vez que el peso hidratado cerraba contra los sensores replicados. La primera había costado descoser once metros ochenta de respaldo; la segunda, cuarenta y ocho horas de ciclo; esta, nada. Vidal firmó la línea y resistió el impulso de pedir una lectura confirmatoria.

---Cinco cincuenta y uno con margen de nueve segundos sobre el peor caso ---dijo Ledda---. Y veintiuno sobre la media. Sirve.

Fue Chiara la que vio primero la hoja que Ledda había dejado sobre la mesa, boca abajo, debajo del cuaderno. La dio vuelta con las dos manos.

---¿Y esto?

---Pronóstico extendido. Llegó esta mañana. ---Ledda lo tomó y lo leyó como leía los manifiestos, renglón por renglón---. Para el veintinueve dan calor fuera de estación. Treinta y uno de máxima, treinta y dos el peor modelo. Récord para abril en la ciudad.

Vidal hizo el cálculo antes de que nadie lo pidiera. El sistema de clima de la teca contra treinta y un grados de nave exterior: compresores al límite, ciclos más largos, deriva en las sondas. Y la revisión de las 16:05, la que abría la ventana, la hacía un técnico con la Comisión mirando: si el sistema llegaba forzado, el técnico podía apurar la intervención para devolverle carga al circuito cuanto antes. La media histórica era seis minutos doce. La mínima registrada, cinco cuarenta.

---La ventana se puede achicar ---dijo Vidal.

---Se puede achicar ---dijo Ledda---. Regla nueva: si la ventana proyectada baja de cinco minutos, se aborta antes de abrir el marco. La tela no sale del carro y vuelve a Collegno esa misma tarde, por la misma ruta. El aborto lo canto yo, por auricular, y no se discute adentro del recinto.

---¿Y quién decide la proyección? ---dijo Ferrero---. Porque cinco minutos no es un dato que esté escrito en la puerta. Alguien mira un reloj y un compresor y dice un número. ¿Quién paga ese número si está mal?

---Lo pago yo. La proyección la hago yo con la telemetría de la mañana y el tiempo real del técnico en las dos revisiones previas del día. Si dudo, aborto. Dudar cuenta como menos de cinco.

Ferrero asintió despacio y se puso los guantes de algodón sin que la conversación lo pidiera.

Vidal miró la pared del fondo. Once hojas pegadas con cinta en las cuatro esquinas, cuatrocientas ocho horas celda por celda, y en la hoja once la celda de las 16:05 con los tres apellidos a lápiz, la S corregida. Después de esa celda quedaban tres más: cierre, desmonte, entrega de llaves. Después de la última celda, el borde de la hoja. El plan entero terminaba en un canto de papel a cuatro centímetros del revoque, y del otro lado del canto estaba el treinta de abril, que en ninguna hoja de ningún cuaderno de nadie tenía un solo renglón escrito para el caso de que la tela volviera esa noche a Collegno. Chiara estaba mirando el mismo borde. Ledda guardó el pronóstico en el cuaderno. Ninguno de los tres dijo la fecha.

---Una más ---dijo Vidal---. Completa.

Corrieron la cuarta a las once menos veinte, sin anuncio previo, con Ledda arrancando el cronómetro a mitad de una frase para robarles el segundo de preparación que el veintinueve nadie les iba a regalar. Cinco cuarenta y cuatro. Vidal cerró el marco, asentó el pestillo con una sola mirada, y sintió las dos lecturas que le faltaban como se siente un escalón que la escalera ya no tiene.

---Alcanza ---dijo.

Nadie le había preguntado. Chiara apagó el cronómetro que Ledda le había dejado en la mano y lo guardó sin mirar la pantalla.

\clearpage

\chapter{Veintinueve de abril}

\lettrine[lines=2, lhang=0.1, nindent=0.2em]{-}{--Nombre} y empresa.


---Vidal, Sebastián. Ambienti Controllati.

El carabiniere pasó la credencial por el lector, esperó, comparó la foto con la cara y después con la pantalla, que hacía lo mismo por su cuenta con la marcha y la silueta desde el momento en que él había bajado de la camioneta. Ocho segundos. En el turno anterior de la mañana habían sido seis y once, según el cronómetro de Ledda desde la esquina. Un contrato marco del catorce de febrero, dos facturas emitidas, aportes al día, alta en el anexo de personal autorizado de la diócesis. Todo cierto, todo verificable, y ninguna de esas verdades tenía relación con lo que iba a hacer esa tarde.

---Fabbri, Chiara. La misma empresa.

---Función.

---Asistencia de manipulación de textil en intervención sobre soporte ---dijo ella, y le dio el papel con las dos manos.

El hombre lo leyó entero. Atrás, en el hueco del acceso, el aire de la calle entraba caliente y quieto; a las nueve menos diez marcaba veintiséis grados y el pronóstico extendido del ensayo había dado treinta y uno, treinta y dos en el peor modelo. Vidal sintió la camisa gris pegada abajo del overol de la empresa. Nada de leyendas, había dicho Ledda en Collegno la primera semana, con la voz plana con que decía todo: personal con historia laboral y domiciliaria verificable, sin coberturas, porque una identidad falsa no aguanta una tarde de 2033. La frase terminaba ahí y la mitad que faltaba la había completado Vidal solo, en su cabeza, sin escribirla en ningún registro: una identidad verdadera con motivos ocultos, sí.

El carabiniere devolvió las dos credenciales y volvió a la tableta.

Ferrero había entrado a las nueve en punto por la puerta que le correspondía, con corbata de lana a veintiséis grados y la carpeta de la Comisión bajo el brazo, y no los había mirado. Lo esperaban para la revisión preostensión. Ese era el eje de la jornada, y las tres inspecciones que monseñor Ceruti había concentrado en un solo día colgaban de él como los renglones de un remito.

El carro estaba en el depósito de servicio, contra la pared, donde lo dejaba la empresa entre intervención e intervención. Acero inoxidable, cuatro ruedas de goma dura, dos bandejas superiores para instrumental y un bastidor inferior cerrado con puerta abatible y llave de cuadradillo. Vidal la abrió. Adentro, el juego de filtros de repuesto, dos manómetros en su caja, el rollo de tubo flexible. Y en el fondo, encajado entre el bastidor y la traviesa, sujeto con dos correas de velcro que Chiara había cosido a medida en Collegno, el tubo de transporte con el lienzo adentro, purgado, plano en la cabeza y a doce grados de la temperatura de la sacristía, que era la diferencia que iban a tener que dejar bajar. Al lado, el contenedor rígido que saldría lleno con la otra tela cuando llegara la ventana.

La donación de 2031 había pagado un sistema de atmósfera de nueva generación, sensores continuos, telemetría auditada. Dentro de ese mismo precio, sin renglón separado, había pagado también una empresa con contrato marco, un carro de acero con bastidor cerrado, dos cámaras nuevas en el acceso de servicio y el derecho de que un hombre con overol empujara todo eso por el pasillo de piedra sin que nadie levantara la cabeza. Vidal empujó. Las ruedas hicieron un ruido parejo sobre las losas, y en la junta de cada losa el tubo se movió dos milímetros dentro de las correas.

Segunda vez. Última.

A las once y veinte el compresor de la unidad entró y se quedó adentro. Vidal lo escuchó desde el pasillo de máquinas, un zumbido bajo con una vibración de tercera armónica que a Chiara le crispaba la mandíbula. En la pantalla de servicio, la temperatura de sala subía cuatro décimas cada hora y la unidad respondía con ciclos cada vez más largos: cincuenta y dos minutos, cincuenta y ocho, sesenta y siete. La banda de la sonda de humedad de la teca era de dos puntos y la lectura llevaba desde las diez pegada al borde superior de la banda, sin cruzarlo.

---Si sigue así lo van a mirar ---dijo Chiara.

---Van a tener que mirarlo. Eso es lo que queremos.

---No te pregunté qué queremos.

Vidal anotó el ciclo en la hoja de servicio con el bolígrafo de la empresa. El papel se le pegaba al canto de la mano.

Ferrero cruzó el crucero a la una menos cuarto, del lado de adentro de la barrera, con dos hombres de la Comisión y una mujer que llevaba el plano de la instalación eléctrica. Al pasar frente al carro se detuvo, se tocó la muñeca izquierda con dos dedos y le habló a Vidal como se le habla a un técnico.

---¿Ustedes tienen prevista la revisión de las cuatro?

---Está en el parte, dottore.

---Bien. ---Miró el número de la pantalla de servicio y después miró hacia el crucero, hacia el paño oscuro y la barrera baja---. Si esto sale como está previsto, mañana usted no tiene nada. Ni muestra, ni archivo propio, ni un papel que pueda mostrarle a nadie. ¿Ya pensó qué hace un hombre con las manos vacías un treinta de abril?

Del otro lado de la reproducción fotográfica había gente arrodillada. Vidal contó: once, doce con la mujer que se levantaba.

---No exactamente nada ---dijo---. Tengo diecinueve páginas de estados.

---Eso es un cuaderno. ---Ferrero se acomodó la carpeta---. Le pregunté por las manos.

Se fue detrás de los otros dos. Vidal volvió a la hoja, anotó el ciclo de las 12:51 y se quedó mirando la columna de horas sin agregar nada.

Contó desde la una hasta las cuatro y cinco. No los minutos del reloj: los ciclos del compresor, los pasos de gente por el pasillo norte, los segundos entre el clic del relé y el arranque del motor ---cuatro, siempre cuatro, salvo una vez que fueron cinco y quedó anotado---. La serie no se interrumpió en ningún momento. Contaba así desde hacía treinta años; en octubre, en el sótano de Har Hotzvim, habían sido trescientas quince baldosas grises. Doscientos cuarenta y un clics del relé, ochenta y siete pasos, tres ciclos completos del compresor. En el auricular, cada veinte minutos, Ledda daba temperatura exterior y nada más.

A las 16:03 llegó Bosio, el técnico de la empresa que hacía la revisión del sistema de clima con la Comisión mirando desde el otro lado de la barrera. Treinta y tantos, transpirado, con la tableta de la teca colgada del cuello.

---Aguanta bien ---dijo, y le mostró la curva a Vidal como se le muestra a un colega---. Mira el ciclado. Ahí adentro tienes a media Comisión y a monseñor con cuatro personas que entran y salen. Cada cuerpo mete cuarenta gramos de agua por hora en el aire del recinto.

---Ochenta con dos ---dijo Vidal---. Ciento sesenta con cuatro.

---Bueno. Si hago la revisión ahora, mido eso. Mido a la gente. ---Bosio pasó el dedo por la pantalla y ordenó la curva por hora---. Le voy a pedir a la Comisión que la corra al cierre. Cuando se vacíe el edificio y baje la carga, con dos horas de estabilización, ahí sirve el número. Antes, no.

---¿A qué hora sería?

---Nueve y media, diez. Depende de cuándo se vayan.

Ferrero, del otro lado de la barrera, se sacó los anteojos y empezó a explicarle a la mujer del plano por qué las tablas tomadas con tránsito de personal habían dado problemas en el sesenta por ciento de las intervenciones que él conocía. Habló tres minutos y dijo la palabra ``cierre'' cuatro veces. Ninguna decisión pareció tomarse; a las 16:22 la Comisión había decidido.

En el auricular hubo un silbido y después Ledda:

---Confirmado. Revisión reprogramada a las 21:40. ¿Qué hace con la tela?

Vidal miró el bastidor cerrado del carro. Cinco horas y dieciocho minutos con cuatro metros y cuarenta de lino del primer siglo dentro de un mueble de acero, en un edificio con dos cámaras nuevas en el acceso de servicio y una revisión de vigilancia que empezaba el lunes.

---¿Estamos en aborto?

---No hay ventana proyectada. La regla es sobre la ventana. Sin ventana no hay número que baje de cinco.

---Entonces la tela se queda ---dijo Vidal---. Si la sacamos ahora, la sacamos con el edificio lleno y con la camioneta rotulada un día en que la prefectura cuenta camionetas. La exposición de salir es mayor que la de esperar.

---De acuerdo. Registro que la decisión es suya.

---Regístrela.

Chiara pasó las cinco horas sentada en el piso del pasillo de máquinas, con la espalda contra la pared fría del tablero, la banda del pelo en la muñeca y el pelo atado con un trozo de hilo de algodón que había sacado del bolsillo del overol. A las siete y media se durmió veinte minutos, con la cabeza caída de costado y la respiración lenta. Vidal contó las respiraciones un rato ---doce por minuto, después once--- y dejó de contar cuando se dio cuenta de que las estaba contando.

Se despertó sola.

---¿Comiste algo? ---dijo.

---A las dos.

---Comiste el pan y dejaste lo demás. ---Se estiró los dedos uno por uno, como hacía antes de tocar---. Después de esto te llevo a comer a un lugar que conozco. No pregunto si querés.

A las 21:31 Bosio bajó al pasillo con la tableta y el permiso firmado, y las luces del crucero ya estaban en el nivel de noche. El edificio olía a cera caliente y a piedra que había estado al sol todo el día. Vidal soltó el freno del carro y empujó; Chiara caminó del lado izquierdo, entre el carro y la pared, con las dos manos apoyadas en la baranda del bastidor. Ochenta y tres pasos hasta la curva del pasillo. En la curva, treinta y uno más hasta la puerta de la sacristía.

En el paso veintidós de los treinta y uno, Vidal frenó el carro con la palma.

Debajo de la puerta de la sacristía había una línea de luz. Amarilla, continua, de las lámparas de techo y no de la piloto de emergencia. A esa hora, con el edificio cerrado y la revisión de Bosio recién a las 21:40, esa sala tenía que estar a oscuras.

El auricular chasqueó dos veces, que era el aviso de Ledda antes de hablar.

---Hay una persona adentro. Es Ceruti.

\clearpage

\chapter{Ceruti}

\lettrine[lines=2, lhang=0.1, nindent=0.2em]{L}{a} línea de luz medía el ancho de la puerta menos los dos burletes. Vidal dejó la mano sobre el freno del carro, que ya estaba puesto, y lo apretó otra vez.


---Entró a las nueve y cuatro por la puerta de la Comisión ---dijo Ledda en el auricular---. Firmó el libro. El sacristán le abrió la sacristía y se fue. Está solo.

Chiara había quedado a su izquierda, entre el carro y la pared, con las dos manos en la baranda del bastidor. A través de la puerta se oía un ventilador pequeño, de los portátiles, y debajo del ventilador un roce de papel a intervalos irregulares.

La aritmética era simple y ninguna de sus salidas cerraba. La revisión de Bosio estaba fijada para las 21:40, dentro de nueve minutos. La ventana existía dentro de la revisión y en ningún otro lugar. Un monseñor a treinta metros de la teca era, según la regla que Ledda había escrito y Vidal había firmado, menos de cinco minutos. La duda contaba como menos de cinco minutos. Y el plan de once hojas de la pared de Collegno tenía tres celdas después de esta y ningún renglón para el treinta de abril.

Se le había secado la boca. Contó las vetas de la baldosa que tenía bajo el pie derecho ---cinco, una bifurcada--- y volvió a empezar.

Entonces la puerta de la sacristía se abrió desde adentro, apenas, y la línea de luz se ensanchó hasta hacerse un rectángulo. Ceruti había salido a buscar aire al pasillo del lado del crucero: Vidal le vio la espalda, la camisa sin alzacuello, arremangada, y una tableta en la mano izquierda con la pantalla encendida. En la pantalla había una tabla. A esa distancia las columnas eran solo columnas, pero una franja horizontal cruzaba la tabla en amarillo, y el amarillo, en cualquier volcado de telemetría que Vidal hubiera visto en su vida, marcaba una sola cosa.

---Proyección: cero ---dijo Ledda---. Abor\ldots{}

Pasos en el crucero. Ledda se calló a mitad de la palabra, cosa que en cuatro semanas nunca había hecho, y por el fondo del pasillo apareció Ferrero.

Venía caminando a su paso de siempre, encorvado, la corbata de lana torcida por el calor, y mientras caminaba se estaba quitando los guantes de algodón. Primero el derecho, después el izquierdo, un tirón por dedo, sin apuro. Los dobló juntos y se los guardó en el bolsillo del saco, y llegó hasta Ceruti con las manos desnudas.

---Monseñor. Qué noche para estar de pie.

---Professore. ---Ceruti se abanicó con la carpeta que llevaba en la otra mano---. Mi departamento da al oeste. Entre dormir ahí y venir a ver los sensores, elegí los sensores. ---Levantó la tableta---. Y me alegro de haber venido, porque el volcado que pedí llegó esta tarde y tiene una cosa que quiero que alguien me explique antes de que la explique otro.

---A ver.

---Doce de abril. Tres y veintiuno de la madrugada. Las cuatro células de carga, pico de dos gramos y una décima. Once segundos fuera de banda, estabilizada al doce. Después queda más quieto que antes. ---El roce de la carpeta contra el aire se detuvo---. Yo soy conservador, professore, no ingeniero. Pero sé leer una tabla, y sé que a las tres y veintiuno de esa madrugada esa sacristía tenía que estar cerrada, con el marco puesto. Yo volví a las siete y cuarenta a bendecir una teca que se suponía que nadie había tocado.

Vidal dejó de sentir la palma sobre el freno. Doce segundos. El pico que él había decidido dejar escrito porque una corrección dejaba firma de corrección, y la decisión había sido correcta, seguía siendo correcta, cada término de ese cálculo resistía revisión. La operación entera colgaba de un renglón amarillo en la tableta de un hombre que había venido a la catedral porque hacía calor.

Al lado suyo, Chiara movía los labios sin voz. Dos frases. La segunda terminó en algo que podía ser un nombre.

---La lectura es real ---dijo Ferrero.

Lo dijo con la voz de acta, la que llegaba hasta el fondo de cualquier sala, y Vidal se encontró contando las afirmaciones que vinieron después como contaba renglones en una cadena de custodia.

---La conozco. La tengo anotada. Y le voy a decir por qué la conozco, monseñor, y le pido que me escuche entero antes de contestar, porque la parte fácil de esto es enojarse y la parte difícil es lo que viene a cinco días de la ostensión. ---Una pausa; el ventilador siguió---. Esa teca la pagó una donación que la Comisión nunca auditó. Las especificaciones vinieron de afuera, más completas de lo que nosotros habríamos sabido pedir, y nadie preguntó por qué. Yo pregunté. Llevo semanas ejecutando verificaciones discretas del instrumental, bajo mi propia autoridad, sin acta, porque un acta a medio hacer, filtrada en la semana equivocada, hace un daño que no lo mide ningún sensor. Usted vio lo que pasa cuando un dato sale sin nadie al lado. Yo también. Los dos éramos jóvenes.

---Professore\ldots{}

---Todavía no terminé. ---Otra pausa, más corta---. Yo firmé el acta de cierre del año tres, dottore. Tercero de cinco firmas. Al año siguiente conté los contenedores y faltaban tres, y reclasifiqué la diferencia y me callé. Llevo veintinueve años callado. No vuelvo a firmar un silencio. Ayúdeme a mirar los logs ahora, conmigo, en la sala de clima. Los dos, esta noche, línea por línea. Si hay algo que la Comisión tiene que saber, lo sabe por nosotros y con las tablas completas, no por un renglón amarillo suelto.

Cinco afirmaciones verificables. Vidal las recontó. Ninguna falsa. Ninguna completa. La misma estructura de las frases de Ginebra, en boca de un hombre de sesenta y tres años con la corbata torcida.

Ceruti tardó en contestar. Bajó la tableta, y cuando habló, la voz institucional se le había corrido a un costado.

---Mi antecesor decía que en esta casa los problemas se dividen en dos: los que se resuelven rezando y los que se resuelven a tiempo. ---Un silencio---. ¿Los tres contenedores existen todavía?

---Esa es otra conversación, y también se la debo. Esta noche, los logs.

---Los logs. ---El monseñor se guardó la carpeta bajo el brazo---. Sepa que voy a anotar todo lo que me diga, professore. Con su nombre.

---Para eso se lo estoy diciendo con mi nombre.

Se fueron por el crucero, hacia el pasillo de máquinas, la tableta encendida entre los dos. Vidal siguió las dos siluetas hasta que la columna las cortó. La sacristía quedó con la luz prendida y vacía, el ventilador girando para nadie.

El auricular chasqueó dos veces.

---Ventana. Ahora.

Del lado de la Comisión, del otro lado de la barrera, quedó la mujer del plano eléctrico ---Ferrero la había retenido antes con la misma queja de calendario que había ensayado--- mirando sin mirar, con la carpeta cerrada bajo el brazo. Bosio esperaba junto a la barrera baja con la tableta apoyada en el tablero del módulo; a las 21:41 puso las células de carga en modo de servicio y liberó los pasadores del marco con la llave de cuadradillo, recitando valores en voz alta para su propia acta, de espaldas al recinto, como recitaba siempre. Chiara ya tenía el tubo fuera del bastidor. Ledda arrancó el cronómetro sin aviso, a mitad de un movimiento, igual que en la cuarta pasada del ensayo.

Vidal levantó la gemela del marco. Pesaba lo previsto, caía como una tela de tres años tratada de mil maneras para caer como una de seiscientos, y por primera vez desde Bélgica le pareció exactamente eso: un objeto correcto en todos sus parámetros y equivocado en algo que no tenía columna. La plegó sobre el alma del contenedor rígido en los tiempos ensayados y cerró las presillas. Dos minutos diez.

---Tres cuarenta de margen ---dijo Ledda.

Las manos de Chiara sacaron la tela del tubo. El lienzo cedió sobre sí mismo con el ruido seco, de papel viejo, que ningún ensayo había imitado, y se desplegó sobre el soporte del marco en un solo movimiento continuo que Vidal había visto nueve veces en Collegno y que seguía sin poder descomponer en fases. Ella lo asentó con los dedos abiertos, sin uñas, leyendo la caída del borde contra las marcas de referencia.

---Posición ---dijo.

Vidal verificó. Tres cotas dentro de banda, limpias. La cuarta, el borde inferior derecho, pegada al límite sin cruzarlo. El impulso de tocarla, de correrla un milímetro hacia el centro de la banda, le subió por el brazo con toda la fuerza de treinta años de segundas lecturas, y se estrelló contra una cláusula que llevaba su propia firma al margen: verificación única por punto de control. Dentro de banda era dentro de banda. La cláusula ganó. Bajó la mano.

---Conforme.

Bosio repuso los pasadores, giró la llave, sacó las células del modo de servicio. Los cuatro indicadores parpadearon, buscaron, y se quedaron quietos.

---Cinco cuarenta ---dijo Ledda---. Cerrado.

Chiara se apartó del cristal un paso y se quedó mirando la tela un segundo entero, uno solo, que el cronómetro ya no medía. Después se agachó a juntar las correas de velcro.

Cargar el carro llevó otros cuatro minutos que ya eran administrativos. Vidal firmó la hoja de servicio del instrumental ---lecturas verificables, documento falso, la mano no le tembló y registró que no le temblara--- y empujó el carro de vuelta por el pasillo, con el contenedor rígido sujeto donde a la ida había viajado el tubo. En algún punto de esos cuatro minutos hizo la cuenta que faltaba, la que ninguna celda del plan tenía prevista: Ferrero acababa de atar su nombre, con testigo y a pedido propio, a una desviación de telemetría que cualquier auditoría futura podía levantar del registro. Vidal dio vuelta el cálculo tres veces buscándole el error y el error se negó a aparecer.

La puerta de la sala de clima estaba entornada y adentro había luz de pantalla. Al pasar con el carro, Vidal aflojó el paso sin decidirlo.

---\ldots{}y esta columna es la deriva térmica del sensor, que se corrige contra la sonda patrón ---decía la voz de Ferrero, la voz de un hombre que explicaba tablas como si fueran las diez de la mañana---. Lo que usted tiene que mirar es la pendiente, dottore, no el valor.

---¿Y si al final de todo esto la Comisión me pregunta por qué lo supe por usted y en abril?

Un silencio del largo de un trago de agua.

---Entonces habrá consecuencias, y las firmo yo. Cuarenta años firmando actas de otros. Ya era hora de que una fuera mía.

Chiara, del otro lado del carro, miró a Vidal por encima del bastidor. Traía el pelo atado con el hilo de algodón y no dijo nada.

Vidal siguió empujando. El carro pesaba lo mismo que a la ida.

\clearpage

\chapter{Devolución}

\lettrine[lines=2, lhang=0.1, nindent=0.2em]{L}{a} consola de la sala de clima tenía una pantalla de once pulgadas y un teclado con el barniz gastado en cuatro teclas. Bosio la había puesto en modo lectura y estaba llenando su acta a mano, apretando el papel contra el borde del mueble, con esa letra chica de los que escriben mucho de pie.


---¿Le sirve la ventana de estabilización? ---dijo---. Se la dejo abierta, si quiere anotar el instrumental.

---Me sirve.

Cuatro columnas. Célula uno, dos, tres, cuatro, muestreo cada segundo desde las 21:47. Vidal leyó los primeros doce segundos y los volvió a leer. Desvío máximo: tres décimas de gramo en la célula tres. Al sexto segundo la pendiente se aplanó, y de ahí en adelante las cuatro columnas hacían lo que hace un objeto que está donde tiene que estar: nada, con dos decimales. Ninguna entrada en el registro de eventos. Ninguna banda tocada.

---Se portó mejor que en abril ---dijo Bosio, sin levantar la vista del acta---. El doce tuve un pico al cierre. Dos gramos y pico, once segundos. Lo dejé asentado como asentamiento del soporte, que es lo que fue.

---Es la explicación razonable.

---Es la única. Estas células son buenas pero les molesta que las toqués.

Vidal copió las cuatro lecturas en la hoja de servicio y firmó abajo la línea de instrumental conforme. Dos meses de telar en Bélgica. Ciento veinte horas de horno a noventa y ocho grados con ciclos de humedad. Once metros ochenta de respaldo descosido, medido y vuelto a coser en tres días porque la batista tratada tomaba seis gramos y cuatro de agua que la otra no tomaba. Tres pruebas contra sensores replicados, la última cerrada apenas dos noches antes del traslado. Todo eso para entrar por una puerta que la tela verdadera cruzó sin que nadie tuviera que abrirle nada. Lo auténtico no pedía calibración: era el patrón. Anotó la frase en el margen de la hoja de servicio y después la tachó, porque la hoja de servicio iba a quedar en el archivo de una empresa que existía de verdad.

Salió al pasillo. Chiara ya se había ido con el carro por el acceso de servicio a las 21:58, tubo vacío y contenedor rígido bien sujeto, con el remito de descartes en la mano y la naturalidad de una operaria que termina su turno. Ledda había fijado catorce minutos de separación entre las dos salidas. Faltaban cinco.

En el crucero la luz estaba baja, la de las noches sin público: la reproducción fotográfica apagada, dos reflectores lejanos y la barrera con su paño oscuro recogido a un costado porque la Comisión seguía de inspección. La teca era una caja horizontal de aristas limpias, y adentro, detrás del cristal, la tela plana con las quemaduras simétricas de 1532 corriendo a lo largo como dos hileras de puntos de sutura.

Ferrero venía de la sala de clima con la tableta de Ceruti bajo el brazo. Se detuvo del lado de la barrera, del lado permitido, y sacó del bolsillo derecho los guantes de algodón. Se los puso sin apuro, primero el izquierdo. Después apoyó dos dedos en el cristal, arriba, cerca del borde, en el punto exacto donde se controla la condensación de junta.

Vidal contó. Uno. Al tres el hombre se inclinó apenas, la nariz a un palmo del vidrio. Al seis retiró la mano, la miró como se mira un guante que podría haber recogido humedad, y anotó una cifra en la libreta que llevaba en el bolsillo del saco: la anotó para que el gesto tuviera acta.

Dieciséis mil placas publicadas entre 2004 y 2019. Cuarenta años de calibración cromática documentada, de notas al pie que decían urdimbre y decían reliquia sin comillas. Seis segundos con dos dedos en un cristal.

Los guantes volvieron al bolsillo. Ferrero se dio media vuelta hacia el pasillo de la Comisión, y desde el fondo la voz de Ceruti preguntó algo sobre la columna de oxígeno residual, y Ferrero contestó que eso se lee contra la sonda patrón, que la pendiente importaba más que el número suelto, y que si querían podían repasar juntos el histórico de febrero, que era el mes en que la teca todavía estaba aprendiendo a respirar.

Se iba a quedar hasta el final. La mejor cobertura del mundo, y la más cara, era seguir trabajando.

Vidal salió a las 22:12. El control de credencial le llevó ocho segundos, los mismos de la mañana. Afuera la piazza San Giovanni conservaba el calor del récord de abril: veinticinco grados a las diez de la noche, asfalto tibio a través de la suela, un olor a piedra recalentada que en la ciudad vieja de Jerusalén había sido idéntico y con otro polvo.


\scenebreak


En Collegno el desmontaje había empezado antes de que él llegara.

Ledda trabajaba con la misma prolijidad con la que había construido todo, en el mismo orden, invertido. La cinta azul de la réplica de la sacristía salía del piso de la nave en tiras enteras, enrolladas sobre sí mismas y metidas en una bolsa negra: cuarenta y un metros cuadrados de geometría transcrita de la pestaña dos, que a las once y media eran una bolsa de basura de doce litros. Las once hojas del plan de las cuatrocientas ocho horas bajaron de la pared del fondo por orden inverso, once, diez, nueve, con la cinta despegada de las cuatro esquinas de a una para no arrancar revoque. Los tres racks de Tamiz quedaron sin cables, con los veinticuatro conectores del frente al aire, y el hueco donde había estado el rack central mostraba en el piso flotante cuatro marcas de rueda y un ovillo de polvo que la sala blanca, con su presión positiva, no había tenido derecho a producir. Ledda barrió el ovillo con la mano. Después llenó una etiqueta, la pegó en un flanco, y anotó en su cuaderno la hora y el apellido de la empresa de transporte, subrayando el renglón con regla.

La gemela estaba en su caja de conservación, plana, sobre la mesa del taller textil, con la sábana de papel libre de ácido doblada hacia atrás en un triángulo.

Chiara la miró un rato largo, de pie, con los antebrazos cruzados.

---¿Y esto?

---Se destruye ---dijo Ledda.

---Te pregunto cuándo.

---Esta noche. Antes de que salga el segundo camión. ---Cerró el cuaderno con el pulgar adentro---. Una réplica perfecta es una acusación esperando dueño.

Nadie lo contradijo. Vidal calculó la afirmación buscándole la falla y la afirmación no tenía falla: mientras esa tela existiera, existía un objeto físico que probaba capacidad, intención y ejecución, y bastaba una fibra bajo microscopio para que un peritaje de cuarenta minutos convirtiera diecisiete días en una carátula.

---¿Cómo la piensan hacer? ---preguntó Chiara.

---Horno industrial. Un taller que ya usamos antes, con acta que nadie va a archivar. Setenta y dos horas.

---Antes de que se vaya, la corto yo.

---No hace falta.

---Hace falta para mí. ---Ya estaba abriendo el cajón donde tenía las tijeras de conservación, las de hoja corta---. Cuatrocientas ochenta y dos perforaciones a mano en el parche de arriba, y las de abajo, y el orillo. Nadie más las hizo, y no las mando enteras a que las mire un desconocido con un soplete. ---Levantó la vista, y por primera vez en cuatro semanas usó con Ledda el tono con el que le hablaba a Vidal---. El resto que arda. Esto lo deshago yo, ahora, con mis manos.

Trabajó desde la medianoche hasta las cuatro menos veinte, sentada, con la lámpara rasante bajada a cuarenta grados sobre la mesa, separando puntada por puntada el parche superior del cuerpo de la tela. Vidal la escuchó desde la sala de máquinas: el ruido de la tijera en el lino seco era corto y limpio, dos golpes por corte, con una pausa de un segundo entre uno y otro, y no se aceleró nunca. En algún momento de la primera hora ella empezó a rezar en voz muy baja, las mismas dos frases de antes de comer, y Vidal advirtió que la segunda ya no terminaba en el nombre que terminaba antes.

Él tenía sus propios cierres. Volcado del acelerómetro a papel, copia única. Verificación de las llaves de la caja de dos cerraduras, una suya, una de Ferrero, y la caja vacía. Recuento del remanente de muestra: cero, en las cuatro vías, sin reserva. A las 4:26 abrió la consola del conmutador para desmontar el puerto cuatro y se quedó cuatro minutos mirando la ventana fija abrirse a las 4:30, como se había abierto todas las madrugadas desde el día uno, declarada en la pestaña nueve, tercera hoja, con sus iniciales al margen desde enero.

Doscientos catorce archivos. Manifiesto con hash y hora, tres por ciento de sobrecarga, y en el campo de estado, a las 4:33: \textit{sincronización completa. Sin pendientes.}

Ahí estaba todo. Los cincuenta y ocho gigabytes de datos crudos y los ochenta y dos barridos de doce micrones, los cinco controles ciegos con su lino copto del siglo IV, la salida de catorce kilobytes de la cuarta pasada con el índice global en 0,97, el consolidado con clasificación \textit{muestra inescrutable} y detalle \textit{no existe población de referencia}, las actas escaneadas del despliegue con la letra de Chiara debajo de la suya, y el borrador de diecinueve páginas con el índice que decía Estados y con la página catorce, escaneada para el anexo, donde su propia mano había escrito que la señal no se descarta como ruido de sistema y se asienta sin clasificar, y donde Ferrero había firmado abajo como testigo estando en contra de que la línea existiera.

Ninguna copia ilegítima. Cada byte había viajado por un canal que él autorizó con sus iniciales cuatro meses antes de saber lo que iba a haber para mandar.

El teléfono fijo sonó a las 5:04 y no se cortó antes del segundo timbre.

---Doctor Vidal. Buenas noches, o buenos días, no sé cuál corresponde a esta hora. Le pido disculpas por el horario y le aclaro de entrada que no llamo por ninguna novedad, porque no hay ninguna.

---Terminó a las 21:47.

---Lo sé, y no lo sé por usted, de modo que no le estoy pidiendo confirmación de nada. Tengo instrucciones antiguas para este momento y quiero cumplirlas con exactitud. El fundador le pidió que le comunicara que lo espera en Cologny cuando esto termine. La frase es textual y es de enero. Yo no estoy autorizado a interpretarla, pero entiendo que usted está en mejores condiciones que yo para saber si el momento llegó.

Vidal miró la pantalla, todavía con el manifiesto abierto.

---Depende de qué llame usted terminar.

---Le concedo que la palabra admite lecturas. Le propongo la más simple: el objeto está en su sitio, la inspección se cerró y nadie del equipo tiene una tarea pendiente distinta de guardar sus cosas. Si esa lectura le parece razonable, el automóvil puede estar en la estación de Ginebra a las diez y media. Si prefiere dormir un día, la misma oferta vale para el jueves y no cambia en nada.

Desde el taller llegaron dos golpes de tijera y la pausa de un segundo.

---Hoy ---dijo Vidal.

---Entonces a las once. Doctor: hay una sola cosa que le pido, y se la pido en nombre propio. No lleve nada escrito.

Colgó él primero, sin cortesía de cierre, y anotó la cita en el cuaderno de tapas de cartón naranja: fecha, hora, lugar. En el campo de entregable, donde su sistema exigía texto desde 2029, escribió \textit{los motivos}, y la punta del lápiz se quedó ahí, en el punto final, hasta que en el taller la tijera dejó de sonar y Chiara preguntó, sin levantar la voz, si quedaba café.

\clearpage

\chapter{Confesión}

\lettrine[lines=2, lhang=0.1, nindent=0.2em]{Q}{uedaba} café para media cafetera y el taller ya devolvía eco. Sin la cinta azul en el piso, sin las once hojas del plan en la pared, la nave había recuperado su acústica de galpón: cada taza apoyada sonaba dos veces. La línea de nitrógeno estaba cerrada desde las tres y el silencio donde antes hacía clic el caudalímetro era un dato nuevo que Vidal registró sin columna donde ponerlo. Ledda despegaba etiquetas de los racks al fondo, una por una, y las pegaba en una hoja de descargo por orden de retiro.


Chiara sirvió las dos tazas y se sentó del otro lado de la mesa del taller textil, la única superficie que quedaba sin embalar.

---¿Tenés diez minutos?

Vidal miró el reloj. El tren de alta velocidad a Ginebra, el que habían abierto el año pasado por el nuevo túnel, salía a las 7:12; el margen era de cincuenta minutos más traslado. Diez cabían.

---Tengo doce.

Ella asintió. Sobre la mesa, a veinte centímetros de su mano derecha, había un resto de hilo de embalar, medio metro, suelto. Chiara puso las manos una sobre la otra, planas, y las dejó ahí. El hilo quedó donde estaba. Vidal llevaba cuatro semanas leyendo a esa mujer por lo que enrollaba antes de hablar, y ahora las manos quietas le decían que lo que venía no era una pregunta difícil. Era otra cosa. Algo ya decidido.

---Yo no entré a esto por la tela ---dijo---. Entré por tres contenedores.

---Cuarenta y uno, cuarenta y dos, cuarenta y tres.

---Los conocés por Ferrero.

---Por Ferrero conozco los números. Por vos, aquella noche, conocí los cuadernos. Nunca uní las dos cosas por escrito.

---Unilas ahora. ---Sacó de su bolso una carpeta fina, de cartón, y la dejó cerrada entre los dos---. Mi maestra anotaba primero el daño y después la causa. En junio de 2002 dibujó a mano, con cotas, las fibras del refuerzo dorsal antes de descoserlo. Torsión, diámetro, un defecto de hilatura cada tantos centímetros. Dibujó también los hilos carbonizados, los sesenta y tres, uno por tubo. Y pesó el polvo del aspirado: veintiún gramos. Veintidós con portamuestras.

---Eso está en los cuadernos.

---Eso está en los cuadernos. Lo otro está acá. ---Tocó la carpeta con un dedo y volvió a poner la mano sobre la otra---. El corpus de referencia de este laboratorio. Las nueve muestras de Judea del siglo primero. Revisé cinco, las noches que me sobraban. Cuatro coinciden con los esquemas de mi maestra. Cota por cota. El defecto de hilatura está, Sebastián. Dibujado en 2002 con lápiz, y abajo del microscopio en 2033, en una muestra con procedencia documentada y cadena de custodia en regla.

Vidal dejó la taza. El eco del galpón la hizo sonar dos veces.

---Una coincidencia de torsión y diámetro admite población. Un defecto de hilatura con cota\ldots{}

---No admite población. Lo sé. Por eso revisé cinco y no una.

Él se levantó a medias, hacia el fondo de la nave, donde los cajones del corpus esperaban flejados con el inventario de la empresa de transporte firmado por Ledda. Se detuvo antes de dar el segundo paso. Abrir un cajón ahora rompía una cadena de custodia que ya no era suya; el material volvía a la Fondation esa mañana, contado y sellado. Volvió a sentarse. Era la primera medición importante de su vida adulta que iba a aceptar sin repetirla, y el instrumento era la letra de una mujer que limpiaba estandartes municipales.

---Los perfiles de carbonización ---dijo.

---También. Los patrones de historia térmica que usaste para calibrar el chamuscado. Comparé tres contra las fichas de los tubos del contenedor cuarenta y dos. Longitud, sección, el punto donde el hilo pasa de tostado a íntegro. Está todo en la carpeta, con la referencia de cuaderno al lado de cada cota.

El cálculo se armó solo, sin que él lo convocara. El corpus superaba a cualquier colección universitaria conocida y seguía siendo pobre: esa frase la había escrito él mismo el día dos, y debajo había dejado una línea sin completar que empezaba \textit{procedencia del corpus:}. Nadie arma en cuatro años una colección de lino de Judea con cadenas limpias. La arma en treinta, o la arma con material que ya tiene. Cassiodore no había comprado los contenedores para guardarlos en un puerto franco. Los había comprado para medir con ellos. Y si alguien medía con fragmentos del propio lienzo desde hacía años, entonces alguien ya había hecho la pregunta antes, en pedazos, con instrumentos de otra década, y había recibido lo que se recibe de un fragmento: ruido, ambigüedad, un resultado que no alcanza. La cláusula octava decía \textit{toda anomalía registrada durante los trabajos}, redactada antes de que existieran los trabajos. Uno escribe esa palabra cuando ya sabe que va a haber una.

---No fue una pregunta ---dijo Vidal. La boca se le había secado y tomó café frío para disimularlo---. Fue una verificación. A escala completa. Nosotros fuimos la segunda corrida del mismo experimento.

---Vos lo decís mejor que yo. Yo lo pensé así: alguien ya había probado con retazos y no le alcanzó. Nosotros éramos la tela entera.

Vidal abrió la carpeta. Adentro había seis hojas escritas a mano, prolijas, cada renglón con su referencia doble: cuaderno, folio, cota; muestra, portaobjetos, medición. La estructura era auditable. Primero el daño, después la causa. Cerró la carpeta y se quedó con la mano encima.

---¿Cuándo lo supiste?

---Con certeza, la segunda semana. La noche que me encontraste con los portamuestras ya tenía tres de las cuatro coincidencias.

---Y no dijiste nada.

---No. ---Las manos seguían quietas. Ni el hilo de embalar, ni el borde de la carpeta, ni la banda del pelo---. Si lo decía, ¿qué hacías vos?

---Detenía el análisis. Exigía trazabilidad del corpus. Sin corpus verificado no hay calibración, sin calibración no hay\ldots{}

---No hay peritaje. Y sin peritaje, ¿la devolución cuándo era?

Vidal no completó la cuenta en voz alta. La cuenta daba nunca. Una operación reventada el día catorce dejaba el original en Collegno y a la gemela en la teca el 4 de mayo, con el papa delante.

---Elegí la tela ---dijo Chiara---. Sobre la verdad, sobre vos, sobre mi maestra. Y quiero que sepas una cosa antes de que te subas al tren: lo volvería a hacer. Mañana, igual. No te lo cuento para que me perdones. Te lo cuento porque ya no hace daño.

Él se quedó mirando el punto de la mesa donde las manos de Chiara no se movían. Conservarla hasta que dejara de hacer daño. Ella lo había hecho con algo más grande y con mejor pulso: custodia, no secreto.

---¿Cómo se llama? ---preguntó.

Chiara parpadeó. Era evidente que esperaba cualquier pregunta menos esa.

---¿Quién?

---Tu maestra. En los cuadernos la letra no tiene firma.

---Bruna. Bruna Ottino. ---Hizo una pausa corta---. Le va a gustar saber que alguien preguntó el nombre antes que el método.

---No exactamente. Pregunté el nombre porque el método ya lo había leído.

Ella se rió una vez, baja, cansada de una noche de tijeras, y empujó la carpeta los últimos diez centímetros.

---Es copia. Los cuadernos se quedan conmigo, eso no se negocia. Vos vas a Cologny.

---¿Quién te dijo que\ldots{}?

---Nadie. Te afeitaste. ---Se puso de pie y juntó las dos tazas---. Yo con esto ya no puedo hacer nada. Denunciarlo es reventar la devolución para atrás; callarlo para siempre es regalárselo. Vos sí podés hacer algo. Podés preguntar. Ahí, en el jardín ese, mirándolo. Es lo único que te pido. Que la pregunta exista en voz alta una vez.

Vidal tomó la carpeta con las dos manos. Advirtió el gesto cuando ya lo había hecho, el peso repartido, el cartón horizontal como si fuera tela, y no lo deshizo. Ansermet le había pedido, en nombre propio, que no llevara nada escrito. La discrepancia quedaba registrada. No había a quién elevarla: la asentó igual.

Antes de guardarla abrió el cuaderno de tapas naranjas sobre la mesa y escribió, con la letra de las cadenas de custodia, la línea que había dejado incompleta el día dos: \textit{Procedencia del corpus: material removido del lienzo en 2002. Contenedores 41, 42, 43. Integrado a mi calibración sin pregunta previa. El peritaje midió la reliquia contra sí misma en cuatro de nueve referencias de Judea. Constancia: Fabbri, comparación adjunta.} Releyó el párrafo buscando un adjetivo que tachar y comprobó que no había puesto ninguno. La medición más limpia de su carrera, con los cimientos hechos de material robado, y la frase aguantaba sin adjetivos. Eso también era un dato.

Ledda cruzó la nave con la hoja de descargo completa y el bolso al hombro. Se detuvo a un metro de la mesa, la distancia de siempre.

---Doctor. El transporte pasa a las ocho. Setenta y dos horas y esta dirección no existió. Ni para nosotros ni para nadie que pregunte bien. ---Extendió la mano y Vidal se la estrechó: un apretón exacto, sin segundo movimiento---. Su tren sale a las 7:12. Con este tránsito, salga 6:25.

Después se volvió hacia Chiara. Vidal esperó el mismo apretón y el mismo horario. En cambio Ledda se quedó un segundo con las manos a los costados.

---Cuando me dieron su perfil, en noviembre, decía historia verificable completa y ninguna observación. Yo agregué una a lápiz: motivo de participación no establecido. La regla de esta operación fue siempre la misma: historia verdadera, motivos ocultos. Es la única combinación que resiste. ---Hizo una pausa que en él equivalía a un discurso---. Sus motivos ocultos fueron los mejores del proyecto.

Era el segundo adjetivo que Vidal le escuchaba en cuatro semanas, y esta vez venía en superlativo. Chiara recibió la frase como recibía los objetos ajenos, sin soltarla enseguida.

---¿Desde cuándo lo sabe?

---Desde noviembre ---dijo Ledda, y se colgó mejor el bolso---. Los motivos ocultos no son un problema de seguridad. Son un problema de dirección, y la de ella cerró sola. ---Miró una vez más la nave vacía, comprobó algo contra su hoja y se fue por la puerta de personal sin despedirse dos veces.

Vidal guardó la carpeta en el compartimiento interno del bolso, en el mismo lugar donde en enero había viajado el expediente de catorce pestañas, entre el pasaporte y el billete. Chiara lo acompañó hasta el portón. Afuera el polígono olía a asfalto tibio; el calor fuera de estación seguía instalado y a las seis y cuarto ya no hacía frío.

---Preguntá ---dijo ella---. Aunque conteste bien. Sobre todo si contesta bien.

---Depende de qué llames contestar bien.

---Sí. Depende. ---Se apoyó en el marco del portón, los antebrazos cruzados, las manos por fin desocupadas de verdad---. Esto de hoy no era. Cuando vuelvas, lo que te quería mostrar sigue en pie. No es de ahora. Es de antes.

---Lo sé. Lo dejaste inconcluso el veintiocho a la noche. Treinta y cuatro segundos de cable.

---Los contaste.

---Conté. Es lo que hago.

Ella sonrió con la mitad de la boca y le hizo un gesto con el mentón hacia el auto, andá. Vidal ya había abierto la puerta cuando la escuchó a su espalda, en el tono exacto de quien repite una fórmula heredada, mitad taller y mitad otra cosa:

---Lavate las manos antes y después.

Subió, puso las manos en el volante y las miró un momento: limpias, secas, con el pulso en banda. La carpeta le marcaba una arista contra el costado del pecho a través del bolso. Arrancó con esa presión ahí, sin acomodarla.

\clearpage

\chapter{El archivo}

\lettrine[lines=2, lhang=0.1, nindent=0.2em]{E}{l} túnel se tragó el paisaje a los veinte minutos de Turín y el vagón quedó reducido a su propio ruido, un zumbido parejo con una vibración de fondo que Vidal estimó por debajo de la banda que le habría preocupado en un transporte de muestras. Bajó la mesita, sacó el informe del bolso y lo apoyó sobre la superficie gris. Cuarenta y una páginas. Las había contado en el andén y volvió a contarlas ahora, pasando las esquinas con el pulgar. Cuarenta y una. Registró la coincidencia con la fila de octubre y le negó valor en el mismo gesto: un número no significa, se repite.


Lo había escrito en los últimos tres días del cronograma, entre una tarea y otra, en la mesa de la sala de reuniones de Collegno; la última media página, la del cierre, la completó esa madrugada mientras la nave se desarmaba alrededor con un orden que era el del montaje leído al revés. Escribió a mano, con la letra de las cadenas de custodia, deteniéndose dos veces por página a tachar adjetivos. La estructura era la de un peritaje clásico porque la estructura clásica era la más difícil de tergiversar: objeto, método, cadena de manipulación, resultados por canal, estados. En el índice, donde cuarenta años de literatura del campo ponían ``Conclusiones'', decía ``Estados'', y los estados eran cuatro y cabían en media página. La datación de 1988 midió una reparación por empalme: intervalo de la intervención 1290--1420, confianza 0,93, el corte de siete centímetros íntegramente dentro de la zona remendada. La formación de la imagen carece de mecanismo representado: once hipótesis evaluadas y descartadas por rasgo medido, clasificación fuera de corpus. La traza de paleta medieval, confianza 0,91, no discrimina entre ejecución y contacto: las copias se consagraban extendidas sobre el original, a veces durante horas. La muestra es inescrutable para todo método presente y para todo método previsible que dependa de una población de referencia, porque la población de referencia no existe.

El apartado histórico ocupaba tres páginas y era el que más le había costado, porque ahí el rigor trabajaba contra la mitad de los lectores posibles y él lo sabía. El memorial de d'Arcis quedaba degradado a indicio, con sus tres defectos en fila: borrador sin sello, confesión de oídas, litigio abierto. Y a continuación, con todo su peso, el argumento que había encontrado en noviembre, en su archivo de lectura, y que ningún dato del laboratorio había logrado mover: trece siglos sin documento. Ningún intervalo estrecho compensa un silencio de esa longitud, había escrito. La frase era suya, era correcta, y en un informe que dejaba la imagen sin mecanismo funcionaba como contrapeso exacto. Un lector honesto saldría de esas tres páginas sin saber qué creer. Ese era el diseño.

El apartado de anomalías tenía una sola entrada. Índice de convergencia 0,97, capa no auditable, herramientas de interpretación con retorno de ruido, ablación por canal documentada en anexo. Rótulo: \textit{señal no auditable --- no clasificada como ruido}. Y debajo, la fórmula de la página catorce del borrador, la que Ferrero había firmado como testigo estando en contra de que la línea existiera: la señal no se descarta como ruido de sistema; se asienta sin clasificar.

Pasó a la última página. El cierre era una línea de inventario, la más neutra del documento: constancia de restitución, fecha, hora, estado conforme al registro adjunto. La leyó completa por procedimiento, como leía todo lo que llevaba su firma, y en la penúltima línea encontró la palabra. \textit{La tela fue restituida el 29 de abril a las 21:46, dentro de la ventana prevista, sin novedad.} La tela. Su sistema de redacción usaba ``el objeto'' desde 2029; el protocolo interno de la operación usaba ``el lienzo'' o ``la unidad'' según la fase; ``la tela'' era vocabulario de Chiara, de dentro de fase, de cuando traducir costaba dos décimas de segundo. La palabra estaba en su letra, con su tinta, y él no recordaba haberla elegido. Sostuvo el lápiz sobre el renglón el tiempo de una corrección. Después lo apoyó en la mesita y miró por la ventanilla el muro del túnel, que a esa velocidad era una sola veta gris continua.

La dejó.

En Ginebra el auto de la Fondation esperaba donde Ansermet había dicho, y el chofer era el mismo de abril, y el trayecto hasta la ciudad vieja tomó los mismos minutos. El estudio olía a papel y a la cera del pasillo de abajo. El radiador crujió a los cuarenta segundos de sentarse; Vidal lo esperó y llegó puntual, como en enero. Sobre el escritorio sin computadora había tres cosas dispuestas en paralelo: una carpeta de tela gris, un recibo en dos ejemplares y una caja de madera clara del tamaño de un diccionario, abierta y vacía.

---Doctor ---dijo Ansermet---. La cláusula octava prevé la entrega en el domicilio de la Fondation de todos los soportes, claves y documentos comprendidos en la definición de archivo. Le propongo el orden siguiente: soportes físicos primero, claves segundo, informe tercero. Si prefiere otro orden, cualquier orden es válido; el acta se redacta al final y refleja el que usted haya elegido.

---Ese orden está bien.

Los soportes eran cuatro discos y entraron en la caja de a uno, con Ansermet leyendo en voz alta el número de serie de cada uno y Vidal confirmando contra su propia lista. El procedimiento tomó nueve minutos y fue indistinguible de cualquier cadena de custodia de su carrera, salvo por la dirección. Cuatrocientas veces había recibido material ajeno con esa liturgia de números cantados y confirmados. Recibir era lo que hacía. La palabra para lo de hoy era otra y el acta la usaba sin pudor: entrega.

Las claves iban en sobre lacrado, escritas por él esa madrugada con la letra más prolija de las cuarenta y una páginas. Las frases de acceso al canal, la del directorio 2033-04, la del consolidado. Ansermet tomó el sobre sin mirarlo y lo puso bajo la carpeta gris.

---El informe ---dijo.

Vidal lo sacó del bolso. Cuarenta y una páginas, cosidas con hilo esa mañana en Collegno con la cosedora que Ledda no había querido llevarse. Lo dejó sobre el escritorio y le pareció que pesaba menos afuera de la mano que adentro, lo cual era un error de percepción y lo registró como tal.

---La Fondation acusa recibo del archivo completo de los trabajos ---dijo Ansermet---, que queda en su custodia a perpetuidad conforme a la cláusula octava, con el régimen de autorización que usted negoció en enero. Sesenta días, denegación motivada, segunda presentación. Le corresponde firmar los dos ejemplares del recibo y quedarse con uno.

La lapicera era la misma de laca negra comida en la zona de agarre. Vidal firmó el primer ejemplar y, con la pluma sobre el segundo, se detuvo.

---Hay dos archivos que retengo ---dijo---. Los declaro. Un contenedor cifrado de marzo de 2029, denominado 2029-03, relativo a un trabajo anterior a todo vínculo con la Fondation. Y una planilla local de once kilobytes, de octubre de 2032, sin relación con los trabajos. Ninguno de los dos contiene datos del período doce al veintinueve de abril. Si la Fondation entiende que la definición de archivo los alcanza, quiero la interpretación por escrito antes de firmar.

Ansermet apoyó las dos manos en el borde del escritorio. Durante dos o tres segundos miró un punto de la pared, a la izquierda de Vidal, donde había un grabado de la rada de Ginebra, y la mirada tenía la calidad de una consulta, aunque en la sala eran dos y el grabado tenía cien años. Después volvió.

---Tengo instrucción previa sobre ese punto exacto ---dijo---. Textual: el fundador dice que lo que es suyo, es suyo. La instrucción es de enero, anterior a que yo supiera que usted retendría algo, y no distingue entre los dos archivos que usted declara ni exige que los declare. Constará en acta que los declaró de todos modos.

Vidal recontó la frase. De enero. Anterior a la declaración. Sin distinción entre los dos archivos, cuando uno de los dos ---el de octubre, la fila, los dos puntos abiertos--- era una entrada que ningún inventario del mundo debía conocer. La instrucción cubría por adelantado una excepción que él acababa de inventar, y la cubría con una fórmula de seis palabras que servía igual para un log de sensor que para dos horas y veintiocho minutos de una tarde en Jerusalén. Buscó la lectura mínima: cortesía contractual, propiedad intelectual previa, nada más. La lectura mínima no explicaba el ``es suyo'' dicho dos veces, con el verbo en presente, como una devolución.

Firmó el segundo ejemplar. Ansermet le entregó el suyo doblado en tres, con el pliegue hecho a regla.

---El automóvil lo lleva a Cologny ---dijo---. El fundador lo espera. Me pidió que le transmita que no hay orden del día.

---Todo lo que él transmite tiene orden del día.

---Es posible. Yo transmito. ---Ansermet se puso de pie y, por primera vez en cuatro reuniones, agregó una frase que ningún tribunal habría podido usar contra nadie y que sin embargo sonó fuera del registro notarial---: Fueron cuatro meses de trabajar con usted, doctor. En mi oficio eso se dice al final, porque antes contamina.

Le tendió la mano. Vidal la estrechó y salió con el recibo en el bolso, del lado donde durante un año habían viajado, sucesivamente, un expediente de catorce pestañas, un contrato de diecinueve cláusulas y un informe de cuarenta y una páginas. El compartimiento iba casi vacío y el bolso colgaba distinto, más liviano de un lado, y el cuerpo tardó dos cuadras en dejar de corregir un peso que ya no estaba.

El auto tomó la costa. El lago venía picado por un viento del norte que no llegaba a mover los árboles de la orilla, y Vidal pasó el trayecto con la vista en el agua, sin sacar el cuaderno naranja. En el campo de entregable de la entrada de hoy había escrito, a las cinco de la mañana, ``los motivos''. Ahora dudaba de la palabra. Los motivos eran lo que Sandoz había comprometido para el final; lo que Vidal llevaba, en cambio, lo sacó del bolso y lo pasó al bolsillo interior del saco, contra el pecho, del lado izquierdo: seis hojas manuscritas dobladas en dos, cada renglón con doble referencia, cuaderno y folio y cota de un lado, muestra y portaobjetos y medición del otro. La comparación de Chiara. Cuatro coincidencias sobre cinco. Un defecto de hilatura dibujado a lápiz en 2002 y medido en Collegno en 2033. El papel era común, de taller, y hacía un ruido mínimo contra el forro cuando él respiraba hondo, de modo que dejó de respirar hondo.

En el patio de Cologny la grava estaba rastrillada en líneas paralelas, con una interrupción de medio metro cerca del sendero sur donde alguien había pasado esa mañana y nadie había vuelto a rastrillar. Vidal la registró al cruzarla. En abril la grava había estado perfecta.

La puerta se abrió antes de que tocara. La mujer que lo guió por el pasillo de piedra clara caminaba más despacio que en abril, o él caminaba más rápido; olía a cera y, al fondo, a tierra húmeda y a hoja verde, el olor del lado sur. La puerta del jardín de invierno estaba entornada y por la ranura salía el aire tibio y cargado del vidrio al sol. Adentro, entre los paneles curvos, la gotera del techo marcaba su intervalo regular contra algo de metal, y entre dos goteos se oyó el agua sirviéndose en una taza. Una sola.

Vidal se llevó la mano al bolsillo interior, comprobó la arista de las seis hojas con dos dedos, y empujó la puerta.

\clearpage

\chapter{El jardín de invierno}

\lettrine[lines=2, lhang=0.1, nindent=0.2em]{E}{l} aire del jardín de invierno estaba dos o tres grados por encima del pasillo y olía a tierra regada esa mañana. Sandoz estaba de pie junto a la mesa de plantines, con el lápiz cruzado sobre una pila de dos libros donde en abril había habido tres, y la bandeja del té ya servida: una tetera, una taza. La segunda taza estaba boca abajo sobre su platito, igual que la otra vez, y Vidal la miró el tiempo que en abril se había ahorrado mirarla. La chaqueta de Sandoz era la misma. Le quedaba grande. Entre abril y hoy había perdido masa donde la ropa cara no disimula: el cuello, el dorso de las manos. La gotera del techo marcó su intervalo contra el balde de zinc y Vidal contó tres goteos antes de entender lo que llevaba dos meses delante de sus instrumentos. Un hombre que sirve a todos y nunca a sí mismo. En abril lo había archivado como un tic. Era una dieta. Era lo que queda cuando el cuerpo ya no admite lo que la mano ofrece.


---Doctor ---dijo Sandoz, y le llenó la taza---. El auto de Ansermet es puntual. El tren de la tarde hacia Zúrich, en cambio, se detiene en Nyon sin motivo publicado.

Vidal se sentó, sacó las seis hojas del bolsillo interior y las puso sobre la mesa, desplegadas, con las cotas hacia Sandoz.

---Antes de los motivos ---dijo---. Esto.

Sandoz las levantó con las dos manos y las leyó enteras. Tardó lo que se tarda: renglón por renglón, las dos columnas, la doble referencia. Diez minutos, quizá doce. En ese lapso la gotera cayó cuarenta y una veces y Vidal perdió la cuenta una sola vez, cuando Sandoz pasó de la cuarta hoja a la quinta sin cambiar de expresión. Al terminar, las alineó por el borde, las dobló por el pliegue original y las devolvió a la mesa, del lado de Vidal.

---Está bien hecha ---dijo---. Dígale a la señorita Fabbri que su maestra la habría firmado. Ottino tenía la mejor mano de su generación y la peor suerte institucional. Las dos cosas constan.

---Usted tiene los contenedores.

---Cuarenta y uno, cuarenta y dos y cuarenta y tres. Desde el año 2003. ---Sirvió un hilo más de té aunque la taza de Vidal seguía llena---. Pregunte en el orden que quiera. Hoy contesto entero.

Vidal preguntó en el orden en que auditaba: adquisición, cadena, uso, resultado.

---¿Por qué vía?

---La asociación de amigos del lienzo. La misma red que veintiocho años después canalizó la donación de la teca. La compré antes de saber para qué la iba a necesitar la segunda vez. Un año después de la intervención, el material estaba clasificado como conservación en curso, que es la categoría que nadie revisa porque revisar implicaría admitir que existe.

---¿Qué hizo con él?

---Medirlo. Veinte años, con lo mejor que hubo en cada momento. Laboratorios privados, de a un canal por vez, sin que ninguno supiera qué tenía en la mesa. ---Se sentó, despacio, administrando el movimiento---. Usted formuló mi resultado hace dos días, doctor, con mejor instrumental y la tela entera. Inescrutable. Yo lo tuve hace diez años, sobre fragmentos. Y un resultado sobre fragmentos es atacable. Contaminación, cadena rota, muestra no representativa. Lo sé con precisión, porque la demostración de que todo certificado es atacable la pagué yo.

La taza de Vidal estaba a medio camino de la boca. La devolvió al platito sin beber.

---Basilea.

---Basilea. ---Sandoz apoyó las manos sobre las rodillas, una sobre otra, manos de jardinero con la tierra todavía en las cutículas---. Usted sostiene que el falsificador entrenó treinta ídolos contra el modelo certificador y que el enigma es cómo obtuvo los pesos. Es la formulación correcta y es la que todos hicieron. La respuesta es que los pesos se compran cuando la infraestructura de datos donde residen te pertenece. El falsificador cobró su parte y desapareció porque desaparecer era su parte. Lo que se probó en Basilea fue una sola cosa: la certeza auditable no muere por sus datos. Muere por su custodia. El mejor certificado de Europa quedó muerto en una semana y ningún dato suyo era falso.

Tres goteos. Vidal se descubrió recontando la cadena. La denegación de diciembre citaba el reglamento de 2029. El reglamento de 2029 salió de Basilea. Basilea era de Sandoz. Tres eslabones. Los recorrió de nuevo. Cerraban.

---El reglamento me dejó afuera de Turín ---dijo Vidal. La voz le salió pareja; la boca, seca---. Tamiz es cerrado. Un modelo cerrado es un perito inexistente. La Custodia me denegó el acceso con un artículo que existe por algo que usted hizo.

---Y en enero le llegó una carta de Ginebra con la única puerta que quedaba abierta. ---Sandoz asintió una sola vez, con la cabeza grande y lenta---. Construí la puerta cerrada, doctor. Dos años más tarde construí la otra. Usted recorrió el pasillo midiendo cada baldosa y creyendo que elegía. Le concedo esto: dentro del pasillo, cada elección fue suya. El pasillo era mío.

Vidal contó las juntas del piso de piedra entre su silla y la mesa de plantines. Once. Las volvió a contar. Once. El número era estable y nada más lo era.

---Los motivos ---dijo---. Usted comprometió los motivos para cuando esto terminara.

---Terminó. ---Sandoz miró la taza boca abajo, un instante, y después a Vidal---. Hay dos hechos y ninguno pide adjetivos. El primero: me estoy muriendo. Los médicos me dieron un intervalo y es más estrecho que cualquiera de los suyos. Los meses cupieron en el calendario de la ostensión con margen razonable; ``cuando esto termine'' tuvo siempre dos lecturas y las dos eran verdaderas. El segundo hecho es de 1988. Mi madre financió una parte de aquella datación convencida de que la ciencia iba a devolverle confirmada la tela que amaba. El resultado se anunció en conferencia de prensa, en una tarde, sin custodia del relato, sin nadie del otro lado del anuncio. Usted leyó mi ejemplar del STURP. La doble raya de la página cuatrocientos ochenta y siete, junto al párrafo que recomendaba coordinar la difusión. La marqué en el ochenta y ocho, antes del anuncio, sin saber todavía que la discusión ya estaba perdida. Mi madre no perdió la fe, doctor; perdió la Iglesia, que es una amputación más lenta, y murió sin\ldots{} ---Se detuvo. Fue la primera frase que Vidal le oía dejar a medias, y duró lo que dura un goteo. Sandoz retomó con la sintaxis intacta---: Murió sin reconciliarse con nada. Un sí a medias o un no a medias, anunciados con su firma o con la mía, matan gente que espera del otro lado. Usted lo oyó de Ferrero. Yo lo leí en el ochenta y ocho, margen por margen.

---El archivo a perpetuidad.

---El archivo a perpetuidad. Yo no compré la respuesta, doctor. La respuesta, usted lo certificó, no existe en ninguna población de referencia. Compré la imposibilidad de que alguien la falsifique. Cada dato de esos diecisiete días, sellado con hash y hora, fuera del alcance de todo tribunal, toda diócesis y todo laboratorio con una tesis que defender. Una tumba para las certezas parciales. Es la obra de conservación más cara que se hizo nunca sobre esa tela, y la única que la tela necesitaba.

Vidal tenía una sola pregunta en pie y la hizo con las palabras contadas.

---¿Custodia o expropiación?

Sandoz sonrió por primera vez en la reunión, sin mostrar los dientes.

---Elija usted la palabra. Es lo único que no puedo comprarle.

La gotera siguió con su intervalo. Vidal miró los plantines: almácigos de algo de hoja ancha, demasiado tiernos para abril, sembrados por un hombre que había firmado un calendario donde él ya no figuraba después del verano. Aquello tampoco pedía adjetivos.

---Queda un punto ---dijo Sandoz---. Usted formuló mejor que nadie la objeción a mi arquitectura: un archivo custodiado por un solo hombre muere con el hombre. Es correcta y la respuesta es la sucesión. Le ofrezco la custodia del archivo a mi muerte. Acceso vitalicio, el corpus completo, los tres contenedores con su expediente, y la salida de la cuarta pasada con todo su contexto. El cero noventa y siete, doctor, con los cincuenta y ocho gigabytes alrededor. Todo lo que fue a buscar a Turín, entregado con recibo. Ansermet tiene el instrumento redactado desde marzo. Falta una firma.

Vidal escuchó el ofrecimiento entero sin mover las manos. Era exacto. Era, ítem por ítem, la lista de lo que había querido desde noviembre, y el hombre del otro lado de la mesa la había redactado antes de conocerle la cara. Esperó la señal interna de siempre, el tirón hacia el dato, la necesidad de tener la serie completa bajo llave propia. Registró, en su lugar, una calma que ya conocía: la de la última noche en Collegno, con un cable ajeno entre las manos.

---No ---dijo.

Sandoz esperó, como quien deja estabilizar una lectura.

---Usted colecciona la pregunta ---dijo Vidal---. Yo aprendí a devolverla.

---Formulado así, es un rechazo que no puedo mejorar. ---Sandoz inclinó la cabeza---. ¿Pide algo en su lugar? Me consta que usted audita hasta las despedidas.

---Dos cosas. El expediente de adquisición de 2003, completo, con fechas y firmas, entregado a Chiara Fabbri por vía que ella pueda mostrar. Bruna Ottino lleva treinta años lavando estandartes municipales con una sospecha encima que fabricó el silencio de otros. El polvo del aspirado faltaba porque usted lo tenía. Ella tiene derecho a la cadena de custodia de su propia inocencia, y a tenerla en vida.

---Concedido. Ansermet lo ejecuta esta semana. ¿La segunda?

---El cero noventa y siete queda en el archivo con su rótulo exacto. Señal no auditable, no clasificada como ruido. A perpetuidad, sin reclasificación posible, ni por usted ni por quien lo suceda. Si su tumba va a guardar algo, que guarde el estado verdadero del dato.

---Es la cláusula que yo habría escrito si usted no la escribía primero. ---Sandoz tomó la tetera y volvió a llenar la taza de Vidal, que estaba a dos tercios---. Concedido también. Tome el té, doctor. Es el último que le sirvo y es el de su casa de Zúrich; verifíquelo si quiere.

Vidal lo tomó. Era el de su casa.

Se levantó cuando la luz del lado sur empezó a bajar de ángulo. En la puerta del jardín de invierno se dio vuelta una vez: Sandoz seguía sentado entre sus plantines, con la taza boca abajo delante y las manos quietas, y levantó una de las dos, apenas, sin ceremonia. En el patio, la interrupción de la grava seguía sin rastrillar, medio metro de piedras revueltas cerca del sendero sur, y ahora Vidal supo leerla: el jardinero ya elegía qué partes del jardín iban a sobrevivirlo.

El auto lo dejó en la estación con veinte minutos de margen. Los usó de pie en el andén, sin contar nada.

En el tren sacó el cuaderno de tapas de cartón naranja y buscó la hoja de abril, la de la primera visita, donde había listado las afirmaciones de Sandoz una por línea con su casillero de verificación y había escrito en el renglón de síntesis: \textit{No dijo nada verificablemente falso.} Los casilleros seguían vacíos. Iban a seguir así: los motivos de un hombre no entregan espectro, y el único método disponible había sido sentarse enfrente y decidir. Escribió debajo, con la letra de las cadenas de custodia:

\textit{Elegí igual.}

Miró la frase el tiempo de tres estaciones de lago. Carecía de mecanismo, de intervalo y de testigo, y era la primera anotación de su vida que estaba completa así. Cerró el cuaderno sin releerla y dejó que el tren hiciera su trabajo, que era volver.

\clearpage

\chapter{Registro, no anuncio}

\lettrine[lines=2, lhang=0.1, nindent=0.2em]{L}{a} llamada entró el lunes a las 8:40, mientras Vidal cotejaba la cuenta del hotel contra su registro de gastos, renglón por renglón, aunque ya no había nada en ese gasto que necesitara esconderse. Reconoció el número. Atendió al segundo timbre.


---Le llegó a Bruna. Por escribano, esta mañana: fechas, cinco firmas, los tres contenedores con número de inventario, y una sola línea de carta, sin membrete: «Lo que es de ella, es de ella».

---¿Puede mostrarlo?

---Puede. Treinta años escribiendo cartas que nadie contestaba, y ahora puede.

---Entonces el papel está donde tenía que estar.

Cortó ella. Cuatro frases en cuarenta segundos. Vidal se quedó un momento con el teléfono contra la oreja, escuchando el tono vacío, y después anotó la llamada en el cuaderno de tapas naranjas: hora, duración, entregable cumplido. En el campo de observaciones dibujó el círculo de no requiere acción y comprobó que la mano, al dibujarlo, había hecho un círculo casi perfecto.

A mediodía lo esperaba Ferrero. El mensaje había llegado la noche anterior, por el conserje, en un sobre de la Comisión con dos líneas manuscritas: la puerta lateral, las doce, sin carpeta. Vidal caminó desde el hotel bajo los pórticos, entre andamios y vallas apiladas todavía sin armar, y llegó con tres minutos de margen. Ferrero salió a las doce y cuatro. Sin guantes. Las manos manchadas de reactivos colgaban de las mangas del saco como herramientas fuera de servicio, y la corbata de lana estaba derecha.

Caminaron hacia el río sin acordarlo.

---¿Cómo terminó? ---dijo Vidal.

---Terminó en un informe. Cuatro páginas, dos firmas. Bosio sostiene el asentamiento del soporte, y su explicación es la única que las tablas admiten sin forzarlas: yo mismo la verifiqué contra los nueve registros de servicio y cierra mejor que cualquier alternativa que a la Comisión se le pudiera ocurrir. Ceruti firmó al lado mío. Conservador, no policía; me lo dijo con esas palabras, mirando la tableta y no a mí, que es como se dicen las cosas que uno decidió antes de decirlas.

---¿Y sus verificaciones?

---Anexadas. Fechas, valores, instrumento, mi firma en cada hoja. Trabajo real, doctor; a usted no tengo que explicarle que la mejor coartada es la que además es cierta. ---Se detuvo en una esquina, esperó que pasara un tranvía---. Lo que el anexo no dice es por qué un miembro de la Comisión verificaba de noche lo que podía verificar de día. Esa pregunta quedó asentada en el expediente, sin respuesta, con mi nombre al lado. El pico también queda: once segundos fuera de banda, estabilizado al doce, doce de abril, tres y veintiuno de la madrugada. El que dentro de veinte años sepa qué buscar, lo va a encontrar, y va a encontrar mi anexo cosido atrás, y va a hacer la cuenta.

---El historial no se puede editar.

---Ni debe. ---Ferrero retomó el paso---. Quedó todo registrado y nada anunciado. La distinción la firmamos usted y yo en Collegno, con testigo; la Comisión la aplicó sin saber que existía. Es la clase de ironía que en mi familia se consideraba una forma de la misericordia.

Cruzaron frente a una iglesia con el atrio en obra. Ferrero miró los andamios el tiempo de dos pasos.

---Comuniqué mi retiro ---dijo---. Efectivo después de la ostensión. La ostensión abre el miércoles; entre una cosa y la otra hay una misa y una firma, y las dos me corresponden.

---¿Se lo pidieron?

---Nadie pide estas cosas. Se calculan. ---Se tocó la muñeca izquierda, un roce breve, y bajó la mano---. Cuarenta años preguntando quién paga si esto sale mal. Resulta que uno duerme mejor cuando la respuesta es «yo».

---¿Consecuencias?

---Consecuencias. Al contado. ---Y por primera vez en toda la operación la palabra sonó a factura saldada y a ninguna otra cosa.

En la esquina siguiente se separaron. Ferrero le tendió la mano desnuda, y antes de soltarla preguntó, con la voz de acta:

---¿Va a venir el miércoles?

---Sí.

Ferrero lo miró un segundo de más.

---Primera respuesta suya sin margen de error, doctor. Guárdela. Salen caras.


\scenebreak


El taller quedaba en el mismo polígono, a cuatrocientos metros de la nave, y la operación ya lo había usado dos veces: tratamiento térmico, hornos de recocido, un dueño con treinta años de facturas en regla y ninguna curiosidad. Vidal llegó a las cinco. Ledda estaba junto al horno grande con el cuaderno abierto y un cronómetro colgado del cuello, y sobre un carro de carga, plegada sobre sí misma en un paquete de un metro por medio, estaba la gemela.

Sin lo que Chiara le había descosido esa madrugada con tijeras de hoja corta, hasta las cuatro menos veinte, la tela pesaba menos de lo que la memoria de Vidal esperaba. Dos meses de telar en Bélgica. Ciento veinte horas de envejecimiento térmico. Cuatrocientas ochenta y dos perforaciones de aguja contadas una por una sobre las placas de Ferrero y reproducidas a razón de cuarenta por hora. Las quemaduras de 1532 ejecutadas con plata fundida y una curva de temperatura de tres páginas. Un objeto que había engañado a cuatro células de carga, a una sonda de humedad, a un sensor de oxígeno, a ocho observadores pagados en efectivo y a una ciudad entera durante diecisiete días, y que ahora ninguna colección del mundo podía recibir, porque su única función había sido existir sin que nadie supiera que existía.

---Ochocientos veinte grados ---dijo Ledda---. Ingreso a las 17:10. Estimado de consumo total, dieciocho a veintidós minutos. El acta la escribo durante.

La boca del horno se abrió con un golpe de aire caliente que Vidal sintió en la cara y en el dorso de las manos. El dueño del taller empujó el carro con una pértiga, la tela entró entera, y por la mirilla Vidal la vio encenderse de los bordes hacia adentro, rápido, con la llama clara de la celulosa limpia. Olía a papel quemado y, debajo, a algo apenas dulce, el apresto o el tinte de igualación. Las quemaduras falsas de 1532 ardieron por segunda vez, y esta vez hasta el final. A los veinte minutos exactos el horno contenía una capa de ceniza gris que el tiro deshacía en espiral.

Ledda terminó de escribir, subrayó el último renglón con la regla y le acercó la hoja: hora de ingreso, temperatura de cámara, duración, hora de cierre, dos firmas previstas.

---¿Destinatario? ---preguntó Vidal.

---Ninguno. Es un acta sin adónde ir. ---Ledda la firmó, esperó la firma de Vidal, la dobló en tres y se la guardó en el bolsillo interior---. La conservo yo. Después corre la misma suerte que la dirección.

Setenta y dos horas, contadas desde la noche del veintinueve. El plazo se cumplía esta misma tarde, y el resto ya estaba hecho: la sala blanca era otra vez un taller de mantenimiento electromecánico con dos camionetas rotuladas, la línea de nitrógeno un remito genérico, las once hojas del plan un rollo en una bolsa negra. La cinta azul de la sacristía a escala ya había salido en tiras enteras. Quedaban las llaves, que Ledda entregó al dueño del taller contándolas en voz alta, y quedaba el trayecto hasta los autos, cruzando el playón en diagonal, los dos callados.

Junto a su auto, Ledda se detuvo. Miró el polígono una vez, de punta a punta, como quien firma un inventario, y después miró a Vidal.

---Si alguna vez necesita una operación limpia, no me llame ---dijo---. Que no la necesite es mi mejor trabajo.

Era la primera vez que Vidal le oía reclamar algo como suyo: mi mejor trabajo, no el mejor. Le tendió la mano.

---Sin pendientes ---dijo.

Algo se movió en la cara de Ledda, una contracción de dos décimas que en otro hombre habría sido una sonrisa, y subió al auto.


\scenebreak


De noche, desde la ventana del hotel, se veía un sector de techos y, entre dos cúpulas, el resplandor intermitente de la plaza donde estaban montando las luces. Vidal dejó la ventana entreabierta. Subía el ruido de las cuadrillas, metal contra metal, las vallas encastrándose de a una en serpentinas todavía vacías.

El archivo completo del peritaje dormía en Ginebra. Cincuenta y ocho gigabytes de datos crudos, ochenta y dos barridos, cinco controles ciegos, un consolidado que decía muestra inescrutable, una salida de catorce kilobytes con un número que ninguna herramienta traducía, y un informe de cuarenta y una páginas cosido con hilo. Custodia perpetua. Ningún tribunal podía citarlo, ninguna iglesia reclamarlo, ningún laboratorio replicarlo, y esa triple imposibilidad era exactamente lo que su dueño había comprado, factura por factura, con veinte años de paciencia y una donación que nadie pudo permitirse rechazar. Un año atrás Vidal habría sabido clasificar eso sin levantar la vista del teclado: derrota del método, victoria de la administración, artefacto y se descarta.

Abrió el cuaderno naranja en una página nueva. Escribió la fecha. Debajo:

\textit{Archivo: custodia perpetua, Ginebra. Estado: fuera de alcance de todo peritaje, incluido el mío. Clasificación: sin datos.}

Y en el renglón siguiente, con la misma letra:

\textit{El punto anterior no impide dormir.}

Le puso el círculo al margen y dejó el cuaderno abierto sobre la mesa, para que la tinta secara.

En la plaza, un reflector se encendió, barrió el frente de la catedral de abajo hacia arriba y se apagó. Prueba de luces. Volvió a encenderse, quedó fijo. Las vallas seguían encastrándose a un ritmo de una cada pocos segundos, y Vidal notó que llevaba un rato oyéndolas sin llevar la cuenta. La fila todavía no existía. Faltaban dos días para que empezara a formarse, y él ya sabía dónde iba a pararse.

\clearpage

\chapter{Ostensión}

\lettrine[lines=2, lhang=0.1, nindent=0.2em]{L}{a} fila empezaba detrás del Palazzo Reale y daba dos vueltas antes de entrar a los jardines. Vidal se sumó a las 8:40 por el extremo abierto, donde una voluntaria con pechera amarilla repartía folletos en cuatro idiomas y decía a cada grupo, con la entonación gastada de la segunda hora, que la espera estimada era de tres horas. Tres horas y cuarto, corrigió otra voluntaria más adelante, sin que nadie se lo pidiera. Vidal aceptó el folleto con las dos manos y se quedó con él sin abrirlo.


Había una acreditación posible. Habría bastado un mensaje a Ferrero, o presentarse en la mesa de la Comisión con su nombre real, que figuraba en dos contratos y una nómina de personal autorizado. Eligió la puerta de las tres horas: bolso abierto para dos carabinieri con guantes azules, documento pasado por un lector de mano, y adentro del cordón un dron quieto a media altura que nadie miraba dos veces. La visita del papa era recién a la tarde, pero el dispositivo ya estaba completo desde temprano.

El sol de mayo caía sesgado sobre las vallas, las mismas que había oído encastrarse de noche desde la ventana del hotel. La fila avanzaba diez o doce pasos, se detenía un minuto, avanzaba otra vez. Olía a protector solar y a piedra caliente. Un chico de unos seis años, dos posiciones más adelante, contaba las vallas en voz alta y perdía la cuenta cada vez que la fila se movía, y volvía a empezar desde uno. Detrás de Vidal, una mujer mayor con un pañuelo al cuello lo estudió un momento.

---¿El del conteo es suyo? ---le preguntó, con el mentón hacia el chico.

---No ---dijo Vidal---. Pero cuenta bien.

La mujer sonrió y volvió a su rosario. La fila avanzó otro tramo. Por los altavoces, una voz pidió en italiano y luego en inglés que se apagaran los flashes dentro de la catedral.

Adentro, el aire cambiaba dos grados y el murmullo bajaba a la mitad, como si el edificio lo descontara en la puerta. El recorrido estaba marcado con cordones y desembocaba en la nave por el costado, y desde el fondo, antes de ver nada, Vidal ya sabía qué iba a ver: la teca elevada sobre el crucero, inclinada dieciocho grados, el cristal con tratamiento antirreflejo que la donación había especificado hasta el índice de refracción. Debajo del soporte, invisibles, las cuatro células de carga que él había medido con hoja de servicio de una empresa que existía sobre todo en papel. En algún registro que ya nadie consultaría, un pico de dos gramos y una décima, once segundos fuera de banda, archivado como asentamiento del soporte con la firma de un técnico que creía lo que firmaba. Encima de los instrumentos, la tela, iluminada por una luz rasante de diseño exacto, pagada por un hombre que se estaba muriendo en un jardín de invierno a doscientos cincuenta kilómetros de ahí.

La imagen se armaba a la distancia correcta. Vidal lo sabía por las placas, por el barrido de doce micrones, por la madrugada del día doce en la sala blanca, y aun así el efecto se produjo igual: a dos metros y medio del cordón, el rostro en negativo salió de la trama como sale un resultado de una consolidación, de golpe y entero, con los ojos cerrados y la asimetría del pómulo en su lugar. Alrededor, la gente se detenía los segundos que los cordones permitían. Algunos se persignaban. Una mujer levantó a un bebé dormido para que mirara algo que el bebé no miró.

Del otro lado de la barrera baja, en el sector reservado, Vidal vio a Ferrero.

Corbata de lana con treinta grados afuera. Estaba de pie junto a monseñor Ceruti, los dos inclinados sobre una tableta, y Ferrero señalaba algo en la pantalla con el índice, una curva o una tabla, y Ceruti asentía. Era el último acto público del cargo; el retiro corría con efecto posterior a la ostensión. En un momento Ferrero le devolvió la tableta a Ceruti, cruzó por delante de la teca camino al otro extremo del sector, y al pasar redujo el paso.

Vidal contó. Seis segundos con los dedos de la mano derecha a la altura del cristal, a dos o tres centímetros, sin tocarlo. Manos desnudas. Ningún guante asomaba del bolsillo del saco, ni de algodón ni de nitrilo, y las manos manchadas de reactivos quedaron ahí, suspendidas cerca de la superficie, el tiempo justo de una verificación que ya no medía nada. Después siguió caminando. En ningún momento miró hacia la fila.

---Rasante. ¿Cuarenta? ---dijo Chiara.

Estaba a su lado. Vidal no la había vuelto a ver desde Collegno, y ahora estaba a su izquierda en la fila como si acabara de volver de buscar una herramienta, el pelo atado con una banda elástica azul, sin explicación de por dónde había entrado ni desde cuándo caminaba junto a él. La fila avanzó un tramo y ella avanzó con él.

---Cuarenta ---confirmó Vidal---. La misma que usábamos con la lámpara.

---Por eso. Alguien especificó bien. ---Chiara miró la teca entrecerrando un poco los ojos, con la cabeza apenas ladeada, el gesto de taller de evaluar una tensión---. ¿Viste cómo cae?

---Está plana. Las dieciséis zonas.

---No te pregunto eso. Te pregunto cómo cae.

Vidal miró. Debajo del cristal, sobre el soporte inclinado, la tela descansaba con ese peso repartido que ninguna sarga nueva imitaba, seiscientos años de plegados durmiendo en los mismos lugares, y la caída era la de siempre, la de las placas, la del inventario al revés sobre el papel del día dos, con la lámpara rasante a cuarenta grados y ochenta y una líneas.

---La tela cae bien ---dijo.

Chiara tardó un segundo de más en contestar. Después dijo:

---Sí. Cae bien.

Avanzaron juntos el tramo siguiente.

---La cota del borde inferior no se aprecia desde acá ---dijo Chiara.

---Con esta carga de gente, la banda de la sonda debe estar comida entera.

---Bosio va a tener trabajo a la noche.

Ninguno de los dos dijo abril. Frente al punto más cercano del cordón, donde el recorrido concedía su pausa más larga, Chiara juntó las manos sin cruzarlas del todo, bajó la cabeza y rezó dos frases en italiano, la voz por debajo del murmullo, sin mirar a nadie. La segunda frase terminó donde ahora terminaba. Él no. Se quedó de pie a su lado, derecho, con el folleto sin abrir en la mano, mirando el rostro que se armaba y se desarmaba según los cuerpos que pasaban por delante. Ninguno de los dos corrigió al otro.

Después el cordón los soltó hacia la salida lateral y el rostro quedó atrás.

En el atrio, el sol pegaba de frente. La plaza estaba llena hasta las vallas del fondo y la fila seguía llegando desde los jardines, y por los altavoces la misma voz repetía la espera estimada, que había subido a tres horas y media. Vidal se detuvo un momento al borde de la escalinata. Dos manos habían devuelto esa tela a su sitio la noche del veintinueve. Las dos estaban paradas en ese escalón. Buscó, por hábito, la columna donde asentar el hecho. No había columna. Se miró las manos. Contó los escalones que le faltaban para bajar del atrio. Seis. Los volvió a contar. Seis.

---¿Tenés tren hoy? ---preguntó Chiara.

---Mañana a las siete.

---Entonces caminá. Yo voy para Porta Palazzo.

Bajaron la escalinata contra el sentido de la fila. A la altura de la segunda serpentina de vallas, el hábito se disparó solo, treinta años de reflejo: la fila, el número, el conteo. Jerusalén, siete meses atrás, cuarenta y una personas, la columna de la rotonda, la segunda pasada para descartar el sesgo de anclaje. Vidal fue a buscar la cifra de esta fila y encontró el lugar vacío. En algún punto de la mañana la serie se había interrumpido, o directamente nunca había empezado; revisó hacia atrás y la última cuenta completa que pudo ubicar era la del chico de seis años, que era de otro y estaba mal. Se quedó un instante con esa comprobación en la mano, esperando el impulso de reiniciar desde uno.

El impulso llegó al ritmo de siempre. Lo dejó pasar.

Chiara se detuvo y lo miró.

---¿Qué?

---Nada ---dijo Vidal---. Estaba contando.

---¿Cuánto te dio?

---Nada. Perdí la cuenta hace rato.

Chiara se rió, corta, y le señaló con el mentón el final de la fila, allá al fondo, donde los últimos llegaban a preguntarle a la voluntaria de amarillo cuánto faltaba.

---Tres horas ---dijo---. No van a ver nada que no esté en las placas de Ferrero.

Vidal caminó unos pasos antes de contestar.

---Depende de qué llames nada ---dijo.

Ella no lo corrigió. Cruzaron la plaza entre la gente, hacia el mercado, dos técnicos sin acreditación entre miles de personas que esperaban al sol, y ninguna de todas ellas, hasta donde Vidal podía medir, estaba cometiendo un error de método.

\clearpage

\chapter{Márgenes}

\lettrine[lines=2, lhang=0.1, nindent=0.2em]{L}{a} esquela salió un jueves, en la página de avisos de la prensa financiera de Ginebra, entre la fusión de dos aseguradoras y un cambio de directorio. Tres líneas. «Emeric Sandoz. Fondation Cassiodore. Los servicios se celebraron en la más estricta intimidad.» La edad y la causa habían quedado afuera; fotografía tampoco había. Vidal la leyó dos veces y le contó las palabras: trece. Después buscó la fecha de la edición y calculó hacia atrás. Los médicos habían dado un intervalo que cabía en el calendario de la ostensión; el hombre lo había excedido en dos semanas y media. Un intervalo bien construido admite ese margen. Igual lo verificó de nuevo, edición en mano, contra el calendario de la pared.


Se hizo té. El de siempre, la misma casa, el mismo grado de tueste. A los dos tercios de la taza se descubrió detenido junto a la ventana, midiendo un nivel que ya nadie iba a completar. Lo terminó de pie. Abajo, el seis dobló la curva de la Universitätstrasse con su chirrido.

La carta de Ansermet llegó el lunes, en el papel y la sintaxis habituales: ninguna frase que un tribunal pudiera usar, ninguna que dejara un hueco. La estructura de custodia perpetua había quedado constituida el veinticuatro de mayo. Las dos condiciones de Vidal figuraban en el instrumento con rango de cláusulas fundacionales: la entrega a Turín, ejecutada y documentada; el rótulo de la señal, intocable a perpetuidad, sin reclasificación posible por sucesor alguno. «El resto de las disposiciones del fundador no le concierne», cerraba el párrafo. «Acompaño la única que sí.»

El sobre interior traía una fotocopia. Era la última página de su propio informe, la del apartado de anomalías con su única entrada: \textit{señal no auditable --- no clasificada como ruido}. Y al margen, a lápiz, con mina blanda, la letra que Vidal conocía de quinientas páginas de actas: la mano que corregía gradientes de chamuscado, que renumeraba hipótesis diez años antes que la literatura, que discutía con los autores y solía ganar. Dos palabras, en español esta vez.

\textit{Bien conservado.}

Vidal cotejó la inclinación de los trazos contra las anotaciones del libro, por oficio, porque cotejar era gratis. La misma mano. Después recorrió de memoria los márgenes del ejemplar: en quinientas y pico de páginas, el lápiz objetaba, precisaba, tachaba, exigía. En ningún margen concedía. Esta era la única vez, y había elegido su idioma para hacerlo. Abrió las actas del STURP en la página 487, donde la doble raya vertical seguía marcada hasta el reverso de la hoja, y deslizó la fotocopia contra esa página. El libro cerró con un milímetro más de lomo. Lo devolvió al cajón, debajo de los manuales de calibración.


\scenebreak


El addendum le llevó una noche, y la noche empezó con once minutos de reconstrucción del respaldo distribuido. El contenedor 2029-03 se abrió como siempre; lo que cambió fue lo que Vidal hizo con él. Leyó el log completo una vez más, de la línea uno a la última, y después redactó en el formato del repositorio donde los datos crudos de Amberes vivían desde 2029, públicos, citados, intactos.

«Addendum al conjunto de datos, panel flamenco, 2029. Líneas 3.407 a 3.418: señal discordante en la capa de imprimación del panel izquierdo, peso agregado 0,04. Clasificada en origen como ruido de sensor por decisión del operador, sin verificación de aislamiento ni repetición en frío. La verificación ya no es posible: el instrumento fue recalibrado en 2030. Reclasificación: señal no auditable, no clasificada como ruido. La conclusión sobre el panel se sostiene por seis vías independientes posteriores a la publicación, ninguna de las cuales pertenece al firmante. La clasificación original estaba fuera de método. Fecha. Firma.»

Lo leyó una vez. Era lo contrario exacto de lo que había hecho a las 2:31 del 14 de marzo de 2029, con las mismas doce líneas y la misma mano. Depositó el archivo. El repositorio devolvió un identificador y una hora, y a partir de esa hora el asunto dejó de ser suyo.

Las consecuencias tardaron nueve días en organizarse. El editor de la revista que había publicado Amberes pidió una nota aclaratoria; el estudio que llevaba el juicio del museo pidió copia certificada, y Vidal mandó las dos cosas el mismo día. La carta de Halevi llegó al décimo, manuscrita, con el membrete de su cátedra.

«Doctor Vidal: leí su addendum. Hace ocho meses, en Jerusalén, mencioné los bronces etruscos por elegancia, sin nombrarlos. Se los nombro ahora: los suyos. Un hombre que a los veintinueve firmó sin pedir los datos crudos, y que diecisiete años después confiesa haber archivado una señal a mano, tiene una pregunta que responder, y la pregunta es mía, de los periodistas que se ocupe su repositorio. ¿Cuántas veces más decidió su mano lo que firmó su máquina? Le pido la lista. Si la lista está vacía, dígalo bajo su firma, que es lo único suyo que todavía cotiza. La dalet estaba bien leída; eso lo mantengo donde lo dije. Hundirlo no me interesa. Cuarenta años de ojo también son un archivo, doctor, y nadie me pidió nunca el mío.»

Vidal la leyó dos veces. La segunda vez se quedó en la última frase más de lo que la frase pedía. Después contestó, a mano, tres oraciones:

«La clasificación de 2029 fue decisión mía, tomada sin las verificaciones que exijo a terceros, y el addendum contiene todo lo que sé del caso. Otros registros míos con decisión de operador sin verificación posterior son dos, ambos de esta primavera, ambos depositados en una custodia que ya no me pertenece y con un rótulo que no se puede reclasificar. Su ojo leyó la dalet antes que mi máquina, y ese dato figura en mi acta de Jerusalén con su apellido.»

Firmó y despachó. Defensa, cero líneas. Halevi contestaría o haría circular la carta o las dos cosas; el costo estaba calculado desde antes de depositar el addendum, con margen, y hasta ahora corría dentro del intervalo previsto. En la planilla de asignación, la noche del addendum ocupaba una celda con código propio, creado esa misma semana. Restitución. Entregable: ninguno.


\scenebreak


La carta de Chiara llegó a fin de mes, en papel de cuaderno cuadriculado, doblada en tres, con la letra ancha de quien escribe poco a mano y no le importa.

«Sebastián: el círculo del Opificio recibió el expediente. No hubo prensa, y estuvo bien que no hubiera: Bruna no necesitaba un titular, necesitaba que dos o tres personas concretas leyeran cinco firmas concretas. Las leyeron. La semana pasada la invitaron a dar el seminario de fibras de otoño, el que dio quince años seguidos hasta 2002. Dijo que sí, y después me corrigió una cota de memoria, así que está bien. Los estandartes municipales los va a seguir lavando, dice, porque ahora los lava por elección y eso cambia la química del agua. Es un chiste de ella. Creo.

»Del resto: la teca marca 51 \% y la ciudad se está olvidando de la fila, que es lo que las ciudades hacen mejor. Ferrero me saluda cuando nos cruzamos y una vez hasta me preguntó por vos, con esa sintaxis de acta que tiene, ``el doctor Vidal, ¿ha retomado actividad pericial?''. Le dije que sí. No me desmientas.

»Ya no tengo cables que enrollar. Preguntá igual.

»C.»

Vidal la leyó de pie, junto a la ventana. Después la leyó otra vez, sentado. La frase final quedaba fuera de todo formato conocido: era una autorización sin objeto, sin plazo, sin condición de ejercicio. La archivó en el cajón de abajo, con la correspondencia que se contesta, y le puso fecha de respuesta: antes del viernes. Escribió el borrador esa misma noche. Salió corto, con una sola pregunta, y la pregunta era por Bruna y no por la tela. Le pareció el orden correcto.


\scenebreak


El sábado ordenó el estudio. Los sábados de orden existían en su calendario desde que Zúrich era Zúrich, cuatro horas con código, entregable «inventario conforme», y esta vez la palabra inventario le rindió más que de costumbre porque el inventario había cambiado de composición.

El disco de marzo de 2029, el que decía «soporte, misceláneo», ya tenía los campos completos. Procedencia: propia. Contenido: log de Amberes, versión íntegra, sin cifrar. Disposición final: repositorio público, depositado, identificador anotado. La entrada había salido del amarillo por primera vez en cuatro años, y el cajón donde vivía pesaba lo mismo y contenía menos deuda.

El libro del STURP tenía ficha nueva. Procedencia: donación del fundador de la Fondation Cassiodore, marzo de 2033. Estado: anotado a lápiz por mano identificada, tres épocas de mina más una. Disposición final: se conserva. En el campo de observaciones Vidal había escrito «contiene concesión única, p. 487 y hoja suelta» y le había puesto el círculo de no requiere acción, porque era exacto: la página fotocopiada iba a estar ahí dentro de treinta años y ninguna acción del mundo la mejoraba.

Las cinco camisas grises colgaban en el placard en su orden de rotación. La primera todavía tenía, contra la luz, el tono apenas corrido que le había dejado el polvo de Horvat Adullam; en octubre iba a hacer un año de eso. Vidal la corrió al final de la fila y la fila quedó igual de gris, con una variación por debajo de cualquier umbral que valiera la pena declarar.

Quedaba la planilla. Once kilobytes, carpeta local de 2029, fuera del alcance del sincronizador de las 20:00, que esa noche pasó a su hora y la salteó como la había salteado dos mil veces. La columna seguía llamándose tiempo perdido. Adentro había una fila de octubre, dos horas y veintiocho minutos, un círculo al margen y una nota de teléfono en otra parte, con dos puntos abiertos desde hacía nueve meses.

Vidal apoyó el cursor sobre el archivo y midió el gesto que seguía. Renombrar la columna eran cuatro segundos. Completar la nota, quizá dos líneas. Los dos gestos estaban disponibles, auditados, sin costo técnico, y ninguno de los dos era todavía verdadero. Una reclasificación antes de tiempo es la forma prolija del falso positivo; eso lo había escrito él, en un informe de 2031, sobre el trabajo de otro.

Retiró la mano. Apagó la pantalla. En la cocina quedaba té para una taza y la hizo durar hasta que el seis pasó dos veces más allá abajo, con su curva y su chirrido, cada siete minutos, verificable. La planilla quedó donde estaba, con su columna mal llamada y su fila sin cerrar. Mantenida.

El lunes, entre el correo de la mañana, llegó un sobre con sellos de Israel y el remitente escrito a máquina: Autoridad de Antigüedades, distrito de Judea.

\clearpage

\chapter{La fila}

\lettrine[lines=2, lhang=0.1, nindent=0.2em]{L}{a} segunda lasca de Horvat Adullam salió del tiesto a las 9:31 y pesó setecientos cuarenta microgramos. Vidal la vio entrar al vial ---HA-33-047, rotulado a mano--- y firmó la cadena de custodia con la letra de siempre, mayúsculas alineadas, sobre la misma mesa del mismo sótano de Har Hotzvim donde un año antes había firmado la primera. Las baldosas grises de treinta por treinta seguían ahí. Trescientas quince, con margen de dos: el dato tenía doce meses de vigencia y esta vez alcanzó con saberlo.


El lote nuevo eran once fragmentos de la campaña de primavera, y uno casaba con el borde inferior del óstracon: la misma pasta, la fractura complementaria, y dos letras más de la segunda línea. Be'eri lo había anunciado como la confirmación que faltaba. Segev lo había anunciado como la prueba de que Be'eri anunciaba demasiado. El sobre de la Autoridad había llegado a Zúrich en junio; el contrato, en agosto; el fragmento, ahora, bajo la lupa binocular, con su etiqueta de excavación intacta y una foto de contexto que esta vez incluía la escala en el encuadre. Alguien había aprendido.

Vidal llevaba la tercera de las cinco camisas grises. El polvo del sitio ya le había corrido el tono a otra de la rotación, como el año anterior, y la variación seguía por debajo de cualquier umbral que valiera la pena declarar.

---Cuarenta y siete series, todas conformes ---dijo Ruti, y le acercó el manifiesto. Esperó la firma y, con el papel ya firmado, agregó---: Mi abuela pregunta si ahora sí es David.

---Es un intervalo. El mismo del año pasado, más angosto.

---Se lo voy a decir así. Ella entiende de intervalos: tiene noventa y uno.

Tamiz corrió la pasada consolidada en dieciséis horas, con dos controles ciegos intercalados que Vidal ubicó en la cola sin mirar dónde. El junte era auténtico: misma vasija, misma cocción, continuidad de pátina a través de la fractura. La datación cerró en 1000 a 930 antes de la era común, confianza 0,95, un intervalo adentro del anterior. Y el canal epigráfico, con las dos letras nuevas, subió la hipótesis del antropónimo regio de 0,61 a 0,74.

El umbral firmado por las dos comisiones seguía siendo 0,90.

Be'eri sostuvo en la reunión que 0,74 con sintaxis de cancillería valía una atribución, que ningún escriba de aldea escribía así, que el país entero esperaba. Segev sostuvo que 0,74 era la palabra ``casi'' con decimales, y que ella llevaba veinte años publicando contra el casi. Los dos tenían el tono de quien repite un argumento que ya perdió una vez y piensa perderlo mejor. Ashkenazi levantó acta, recordó el umbral, cerró la sesión a las 12:15. Después le pidió a Vidal que se quedara.

La oficina daba a un estacionamiento. Ashkenazi apoyó las dos manos sobre el escritorio, sin sentarse.

---Doce personas van a leer su informe, doctor. El intervalo lo va a leer el mundo. Y la gente que espera afuera de un anuncio no vive en un intervalo. En algún punto, todos eligen.

---Lo sé. ---Vidal cerró la carpeta que tenía sobre las rodillas---. Por eso entrego el intervalo: para que la elección sea de ellos y no mía. Eso es lo que vendo.

Ashkenazi lo miró dos segundos completos.

---El año pasado me habría contestado que el intervalo no es una opinión.

---Le contesté eso.

---Sí. ---Firmó la recepción del preliminar sin releerla---. Por eso me acuerdo.

Lo acompañó hasta la puerta y ahí, con la mano ya en el picaporte, asintió una sola vez, despacio, con el gesto de quien da por buena una lectura después de mucho servicio del instrumento, y pasó a lo siguiente.


\scenebreak


El jueves, víspera del vuelo, quedó como una celda en blanco en la planilla de asignación. Vidal la dejó así. Era la tercera desde 2029 y la primera que quedaba vacía a propósito, aunque el archivo carecía de campo para registrar el propósito y él dejó también eso en paz.

Entró a la ciudad vieja por la puerta de Damasco, que quedaba a trasmano del hotel, y bajó por el mercado sin plan de recorrido. El año anterior una caminata así le había producido veintiséis entradas de observaciones, cotas de hiladas, tiempos de marcha. Esta vez el teléfono fue apagado en el bolsillo, un peso constante contra el muslo, y las hiladas pasaron sin medirse. Olía a cardamomo y a fritura, y más abajo, donde el pavimento se volvía piedra pulida por pies, a piedra mojada aunque hiciera semanas que el cielo estaba seco. Un carro de reparto bajó los escalones a los saltos y todo el mundo se apretó contra los puestos con la coordinación de una maniobra ensayada mil años.

Llegó al patio del Sepulcro por el camino largo, dando un rodeo por el barrio cristiano que un topógrafo habría tachado entero.

La fila salía del portal, cruzaba el patio en diagonal y se doblaba contra el muro. Avanzaba, a ojo, un metro cada tres minutos. Adentro, decía el panel de la entrada en tres idiomas, no había nada.

Se puso al final. No la contó.

El sol daba en la mitad del patio y la fila entraba y salía de la sombra según avanzaba. Delante de él, una mujer con un bolso lleno de velas envueltas en papel de diario le explicaba algo a una adolescente que miraba el techo. Más adelante, un grupo con gorras iguales. El hombre inmediatamente delante se dio vuelta a medias y señaló la cúpula con el mentón.

---First time? ---preguntó.

---Second ---dijo Vidal.

El hombre asintió como si eso explicara algo y volvió a su teléfono.

Adentro, la rotonda tenía el mismo aire espeso de cera y piedra fría, y las lámparas colgaban a la misma altura sin fecha. Vidal avanzó con la fila, un metro, una espera, otro metro. En algún momento pasó junto a la columna donde un año atrás había apoyado la espalda para contar mejor, y la columna quedó atrás como cualquier otra.

El monje del edículo administraba el ingreso con dos dedos en alto y una consigna repetida en inglés: two minutes. Hizo pasar a la mujer de las velas con la adolescente, y a Vidal detrás.

Adentro estaba la antecámara, y después el hueco bajo, y la losa. Lo mismo que la primera vez. Exactamente lo mismo: la piedra gastada por manos, las lámparas, el espacio para tres personas donde apenas cabían tres personas, y en el centro de todo eso, nada. Eso ya no era un dato que faltaba.

La mujer apoyó la frente contra la losa. La adolescente miró a Vidal, incómoda, y Vidal miró la piedra. Las manos le quedaron quietas a los costados. Ninguna serie se puso en marcha. Duró lo que duró, y el golpe en la puerta llegó a los dos minutos, como estaba anunciado, y los tres salieron agachando la cabeza bajo el dintel.

Afuera, en los candelabros de arena, ardían velas a tres alturas. Vidal compró una en el puesto del patio interior, la inclinó contra otra encendida hasta que la mecha tomó, y la plantó derecha entre las demás. La cera tardó en agarrar. Esperó a que agarrara.

Después cruzó el patio contra la fila, que seguía avanzando a su metro cada tres minutos, con la regularidad de las cosas que nadie administra, y salió a la calle por donde había venido, o por otra; a esa altura del día ya daba igual.


\scenebreak


El hotel de la calle Jaffa era el mismo del año anterior, dos pisos más arriba. Vidal cenó algo en el cuarto, lavó la taza en el baño y la dejó boca abajo sobre una servilleta. Por la ventana, abajo, pasó el tranvía con su campanilla, lleno de gente a las nueve de la noche.

Abrió la computadora. La carpeta local de 2029 seguía donde siempre, fuera del alcance del sincronizador, que había pasado a las 20:00 y la había salteado una vez más. Adentro, la planilla: once kilobytes. La columna todavía se llamaba tiempo perdido y tenía una sola entrada, del octubre anterior, con su círculo al margen.

Escribió la segunda entrada debajo de la primera, en la columna de siempre: la fecha, basílica del Santo Sepulcro, 15:20--17:55. Revisó la hora de inicio contra su memoria, una sola vez, y la dio por buena.

Después seleccionó el título de la columna y lo borró.

El encabezado quedó en blanco. Las dos entradas quedaron: fecha, lugar, duración. La de arriba con su círculo, la de abajo sin nada. Una columna sin nombre encima de dos filas ciertas, en un archivo que pesaba ahora un poco más de once kilobytes y que conservó el nombre que tenía, porque cualquier nombre nuevo habría sido una clasificación y la clasificación seguía sin existir.

Guardó. Cerró.

En la nota del teléfono, los dos puntos de octubre seguían abiertos. Los dejó abiertos también; ya eran menos una deuda que una entrada del inventario, procedencia propia, disposición final: se conserva.

Abajo pasó otro tranvía, en el sentido contrario. Vidal se quedó junto a la ventana con la luz apagada, mirando la calle, y dejó pasar el siguiente sin medir el intervalo, y el de después, hasta que la calle estuvo un rato larga y vacía y él seguía ahí, de pie, sin nada pendiente entre las manos.
